% !TEX root = main.tex
\documentclass[12pt,a4paper]{ctexbook}
\usepackage{geometry}
\usepackage{tikz}
\usepackage{xcolor}
\usepackage{amsmath}
\usepackage{fancyhdr}
\usepackage{titlesec}
\usepackage{amsmath,amssymb,amsfonts}
\usepackage{geometry}
\usepackage{esint}
\usepackage{lastpage}
\usepackage{datetime2}
\usepackage{tcolorbox}
\usepackage{hyperref}
\usepackage{tikz-3dplot}
\usepackage{pgfplots}
\usetikzlibrary{calc}
\usetikzlibrary{matrix,3d}
\usetikzlibrary{tikzmark,backgrounds}
\usetikzlibrary{shapes.geometric}
\pgfdeclarelayer{background}
\pgfsetlayers{background,main}
\pgfplotsset{compat=1.13}
\DTMsetdatestyle{iso}
\setcounter{chapter}{-1}
\tcbuselibrary{theorems, skins, breakable}
\everymath{\displaystyle}
\geometry{margin=2.5cm}

% 定义统一的方框样式
\newtcbtheorem[number within=chapter]{definition}{定义}{%
    enhanced,
    colback=blue!5,      % 背景浅蓝色
    colframe=blue!75!black,
    coltitle=white,
    fonttitle=\bfseries,
    title={#2},          % 方框标题
    attach boxed title to top left={yshift=-2mm, xshift=5mm},
    boxed title style={sharp corners, rounded corners, colback=blue!75!black},
    separator sign none, % 不显示默认的冒号
    before skip=10pt,
    after skip=10pt,
}{def}

\newtcbtheorem[number within=chapter]{theorem}{定理}{%
    enhanced,
    colback=green!5,     % 背景浅绿色
    colframe=green!50!black,
    coltitle=white,
    fonttitle=\bfseries,
    title={#2},
    attach boxed title to top left={yshift=-2mm, xshift=5mm},
    boxed title style={sharp corners, rounded corners, colback=green!50!black},
    separator sign none,
    before skip=10pt,
    after skip=10pt,
    fontupper=\itshape
}{thm}

% 手动计数器，随章编号
\newcounter{myexample}[chapter]
\renewcommand{\themyexample}{\thechapter.\arabic{myexample}}

% 例题环境：椭圆标签在左，与文字同行
\newenvironment{myexample}
  {%
    \par\noindent                % 新段落，无缩进
    \refstepcounter{myexample}   % 计数器加1
    \tikz[baseline=(ell.base)]   % 对齐椭圆文字基线
      \node[ellipse, draw=black!40, fill=black!8,
            inner sep=1pt, font=\small,
            anchor=base] (ell) {例\themyexample};%
    \quad                        % 留一点间距
  }
  {%
    \par                         % 结束段落
  }


% 封面配色
\definecolor{titleblue}{HTML}{1B3A5C}
\definecolor{gold}{HTML}{C5A54E}

% 正文章节标题格式
\titleformat{\section}{\Large\bfseries\color{titleblue}\centering}{\thesection}{1em}{}
\titleformat{\subsection}{\large\bfseries\color{titleblue!80}}{\thesubsection}{1em}{}


\begin{document}

% ================= 封面页 =================
\thispagestyle{empty}
\pagestyle{empty}

% ---------- 背景水印装饰 ----------
\begin{tikzpicture}[remember picture, overlay]
  % 中央大积分号
  \node[opacity=0.06, scale=7] at (current page.center) {$\displaystyle\iint_Df(x,y)dxdy$};

  % 散布的小积分号（使用 shift 精确定位，无 calc 库也可）
  \foreach \x/\y/\sym in {
      0.2/0.8/\frac{\partial y}{\partial x},
      0.8/0.8/\iiint_\Omega,
      0.15/0.15/\oint_L,
      0.85/0.15/\oiint,
      0.5/0.85/\int_a^b,
      0.5/0.2/\lim_{x \to \infty}f(x)
  } {
    \node[opacity=0.07, scale=4]
      at ([shift={(\x\paperwidth, \y\paperheight)}]current page.south west)
      {$\displaystyle\sym$};
  }

  % 顶部金色装饰线
  \draw[line width=2pt, gold]
    ([yshift=-1.5cm]current page.north west) --
    ([yshift=-1.5cm]current page.north east);
  % 底部金色装饰线
  \draw[line width=1.5pt, gold]
    ([yshift=1.5cm]current page.south west) --
    ([yshift=1.5cm]current page.south east);
\end{tikzpicture}

\vspace*{-1cm}
\begin{center}
  % 中文主标题
  {\fontsize{48}{52}\selectfont\bfseries\color{titleblue} 从极限

  到广义斯托克斯定理} \\[0.6cm]

  % 英文标题
  {\fontsize{22}{26}\selectfont\bfseries\color{titleblue} From Limits

  to the Generalized Stokes' Theorem} \\[1cm]

  % 分隔装饰
  \textcolor{gold}{\rule{0.6\textwidth}{1.5pt}} \\[1.5cm]

  % 副标题
  {\Large\color{titleblue} \pageref{LastPage} 页了解经典微积分} \\[1.5cm]

  % 日期
  {\large\color{titleblue} \DTMnow}
\end{center}

\pagebreak

\begin{quote}\footnotesize
	\thispagestyle{empty}
		Copyright \copyright{} {\the\year} From Limits to the Generalized Stokes' Theorem

		Permission is hereby granted, free of charge, to any person obtaining a copy of this document and associated documentation files (the "Document"), to deal in the Document without restriction, including without limitation the rights to use, copy, modify, merge, publish, distribute, sublicense, and/or sell copies of the Document, and to permit persons to whom the Document is furnished to do so, subject to the following conditions:

		\vspace{1ex}
		The above copyright notice and this permission notice shall be included in all copies or substantial portions of the Document.
		\vspace{1ex}

		THE DOCUMENT IS PROVIDED "AS IS", WITHOUT WARRANTY OF ANY KIND, EXPRESS OR IMPLIED, INCLUDING BUT NOT LIMITED TO THE WARRANTIES OF MERCHANTABILITY, FITNESS FOR A PARTICULAR PURPOSE AND NONINFRINGEMENT. IN NO EVENT SHALL THE AUTHORS OR COPYRIGHT HOLDERS BE LIABLE FOR ANY CLAIM, DAMAGES OR OTHER LIABILITY, WHETHER IN AN ACTION OF CONTRACT, TORT OR OTHERWISE, ARISING FROM, OUT OF OR IN CONNECTION WITH THE DOCUMENT OR THE USE OR OTHER DEALINGS IN THE DOCUMENT.
	\end{quote}


% ================= 正文区域 =================
\newpage
\pagestyle{fancy}
\fancyhf{}
\setlength{\headheight}{14pt}
\fancyhead[LE,RO]{\color{titleblue!70}\hyperlink{toc}{从极限到广义斯托克斯定理}}
\fancyhead[RE,LO]{\color{titleblue!70}\leftmark}
\fancyfoot[C]{\thepage}
\renewcommand{\headrulewidth}{0.4pt}
\renewcommand{\headrule}{\hrule width\headwidth height\headrulewidth \vskip-2pt {\color{gold}\hrule width\headwidth height 0.5pt}}


\frontmatter
\chapter{前言}

本书是为了让已经学过经典微积分学的人，或者将要学习经典微积分学的人，例如高三或者大一学生，用最短的时间和最少的额外工具，把微积分基本定理、格林公式、高斯公式、斯托克斯公式统一成一个公式——广义斯托克斯定理。全书不引入复杂的线性代数概念，所有推导都用极限、偏导数、二阶三阶行列式和向量运算完成。

本书假定读者具备高中数学的基本知识，尤其是导数和向量的直观概念与初步计算。前置的线性代数知识包含在预备章中。

本书的结构是直线型的。每一章都在为下一章做铺垫，最后一章是全书的终点，前面所有的内容都是为了它服务的。你可以选择性地跳过某些章节，但建议按章节顺序阅读，以获得最佳阅读体验。相对于同济版高等数学，本书着重于说清楚黎曼积分的计算，而不是说清楚它的概念和证明。对那些内容感兴趣的读者，书中也会提供一些相关的讨论和说明，但不会过于深入。因此本书会忽略大部分细节，只讲最核心的主干。

由于精力有限，本书可能存在一些瑕疵，欢迎各位读者不吝赐教，批评指正。

以上，让我们开始吧。


\tableofcontents
\hypertarget{toc}{}


\mainmatter

% !TEX root = ../main.tex
\chapter{前置知识}
在学习微积分之前，我们需要先掌握几个工具。它们分别是行列式，向量，空间解析几何。有了这些，我们才能开启微积分的大门。我们只作最简单的介绍，掌握运算即可，不涉及过多的理论和证明。

\section{行列式}
\textbf{行列式}不仅是判断线性方程组有唯一解的关键量，也是后续计算向量叉乘的基础。本章将从二阶行列式出发，并将其推广到三阶。我们主要学习对角线法则和展开降阶两种计算方法，介绍行列式最基本的几条性质，如转置不变，交换变号等，并展示如何利用这些性质简化计算。本章内容限于二阶与三阶行列式，高阶行列式以及其他内容不作讨论。

\subsection{二阶行列式}

\begin{definition}{二阶行列式}{}
  由四个数排成两行两列，记作
  \[
  A = \begin{vmatrix}
    a & b \\
    c & d
  \end{vmatrix} = ad - bc
  \]
\end{definition}

我们可以这样理解：行列式 $A$ 的值等于主对角线（红色）的积减去副对角线（蓝色）的积。简记为“左上$\times$右下$-$右上$\times$左下”。

\[
A = \begin{vmatrix}
    \tikzmarknode{a}{a} & \tikzmarknode{b}{b} \\
    \tikzmarknode{c}{c} & \tikzmarknode{d}{d}
\end{vmatrix}
= \textcolor{red}{ad} - \textcolor{blue}{bc}
\begin{tikzpicture}[remember picture, overlay]
  \scoped[on background layer]{
    \draw[red, thick, opacity=0.5] (a.center) -- (d.center);
    \draw[blue, thick, opacity=0.5] (b.center) -- (c.center);
  }
\end{tikzpicture}
\]

\begin{myexample}
举例说明：设 $A = \begin{vmatrix}
    1 & 2 \\
    3 & 4
\end{vmatrix}$，求$A$。
$$
A = \begin{vmatrix}
    \tikzmarknode{a}{1} & \tikzmarknode{b}{2} \\
    \tikzmarknode{c}{3} & \tikzmarknode{d}{4}
\end{vmatrix}
= \textcolor{red}{1 \cdot 4} - \textcolor{blue}{2 \cdot 3} =\textcolor{red}{4}-\textcolor{blue}{6}=-2
\begin{tikzpicture}[remember picture, overlay]
  \scoped[on background layer]{
    \draw[red, thick, opacity=0.5] (a.center) -- (d.center);
    \draw[blue, thick, opacity=0.5] (b.center) -- (c.center);
  }
\end{tikzpicture}
$$
\end{myexample}

这就是二阶行列式的定义及其计算方法。

\subsection{三阶行列式}
与二阶行列式类似，我们给出三阶行列式的定义：
\begin{definition}{三阶行列式}{}
  由九个数排成三行三列，记作
  $$
  A = \begin{vmatrix}
    a & b & c \\
    d & e & f \\
    g & h & i
  \end{vmatrix} = aei + bfg + cdh - ceg - bdi - afh
  $$
\end{definition}
我们可以这样理解：行列式 $A$ 的值等于三条主对角线（红色）的积之和减去三条副对角线（蓝色）的积之和。

\[
A = \begin{vmatrix}
    \tikzmarknode{a}{a} & \tikzmarknode{b}{b} & \tikzmarknode{c}{c} \\
    \tikzmarknode{d}{d} & \tikzmarknode{e}{e} & \tikzmarknode{f}{f} \\
    \tikzmarknode{g}{g} & \tikzmarknode{h}{h} & \tikzmarknode{i}{i}
\end{vmatrix}
= \textcolor{red}{aei} + \textcolor{red}{bfg} + \textcolor{red}{cdh} - \textcolor{blue}{ceg} - \textcolor{blue}{afh} - \textcolor{blue}{bdi}
\begin{tikzpicture}[remember picture, overlay]
  \scoped[on background layer]{
    % 三条主对角线（红色，半透明）
    \draw[red, thick, opacity=0.5]
      (a.center) -- (e.center) -- (i.center);
    \draw[red, thick, opacity=0.5]
      (b.center) -- (f.center) -- (g.center);
    \draw[red, thick, opacity=0.5]
      (c.center) -- (d.center) -- (h.center);
    % 三条副对角线（蓝色，半透明）
    \draw[blue, thick, opacity=0.5]
      (c.center) -- (e.center) -- (g.center);
    \draw[blue, thick, opacity=0.5]
      (a.center) -- (f.center) -- (h.center);
    \draw[blue, thick, opacity=0.5]
      (b.center) -- (d.center) -- (i.center);
  }
\end{tikzpicture}
\]

显然，这种方法是非常难记住的，因此我们有一个相对简单的方法来计算三阶行列式，叫做展开降阶法：
$$A = a \begin{vmatrix}
    e & f \\
    h & i
  \end{vmatrix} - b \begin{vmatrix}
    d & f \\
    g & i
  \end{vmatrix} + c \begin{vmatrix}
    d & e \\
    g & h
  \end{vmatrix}$$

其中，$a$ 后面的二阶行列式叫做 $a$ 的余子式，$b$ 后面的二阶行列式叫做 $b$ 的余子式，$c$ 后面的二阶行列式叫做 $c$ 的余子式。我们可以选择任意一行或者一列来展开，结果都是一样的。

怎么理解呢？我们还是先看一下这个三阶行列式：
\[
A = \begin{vmatrix}
    a & b & c \\
    d & e & f \\
    g & h & i
  \end{vmatrix}
\]

先固定$a$所在的行和列不动，暂时忽略。剩下的数看作一个二阶行列式（ $a$ 的余子式），计算这个二阶行列式，结果再与$a$相乘。

\[
A = \begin{vmatrix}
    a & \blacksquare & \blacksquare \\
    \blacksquare & e & f \\
    \blacksquare & h & i
  \end{vmatrix}
  = a \begin{vmatrix}
    e & f \\
    h & i
  \end{vmatrix}
  \cdots
\]

同理，再固定$b$所在的行和列不动，暂时忽略。剩下的数看作一个二阶行列式（ $b$ 的余子式），计算这个二阶行列式，结果再与$b$相乘。\textbf{注意：}这一项的符号是负的。

\[
A = \begin{vmatrix}
    \blacksquare & b & \blacksquare \\
    d & \blacksquare & f \\
    g & \blacksquare & i
  \end{vmatrix}
  = a \begin{vmatrix}
    e & f \\
    h & i
  \end{vmatrix}
  - b \begin{vmatrix}
    d & f \\
    g & i
    \end{vmatrix}
    \cdots
\]

最后，固定$c$所在的行和列不动，暂时忽略。剩下的数看作一个二阶行列式（ $c$ 的余子式），计算这个二阶行列式，结果再与$c$相乘。\textbf{注意：}这一项的符号是正的。

\[
A = \begin{vmatrix}
    \blacksquare & \blacksquare & c \\
    d & e & \blacksquare \\
    g & h & \blacksquare
  \end{vmatrix}
  = a \begin{vmatrix}
    e & f \\
    h & i
  \end{vmatrix}
  - b \begin{vmatrix}
    d & f \\
    g & i
    \end{vmatrix}
  + c \begin{vmatrix}
    d & e \\
    g & h
  \end{vmatrix}
\]

由此，我们得到了三阶行列式的展开降阶法：
\begin{theorem}{三阶行列式的展开降阶法}{}
  $$A =  \begin{vmatrix}
    a & b & c \\
    d & e & f \\
    g & h & i
  \end{vmatrix} =
  a \begin{vmatrix}
    e & f \\
    h & i
  \end{vmatrix} - b \begin{vmatrix}
    d & f \\
    g & i
  \end{vmatrix} + c \begin{vmatrix}
    d & e \\
    g & h
  \end{vmatrix}$$
\end{theorem}

\begin{myexample}
  举例说明：设 $A = \begin{vmatrix}
    1 & 2 & 3 \\
    4 & 5 & 6 \\
    7 & 8 & 9
  \end{vmatrix}$，求$A$。

  \begin{align*}
    A &= 1 \cdot \begin{vmatrix} 5 & 6 \\ 8 & 9 \end{vmatrix}
        - 2 \cdot \begin{vmatrix} 4 & 6 \\ 7 & 9 \end{vmatrix}
        + 3 \cdot \begin{vmatrix} 4 & 5 \\ 7 & 8 \end{vmatrix} \\
      &= 1 \cdot (5\cdot 9 - 6\cdot 8) - 2 \cdot (4\cdot 9 - 6\cdot 7) + 3 \cdot (4\cdot 8 - 5\cdot 7) \\
      &= 1 \cdot (45 - 48) - 2 \cdot (36 - 42) + 3 \cdot (32 - 35) \\
      &= 1 \cdot (-3) - 2 \cdot (-6) + 3 \cdot (-3) \\
      &= -3 + 12 - 9 \\
      &= 0
  \end{align*}
\end{myexample}

实际上，由降阶法，我们可以计算更高阶的行列式，但本章不作讨论。

\subsection{行列式的简单性质}
行列式有很多性质，这里我们只列举几个最基本的性质，计算的时候可能会用到它们。

\begin{theorem}{行列转置}{}
\[
\begin{vmatrix}
a_1 & a_2 & a_3 \\
b_1 & b_2 & b_3 \\
c_1 & c_2 & c_3
\end{vmatrix}
=
\begin{vmatrix}
a_1 & b_1 & c_1 \\
a_2 & b_2 & c_2 \\
a_3 & b_3 & c_3
\end{vmatrix}
\]
\textbf{推论}：若对行有成立的性质，对列也必成立。因为转置后行列式的值不变，这意味着\textbf{行与列的地位完全对称}。
\end{theorem}

\begin{theorem}{交换变号}{}
\[
\begin{vmatrix}
a_1 & a_2 & a_3 \\
b_1 & b_2 & b_3 \\
c_1 & c_2 & c_3
\end{vmatrix}
= -
\begin{vmatrix}
b_1 & b_2 & b_3 \\
a_1 & a_2 & a_3 \\
c_1 & c_2 & c_3
\end{vmatrix}
\]
\textbf{推论：}若有两行（列）相同，行列式为$0$。因为交换相同的两行，行列式既要等于自己，又要变号，只能是$0$。
\end{theorem}

\begin{theorem}{提取公因子}{}
\[
\begin{vmatrix}
ka_1 & ka_2 & ka_3 \\
b_1 & b_2 & b_3 \\
c_1 & c_2 & c_3
\end{vmatrix}
= k
\begin{vmatrix}
a_1 & a_2 & a_3 \\
b_1 & b_2 & b_3 \\
c_1 & c_2 & c_3
\end{vmatrix}
\]
\textbf{推论：}若某一行（列）全为$0$，行列式为$0$（相当于提出因子$0$）。若两行（列）成比例，提出比例因子后剩下两行相同，行列式也为$0$。
\end{theorem}

这里我们不给出例子。具体计算的时候自然会用到。

\section{空间向量}
这个章节中，我们假定读者有基础的高中知识，已经学习过平面向量的数乘，点积，坐标表示等内容，因此我们直接从向量的叉积开始。

\subsection{向量的叉积}
\begin{definition}{向量的叉积}{}
设 $\vec{a} = (a_1, a_2, a_3)$ 和 $\vec{b} = (b_1, b_2, b_3)$ 是三维空间中的两个向量，夹角为 $\theta (0 \leq \theta \leq \pi)$，记它们的叉积（又称向量积）为$\vec{c}$，则$\vec{c}$定义为：
\[\vec{c}=
\vec{a} \times \vec{b}
= \begin{vmatrix}
\vec{i} & \vec{j} & \vec{k} \\
a_1 & a_2 & a_3 \\
b_1 & b_2 & b_3
\end{vmatrix}
= (a_2 b_3 - a_3 b_2,\; a_3 b_1 - a_1 b_3,\; a_1 b_2 - a_2 b_1)
\]
其中 $\vec{i}, \vec{j}, \vec{k}$ 分别是沿 $x,y,z$ 轴正方向的单位向量。

$\vec{c}$的模长为$|\vec{c}|=|\vec{a}||\vec{b}|\sin\theta$，方向垂直于 $\vec{a}$ 和 $\vec{b}$ 所在的平面。具体指向由右手定则确定。
\end{definition}

\begin{definition}{右手定则}{}
  伸出右手，四指指向第一个向量 $\vec{a}$ 的方向，然后向第二个向量 $\vec{b}$沿较小的夹角（$\leq \pi$） 弯曲握拢，此时大拇指所指的方向就是 $\vec{a} \times \vec{b}$ 的方向。
\end{definition}

应当注意，向量的叉积是一个向量，而不是一个数。叉积的模长等于由这两个向量构成的平行四边形的面积。特殊地，若$\vec{a} \times \vec{b} = \vec{0}$，则表示两向量共线。

\begin{myexample}
  举例说明：设 $\vec{a} = (0, 1, 0)$ 和 $\vec{b} = (0, 0, 1)$，记 $\vec{c} = \vec{a} \times \vec{b}$，求 $\vec{c}$ 的方向。

  我们画一个草图展示 $\vec{a}$（红色）、$\vec{b}$（蓝色）以及它们的叉积 $\vec{c}$（绿色）的关系。

    \begin{center}
    \begin{tikzpicture}[
      scale=1.2,
      x={(-0.5cm,-0.5cm)},
      y={(1cm,0cm)},
      z={(0cm,1cm)},
      >=stealth
    ]
      % 坐标轴
      \draw[->] (-0.5,0,0) -- (2.5,0,0) node[below right] {$x$};
      \draw[->] (0,-0.5,0) -- (0,2.5,0) node[above left] {$y$};
      \draw[->] (0,0,0.5) -- (0,0,2.5) node[above] {$z$};

      \draw[red, very thick, ->] (0,0,0) -- (0,1.8,0) node[below right] {$\vec{a}=(0,1,0)$};

      \draw[blue, very thick, ->] (0,0,0) -- (0,0,1.8) node[above right] {$\vec{b}=(0,0,1)$};

      \draw[green!60!black, very thick, ->] (0,0,0) -- (1.8,0,0) node[below right] {$\vec{c}=(1,0,0)$};

      % 右手定则弧线
      \begin{scope}[canvas is yz plane at x=0]
        \draw[->, gray, thick]
          (0.6,0) arc[start angle=0, end angle=90, radius=0.6cm]
          node[midway, above right] {右手四指};
      \end{scope}
    \end{tikzpicture}
  \end{center}

  从图中可看出，$\vec{a}$ 沿 $y$ 轴，$\vec{b}$ 沿 $z$ 轴，右手四指从 $y$ 转向 $z$ 时，大拇指指向 $x$ 轴正方向，即 $\vec{c}$ 的方向。
\end{myexample}

我们很快发现，如果令$\vec{c} = \vec{b} \times \vec{a}$，此时$\vec{c}$ 的方向与 $\vec{a} \times \vec{b}$ 的方向相反。这是因为交换了 $\vec{a}$ 和 $\vec{b}$ 的位置，右手四指从 $z$ 转向 $y$ 时，大拇指指向 $x$ 轴负方向。由此我们可以得出叉积的一条重要性质：
\begin{theorem}{叉积的反对称性}{}
\[
\vec{a} \times \vec{b} = -(\vec{b} \times \vec{a})
\]
\end{theorem}

\begin{myexample}
  举例说明：设 $\vec{a} = (1, 2, 3)$ 和 $\vec{b} = (4, 5, 6)$，求 $\vec{a} \times \vec{b}$和$\vec{b} \times \vec{a}$。

  \begin{align*}
    \vec{a} \times \vec{b}
    &= \begin{vmatrix}
       \vec{i} & \vec{j} & \vec{k} \\
       1 & 2 & 3 \\
       4 & 5 & 6
       \end{vmatrix} \\
    &= \vec{i} \begin{vmatrix} 2 & 3 \\ 5 & 6 \end{vmatrix}
       - \vec{j} \begin{vmatrix} 1 & 3 \\ 4 & 6 \end{vmatrix}
       + \vec{k} \begin{vmatrix} 1 & 2 \\ 4 & 5 \end{vmatrix} \\
    &= \vec{i}(2\cdot6 - 3\cdot5) - \vec{j}(1\cdot6 - 3\cdot4) + \vec{k}(1\cdot5 - 2\cdot4) \\
    &= \vec{i}(12 - 15) - \vec{j}(6 - 12) + \vec{k}(5 - 8) \\
    &= -3\vec{i} + 6\vec{j} - 3\vec{k} \\
    &= (-3,\; 6,\; -3)
    \end{align*}

    \begin{align*}
    \vec{b} \times \vec{a}
    &= \begin{vmatrix}
       \vec{i} & \vec{j} & \vec{k} \\
       4 & 5 & 6 \\
       1 & 2 & 3
       \end{vmatrix} \\
    &= \vec{i} \begin{vmatrix} 5 & 6 \\ 2 & 3 \end{vmatrix}
       - \vec{j} \begin{vmatrix} 4 & 6 \\ 1 & 3 \end{vmatrix}
       + \vec{k} \begin{vmatrix} 4 & 5 \\ 1 & 2 \end{vmatrix} \\
    &= \vec{i}(5\cdot3 - 6\cdot2) - \vec{j}(4\cdot3 - 6\cdot1) + \vec{k}(4\cdot2 - 5\cdot1) \\
    &= \vec{i}(15 - 12) - \vec{j}(12 - 6) + \vec{k}(8 - 5) \\
    &= 3\vec{i} - 6\vec{j} + 3\vec{k} \\
    &= (3,\; -6,\; 3)
    \end{align*}
\end{myexample}

计算过程中，根据行列式的交换变号性，我们也能看出叉积的反对称性。

\subsection{向量的混合积}
向量的混合积并不常用，我们简单带过。

\begin{definition}{向量的混合积}{}
设 $\vec{a} = (a_1, a_2, a_3)$，$\vec{b} = (b_1, b_2, b_3)$ 和 $\vec{c} = (c_1, c_2, c_3)$ 是三维空间中的三个向量，记它们的混合积为 $[\vec{a}, \vec{b}, \vec{c}]$，则 $[\vec{a}, \vec{b}, \vec{c}]$ 定义为：
\[ [\vec{a}, \vec{b}, \vec{c}] = \vec{a} \cdot (\vec{b} \times \vec{c}) = \begin{vmatrix}
a_1 & a_2 & a_3 \\
b_1 & b_2 & b_3 \\
c_1 & c_2 & c_3
\end{vmatrix} \]
\end{definition}

应当注意，向量的混合积是一个数，而不是一个向量。混合积的绝对值等于由这三个向量构成的平行六面体的体积。特殊地，若 $[\vec{a}, \vec{b}, \vec{c}] = 0$，则表示这三个向量共面。

\section{空间解析几何}
这一节开始难度会升高。我们将介绍空间中的曲线和曲面。我们依旧假定读者有基础的高中知识，已经学习过平面解析几何和空间立体几何的基本内容，因此我们直接从空间解析几何的内容开始。

\subsection{平面方程}

在空间中，确定一个平面最自然的方式是：给定平面上的一个点，以及垂直于平面的方向，这个方向由一个法向量给出。这种确定平面的方法只需要一个点和一个法向量。因此，我们可以给出这样的定义：

\begin{definition}{平面的点法式方程}{}
设 $\Sigma$ 是空间中的一个平面，$P_0(x_0,y_0,z_0)$ 是 $\Sigma$ 上的一个点，$\vec{n} = (A,B,C)$ 是 $\Sigma$ 的一个法向量，则 $\Sigma$ 的平面方程为：
$$A(x-x_0)+B(y-y_0)+C(z-z_0)=0$$
\end{definition}

如果将点法式方程展开，写成$Ax + By + Cz + D = 0$的形式，我们便得到了一般式方程：

\begin{definition}{平面的一般式方程}{}
设 $\Sigma$ 是空间中的一个平面，则 $\Sigma$ 的一般式方程为：
$$Ax + By + Cz + D = 0$$
\end{definition}

由一般式可以看出，$\vec{n} = (A,B,C)$仍然是平面的一个法向量。因此，确定了法向量，就确定了平面的朝向。

\begin{myexample}
举例说明：求过点$P_0=(1,2,3)$且法向量为$\vec{n}=(4,5,6)$的平面方程。

$P_0=(1,2,3)$，$\vec{n}=(4,5,6)$，我们得到：
\[
4(x-1)+5(y-2)+6(z-3)=0
\]

因此，平面 $\Sigma$ 的方程为：
\[
4x+5y+6z-32=0
\]
\end{myexample}

\begin{myexample}
举例说明：已知$A(1,1,1),B(1,2,3),C(3,3,3)$，求过线段$AB,BC$的平面方程。

先计算向量$\overrightarrow{AB},\overrightarrow{BC}$：
\[
\overrightarrow{AB} = (1-1,\,2-1,\,3-1) = (0,\,1,\,2)
\]
\[
\overrightarrow{BC} = (3-1,\,3-2,\,3-3) = (2,\,1,\,0)
\]

两向量不平行，因此 $A,B,C$ 三点不共线，可唯一确定一个平面。

平面的法向量 $\vec{n}$ 可用向量叉积求得：
\[
\vec{n} = \overrightarrow{AB} \times \overrightarrow{BC}
= \begin{vmatrix}
\vec{i} & \vec{j} & \vec{k} \\
0 & 1 & 2 \\
2 & 1 & 0
\end{vmatrix}
= \vec{i}(1\cdot0 - 2\cdot1) - \vec{j}(0\cdot0 - 2\cdot2) + \vec{k}(0\cdot1 - 1\cdot2)
= (-2,\,4,\,-2)
\]

取简化形式 $\vec{n} = (1,\,-2,\,1)$。

由点法式，代入点 $A(1,1,1)$ 得：
\[
1\cdot(x-1) + (-2)\cdot(y-1) + 1\cdot(z-1) = 0
\]

展开并整理：
\[
x - 2y + z = 0
\]

\begin{center}
  \begin{tikzpicture}[scale=1.4, line join=round]
  \tdplotsetmaincoords{60}{110}
  \begin{scope}[tdplot_main_coords]

    % 坐标轴（略大于三角形范围）
    \draw[->] (0,0,0) -- (3,0,0) node[above] {$x$};
    \draw[->] (0,0,0) -- (0,3,0) node[above] {$y$};
    \draw[->] (0,0,0) -- (0,0,3) node[left] {$z$};

    % 三角形 ABC 填充
    \fill[blue!15, opacity=0.7] (1,1,1) -- (1,2,3) -- (3,3,3) -- cycle;
    \draw[thick] (1,1,1) -- (1,2,3) -- (3,3,3) -- cycle;

    % 三个顶点
    \filldraw (1,1,1) circle (2pt) node[below left]  {$A(1,1,1)$};
    \filldraw (1,2,3) circle (2pt) node[above left]  {$B(1,2,3)$};
    \filldraw (3,3,3) circle (2pt) node[right]        {$C(3,3,3)$};

    % 平面方程
    \node at (3,4,4) [below] {$x-2y+z=0$};

  \end{scope}
\end{tikzpicture}
\end{center}

\end{myexample}

\subsection{空间直线方程}
空间中的直线可以看作两个不平行的平面的交线，因此我们可以用两个平面方程联立来表示一条直线。但更方便的方式是使用点向式方程。

\begin{definition}{直线的方向向量}{}
在空间（或平面）中，如果一条直线$L$已经确定，那么任何一个与$L$平行的非零向量，都称为这条直线的一个方向向量。
\end{definition}

显然，一条直线有无数个方向向量。我们可以通过一点和一个方向向量确定一条空间直线。


\begin{definition}{直线的点向式方程}{}
设直线 $L$ 经过点 $P_0(x_0,y_0,z_0)$，且方向向量为 $\vec{s}=(m,n,p)$（其中 $m,n,p$ 不全为零），则 $L$ 的方程为：
\[
\frac{x-x_0}{m} = \frac{y-y_0}{n} = \frac{z-z_0}{p}
\]
特别的，如果某个分母为零，则相应的分子也为零，表示该坐标恒定。例如$m=0$时，直线方程变为：

$$L: \begin{cases}\frac{y-y_0}{n} = \frac{z-z_0}{p}\\x=x_0\end{cases}$$

此时直线垂直于$x$轴。

\end{definition}

因此，已知直线上一点以及直线的一个方向向量，就可以求出空间直线的方程。已知直线的点向式方程，我们可以读出直线上一点与直线的一个方向向量。

\begin{myexample}
  举例说明：已知直线$L$过$P_0(1,2,3)$，其中一个方向向量为$\vec{s}=(4,5,6)$，求$L$的方程。

  由点向式方程可得：
  \[
  \frac{x-1}{4} = \frac{y-2}{5} = \frac{z-3}{6}
  \]
\end{myexample}

由点向式方程我们可以看出，$\frac{x-x_0}{m} , \frac{y-y_0}{n} , \frac{z-z_0}{p}$这三个分数都相等，那么我们可以引入一个参数$t$（$t\in\mathbb{R}$），令$\frac{x-x_0}{m} = \frac{y-y_0}{n} = \frac{z-z_0}{p} = t$，将$m,n,p$都乘到等式右边，则有：

$$\begin{cases}
  x-x_0=mt\\
  y-y_0=nt\\
  z-z_0=pt
\end{cases}$$

这就得到了直线的参数式方程：

\begin{definition}{直线的参数式方程}{}
设直线 $L$ 经过点 $P_0(x_0,y_0,z_0)$，且方向向量为 $\vec{s}=(m,n,p)$（其中 $m,n,p$ 不全为零），则 $L$ 的方程为：
$$L: \begin{cases}
  x=mt+x_0\\
  y=nt+y_0\\
  z=pt+z_0
\end{cases}$$

其中$t\in\mathbb{R}$。

特别的，如果方向向量某个分量为零，则该坐标恒定。例如$m=0$时，直线方程变为：

$$L: \begin{cases}
  x=x_0\\
  y=nt+y_0\\
  z=pt+z_0
\end{cases}$$

此时直线垂直于$x$轴。
\end{definition}

\begin{myexample}
  举例说明：已知空间直线 \(L\) 经过两点 \(A(2, -1, 5)\) 和 \(B(8, -1, 1)\)，求直线 \(L\) 的点向式方程和参数式方程。

  先求\(\overrightarrow{AB}\)：\(\overrightarrow{AB} = (6,\ 0,\ -4)\)。取简化形式 \(\overrightarrow{AB} = (3,0,-2)\)。

  则已知 \(A(2, -1, 5)\)为直线上一点，\(\vec{s} = (3,0,-2)\) 为直线的一个方向向量。

  由点向式方程可得：

  \[
  L: \begin{cases}
  \frac{x-2}{3} = \frac{z-5}{-2}\\
  y = -1\end{cases}
  \]

  由参数式方程可得：

  \[
  L: \begin{cases}
    x = 2 + 3t \\
    y = -1 \\
    z = 5 - 2t
  \end{cases}
  \]

\end{myexample}

这个例子也说明了两点确定一条直线。

\subsection{曲面方程}
类比平面中的曲线方程，我们可以推断曲面方程也由类似的结构组成。本节中仅讨论二次曲面，即方程中坐标的最高次数是$2$的曲面。因此，本节需要一定的解析几何中二次曲线的基础。

面对一个复杂的二次曲面方程，直接想象它的空间形状并不容易。我们通常采用截痕法讨论：

\begin{definition}{截痕法}{}
用一族\textbf{相互平行}的平面与曲面相交，得到一系列平面截线。通过分析这些截线连续变化时的形状、大小、位置和变化趋势，从而分析出该曲面的整体空间形态与结构特征。
\end{definition}

怎么理解呢？假设你现在是脑科医生，为病人做了一个CT检查。你将CT检查得到的图片堆叠起来，终于把病人的脑子看清楚了——CT检查和截痕法的底层逻辑是一样的，二次曲面就像病人的脑子。CT检查是一层层扫描得到图片，然后再拼起来，截痕法是用平行平面一层层与曲面相交，得到一条条交线，然后再将这些交线拼起来，就能大致分析出曲面的形状。

\begin{myexample}
  举例说明：描述曲面$\Sigma: x^2+y^2+z^2=9$的形状。

  我们先用平面$z=0$截这个曲面，那么交线就必须同时满足$z=0$和$x^2+y^2+z^2=9$这两个方程。我们把$z=0$带进$\Sigma$的方程里，得到$x^2+y^2=9$，也就是说：

  在$z=0$时，$\Sigma$是以原点为圆心，半径为$3$的圆。

  我们再用平行的平面$z=1$截这个曲面，那么交线就必须同时满足$z=1$和$x^2+y^2+z^2=9$这两个方程。我们再把$z=1$带进$\Sigma$的方程里，得到$x^2+y^2=8$，也就是说：

  在$z=1$时，$\Sigma$是以$(0,0,1)$为圆心，半径为$2\sqrt{2}$的圆。

  同理：在$z=2$时，$\Sigma$是以$(0,0,2)$为圆心，半径为$\sqrt{5}$的圆。

  在$z=3$时，$\Sigma$是一个点，点的坐标为$(0,0,3)$。

  在$z=4$时，$\Sigma$不存在。

  同理，在$z$轴下方，我们能得到：

  在$z=-1$时，$\Sigma$是以$(0,0,-1)$为圆心，半径为$2\sqrt{2}$的圆。

  在$z=-2$时，$\Sigma$是以$(0,0,-2)$为圆心，半径为$\sqrt{5}$的圆。

  在$z=-3$时，$\Sigma$是一个点，点的坐标为$(0,0,-3)$。

  在$z=-4$时，$\Sigma$不存在。

\begin{center}
  \begin{tikzpicture}[scale=1.2]
  \tdplotsetmaincoords{54.7356}{135}
  \begin{scope}[tdplot_main_coords]
    \def\R{3}
    \def\rI{2.828}   % sqrt(8)
    \def\rII{2.236}  % sqrt(5)

    % ---- 虚线区间参数----
    \def\zMinusTwoDashStart{125.90}    % z=-2
    \def\zMinusTwoDashEnd{324.01}
    \def\zMinusOneDashStart{142.18}  % z=-1
    \def\zMinusOneDashEnd{307.82}
    \def\zZeroDashStart{155.70}        % z=0
    \def\zZeroDashEnd{294.30}
    \def\zOneDashStart{173.68}
    \def\zOneDashEnd{276.32}

    % 坐标轴
    \draw[->, black] (0,0,0) -- (4,0,0) node[below left] {$x$};
    \draw[->, black] (0,0,0) -- (0,4,0) node[below right] {$y$};
    \draw[->, black] (0,0,3) -- (0,0,4) node[above] {$z$};
    \draw[black, dashed] (0,0,0) -- (0,0,-3);
    \draw[black, dashed] (0,0,0) -- (0,0,3);

    % 球轮廓
    \draw[black] (0,0,3) -- (0,0,4);
    \draw[black] (0,0,-3) -- (0,0,-4);

    % ---------- 各截面圆 ----------

    % z=-3 顶点
    \filldraw[black] (0,0,-3) circle (1pt) node[below] {$(0,0,-3)$};

    % z = -2
    \fill[green!20, opacity=0.5] plot[variable=\t, domain=0:360, samples=72]
      ({\rII*cos(\t)},{\rII*sin(\t)},-2);
    \draw[green!50!black, thick, dashed] plot[variable=\t, domain=\zMinusTwoDashStart:\zMinusTwoDashEnd, samples=80]
      ({\rII*cos(\t)},{\rII*sin(\t)},-2);
    \draw[green!50!black, thick] plot[variable=\t, domain=\zMinusTwoDashEnd:360, samples=40]
      ({\rII*cos(\t)},{\rII*sin(\t)},-2);
    \draw[green!50!black, thick] plot[variable=\t, domain=0:\zMinusTwoDashStart, samples=30]
      ({\rII*cos(\t)},{\rII*sin(\t)},-2);
    \node[green!50!black] at ({\rII*cos(90)}, {\rII*sin(90)},-2) [below right] {$z=-2$};
    \filldraw[black] (0,0,-2) circle (1pt) node[below] {$(0,0,-2)$};

    % z = -1
    \fill[blue!20, opacity=0.5] plot[variable=\t, domain=0:360, samples=72]
      ({\rI*cos(\t)},{\rI*sin(\t)},-1);
    \draw[blue, thick, dashed] plot[variable=\t, domain=\zMinusOneDashStart:\zMinusOneDashEnd, samples=90]
      ({\rI*cos(\t)},{\rI*sin(\t)},-1);
    \draw[blue, thick] plot[variable=\t, domain=\zMinusOneDashEnd:360, samples=40]
      ({\rI*cos(\t)},{\rI*sin(\t)},-1);
    \draw[blue, thick] plot[variable=\t, domain=0:\zMinusOneDashStart, samples=30]
      ({\rI*cos(\t)},{\rI*sin(\t)},-1);
    \node[blue] at ({\rI*cos(60)}, {\rI*sin(60)},-1) [above] {$z=-1$};
    \filldraw[black] (0,0,-1) circle (1pt) node[below] {$(0,0,-1)$};

    % z = 0
    \fill[red!20, opacity=0.5] plot[variable=\t, domain=0:360, samples=72]
      ({\R*cos(\t)},{\R*sin(\t)},0);
    \draw[red, thick, dashed] plot[variable=\t, domain=\zZeroDashStart:\zZeroDashEnd, samples=80]
      ({\R*cos(\t)},{\R*sin(\t)},0);
    \draw[red, thick] plot[variable=\t, domain=\zZeroDashEnd:360, samples=40]
      ({\R*cos(\t)},{\R*sin(\t)},0);
    \draw[red, thick] plot[variable=\t, domain=0:\zZeroDashStart, samples=30]
      ({\R*cos(\t)},{\R*sin(\t)},0);
    \node[red] at ({\R*cos(0)}, {\R*sin(0)},0) [above right] {$z=0$};
    \filldraw[black] (0,0,0) circle (1pt) node[below] {$(0,0,0)$};

    % z = 1
    \fill[blue!20, opacity=0.5] plot[variable=\t, domain=0:360, samples=72]
      ({\rI*cos(\t)},{\rI*sin(\t)},1);
    \draw[blue, thick, dashed] plot[variable=\t, domain=\zOneDashStart:\zOneDashEnd, samples=80]
      ({\rI*cos(\t)},{\rI*sin(\t)},1);
    \draw[blue, thick] plot[variable=\t, domain=\zOneDashEnd:360, samples=40]
      ({\rI*cos(\t)},{\rI*sin(\t)},1);
    \draw[blue, thick] plot[variable=\t, domain=0:\zOneDashStart, samples=30]
      ({\rI*cos(\t)},{\rI*sin(\t)},1);
    \node[blue] at ({\rI*cos(90)}, {\rI*sin(90)},1) [above left] {$z=1$};
    \filldraw[black] (0,0,1) circle (1pt) node[below] {$(0,0,1)$};

    % z = 2
    \fill[green!20, opacity=0.5] plot[variable=\t, domain=0:360, samples=72]
      ({\rII*cos(\t)},{\rII*sin(\t)},2);
    \draw[green!50!black, thick] plot[variable=\t, domain=0:360, samples=40]
      ({\rII*cos(\t)},{\rII*sin(\t)},2);
    \node[green!50!black] at ({\rII*cos(120)}, {\rII*sin(120)},2) [left] {$z=2$};
    \filldraw[black] (0,0,2) circle (1pt) node[below] {$(0,0,2)$};

    % 顶点
    \filldraw[black] (0,0,3) circle (1pt) node[below] {$(0,0,3)$};
  \end{scope}
\end{tikzpicture}
\begin{tikzpicture}[scale=1.2]
  \tdplotsetmaincoords{87}{135}
  \begin{scope}[tdplot_main_coords]
    \def\R{3}
    \def\rI{2.828}   % sqrt(8)
    \def\rII{2.236}  % sqrt(5)

    % ---- 虚线区间参数----
    \def\zMinusTwoDashStart{0}    % z=-2
    \def\zMinusTwoDashEnd{0}
    \def\zMinusOneDashStart{0}  % z=-1
    \def\zMinusOneDashEnd{0}
    \def\zZeroDashStart{0}        % z=0
    \def\zZeroDashEnd{0}
    \def\zOneDashStart{0}
    \def\zOneDashEnd{0}

    % 坐标轴
    \draw[->, black] (0,0,0) -- (4,0,0) node[below left] {$x$};
    \draw[->, black] (0,0,0) -- (0,4,0) node[below right] {$y$};
    \draw[->, black] (0,0,0) -- (0,0,4) node[above] {$z$};
    \draw[black] (0,0,0) -- (0,0,-4);

    % 球轮廓
    \draw[black] (0,0,3) -- (0,0,4);
    \draw[black] (0,0,-3) -- (0,0,-4);

    % ---------- 各截面圆 ----------

    % z=-3 顶点
    \filldraw[black] (0,0,-3) circle (1pt) node[below] {$(0,0,-3)$};

    % z = -2
    \fill[green!20, opacity=0.5] plot[variable=\t, domain=0:360, samples=72]
      ({\rII*cos(\t)},{\rII*sin(\t)},-2);
    \draw[green!50!black, thick, dashed] plot[variable=\t, domain=\zMinusTwoDashStart:\zMinusTwoDashEnd, samples=80]
      ({\rII*cos(\t)},{\rII*sin(\t)},-2);
    \draw[green!50!black, thick] plot[variable=\t, domain=\zMinusTwoDashEnd:360, samples=40]
      ({\rII*cos(\t)},{\rII*sin(\t)},-2);
    \draw[green!50!black, thick] plot[variable=\t, domain=0:\zMinusTwoDashStart, samples=30]
      ({\rII*cos(\t)},{\rII*sin(\t)},-2);
    \node[green!50!black] at ({\rII*cos(180)}, {\rII*sin(180)},-2) [above] {$z=-2$};
    \filldraw[black] (0,0,-2) circle (1pt) node[below] {$(0,0,-2)$};

    % z = -1
    \fill[blue!20, opacity=0.5] plot[variable=\t, domain=0:360, samples=72]
      ({\rI*cos(\t)},{\rI*sin(\t)},-1);
    \draw[blue, thick, dashed] plot[variable=\t, domain=\zMinusOneDashStart:\zMinusOneDashEnd, samples=90]
      ({\rI*cos(\t)},{\rI*sin(\t)},-1);
    \draw[blue, thick] plot[variable=\t, domain=\zMinusOneDashEnd:360, samples=40]
      ({\rI*cos(\t)},{\rI*sin(\t)},-1);
    \draw[blue, thick] plot[variable=\t, domain=0:\zMinusOneDashStart, samples=30]
      ({\rI*cos(\t)},{\rI*sin(\t)},-1);
    \node[blue] at ({\rI*cos(180)}, {\rI*sin(180)},-1) [above] {$z=-1$};
    \filldraw[black] (0,0,-1) circle (1pt) node[below] {$(0,0,-1)$};

    % z = 0
    \fill[red!20, opacity=0.5] plot[variable=\t, domain=0:360, samples=72]
      ({\R*cos(\t)},{\R*sin(\t)},0);
    \draw[red, thick, dashed] plot[variable=\t, domain=\zZeroDashStart:\zZeroDashEnd, samples=80]
      ({\R*cos(\t)},{\R*sin(\t)},0);
    \draw[red, thick] plot[variable=\t, domain=\zZeroDashEnd:360, samples=40]
      ({\R*cos(\t)},{\R*sin(\t)},0);
    \draw[red, thick] plot[variable=\t, domain=0:\zZeroDashStart, samples=30]
      ({\R*cos(\t)},{\R*sin(\t)},0);
    \node[red] at ({\R*cos(180)}, {\R*sin(180)},0) [above] {$z=0$};
    \filldraw[black] (0,0,0) circle (1pt) node[below] {$(0,0,0)$};

    % z = 1
    \fill[blue!20, opacity=0.5] plot[variable=\t, domain=0:360, samples=72]
      ({\rI*cos(\t)},{\rI*sin(\t)},1);
    \draw[blue, thick, dashed] plot[variable=\t, domain=\zOneDashStart:\zOneDashEnd, samples=80]
      ({\rI*cos(\t)},{\rI*sin(\t)},1);
    \draw[blue, thick] plot[variable=\t, domain=\zOneDashEnd:360, samples=40]
      ({\rI*cos(\t)},{\rI*sin(\t)},1);
    \draw[blue, thick] plot[variable=\t, domain=0:\zOneDashStart, samples=30]
      ({\rI*cos(\t)},{\rI*sin(\t)},1);
    \node[blue] at ({\rI*cos(180)}, {\rI*sin(180)},1) [above] {$z=1$};
    \filldraw[black] (0,0,1) circle (1pt) node[below] {$(0,0,1)$};

    % z = 2
    \fill[green!20, opacity=0.5] plot[variable=\t, domain=0:360, samples=72]
      ({\rII*cos(\t)},{\rII*sin(\t)},2);
    \draw[green!50!black, thick] plot[variable=\t, domain=0:360, samples=40]
      ({\rII*cos(\t)},{\rII*sin(\t)},2);
    \node[green!50!black] at ({\rII*cos(180)}, {\rII*sin(180)},2) [above] {$z=2$};
    \filldraw[black] (0,0,2) circle (1pt) node[below] {$(0,0,2)$};

    % 顶点
    \filldraw[black] (0,0,3) circle (1pt) node[below] {$(0,0,3)$};
  \end{scope}
\end{tikzpicture}
\end{center}

结合截面和截痕分析，我们就能大致判断出这是一个以原点为球心，半径为$3$的球面。如果我们取更多的平行截面，得到更多的截痕，那么我们要研究的曲面的形状就更明显。

\end{myexample}

也就是说，只要我们截面足够多，就能得到足够多的截痕，就能刻画出曲面的形状。

总结一下使用截痕法的步骤：

首先先确定要截曲面的一族平行平面，通常取垂直于坐标轴的平面，如$z=0$、$z=1$、$z=2$，或者$y=0$、$y=1$、$y=2$等等。

再观察当平行平面连续变化时截痕的形状，得出曲面的大致形状和结构。

截痕法是研究曲面的通用方法，接下来我们介绍由二次曲线出发，在空间直角坐标系中旋转得到的二次曲面。

\begin{theorem}{由二次曲线旋转得到的二次曲面}{}
若在 \(xOy\) 平面上有一条已知曲线 \(C\)，方程为：
\[
C: f(x,y)=0
\]

绕\(x\) 轴旋转时，所得曲面$\Sigma $方程为：

\[
\Sigma: f(x,\ \pm\sqrt{y^2+z^2})=0
\]



绕\(y\) 轴旋转时，所得曲面$\Sigma $方程为：

\[
\Sigma: f(\pm\sqrt{x^2+z^2},\ y)=0
\]
\end{theorem}

这个定理对于其他平面上的曲线 \(C\)同样适用。

你可能会说这还是太难记住了。那么这个怎么理解呢？想象一下一条平面曲线绕某条轴旋转，就扫出了一个立体曲面，旋转轴可以认为没有运动，所以旋转轴上的坐标没有变，而由于旋转，非旋转轴上的坐标到旋转轴的距离是相等的，距离就是另外两个坐标的平方和的平方根，这也就是根号的由来。

总结就是：\textbf{绕着哪根轴旋转，哪根轴不变。另外两个坐标的平方和开根号，替换掉曲线的另一个坐标。}

因为有根号，所以我们一般使用平方和形式。即：

\[
\Sigma: f(x^2,\ y^2+z^2)=0 \quad \text{或} \quad \Sigma: f(x^2+z^2,\ y^2)=0
\]

\begin{myexample}
  举例说明：求曲线$C: y=kx$绕$y$轴旋转一周的曲面。

  我们用之前的总结的方法：

  首先，确定绕哪根轴旋转：$y$轴。那么曲线$C: y=kx$中的$y$不变。

  然后，另外两个坐标的平方和开根号，也就是$\sqrt{x^2+z^2}$，替换掉曲线的另一个坐标，也就是替换掉$x$。

  这个时候我们再写一下曲面方程：$y=k\sqrt{x^2+z^2}$，然后写成平方和形式，也就是$y^2=k^2(x^2+z^2)$，这个时候我们已经得到了曲面方程了。

  我们再用截痕法研究一下这个曲面：

  在$y=0$时，曲面是一个点，点的坐标为$(0,0,0)$。

  在$y=1$时，曲面是以$(0,1,0)$为圆心，半径为$\frac{1}{|k|}$的圆。

  在$y=2$时，曲面是以$(0,2,0)$为圆心，半径为$\frac{2}{|k|}$的圆。

  在$y=3$时，曲面是以$(0,3,0)$为圆心，半径为$\frac{3}{|k|}$的圆。

  以此类推，我们最后可以得到这个曲面是一个圆锥面。
\end{myexample}

\begin{center}
  \begin{tikzpicture}[scale=1.1]
    \tdplotsetmaincoords{60}{130}
    \begin{scope}[tdplot_main_coords]
      % ---- 参数 ----
      \def\maxY{3.2}
      \def\stepAngle{22.5}          % 经线间隔（度）
      \def\stepY{0.5}               % 纬线间隔

      % ---- 坐标轴 ----
      \draw[->, black] (0,0,0) -- (4.5,0,0) node[below left] {$x$};
      \draw[->, black] (0,0,0) -- (0,4.5,0) node[below right] {$y$};
      \draw[->, black] (0,0,0) -- (0,0,4.5) node[above] {$z$};
      \draw[black, dashed] (0,0,0) -- (-3.5,0,0);
      \draw[black, dashed] (0,0,0) -- (0,-3.5,0);
      \draw[black, dashed] (0,0,0) -- (0,0,-3.5);

      % ---- 经线（母线）----
      \foreach \ang in {0, \stepAngle, ..., 360} {
        \pgfmathsetmacro{\x}{\maxY*cos(\ang)}
        \pgfmathsetmacro{\z}{\maxY*sin(\ang)}
        \draw[black!50, very thin] (0,0,0) -- (\x,\maxY,\z);
        \draw[black!50, very thin] (0,0,0) -- (\x,-\maxY,\z);
      }

      % ---- 纬线（水平圆）----
      \foreach \y in {-3.0, -2.5, ..., 3.0} {
        \pgfmathsetmacro{\r}{abs(\y)}
        \draw[black!50, very thin] plot[variable=\t, domain=0:360, samples=40]
          ({\r*cos(\t)}, \y, {\r*sin(\t)});
      }



      % ---- 顶点 ----
      \filldraw[red] (0,0,0) circle (2pt);

    \end{scope}
  \end{tikzpicture}
\end{center}

我们给出所有常见二次曲面的标准方程，其他二次曲面都可以由它们通过伸缩旋转平移等变换得到。

\begin{definition}{常见二次曲面的标准方程}{}
  椭球面：
  \(
  \frac{x^2}{a^2} + \frac{y^2}{b^2} + \frac{z^2}{c^2} = 1
  \)，其中$a,b,c > 0$。

  椭圆抛物面：
\(
z = \frac{x^2}{a^2} + \frac{y^2}{b^2}
\)，其中$a,b > 0$。

旋转抛物面：
\(z = \frac{x^2+y^2}{a^2}\)，其中$a > 0$，此为椭圆抛物面中$a=b$的特殊情形。

双曲抛物面（马鞍面）：
\(
z = \frac{x^2}{a^2} - \frac{y^2}{b^2} \quad \text{或} \quad z = \frac{y^2}{b^2} - \frac{x^2}{a^2}
\)，此曲面\textbf{不可能由平面曲线旋转生成}。

单叶双曲面：
\(
\frac{x^2}{a^2} + \frac{y^2}{b^2} - \frac{z^2}{c^2} = 1
\)，其中$a,b,c > 0$。

双叶双曲面：
\(
\frac{x^2}{a^2} - \frac{y^2}{b^2} - \frac{z^2}{c^2} = 1
\)，其中$a,b,c > 0$。

二次锥面：
\(
\frac{x^2}{a^2} + \frac{y^2}{b^2} - \frac{z^2}{c^2} = 0
\)，其中$a,b,c > 0$。

圆锥面： \(x^2+y^2 = k^2z^2\)。

椭圆柱面： \(\frac{x^2}{a^2} + \frac{y^2}{b^2} = 1\)等，其中$a,b > 0$。

圆柱面：\(x^2+y^2=R^2\)等。

双曲柱面： \(\frac{x^2}{a^2} - \frac{y^2}{b^2} = 1\)等，其中$a,b > 0$。

抛物柱面： \(y = x^2\) 或 \(x = y^2\)等。

球面：\(x^2+y^2+z^2=R^2\)。
\end{definition}

其实这些曲面方程记不住也没关系，因为我们可以通过截痕法分析，最终也能得到曲面的形状。

\subsection{空间曲线方程}

和空间直线类似，空间曲线也可以通过不同的方式来描述。空间直线可以看作两个平面的交线，空间曲线则可以看作两个曲面的交线。

\begin{definition}{空间曲线的一般方程}{}
设 $\Sigma_1: F(x,y,z)=0$ 和 $\Sigma_2: G(x,y,z)=0$ 是空间中的两个曲面，则它们的交线 $C$ 的方程为：
\[
C: \begin{cases}
F(x,y,z)=0 \\
G(x,y,z)=0
\end{cases}
\]
这称为\textbf{空间曲线的一般方程}。
\end{definition}

然而，一般方程使用起来并不方便。因为这样会让我们在想象空间曲线时，先作出两个空间曲面。与空间直线类似，我们更常用的是参数方程。

\begin{definition}{空间曲线的参数方程}{}
设 $C$ 是空间中的一条曲线，其中 $C$ 上任意一点的坐标 $(x,y,z)$ 都可以表示为某个参数 $t$ 的函数，则空间曲线 $C$ 的方程为：
\[
C: \begin{cases}
x = x(t) \\
y = y(t) \\
z = z(t)
\end{cases}
\]
这称为\textbf{空间曲线的参数方程}。
\end{definition}

我们可以这样理解：参数 $t$ 就像一个“时间”，随着 $t$ 的变化，空间中都有对应的点 $(x(t), y(t), z(t))$，点动成线，无数多个点在空间中组成一条轨迹，这条轨迹就是曲线 $C$。

最常见的空间曲线参数方程是螺旋线。

\begin{myexample}
  举例说明：描述曲线
  \(C :
  \begin{cases}
    x = \cos t \\
    y = \sin t \\
    z = \frac{t}{10}
  \end{cases}
  \)的形状。

我们可以这样分析：

对$x,y$来说，曲线在$xOy$ 面上的运动是
\(
  \begin{cases}
    x = \cos t \\
    y = \sin t
  \end{cases}
  \)，所以当 $t$ 增大时，点 $(x,y)$ 在 $xOy$ 平面上画圆。

对$z$来说，曲线在$z$ 轴上的运动是\(z = \frac{t}{10}\)，所以当 $t$ 增大时，点 $(x,y,z)$ 在$z$ 轴上以\(\frac{t}{10}\)的速率匀速上升。

所以，整条曲线既在 $xOy$ 平面上画圆，又在$z$ 轴上以\(\frac{t}{10}\)的速率匀速上升，说明曲线实际上在一边画圆一边上升，也就是螺旋上升，就像一根绕在圆柱面上的弹簧。

这其实就是螺旋线。

  \begin{center}
\begin{tikzpicture}[scale=2]
  \tdplotsetmaincoords{55}{135}
  \begin{scope}[tdplot_main_coords]

    % 坐标轴
    \draw[->,thick] (0,0,0) -- (1.5,0,0) node[below left]{$x$};
    \draw[->,thick] (0,0,0) -- (0,1.5,0) node[below right]{$y$};
    \draw[->,thick] (0,0,0) -- (0,0,2.5)   node[above]{$z$};

    % 螺旋线: x=cos(t), y=sin(t), z=t/10
    \draw[
      domain=0:6*pi,
      samples=400,
      smooth,
      variable=\t,
      red,
      thick
    ] plot ({cos(\t r)}, {sin(\t r)}, {\t/10}) node[black, right] {$C$};

  \end{scope}
\end{tikzpicture}
  \end{center}

\end{myexample}

螺旋线在工程中非常常见，例如螺丝钉的螺纹、弹簧的形状，甚至 DNA 的双螺旋结构都可以用螺旋线来描述。

然而有时在实际问题中，曲线常以一般方程的形式给出。要写出它的参数方程，本质就是\textbf{引入一个参数 \(t\)，将 \(x,y,z\) 用 \(t\) 统一表示}。

如果曲线在某坐标轴上的投影比较简单，可以直接令该坐标为 \(t\)。

\begin{myexample}
  设曲线 \(C\) 是旋转抛物面 \(z = x^2 + y^2\) 与平面 \(y = x\) 的交线。写出它的参数方程。

  平面 \(y = x\) 较为简单，所以我们可以令 \(x = t\)，再代入 \(y=x\) 得 \(y = t\)。

  再将 \(x = t\)，\(y = t\) 代入旋转抛物面方程：\(z = x^2 + y^2\)，得到 \(z = 2t^2\)，其中 \(t \in \mathbb{R}\)。
  
  因此参数方程为：
  \[
  C: \begin{cases}
  x = t \\
  y = t \\
  z = 2t^2
  \end{cases},
  \quad t \in \mathbb{R}
  \]
\end{myexample}

这样做的原理是空间曲线可以看作两个曲面的交线，那么它就必须同时满足这两个曲面的方程。

\begin{center}
  \begin{tikzpicture}[scale=0.6]
  \tdplotsetmaincoords{70}{100}
  \begin{scope}[tdplot_main_coords]

        % 坐标轴
    \draw[->] (-2.5,0,0) -- (2.5,0,0) node[left] {$x$};
    \draw[->] (0,-2.5,0) -- (0,2.5,0) node[above] {$y$};
    \draw[->] (0,0,0) -- (0,0,9) node[left] {$z$};

    % 淡绿色平面 y = x（直接填充四边形）
    \fill[green!15, opacity=0.5] (-2,-2,0) -- (2,2,0) -- (2,2,8) -- (-2,-2,8) -- cycle;

    % 淡蓝色抛物面 z = x^2 + y^2（密致条带填充，步长 0.01 实现类曲面效果）
    \foreach \y [evaluate=\y as \nexty using \y+0.01] in {-2,-1.99,...,1.99,2} {
      \fill[blue!15, opacity=0.5]
        plot[domain=-2:2, smooth, variable=\x] (\x, \y, \x*\x + \y*\y) --
        plot[domain=2:-2, smooth, variable=\x] (\x, \nexty, \x*\x + \nexty*\nexty) -- cycle;
    }

    % 红色交线 z = 2x^2
    \draw[red, very thick, domain=-2:2, samples=100, variable=\t]
         plot (\t, \t, 2*\t*\t) node[right] {交线};

  \end{scope}
\end{tikzpicture}
\end{center}

\section{三角函数初步}


\subsection{三角函数的定义}
我们假定你有基础的高中三角函数知识，认识\(\sin\)，\(\cos\)，\(\tan\)这三个函数。然而它们只是三角函数的一部分，还有另外几个函数高中几乎不涉及，但是也必须了解。所以我们先从单位圆定义出发，定义所有三角函数：

\begin{definition}{三角函数}{}
在平面直角坐标系中，以原点 $O$ 为圆心、$1$ 为半径作单位圆。设角 $\theta$ 的终边与单位圆交于点 $P(x, y)$，定义：

\begin{center}
正弦：\(\sin\theta = y\)，\quad 余弦：\(\cos\theta = x\)
\end{center}

\begin{center}
正切：\(\tan\theta = \dfrac{y}{x}\ (x \neq 0)\)，\quad 余切：\(\cot\theta = \dfrac{x}{y}\ (y \neq 0)\)
\end{center}

\begin{center}
正割：\(\sec\theta = \dfrac{1}{x}\ (x \neq 0)\)，\quad 余割：\(\csc\theta = \dfrac{1}{y}\ (y \neq 0)\)
\end{center}

\end{definition}

其中，\(\sin\)，\(\cos\)，\(\tan\)是高中就学过的，但是\(\cot\)，\(\sec\)，\(\csc\)好像就没有那么熟悉了。我们依旧先从几何意义出发，再推广到全体实数：

\begin{center}
  \begin{tikzpicture}[
  scale=2.4,
  line cap=round,
  font=\small
]
  % 坐标轴
  \draw[->] (-1.5,0) -- (1.5,0) node[right] {$x$};
  \draw[->] (0,-1.5) -- (0,2) node[above] {$y$};

  % 单位圆
  \draw[gray, thin] (0,0) circle(1);

  % 设定角度
  \def\theta{60}
  \pgfmathsetmacro\cs{cos(\theta)}
  \pgfmathsetmacro\sn{sin(\theta)}
  \pgfmathsetmacro\tg{tan(\theta)}
  \pgfmathsetmacro\ctg{1/tan(\theta)}

  % 关键点
  \coordinate (O) at (0,0);
  \coordinate (P) at (\cs,\sn);
  \coordinate (M) at (\cs,0);
  \coordinate (T) at (1,\tg);
  \coordinate (S) at (\ctg,1);

  % 终边
  \draw[gray, thick] (O) -- (1.2,{1.2*tan(\theta)});

  % 角度弧
  \draw[->] (0.2,0) arc (0:\theta:0.2);
  \node at (0.32,0.26) {$\frac{\pi}{3}$};

  % 原点与圆上点
  \fill (O) circle(0.03) node[below left] {$O$};
  \fill (P) circle(0.03);

  % ---- 六个三角函数 ----

  % 1. 正弦 sin(π/3) = MP (红)
  \draw[dashed, gray] (P) -- (M);
  \draw[red, thick, ->] (M) -- (P)
    node[pos=0.5, right] {$\sin\frac{\pi}{3}$};

  % 2. 余弦 cos(π/3) = OM (蓝)
  \draw[blue, thick, ->] (O) -- (M)
    node[midway, below] {$\cos\frac{\pi}{3}$};

  % 3. 正切 tan(π/3) = AT (橙)
  \draw[dashed, gray] (1,0) -- (T);
  \draw[orange, thick, ->] (1,0) -- (T)
    node[pos=0.4, right] {$\tan\frac{\pi}{3}$};

  % 4. 余切 cot(π/3) = BS (紫)
  \draw[dashed, gray] (0,1) -- (S);
  \draw[purple, thick, ->] (0,1) -- (S)
    node[pos=0.5, above] {$\cot\frac{\pi}{3}$};

  % 5. 余割 csc(π/3) = OS (绿)
  \draw[green!60!black, thick, ->] (O) -- (S)
    node[pos=0.5, sloped, above] {$\csc\frac{\pi}{3}$};

  % 6. 正割 sec(π/3) = OT (青)
  \draw[cyan, thick, ->] (O) -- (T)
    node[pos=0.6, sloped, below] {$\sec\frac{\pi}{3}$};

\end{tikzpicture}
\end{center}

我们知道，\(\sin\)，\(\cos\)，\(\tan\)在直角三角形里有它们的定义。相应的，\(\cot\)，\(\sec\)，\(\csc\)也有它们的定义。一般地，在直角三角形中，我们定义\(\cot=\frac{\text{邻边}}{\text{对边}}\)，\(\sec=\frac{\text{斜边}}{\text{邻边}}\)，\(\csc=\frac{\text{斜边}}{\text{对边}}\)。细心的你也许已经发现了，其实\(\cot\)就是\(\tan\)的倒数，\(\sec\)就是\(\cos\)的倒数，\(\csc\)就是\(\sin\)的倒数。也就是说，我们只需要用高中计算\(\sin\)，\(\cos\)，\(\tan\)的方法，再记住它们的倒数关系，就能很快的算出\(\cot\)，\(\sec\)，\(\csc\)了。同样的，我们从直角三角形推广到全体实数，就引入了单位圆定义。高中的单位圆定义了\(\sin\)，\(\cos\)，\(\tan\)，这次我们把\(\cot\)，\(\sec\)，\(\csc\)的定义一起加上了。

不过，我们主要使用的还是 $\sin$，$\cos$，$\tan$ 这三个，$\sec$，$\cot$，$\csc$ 不能说不用，只是出现得比较少。

\begin{myexample}
  利用单位圆定义，求 $\sec\frac{\pi}{3}$、$\cot\frac{\pi}{3}$、$\csc\frac{\pi}{3}$ 的值。

角 $\frac{\pi}{3}$ 的终边与单位圆的交点为 $P\!(\frac{1}{2}, \frac{\sqrt{3}}{2})$。因此：

\[
\sin\frac{\pi}{3} = \frac{\sqrt{3}}{2},\quad 
\cos\frac{\pi}{3} = \frac{1}{2},\quad 
\tan\frac{\pi}{3} = \frac{\sin\frac{\pi}{3}}{\cos\frac{\pi}{3}} = \sqrt{3}
\]

利用倒数关系可得：
\[
\sec\frac{\pi}{3} = \frac{1}{\cos\frac{\pi}{3}} = 2,\quad 
\csc\frac{\pi}{3} = \frac{1}{\sin\frac{\pi}{3}} = \frac{2\sqrt{3}}{3},\quad 
\cot\frac{\pi}{3} = \frac{1}{\tan\frac{\pi}{3}} = \frac{\sqrt{3}}{3}
\]

\end{myexample}

这正是图上各个三角函数的值。在单位圆定义里，它们表示的是线段的长度。

常见的特殊角的三角函数值需要牢记。这些值不需要死记硬背。画一个单位圆，标出各特殊角对应的点坐标，$\sin$ 读纵坐标、$\cos$ 读横坐标、$\tan$ 读纵横之比，剩下的就能全部推出来。


\subsection{三角函数公式}

三角函数有很多公式，例如刚刚说到的倒数关系：
\begin{theorem}{三角函数的倒数关系}{}
\[
\sin\theta \cdot \csc\theta = 1,\quad
\cos\theta \cdot \sec\theta = 1,\quad
\tan\theta \cdot \cot\theta = 1
\]
\end{theorem}

还有平方关系：
\begin{theorem}{三角函数的平方关系}{}
\[
\sin^2\theta + \cos^2\theta = 1,\quad
1 + \tan^2\theta = \sec^2\theta,\quad
1 + \cot^2\theta = \csc^2\theta
\]
\end{theorem}

第一个公式是勾股定理的三角版本，后面两个是第一个两边分别除以 $\cos^2\theta$ 和 $\sin^2\theta$ 得到的。

接着是诱导公式。诱导公式用来把任意角化简为 $0$ 到 $\frac{\pi}{2}$ 之间的角。核心规律是“奇变偶不变，符号看象限”。具体的操作我们假定读者高中已经掌握，所以我们不再赘述。

\begin{theorem}{三角函数的诱导公式}{}
\[
\begin{aligned}
\sin(-\theta) &= -\sin\theta, & \sin(\frac{\pi}{2} \pm \theta) &= \cos\theta, & \sin(\pi \pm \theta) &= \mp\sin\theta \\[6pt]
\cos(-\theta) &= \cos\theta, & \cos(\frac{\pi}{2} \pm \theta) &= \mp\sin\theta, & \cos(\pi \pm \theta) &= -\cos\theta \\[6pt]
\tan(-\theta) &= -\tan\theta, & \tan(\frac{\pi}{2} \pm \theta) &= \mp\cot\theta, & \tan(\pi \pm \theta) &= \pm\tan\theta \\[6pt]
\cot(-\theta) &= -\cot\theta, & \cot(\frac{\pi}{2} \pm \theta) &= \mp\tan\theta, & \cot(\pi \pm \theta) &= \pm\cot\theta \\[6pt]
\sec(-\theta) &= \sec\theta, & \sec(\frac{\pi}{2} \pm \theta) &= \mp\csc\theta, & \sec(\pi \pm \theta) &= -\sec\theta \\[6pt]
\csc(-\theta) &= -\csc\theta, & \csc(\frac{\pi}{2} \pm \theta) &= \sec\theta, & \csc(\pi \pm \theta) &= \mp\csc\theta
\end{aligned}
\]
\end{theorem}



然后是和差公式：
\begin{theorem}{三角函数的和差公式}{}
\[
\begin{aligned}
\sin(\alpha \pm \beta) &= \sin\alpha\cos\beta \pm \cos\alpha\sin\beta \\
\cos(\alpha \pm \beta) &= \cos\alpha\cos\beta \mp \sin\alpha\sin\beta \\
\tan(\alpha \pm \beta) &= \frac{\tan\alpha \pm \tan\beta}{1 \mp \tan\alpha\tan\beta} \\
\cot(\alpha \pm \beta) &= \frac{\cot\alpha\cot\beta \mp 1}{\cot\beta \pm \cot\alpha} \\
\sec(\alpha \pm \beta) &= \frac{\sec\alpha\sec\beta}{1 \mp \tan\alpha\tan\beta} \\
\csc(\alpha \pm \beta) &= \frac{\csc\alpha\csc\beta}{\cot\beta \pm \cot\alpha}
\end{aligned}
\]
\end{theorem}

然后是倍角公式：
\begin{theorem}{三角函数的倍角公式}{}
\[
\begin{aligned}
\sin 2\theta &= 2\sin\theta\cos\theta\\
\cos 2\theta &= \cos^2\theta - \sin^2\theta = 2\cos^2\theta - 1 = 1 - 2\sin^2\theta\\
\tan 2\theta &= \frac{2\tan\theta}{1 - \tan^2\theta}\\
\cot 2\theta &= \frac{\cot^2\theta - 1}{2\cot\theta}\\
\sec 2\theta &= \frac{\sec^2\theta}{1 - \tan^2\theta}\\
\csc 2\theta &= \frac{\csc^2\theta}{2\cot\theta}
\end{aligned}
\]
\end{theorem}

它们可以由和角公式直接得到。

然后是半角公式：
\begin{theorem}{三角函数的半角公式}{}
\[
\begin{aligned}
\sin\frac{\theta}{2} &= \pm \sqrt{\frac{1 - \cos\theta}{2}}\\
\cos\frac{\theta}{2} &= \pm \sqrt{\frac{1 + \cos\theta}{2}}\\
\tan\frac{\theta}{2} &= \pm \sqrt{\frac{1 - \cos\theta}{1 + \cos\theta}} = \frac{\sin\theta}{1 + \cos\theta} = \frac{1 - \cos\theta}{\sin\theta}\\
\cot\frac{\theta}{2} &= \pm \sqrt{\frac{1 + \cos\theta}{1 - \cos\theta}} = \frac{\sin\theta}{1 - \cos\theta} = \frac{1 + \cos\theta}{\sin\theta}\\
\sec\frac{\theta}{2} &= \pm \sqrt{\frac{2\sec\theta}{\sec\theta + 1}} = \pm \sqrt{\frac{2}{1 + \cos\theta}}\\
\csc\frac{\theta}{2} &= \pm \sqrt{\frac{2\csc\theta}{\csc\theta - \cot\theta}} = \pm \sqrt{\frac{2}{1 - \cos\theta}}
\end{aligned}
\]
\end{theorem}
它们是倍角公式的逆用。其中正负由\(\frac{\theta }{2}\)所在的象限决定。

然后是降幂公式：
\begin{theorem}{三角函数的降幂公式}{}
\[
\begin{aligned}
\sin^2\theta &= \frac{1 - \cos 2\theta}{2}, & \cos^2\theta &= \frac{1 + \cos 2\theta}{2} \\[6pt]
\tan^2\theta &= \frac{1 - \cos 2\theta}{1 + \cos 2\theta}, & \cot^2\theta &= \frac{1 + \cos 2\theta}{1 - \cos 2\theta} \\[6pt]
\sec^2\theta &= \frac{2}{1 + \cos 2\theta}, & \csc^2\theta &= \frac{2}{1 - \cos 2\theta}
\end{aligned}
\]
\end{theorem}

半角公式移项就是降幂公式，常用于积分中把高次三角幂降为低次。

其实这些公式不用特别记忆，需要的时候回来看一下，多用一下就记住了。如果有些公式记不住，那就说明不常用，不记也罢。

\subsection{反三角函数}

反三角函数，顾名思义，就是三角函数反过来。三角函数是已知角，求对应比值，反三角函数就是已知比值，求对应角度。例如 $\sin\theta = \frac{1}{2}$，求 $\theta$ 等于多少？你马上能说出来 $\theta=\frac{\pi}{6}$。但注意，$\sin\frac{5\pi}{6}$ 也等于 $\frac{1}{2}$，$\sin\frac{13\pi}{6}$ 也等于 $\frac{1}{2}$。这就产生了一点问题，我们知道，在函数的概念里，一个$x$对应一个$y$。现在的$x$是$\frac{1}{2}$，但是对应了$\frac{\pi}{6}$，$\frac{5\pi}{6}$，$\frac{13\pi}{6}$，甚至更多的$y$。因此，反三角函数必须把值域限制在一个特定区间内，使得一个$x$必须对应一个$y$，否则一个比值会对应无穷多个角，就不是“函数”了。

最常用的三个反三角函数是反正弦、反余弦、反正切。

\begin{definition}{反正弦函数}{}
设 $x \in [-1, 1]$，若存在唯一的 $y \in [-\frac{\pi}{2}, \frac{\pi}{2}]$ 使得 $\sin y = x$，
则称 $y$ 为 $x$ 的反正弦，记作
\[
y = \arcsin x
\]
其中，函数 $y = \arcsin x$ 的定义域为 $[-1, 1]$，值域为 $[-\frac{\pi}{2}, \frac{\pi}{2}]$。
\end{definition}

例如 $\arcsin\frac{1}{2} = \frac{\pi}{6}$，$\arcsin(-1) = -\frac{\pi}{2}$，$\arcsin 0 = 0$。我们这里值域取的是 $\sin$ 在 $[-\frac{\pi}{2}, \frac{\pi}{2}]$ 上的那一段，这恰好是 $\sin$ 单调递增的区间，保证了一一对应。

同理，我们有：

\begin{definition}{反余弦函数}{}
设 $x \in [-1, 1]$，若存在唯一的 $y \in [0, \pi]$ 使得 $\cos y = x$，
则称 $y$ 为 $x$ 的反余弦，记作
\[
y = \arccos x
\]
其中，函数 $y = \arccos x$ 的定义域为 $[-1, 1]$，值域为 $[0, \pi]$。
\end{definition}

例如 $\arccos\frac{1}{2} = \frac{\pi}{3}$，$\arccos(-1) = \pi$，$\arccos 0 = \frac{\pi}{2}$。$\cos$ 在 $[0, \pi]$ 上单调递减。

同理，我们有：

\begin{definition}{反正切函数}{}
设 $x \in \mathbb{R}$，若存在唯一的 $y \in (-\frac{\pi}{2}, \frac{\pi}{2})$ 使得 $\tan y = x$，
则称 $y$ 为 $x$ 的反正切，记作
\[
y = \arctan x
\]
其中，函数 $y = \arctan x$ 的定义域为 $\mathbb{R}$，值域为 $(-\frac{\pi}{2}, \frac{\pi}{2})$。
\end{definition}

例如 $\arctan 1 = \frac{\pi}{4}$，$\arctan\sqrt{3} = \frac{\pi}{3}$，$\arctan 0 = 0$。剩下三个反三角函数也是类似的定义，不过它们实在是出现得太少，我们省略了，需要时查表即可。

同样的，反三角函数也有它们的公式。不过实际情况几乎用不到，而且也不在我们微积分的讨论范围内，所以我们此处也省略不讲了。


% !TEX root = ../main.tex
\chapter{极限}

微积分之所以叫微积分，是因为它建立在两个核心概念之上：微分和积分。但是这两个概念都依赖于一个更根本的东西——\textbf{极限}。本章将正式引入极限的 $\epsilon-\delta$ 定义，这一章之后，希望你能用数学语言精确地描述“无限接近”这个概念。

\section{极限的 $\epsilon-\delta$ 定义}

极限的讨论对象是谁？是函数。所以我们必须先统一函数的概念。

\subsection{函数}

我们假定你在高中数学中学过函数。这里我们给出一个更精确的表述。

\begin{definition}{函数}{}
设有两个非空集合 $D$ 与 $Y$。若对于 $D$ 中的每一个元素 $x$，按照某种确定的对应法则 $f$，都有 $Y$ 中唯一确定的元素 $y$ 与之对应，则称 $f$ 是定义在 $D$ 上的一个\textbf{函数}，记作
\[
f: D \to Y, \quad y = f(x), \quad x \in D
\]
其中 $D$ 称为函数的\textbf{定义域}，全体函数值 $f(x)$ 构成的集合 $\{f(x) \mid x \in D\}$ 称为函数的\textbf{值域}。
\end{definition}

由此我们可以定义基本初等函数：
\begin{definition}{基本初等函数}{}
  以下六种函数统称为\textbf{基本初等函数}：

\textbf{常数函数}：\(f(x) = C\)（$C$为实常数）

\textbf{幂函数}：\(f(x) = x^{\mu} \quad (\mu \in \mathbb{R})\)

\textbf{指数函数}：\(f(x) = a^{x} \quad (a > 0,\ a \neq 1)\)

\textbf{对数函数}：\(f(x) = \log_{a} x \quad (a > 0,\ a \neq 1)\)

\textbf{三角函数}：\(\sin x,\ \cos x,\ \tan x,\ \cot x,\ \sec x,\ \csc x\)

\textbf{反三角函数}：\(\arcsin x,\ \arccos x,\ \arctan x,\ \operatorname{arccot} x,\ \operatorname{arcsec} x,\ \operatorname{arccsc} x\)

\end{definition}

更进一步，我们定义初等函数：

\begin{definition}{初等函数}{}
由基本初等函数经过有限次的四则运算和函数复合所构成的，并且可以用一个解析式表示的函数，称为\textbf{初等函数}。
\end{definition}

微积分讨论的都是初等函数，但实际上函数可以多种多样。例如著名的“狄利克雷函数”：有理数处取值为$1$，无理数处取值为$0$。

\[
D(x)=
\begin{cases}
  1,x\in \mathbb{Q} \\
  0,x\notin \mathbb{Q} 
\end{cases}
\]

它是一个极端的“病态”函数，处处不连续，处处不可导，处处不可积，甚至连图象都画不出来。但它确实是一个函数：$D=\mathbb{R} $，$Y=\{0,1\}$，对于 $D$ 中的每一个元素 $x$，都有 $Y$ 中唯一确定的元素 $y$ 与之对应。

1734 年，欧拉在著作中首次用 $f(x)$ 表示一个依赖于 $x$ 的量，这就是我们今天记号的来源。不过那时数学界理解的“函数”基本上就是“一条公式”，也就是必须有解析表达式才算函数。直到 1837 年，狄利克雷提出了现代函数的定义，也就是\textbf{只要每个 $x$ 对应唯一的 $y$，就是函数}。

接下来可以讨论极限了。在极限的讨论中，我们只关心 $x$ 无限接近某个 $x_0$ 时 $f(x)$ 的变化趋势。至于 $f(x_0)$ 本身是否存在、值是多少，我们并不在意。也就是说，极限只看重“趋近”的过程，不看“到达”的结果。

为了精确描述“接近某个点”这件事，我们需要先引入邻域的概念：

\begin{definition}{邻域}{}
设 $x_0$ 是数轴上的一个点，$\delta > 0$。则将集合
\[
U(x_0, \delta) = \{ x \mid |x - x_0| < \delta \} = (x_0 - \delta,\ x_0 + \delta)
\]
称为点 $x_0$ 的\textbf{$\delta$ 邻域}，其中$\delta$ 称为邻域的半径。
\end{definition}

简单来说，$x_0$ 的 $\delta$ 邻域就是以 $x_0$ 为中心、左右各延伸 $\delta$ 的一个开区间。

\begin{center}
  \begin{tikzpicture}[>=stealth, scale=1.2]

  % 数轴
  \draw[->] (-4.5,0) -- (4.5,0) node[right] {$x$};

  % 刻度线 (只画部分)
  \foreach \x in {-3,-2,-1,0,1,2,3}
    \draw (\x,0.15) -- (\x,-0.15);

  % 邻域区间的高亮 (可选)
  \draw[very thick, blue!30, line width=4pt, line cap=round] (-2,0) -- (2,0);

  % 中心点 x0 (实心)
  \filldraw (0,0) circle (2pt) node[above=2pt] {$x_0$};

  % 区间端点 (空心)
  \draw[fill=white, thick] (-2,0) circle (2pt) node[above=2pt] {$x_0-\delta$};
  \draw[fill=white, thick] (2,0) circle (2pt) node[above=2pt] {$x_0+\delta$};

  % 半径标注
  \draw[<->, red!70!black, thick] (0,-0.4) -- (2,-0.4) node[midway, below] {$\delta$};
  \draw[<->, red!70!black, thick] (-2,-0.4) -- (0,-0.4) node[midway, below] {$\delta$};

  % 邻域名称
  \node[below=0.9cm] at (0,0) {$U(x_0,\delta)$};

\end{tikzpicture}
\end{center}


邻域刻画的是“以 $x_0$ 为中心的一小段范围”。但在极限的讨论中，我们通常研究的是“趋近于 $x_0$ 时”会怎样，而不是“在 $x_0$ 处”怎样。所以我们反而并不关心 $x_0$ 这一点本身，于是我们去掉点 $x_0$，这就得到了去心邻域：

\begin{definition}{去心邻域}{}
设 $x_0$ 是数轴上的一个点，$\delta > 0$。则将集合
\[
\mathring{U}(x_0, \delta) = \{ x \mid 0 < |x - x_0| < \delta \} = (x_0 - \delta,\ x_0) \cup (x_0,\ x_0 + \delta)
\]
称为点 $x_0$ 的\textbf{$\delta$ 去心邻域}，其中$\delta$ 称为去心邻域的半径。

\end{definition}

去心邻域和邻域的唯一区别就是把 $x_0$ 本身挖掉了。

\begin{center}
  \begin{tikzpicture}[>=stealth, scale=1.2]
  
    % 数轴
    \draw[->] (-4.5,0) -- (4.5,0) node[right] {$x$};
  
    % 刻度线
    \foreach \x in {-3,-2,-1,0,1,2,3}
      \draw (\x,0.15) -- (\x,-0.15);
  
    % 邻域高亮（先画出整个区间，再通过覆盖实现挖空效果）
    \draw[very thick, blue!20, line width=4pt, line cap=round] (-2,0) -- (2,0);
  
    % 左端点 x0-δ（空心）
    \filldraw[fill=white, draw=black, thick] (-2,0) circle (2pt) node[above=2pt] {$x_0-\delta$};
    % 右端点 x0+δ（空心）
    \filldraw[fill=white, draw=black, thick] (2,0) circle (2pt) node[above=2pt] {$x_0+\delta$};
  
    % 中心点 x0（挖空，用稍大的空心圆表示）
    \filldraw[fill=white, draw=black, thick] (0,0) circle (2pt) node[above=2pt] {$x_0$};
  
    % 半径标注（左右两段长度均为 δ）
    \draw[<->, red!70!black, thick] (0,-0.4) -- (2,-0.4) node[midway, below] {$\delta$};
    \draw[<->, red!70!black, thick] (-2,-0.4) -- (0,-0.4) node[midway, below] {$\delta$};
  
    % 去心邻域符号
    \node[below=0.9cm] at (0,0) {$\mathring{U}(x_0,\delta)$};
  
  \end{tikzpicture}
\end{center}

有了邻域的语言，接下来我们可以用$\epsilon-\delta$语言来定义极限了。

\subsection{函数极限的 $\epsilon-\delta$ 定义}

先举一个简单的例子：设函数$f(x)=2x$，当 $x$ 越来越接近 $1$ 的时候，函数值越来越接近多少？显然是 $2$。也就是说，直觉上，当$x$ 无限接近 $1$ 时，$f(x)$ 无限接近 $2$。但问题来了，“无限接近”是个很模糊的词。到底是多近？用数学语言怎么描述？

我们先来玩一个游戏。假设你是一个水平极高的弓箭手，拥有一项技能，在相同位置射箭总能命中相同的位置，且已知$x_0$，在$x_0$射箭时总能精确地射中靶心 $y = L$，但在$x_0$附近射箭时会有偏差，不能精确射中靶心。然后我和你说，我指定一个$x_0$，你必须在距离$x_0$的$\delta$范围内射箭，而不是$x_0$，要求你射出的箭必须命中靶心区域，误差小于$\epsilon$，其中$\epsilon$是任意的，但是$x_0$不一定。我可以把$x_0$设在天上，地里，总之让你在$x_0$时无法射中靶心 $y = L$。虽然你水平很高，但是$x_0$的情况不清楚。这时你会怎么做？

让我们仔细想想。在距离$x_0$的$\delta$范围内，$\delta$是任意的。那么你可以说，没问题，但是我有个要求，$\delta$必须是我自己取的，而且$x_0$附近一定有我能射箭的地方。只要我找到一个$\delta$，能保证射出的箭就一定在$y = L$附近，且误差小于$\epsilon$，那么我就成功了。

按照这样的规定，游戏开始了。

我要求你在$x_0=1$处射箭，$\epsilon=1$。你说行，然后取了个$\delta=0.1$，射箭命中靶心附近，误差$0.99$，完美完成任务。我说不行，继续要求你在$x_0=999$处射箭，$\epsilon=0.0001$，你说没问题，然后取了个$\delta=0.00001$，射箭命中靶心附近，误差$0.00001$，完美完成任务。我说不行，我这次要求你在一万米高空中射箭，$\epsilon=0.001$，你说可以，然后找了一栋高楼，高度在$9999$米，接着你在天台射箭，命中靶心附近，误差$0.0009$，完美完成任务。

也就是说，无论我给什么$x_0$和$\epsilon $，只要能找到一个$\delta$，使得射出的箭距离靶心永远在$\epsilon$误差范围内，那么你就赢了。

我们把概念抽象一下。我们把“必须在距离$x_0$的$\delta$范围内射箭，而不是$x_0$”叫做“点 $x_0$ 的某个去心邻域内”，把“$x_0$附近一定有我能射箭的地方”叫做“函数 $f(x)$ 在点 $x_0$ 的某个去心邻域内有定义”，那么我们就可以给出极限的定义：

\begin{definition}{函数极限的 $\epsilon-\delta$ 定义}{}
设函数 $f(x)$ 在点 $x_0$ 的某个去心邻域内有定义，$L$ 是一个实数。若对于任意给定的 $\epsilon > 0$，总存在 $\delta > 0$，使得当
\[
0 < |x - x_0| < \delta
\]
时，一定有
\[
|f(x) - L| < \epsilon
\]
则称当 $x$ 趋近于 $x_0$ 时，$f(x)$ 的\textbf{极限}为 $L$ ，记作
\[
\lim_{x \to x_0} f(x) = L
\]
\end{definition}

也就是说，无论$x_0$取什么值，只要$x_0$ 的某个去心邻域内有定义，即使函数在$x_0$ 处无定义，只要满足上面的条件，那么我们也能求出函数的极限。

从图中可以看出，$f(x)$ 在 $\mathring{U}(x_0, \delta)$ 内的所有函数值都落在 $U(L, \epsilon)$ 内。这就是极限的几何表示。


\begin{center}
  \begin{tikzpicture}[>=stealth, scale=1.2]
  
  % --- 坐标轴 ---
  \draw[->] (-0.5,0) -- (6,0) node[below] {$x$};
  \draw[->] (0,-0.5) -- (0,5) node[left] {$y$};
  
  % --- 函数曲线 (例如 f(x)=0.2*(x-3)^2+2) ---
  \draw[domain=0.5:5.5, samples=100, thick, blue!60!black]
      plot (\x, {0.2*(\x-3)^2 + 2}) node[right] {$f(x)$};
  
  % --- 极限点 x0, L ---
  \def\xo{3}
  \def\L{2}
  
  % x0 处画虚线确定 L
  \draw[dashed, gray] (\xo, 0) node[below] {$x_0$} -- (\xo, \L);
  
  % 在 x0 处画空心圆 (去心邻域)
  \filldraw[fill=white, draw=blue!60!black, thick] (\xo, \L) circle (0.06);
  
  % y=L 的标记
  \draw[dashed, gray] (0, \L) node[left] {$L$} -- (\xo, \L);
  
  % --- 定义 epsilon = 0.8, 则 L ± epsilon ---
  \def\eps{0.8}
  \draw[dashed, red!70!black] (0, \L+\eps) node[left] {$L+\epsilon$} -- (5.5, \L+\eps);
  \draw[dashed, red!70!black] (0, \L-\eps) node[left] {$L-\epsilon$} -- (5.5, \L-\eps);
  
  % 标注 epsilon 的双向箭头
  \draw[<->, red!70!black] (5.3, \L) -- (5.3, \L+\eps) node[midway, right] {$\epsilon$};
  \draw[<->, red!70!black] (5.3, \L) -- (5.3, \L-\eps) node[midway, right] {$\epsilon$};
  
  \def\del{2} % delta=2
  % 竖直线 x0±delta
  \draw[dashed, green!50!black] (\xo-\del, 0) -- (\xo-\del, 4.5);
  \draw[dashed, green!50!black] (\xo+\del, 0) -- (\xo+\del, 4.5);
  
  % 标注 delta 的双向箭头
  \draw[<->, green!50!black] (\xo-\del, 1.2) -- (\xo, 1.2) node[midway, above] {$\delta$};
  \draw[<->, green!50!black] (\xo, 1.2) -- (\xo+\del, 1.2) node[midway, above] {$\delta$};
  
  % 在 x 轴上标出区间
  \draw[green!50!black, thick] (\xo-\del, 0.05) -- (\xo-\del, -0.05);
  \draw[green!50!black, thick] (\xo+\del, 0.05) -- (\xo+\del, -0.05);
  \node[below, green!50!black] at (\xo-\del, 0) {$x_0\!-\!\delta$};
  \node[below, green!50!black] at (\xo+\del, 0) {$x_0\!+\!\delta$};
  
  % --- 填充区间背景，强调带子 ---
  % 水平 epsilon 带
  \fill[red!5] (0, \L-\eps) rectangle (5.8, \L+\eps);
  % 竖直 delta 带（去心，但这里无法方便地打洞，我们覆盖上去）
  \fill[green!5] (\xo-\del, 0) rectangle (\xo+\del, 4.8);
  
  % 重新绘制轴和曲线，以确保它们在最上层
  \draw[->] (-0.5,0) -- (6,0) node[below] {$x$};
  \draw[->] (0,-0.5) -- (0,5) node[left] {$y$};
  \draw[domain=0.5:5.5, samples=100, thick, blue!60!black]
      plot (\x, {0.2*(\x-3)^2 + 2});
  
  % 再次画出 L±eps 的虚线, x0±delta 虚线等，以免被填充覆盖
  \draw[dashed, red!70!black] (0, \L+\eps) -- (5.5, \L+\eps);
  \draw[dashed, red!70!black] (0, \L-\eps) -- (5.5, \L-\eps);
  \draw[dashed, green!50!black] (\xo-\del, 0) -- (\xo-\del, 4.5);
  \draw[dashed, green!50!black] (\xo+\del, 0) -- (\xo+\del, 4.5);
  
  % 重描空心点
  \filldraw[fill=white, draw=blue!60!black, thick] (\xo, \L) circle (0.06);
  
  \end{tikzpicture}
\end{center}

我们理解一下这个定义：

$\epsilon$是我指定的误差范围，也就是我要求你的箭必须落在离靶心 $L$ 小于 $\epsilon$ 的范围内。$\epsilon$ 越小，我对你精度的要求就越苛刻。

$\delta$是你找的射箭范围，只要你在离 $x_0$ 小于 $\delta$ 的范围内选一个位置射箭，你就保证命中。$\epsilon$ 越小，你选的 $\delta$ 通常也得越小。

$0 < |x - x_0|$意味着不能在 $x_0$ 本身射箭，$x_0$ 可能在天上、在地里、在墙里面，总之你没法站在 $x_0$ 那个点上射箭。你只能在 $x_0$ 附近选位置。翻译成数学语言就是 $x \neq x_0$，极限只看附近的表现，不管 $x_0$ 本身什么情况，即使函数在$x_0$处无定义。

$|x - x_0| < \delta$就是说你的射箭位置必须在 $\delta$ 范围内。

$|f(x) - L| < \epsilon$是射出去的箭必须命中的靶心范围，$f(x)$ 是落点的坐标。这个落点必须距离 $L$ 小于 $\epsilon$。


这个游戏的胜负标准就是：\textbf{无论我给出多小的 $\epsilon$，你都能找到一个 $\delta$，保证在距离$x_0$的$\delta$ 范围内的每一次射箭，最终都落在距离靶心$L$的$\epsilon$ 范围内。}一旦你找到了这样的 $\delta$，我就承认 $\lim_{x \to x_0} f(x) = L$。

用数学符号，极限的定义可以写成：

\begin{center}
  \(\forall \epsilon > 0\)，
  \(\exists \delta > 0\)，使得\(0 < |x - x_0| < \delta \)，\(|f(x) - L| < \epsilon\)恒成立，则有\(\lim_{x \to x_0} f(x) = L\)
\end{center}

只需要对每个 $\epsilon$ 分别找到一个 $\delta$即可，$\delta$ 可以随着 $\epsilon$ 的改变而改变。

接下来我们再看看我们最开始举的例子：

\begin{myexample}
  举例说明：用 $\epsilon-\delta$ 定义证明 $\lim\limits_{x \to 1}2x = 2$。

  我们先回忆一下定义：对任意给定的 $\epsilon > 0$，都存在 $\delta > 0$，使得当 $0 < |x - x_0| < \delta$ 时，$|f(x) - L| < \epsilon$。

  也就是说，我们要证明的是，对任意给定的 $\epsilon > 0$，都存在 $\delta > 0$，使得当 $0 < |x - 1| < \delta$ 时，$|2x - 2| < \epsilon$。

  那么假设我们现在就是弓箭手，现在我们的目标就是找到一个$\delta > 0$，使得当 $0 < |x - 1| < \delta$ 时，$|2x - 2| < \epsilon$。

  要让 $|2x - 2| < \epsilon$，也就是$2|x - 1| < \epsilon$，只需 $|x - 1| < \frac{\epsilon}{2}$。因此取 $\delta < \frac{\epsilon}{2}$ 即可。

  那么我们不妨取$\delta = \frac{\epsilon}{3}$，那么我们可以得到这样的证明：

对任意的 $\epsilon > 0$，取 $\delta = \frac{\epsilon}{3}$。则当 $0 < |x - 1| < \delta$ 时，有：
\[
|2x - 2| = 2|x - 1| < 2\delta = \frac{2\epsilon}{3} < \epsilon
\]

由 $\epsilon-\delta$ 定义，$\lim\limits_{x \to 1} 2x = 2$。

\end{myexample}

我们用射箭游戏说明，就是我给出了任意的 $\epsilon$，你算出了 $\delta < \frac{\epsilon}{2}$，并取了$\delta = \frac{\epsilon}{3}$，然后在 $x_0=1$ 的 $\delta$ 去心邻域内任意位置射箭，每一箭的落点 $f(x)$ 都落在距离 $L=2$ 小于 $\epsilon$ 的范围内。游戏获胜。

总结一下：

先根据$|f(x) - L| < \epsilon$求出$\delta$的范围。

然后在$\delta$的范围内取一个合适的$\delta$。

代入定义，得到证明。

当然这只是简单的极限证明。我们了解即可，具体的使用那是数学分析才做的事情，更深入的我们这里不作讨论。

\section{极限的性质}

有了 $\epsilon-\delta$ 定义这个工具，我们可以严格地讨论极限的各种性质了。这些性质会在后续计算极限时用到，但是计算极限本身不是我们的重点，因此这里简单带过。

\subsection{单侧极限}

有时，函数在 $x_0$ 的两侧表现不同，比如分段函数在分段点处，如果还用$x_0$的$\delta $邻域定义极限，会发现两边的定义完全不同。这时我们就要区分“从左边趋近”和“从右边趋近”。

因此，我们就引出了左极限与右极限：

\begin{definition}{左极限与右极限}{}
设函数 $f(x)$ 在 $x_0$ 的左侧某个区间 $(x_0 - \delta_0,\ x_0)$ 内有定义，$L$ 为实数。

\textbf{左极限}：若对任意 $\epsilon > 0$，存在 $\delta > 0$，使得当 $x_0 - \delta < x < x_0$ 时，$|f(x) - L| < \epsilon$，则称 $L$ 为 $f(x)$ 在 $x_0$ 处的\textbf{左极限}，记作
\[
\lim_{x \to x_0^-} f(x) = L \quad \text{或} \quad f(x_0^-) = L
\]

设函数 $f(x)$ 在 $x_0$ 的右侧某个区间 $(x_0,\ x_0 + \delta_0)$ 内有定义，$L$ 为实数。

\textbf{右极限}：若对任意 $\epsilon > 0$，存在 $\delta > 0$，使得当 $x_0 < x < x_0 + \delta$ 时，$|f(x) - L| < \epsilon$，则称 $L$ 为 $f(x)$ 在 $x_0$ 处的\textbf{右极限}，记作
\[
\lim_{x \to x_0^+} f(x) = L \quad \text{或} \quad f(x_0^+) = L
\]
\end{definition}

左极限就是只看 $x_0$ 左边的点，右极限就是只看 $x_0$ 右边的点。用射箭游戏的比喻就是，左极限是只允许你在 $x_0$ 的左侧射箭，右极限是只允许你在 $x_0$ 的右侧射箭。

单侧极限最常见的应用场景就是分段函数在分段点处的极限问题。先分别求左右极限，再看它们是否相等，判断函数是否在该点处有极限。

\begin{myexample}
  设 $f(x) = |x|$，求 $\lim\limits_{x \to 0} f(x)$。

  当 $x < 0$ 时， $f(x) = -x$，因此 $\lim\limits_{x \to 0^-} |x| = \lim\limits_{x \to 0^-} -x = 0$。

  当 $x > 0$ 时， $f(x) = x$，因此 $\lim\limits_{x \to 0^+} |x| = \lim\limits_{x \to 0^+} x = 0$。

  左右极限都等于 $0$，故 $\lim\limits_{x \to 0} |x| = 0$。
\end{myexample}

\begin{myexample}
  设 $f(x) = \begin{cases} x+1, & x \geq 0 \\ x-1, & x < 0 \end{cases}$，讨论 $\lim\limits_{x \to 0} f(x)$。

  左极限：$\lim\limits_{x \to 0^-} f(x) = \lim\limits_{x \to 0^-} x-1 = -1$。

  右极限：$\lim\limits_{x \to 0^+} f(x) = \lim\limits_{x \to 0^+} x+1 = 1$。

  然而$-1 \neq 1$，即左右极限不相等，故 $\lim\limits_{x \to 0} f(x)$ 不存在。
\end{myexample}

也就是说，函数在一点有极限，就是在该点的左右极限都存在且相等，否则就说函数在这一点极限不存在。

\begin{theorem}{极限与单侧极限的关系}{}
函数在一点处极限存在的充要条件是该点左右极限都存在且相等。即：

  \begin{center}
    若\(\lim_{x \to x_0^-} f(x) = \lim_{x \to x_0^+} f(x) = L\)，则\(\lim_{x \to x_0} f(x) = L\)
  \end{center}

  反之亦成立。

\end{theorem}

\subsection{无穷远处的极限}

极限不仅关心 $x$ 趋近于某个有限点 $x_0$，也关心 $x$ 无限增大或无限减小时 $f(x)$ 的趋势。例如 $\lim\limits_{x \to \infty} \frac{1}{x} = 0$，当 $x$ 越来越大时，$\frac{1}{x}$ 越来越接近 $0$。那么我们就可以引出无穷远处的极限：

\begin{definition}{自变量趋于无穷时的极限}{}
设函数 $f(x)$ 当 $|x|$ 大于某个正数时有定义，$L$ 为实数。若对任意 $\epsilon > 0$，存在 $X > 0$，使得当 $|x| > X$ 时，$|f(x) - L| < \epsilon$，则称当 $x$ 趋于无穷大时，$f(x)$ 的极限为 $L$，记作
\[
\lim_{x \to \infty} f(x) = L
\]

其中，把 $|x| > X$ 换成 $x > X$，则得到 $x \to +\infty$ 时的极限；换成 $x < -X$，则得到 $x \to -\infty$ 时的极限。
\end{definition}

这里的结构和 $\epsilon-\delta$ 定义完全平行：$\epsilon$ 仍然控制函数值的精度，而 $\delta$ 的角色被 $X$ 取代。$X$ 控制的是自变量“离开原点”的距离。$\epsilon$ 越小，$X$ 通常需要越大。

在符号表示上，用箭头方向加以区分：
\[
\lim_{x \to +\infty} f(x) = L, \qquad \lim_{x \to -\infty} f(x) = L, \qquad \lim_{x \to \infty} f(x) = L
\]

和单侧极限一样，只有当\(\lim_{x \to +\infty} f(x) = \lim_{x \to -\infty} f(x) = L\)时，才能说\(\lim_{x \to \infty} f(x) = L\)。

让我们回到一开始的那个$\lim\limits_{x \to \infty} \frac{1}{x} = 0$的例子：

\begin{myexample}
  证明$\lim\limits_{x \to \infty} \frac{1}{x} = 0$。

  对任意 $ \epsilon>0$，要使 $| \frac {1}{x} -0|< \epsilon $ ，只需$|x|> \frac {1}{\epsilon } $ ，因为$|x| > X$，取$X= \frac {2}{\epsilon } $ 即可。

  所以对任意的$\epsilon >0$，取$ X= \frac {2}{\epsilon }$。则当$|x|>X$时，有：
  \[|\frac{1}{x}-0|=\frac{1}{|x|}<\frac{1}{X}=\frac {\epsilon }{2}<\epsilon \]

  由$\epsilon-X$ 定义，$\lim\limits_{x \to \infty} \frac{1}{x} = 0$。
\end{myexample}

\subsection{运算律}

如果每次求极限都要从 $\epsilon-\delta$ 定义出发，那就太麻烦了。幸好，极限满足一系列运算律，使得我们可以像做加减乘除一样计算极限。

\begin{theorem}{极限的四则运算法则}{}
若 $\lim\limits_{x \to x_0} f(x) = A$，$\lim\limits_{x \to x_0} g(x) = B$，其中$A$，$B$ 均为有限实数，则
\[
\begin{aligned}
&\text{(1)}\ \lim_{x \to x_0} [f(x) \pm g(x)] = A \pm B
&&\text{(2)}\ \lim_{x \to x_0} [f(x) \cdot g(x)] = A \cdot B \\[4pt]
&\text{(3)}\ \lim_{x \to x_0} [c\,f(x)] = c\,A \quad (c\text{ 为常数})
&&\text{(4)}\ \lim_{x \to x_0} \frac{f(x)}{g(x)} = \frac{A}{B} \quad (B \neq 0)
\end{aligned}
\]
\end{theorem}

简单来说就是：\textbf{极限可以加减乘除}，和的极限等于极限的和，积的极限等于极限的积，商的极限等于极限的商（分母极限不为零）。这些法则对 $x \to \infty$ 以及单侧极限也同样成立。

还有一个非常实用的工具叫夹逼定理：

\begin{theorem}{夹逼定理}{}
若在 $x_0$ 的某个去心邻域内，$g(x) \leq f(x) \leq h(x)$，且
\(\lim_{x \to x_0} g(x) = \lim_{x \to x_0} h(x) = L\)，
则\(\lim_{x \to x_0} f(x) = L\)。
\end{theorem}

夹逼定理的思想很简单：$f(x)$ 被 $g(x)$ 和 $h(x)$ 两个函数“夹”在中间，而这两个极限都趋向同一个值 $L$，那么中间的 $f(x)$ 别无选择，也只能趋向 $L$。不过本章中，我们不会对如何求极限作详细说明，因为我们的重点是微积分。

\subsection{无穷小量与无穷大量}

有了极限的定义，我们可以严格地讨论“小到几乎为零的量”和“大到没有上界的量”。

\begin{definition}{无穷小量}{}
若 $\lim\limits_{x \to x_0} f(x) = 0$（或 $x \to \infty$ 时），则称 $f(x)$ 为该极限过程中的一个\textbf{无穷小量}。
\end{definition}

类似地，我们有：

\begin{definition}{无穷大量}{}
若对任意 $M > 0$，存在 $\delta > 0$（或 $X > 0$），使得当 $0 < |x - x_0| < \delta$（或 $|x| > X$）时，$|f(x)| > M$，则称 $f(x)$ 为该极限过程中的一个\textbf{无穷大量}，记作
\[
\lim_{x \to x_0} f(x) = \infty
\]
\end{definition}

需要注意的是，无穷小量不是“一个很小的数”，而是在极限过程中，绝对值可以小于任何预先给定的正数 $\epsilon$ 的量。$0 $是唯一可以称为无穷小量的常数。因此，趋于无穷小也说成趋于零。同理，无穷大量不是“一个很大的数”，也不是一个确定的数，所以当我们说极限趋于无穷大时，这个极限其实是不存在的。

我们举个例子：

\begin{myexample}
  验证 $\lim\limits_{x \to +\infty} x^2 = +\infty$。

  对任意 $M > 0$，取 $X = \sqrt{M}$。当 $x > X$ 时，$x^2 > X^2 = M$。所以 $\lim_{x \to +\infty} x^2 = +\infty$。
\end{myexample}

也就是说，当$x \to +\infty$时，$x^2 \to +\infty$，极限不存在。

无穷大量和无穷小量之间有一个简洁的倒数关系：

\begin{theorem}{无穷大与无穷小的关系}{}
在同一极限过程中：

1.若 $f(x)$ 是无穷小量且 $f(x) \neq 0$，则 $\frac{1}{f(x)}$ 是无穷大量。

2.若 $f(x)$ 是无穷大量，则 $\frac{1}{f(x)}$ 是无穷小量。

\end{theorem}

这个关系非常直观，把几乎为零的数倒过来，就变成大到没边的数，反之亦然。

有时我们需要比较两个无穷小量“谁小得更快”，这就引出了无穷小量的阶的概念：

\begin{definition}{无穷小量的阶}{}
设 $\alpha(x)$ 和 $\beta(x)$ 都是 $x \to x_0$ 时的无穷小量，且当 $x \neq x_0$ 时， $\beta(x) \neq 0$。定义极限
\(
\lim_{x \to x_0} \frac{\alpha(x)}{\beta(x)}
\)：

1.若$\lim_{x \to x_0} \frac{\alpha(x)}{\beta(x)}
=0$，则称 $\alpha$ 是比 $\beta$ \textbf{高阶}的无穷小，记作 $\alpha = o(\beta)$。

2.若$\lim_{x \to x_0} \frac{\alpha(x)}{\beta(x)}
=c$，其中$c \neq 0$且为常数，则称 $\alpha$ 与 $\beta$ 是\textbf{同阶}无穷小。

3.若$\lim_{x \to x_0} \frac{\alpha(x)}{\beta(x)}
=1$，则称 $\alpha$ 与 $\beta$ 是\textbf{等价}无穷小，记作 $\alpha \sim \beta$。

4.若$\lim_{x \to x_0} \frac{\alpha(x)}{\beta(x)}
=\infty$，则称 $\alpha$ 是比 $\beta$ \textbf{低阶}的无穷小。此时 $\beta$ 就是比 $\alpha$ 高阶的无穷小，即 $\beta = o(\alpha)$。

5.若$\lim_{x \to x_0} \frac{\alpha(x)}{[\beta(x)]^k}=c$，其中$c \neq 0$且为常数，则称 $\alpha$ 是关于$\beta$的 \textbf{$k$ 阶}无穷小。通常以$\beta (x) = x-x_0$为基准，称 $\alpha$ 是$x$的\textbf{$k$ 阶}无穷小。

\end{definition}

越高阶的无穷小，趋向于$0 $的速度就越快。$0 $通常被视为“最高阶”的无穷小，没有比$ 0 $更高阶的无穷小了。

其中，最常用的是等价无穷小。等价无穷小在求极限时非常有用，只要是乘除运算，就可以把一个等价无穷小替换成另一个，大大简化计算。

\section{函数的连续性}

极限只关心 $x$ 趋近于 $x_0$ 时 $f(x)$ 的趋势，不关心 $f(x_0)$ 本身。现在我们把视线拉回来，同时看“趋近”和“在点处”，如果这两者能对上，我们就说函数在这一点是连续的。连续性是比极限存在更“强”的性质。

\subsection{连续的不同定义}

历史上数学家从多个角度描述“连续”，最终它们被证明是等价的。我们先给出最标准的 $\epsilon-\delta$ 定义。

\begin{definition}{函数在一点处的连续性}{}
设函数 $f(x)$ 在 $x_0$ 的某个邻域内有定义。若对任意 $\epsilon > 0$，存在 $\delta > 0$，使得当 $|x - x_0| < \delta$ 时，有
\[
|f(x) - f(x_0)| < \epsilon
\]
则称 $f(x)$ 在点 $x_0$ 处\textbf{连续}。
\end{definition}

对比极限的定义，连续性只改了两处：不再挖掉 $x_0$，把极限值 $L$ 换成了函数值 $f(x_0)$。换句话说：

\begin{theorem}{连续与极限的关系}{}
$f(x)$ 在 $x_0$ 处连续是$\lim\limits_{x \to x_0} f(x) = f(x_0)$的充要条件。

\end{theorem}

这个关系暗含了函数在某一点连续的三个条件：极限存在、函数值存在、两者相等。这三者缺一不可。任何一条不满足，我们就说函数在 $x_0$ 处间断。由此引出间断点的定义：

\begin{definition}{间断点}{}
  若函数$f(x)$在点$x_0$不连续，则称$x_0$为该函数的间断点。
\end{definition}

用$\epsilon-\delta$ 定义函数的连续是现代数学分析的基石。如果我们在直观上理解，还可以用“增量”定义连续：令 $\Delta x = x - x_0$（自变量的微小变化），$\Delta y = f(x_0 + \Delta x) - f(x_0)$（函数值的相应变化），则函数的连续性等价于：
\[
\lim_{\Delta x \to 0} \Delta y = 0
\]

也就是说，如果自变量的微小变化只引起函数值的微小变化，我们就说这个函数连续。这个定义实际上与$\epsilon-\delta$ 定义是等价的。

和单侧极限类似，我们有单侧连续：

\begin{definition}{左连续与右连续}{}
若 $\lim\limits_{x \to x_0^-} f(x) = f(x_0)$，则称 $f(x)$ 在 $x_0$ 处\textbf{左连续}。

若 $\lim\limits_{x \to x_0^+} f(x) = f(x_0)$，则称 $f(x)$ 在 $x_0$ 处\textbf{右连续}。

\end{definition}

同理可得，$f(x)$ 在 $x_0$ 处连续的充要条件是$f(x)$在 $x_0$ 处既左连续又右连续。

\begin{theorem}{连续与单侧连续的关系}{}
  函数在一点处连续的充要条件是函数在该点左右都连续。即：

  \begin{center}
    若\(\lim_{x \to x_0^-} f(x) = \lim_{x \to x_0^+} f(x) = f(x_0)\)，则\(f(x)\)在\(x_0\)处连续。
  \end{center}

  反之亦成立。

\end{theorem}

根据“哪里不连续”，我们可以对间断点进行分类：

\begin{definition}{间断点的分类}{}
设 $x_0$ 是 $f(x)$ 的间断点。

\textbf{第一类间断点}：左右极限都存在。其中：

若左右极限相等，则称 $x_0$ 为\textbf{可去间断点}。

若左右极限不相等，则称 $x_0$ 为\textbf{跳跃间断点}，差值 $|f(x_0^+) - f(x_0^-)|$ 称为\textbf{跃度}。

\textbf{第二类间断点}：左右极限中至少有一个不存在。
\end{definition}

\begin{myexample}
  判断$f(x) = |x|$ 在 $x=0$ 处的连续性。

  我们先回忆一下连续的三个条件：极限存在、函数值存在、两者相等。

  先看极限存不存在。计算左右极限：

  左极限：
\[
\lim_{x \to 0^-} |x| = \lim_{x \to 0^-}-x= 0
\]

右极限：
\[
\lim_{x \to 0^+} |x| = \lim_{x \to 0^+}x= 0
\]

左极限等于右极限，因此函数在$x=0$ 处极限存在，也就是：
\[
\lim_{x \to 0} |x| =0
\]

再看函数值：
\[
f(0) = 0
\]

极限和函数值相等。三个条件都满足，因此$f(x) = |x|$ 在 $x=0$ 处连续。

\end{myexample}

\begin{myexample}
    判断分段函数 $f(x) = \begin{cases} x+1, & x \geq 0 \\ x-1, & x < 0 \end{cases}$ 在 $x=0$ 处的连续性。

    我们同样先计算极限：

    左极限：
\[
\lim_{x \to 0^-} f(x) = \lim_{x \to 0^-} (x-1) = -1
\]

右极限：
\[
\lim_{x \to 0^+} f(x) = \lim_{x \to 0^+} (x+1) = 1
\]

左右极限不相等，因此$f(x)$ 在 $x=0$ 处不连续。

更进一步的，$x=0$是函数 $f(x)$ 的跳跃间断点，跃度为$2$。
\end{myexample}

在这一章最开始，我们提到了狄利克雷函数 $D(x)$ ，它就是一个处处不连续的函数，在任何点的任何邻域内，函数值在 $0$ 和 $1$ 之间不断跳变，极限根本不存在。这是第二类间断点的极端例子。

\subsection{闭区间上连续函数的性质}

闭区间上的连续函数具有几条非常“舒服”的性质。它们意义重大，但我们不给出严格证明，直观接受即可。

\begin{theorem}{有界性定理}{}
若 $f(x)$ 在闭区间 $[a,b]$ 上连续，则 $f(x)$ 在 $[a,b]$ 上有界。
\end{theorem}

\begin{theorem}{最值定理}{}
若 $f(x)$ 在闭区间 $[a,b]$ 上连续，则 $f(x)$ 在 $[a,b]$ 上一定能取到最大值和最小值。
\end{theorem}

以上两个定理，都要求 $f(x)$ 在闭区间 $[a,b]$ 上连续。开区间上的连续函数可能无界，闭区间上的不连续函数也可能取不到最值。

\begin{theorem}{介值定理}{}
若 $f(x)$ 在闭区间 $[a,b]$ 上连续，且 $f(a) \neq f(b)$，则对于 $f(a)$ 与 $f(b)$ 之间的任意一个数 $\mu$，都存在至少一个 $\xi \in (a,b)$，使得 $f(\xi) = \mu$。

\end{theorem}

也就是说，连续函数在闭区间上能取到两个端点值之间的所有值。

介值定理的直接推论是零点定理：

\begin{theorem}{零点定理}{}
若 $f(x)$ 在闭区间 $[a,b]$ 上连续，且 $f(a) \cdot f(b) < 0$，则至少存在一个 $\xi \in (a,b)$，使得 $f(\xi) = 0$。
\end{theorem}

零点定理的几何意义非常直观：一条连续的曲线，如果它从 $x$ 轴下方出发到达上方（或反过来），中间必然穿过 $x$ 轴至少一次。

还有一个非常重要的定理，它是我们后面讨论可积和可微的基础：

\begin{theorem}{初等函数的连续性}{}
  一切初等函数在其定义域内处处连续。
\end{theorem}

这个定理是初等函数的一个重要性质。由于这些函数在其定义域内处处连续，因此我们可以放心地在这些函数上进行极限运算。
% !TEX root = ../main.tex
\chapter{一元函数微分学}
这一章节，我们将进入微分学，从高中的导数出发，我们假定你有良好的导数基础，因此我们会跳过大量前置知识，并默认你已了解。


\section{微分}

当自变量 $x$ 发生一个微小变化 $\Delta x$ 时，函数值 $y = f(x)$ 的变化量 $\Delta y$ 能否用 $\Delta x$ 的线性表达式来近似？也就是说，是否可以构造一个形似$y=kx$的表达式来近似$\Delta y$？更关键的是，这个近似的误差是否比 $\Delta x$ 本身小得多？由此我们引入可微的概念：

\subsection{可微与微分}

\begin{definition}{可微的定义}{}
  设函数 $y = f(x)$ 在点 $x_0$ 的某个邻域内有定义。记自变量的增量为 $\Delta x$，函数值的增量为$\Delta y$，则：
  \[
  \Delta y = f(x_0 + \Delta x) - f(x_0)
  \]
  若存在一个与 $\Delta x$ 无关的常数 $A$，使得：
\begin{center}
    当$\Delta x \to 0$时，$\Delta y = A \cdot \Delta x + o(\Delta x)$
\end{center}
  我们就说 $f(x)$ 在 $x_0$ 处\textbf{可微}。
\end{definition}

这个式子的意思是，函数值的变化量 $\Delta y$ 可以分解为两部分。第一部分 $A \cdot \Delta x$ 是 $\Delta x$ 的线性函数，称为\textbf{线性主部}。第二部分 $o(\Delta x)$ 是比 $\Delta x$ 更高阶的无穷小，它除以 $\Delta x$ 的商趋于 $0$。也就是说，当 $\Delta x$ 足够小时，$\Delta y$ 的主体就是 $A \Delta x$，$o(\Delta x)$比$\Delta x$ 趋于$0$的速度更快，误差可以忽略不计。

由此，我们可以引出微分的定义：

\begin{definition}{微分}{}
若 $y=f(x)$ 在 $x_0$ 处可微，则称线性主部 $A \Delta x$ 为 $y$ 在点 $x_0$ 处的\textbf{微分}，记作
\[
dy = A \Delta x
\]
对于自变量 $x$ 本身，$dx = \Delta x$。因此$y$ 对 $x$ 的微分也常记作
\[
dy = A dx
\]
\end{definition}

这里的常数 $A$ 是什么？我们把 $\Delta y = A \cdot \Delta x + o(\Delta x)$ 两边同时除以 $\Delta x$：
\[
\frac{\Delta y}{\Delta x} = A + \frac{o(\Delta x)}{\Delta x}
\]

又因为 $\Delta x \to 0$，所以$\frac{o(\Delta x)}{\Delta x}$趋于$0$，于是：
\[
\lim_{\Delta x \to 0} \frac{\Delta y}{\Delta x} = A
\]

也就是说，$A$ 不是别的，正是我们熟悉的导数。

所以可微等价于可导，且 $A = f'(x_0)$。微分和导数的关系可以写成：

\begin{theorem}{微分和导数的关系}{}
  若 $y=f(x)$ 在 $x_0$ 处可微，则：
  \[
  dy = f'(x_0)\, dx
  \]

  反之亦成立。
\end{theorem}

这个关系实则暗含了两层意思：可微必可导，可导必可微。

因此，莱布尼茨用记号 $\frac{dy}{dx}$ 表示 $f'(x)$，这表示导数是两个微小量 $dy$ 和 $dx$ 的比值，只不过这个比值在取极限后从近似变成了精确。所以，导数也称作“微商”（不是微信电商的意思）。

从几何上看，$dy = f'(x_0)\,\Delta x$ 是切线上的增量，也就是沿切线方向走 $\Delta x$ 所上升的高度，而 $\Delta y$ 是曲线上的真实增量。当 $\Delta x \to 0$ 时，两者的差 $|\Delta y - dy|$ 是比 $\Delta x$ 更高阶的无穷小。这也体现了微分“以直代曲”的数学思想，在极小范围内，函数的曲线和它的切线几乎重合。

\begin{myexample}
  设 $y = x^3$，讨论函数在 $x = 1$ 处的可微性，并求 $\Delta x = 0.01$ 时的微分和真实增量。

  先求 $A$。计算极限：
  \[
  \lim_{\Delta x \to 0} \frac{(1+\Delta x)^3 - 1^3}{\Delta x}
  = \lim_{\Delta x \to 0} \frac{3\Delta x + 3(\Delta x)^2 + (\Delta x)^3}{\Delta x}
  = 3
  \]

  因此 $f(x) = x^3$ 在 $x=1$ 处可微，且 $A = 3$。

  微分：$dy = A \Delta x = 3 \cdot 0.01 = 0.03$。

  真实增量：$\Delta y = (1.01)^3 - 1^3 = 1.030301 - 1 = 0.030301$。

  误差： $|\Delta y - dy| = 0.000301$，是 $(\Delta x)^2$ 量级的，确实比 $\Delta x$ 小得多。随着 $\Delta x$ 减小，误差会以更快的速度趋近于$0$。
\end{myexample}

\subsection{微分运算法则}

我们已经知道微分与导数的关系 $dy = f'(x)\,dx$，因此求微分的本质仍然是求导。但微分有自己独立的运算法则，它们在形式上比导数法则更对称，在积分学中尤其方便。我们默认你已经掌握了基本初等函数的求导公式，此处将它们转写成微分形式：

\begin{theorem}{微分公式}{}
  \[
  \begin{aligned}
  &\text{(1)}\ d(C) = 0
  &&\text{(2)}\ d(x^\mu) = \mu x^{\mu-1}\, dx \\[4pt]
  &\text{(3)}\ d(\sin x) = \cos x\, dx
  &&\text{(4)}\ d(\cos x) = -\sin x\, dx \\[4pt]
  &\text{(5)}\ d(\tan x) = \sec^2 x\, dx
  &&\text{(6)}\ d(\cot x) = -\csc^2 x\, dx \\[4pt]
  &\text{(7)}\ d(\sec x) = \sec x \tan x\, dx
  &&\text{(8)}\ d(\csc x) = -\csc x \cot x\, dx \\[4pt]
  &\text{(9)}\ d(e^x) = e^x\, dx
  &&\text{(10)}\ d(a^x) = a^x \ln a\, dx \\[4pt]
  &\text{(11)}\ d(\ln x) = \frac{1}{x}\, dx
  &&\text{(12)}\ d(\log_a x) = \frac{1}{x \ln a}\, dx \\[4pt]
  &\text{(13)}\ d(\arcsin x) = \frac{1}{\sqrt{1-x^2}}\, dx
  &&\text{(14)}\ d(\arccos x) = -\frac{1}{\sqrt{1-x^2}}\, dx \\[4pt]
  &\text{(15)}\ d(\arctan x) = \frac{1}{1+x^2}\, dx
  \end{aligned}
  \]
\end{theorem}

这些公式不需要死记，只需要将任何一个求导公式两边乘以 $dx$，就得到了对应的微分公式。

接下来是微分的四则运算法则。设 $u = u(x)$，$v = v(x)$ 都可微，则有：

\begin{theorem}{微分的四则运算法则}{}
\[
\begin{aligned}
&\text{(1)}\ d(u \pm v) = du \pm dv \\[4pt]
&\text{(2)}\ d(cu) = c\, du \quad (c\text{ 为常数}) \\[4pt]
&\text{(3)}\ d(uv) = u\, dv + v\, du \\[4pt]
&\text{(4)}\ d(\frac{u}{v}) = \frac{v\, du - u\, dv}{v^2} \quad (v \neq 0)
\end{aligned}
\]
\end{theorem}

这些法则与导数法则一一对应，只是换了一个写法。

也就是说，我们怎么求导，就怎么求微分，最后别忘记加上$dx$和$dy$就行。

\begin{myexample}
  设 $y = x^2 \sin x$，求 $dy$。

  用乘法法则：令 $u = x^2$，$v = \sin x$。则
  \(du = 2x\, dx\)，\( dv = \cos x\, dx\)。

  所以：
  \[
  dy = u\, dv + v\, du = x^2 \cdot \cos x\, dx + \sin x \cdot 2x\, dx = (x^2 \cos x + 2x \sin x)\, dx
  \]
\end{myexample}

\begin{myexample}
  设 $y = \dfrac{x}{1+x^2}$，求 $dy$。

  用除法法则：令 $u = x$，$v = 1+x^2$。则 $du = dx$，$dv = 2x\, dx$。

  所以：
  \[
  dy = \frac{v\, du - u\, dv}{v^2}
  = \frac{(1+x^2) \cdot dx - x \cdot 2x\, dx}{(1+x^2)^2}
  = \frac{1 + x^2 - 2x^2}{(1+x^2)^2}\, dx
  = \frac{1 - x^2}{(1+x^2)^2}\, dx
  \]
\end{myexample}

\section{导数}

导数在高中时期可谓是令人闻风丧胆。不过在这里，我们不会讨论导数难题，我们也不需要。和高中不同，高中导数学得深，而微分学的导数学得广。为什么广？往下看就知道了。

\subsection{可导与连续的关系}
我们假设你还记得导数的定义。不过不记得也没关系，现在帮你回忆一下：

\begin{definition}{导数}{}
设函数 $y = f(x)$ 在点 $x_0$ 的某个邻域内有定义。若极限
\[
\lim_{\Delta x \to 0} \frac{f(x_0 + \Delta x) - f(x_0)}{\Delta x}
\]
存在，则称此极限值为 $f(x)$ 在点 $x_0$ 处的\textbf{导数}，记作
\[
f'(x_0),\quad y'\big|_{x=x_0},\quad \frac{dy}{dx}\Big|_{x=x_0},\quad \frac{df}{dx}\Big|_{x=x_0}
\]
此时称 $f(x)$ 在点 $x_0$ 处\textbf{可导}。
\end{definition}

导数的几何意义是曲线在 $(x_0, f(x_0))$ 处的切线斜率，物理意义是瞬时变化率。

关于导数和连续两个概念，我们有这样的关系：

\begin{theorem}{可导必连续}{}
若 $f(x)$ 在点 $x_0$ 处可导，则 $f(x)$ 在点 $x_0$ 处必连续。
\end{theorem}

要证明这个结论并不难，读者可尝试自行证明。这里给出证明的大致思路：

我们知道，若 $f(x)$ 在点 $x_0$ 处可导，则 $f'(x_0)$ 存在，那么有：
\[
\lim_{\Delta x \to 0} \frac{f(x_0 + \Delta x) - f(x_0)}{\Delta x} = f'(x_0)
\]

因为 $\Delta y = f(x_0 + \Delta x) - f(x_0)$，那么：
\[
\lim_{\Delta x \to 0} \Delta y
= \lim_{\Delta x \to 0} ( \frac{\Delta y}{\Delta x} \cdot \Delta x )
= f'(x_0) \cdot 0 = 0
\]

而这正是连续性的增量语言 $\lim_{\Delta x \to 0} \Delta y = 0$，所以 $f(x)$ 在 $x_0$ 处连续。

直观上理解，导数存在意味着 $\frac{\Delta y}{\Delta x}$ 趋近于一个有限值，而当 $\Delta x$趋近于$0$时，$\Delta y$ 也必须趋近于 $0$，否则此时的导数就是$\frac{\text{常数}}{0}=\infty$，这不是一个有限值。

如果函数在某点不连续，那么它可能连极限都不存在，导数就更无从谈起了。所以\textbf{不连续一定不可导}。

我们能证明，\textbf{初等函数在其定义区间内都是连续的}。

$f(x)$ 在点 $x_0$ 处可导，则 $f(x)$ 在点 $x_0$ 处必连续。反过来成立吗？不成立，也就是说，函数连续却可能没有导数。

与极限和连续类似，我们先定义左导数和右导数：

\begin{definition}{左导数与右导数}{}
设函数 $y = f(x)$ 在点 $x_0$ 的某个左邻域或右邻域内有定义。

\textbf{左导数}：若极限
\[
\lim_{\Delta x \to 0^-} \frac{f(x_0 + \Delta x) - f(x_0)}{\Delta x}
\]
存在，则称此极限值为 $f(x)$ 在点 $x_0$ 处的\textbf{左导数}，记作 $f'_-(x_0)$。

\textbf{右导数}：若极限
\[
\lim_{\Delta x \to 0^+} \frac{f(x_0 + \Delta x) - f(x_0)}{\Delta x}
\]
存在，则称此极限值为 $f(x)$ 在点 $x_0$ 处的\textbf{右导数}，记作 $f'_+(x_0)$。

\end{definition}

左导数与右导数统称为单侧导数。与极限和连续类似，$f(x)$ 在 $x_0$ 处可导的充要条件是$f(x)$ 在 $x_0$ 处的左右导数都存在且相等。

\begin{theorem}{导数与单侧导数的关系}{}
  函数在一点可导的充要条件是函数在该点处的左右导数都存在且相等。即：
  \begin{center}
    若$f'_-(x_0)=f'_+(x_0)=f'(x_0)$，则$f(x)$在$x_0$处可导。
  \end{center}
  反之亦成立。
\end{theorem}

所以，即使函数连续，也无法推出左右导数相等，因此函数连续不一定可导。我们来看一个例子。

\begin{myexample}
  判断$f(x) = |x|$ 在 $x = 0$ 处的连续性和是否可导。

  连续性：由第一章已知 $\lim_{x \to 0} |x| = 0 = f(0)$，连续。

  可导性：先计算左右导数。
  
    左导数：$\quad \lim_{\Delta x \to 0^-} \frac{|0 + \Delta x| - 0}{\Delta x}
    = \lim_{\Delta x \to 0^-} \frac{-\Delta x}{\Delta x} = -1$

    右导数：$\quad \lim_{\Delta x \to 0^+} \frac{|0 + \Delta x| - 0}{\Delta x}
    = \lim_{\Delta x \to 0^+} \frac{\Delta x}{\Delta x} = 1$
  
  左右导数不相等，故导数不存在，因此$f(x) = |x|$ 在 $x = 0$ 处连续但不可导。

\end{myexample}

那么$f(x) = |x|$ 在 $x = 0$ 处可微吗？刚刚我们说到，可微必可导，可导必可微，所以$f(x) = |x|$ 在 $x = 0$ 处不可微。

几何上看，$f(x)=|x|$ 在 $x=0$ 处有一个“尖角”，左边的切线斜率和右边的切线斜率不一样，画不出一条唯一的切线。

总结一下：

\textbf{连续是可导的必要条件，但不是充分条件}。可导必连续，连续不一定可导。具体是否可导，看左右导数是否相等。

\subsection{高阶导数}

导数的导数称为二阶导数，二阶导数的导数称为三阶导数，以此类推。一般地，$n$ 阶导数就是求导 $n$ 次的结果。记号上，$y = f(x)$ 的高阶导数有两种主流写法：
\begin{definition}{高阶导数的记号}{}
  若函数$y = f(x)$ 有任意阶导数，则其高阶导数记为：

  拉格朗日记号：
  \[
  f'(x),\quad f''(x),\quad f'''(x),\quad f^{(4)}(x),\quad \dots,\quad f^{(n)}(x)
  \]
  
  莱布尼茨记号：
  \[
  \frac{dy}{dx},\quad \frac{d^2y}{dx^2},\quad \frac{d^3y}{dx^3},\quad \dots,\quad \frac{d^ny}{dx^n}
  \]
  
\end{definition}

莱布尼茨记号的写法值得解释一下。$\frac{d^2y}{dx^2}$ 的意思是“对 $y$ 取两次微分除以 $dx$ 的平方”：
\[
\frac{d^2y}{dx^2} = \frac{d}{dx} (\frac{dy}{dx})
\]

分子上的 $d^2y$ 表示微分运算执行了两次，分母上的 $dx^2$ 实际上是 $(dx)^2$ 的简写。这个记号的优点是操作感极强：求两次导就是分子上两个 $d$，分母上两个 $dx$。

看几个简单例子：

\begin{myexample}
  求 $y = x^4$ 的各阶导数。

  $y' = 4x^3$，$y'' = 12x^2$，$y''' = 24x$，$y^{(4)} = 24$，$y^{(5)} = y^{(6)} = \cdots = 0$。

\end{myexample}

由此，我们有：
\begin{theorem}{$n$ 次多项式的高阶导数}{}
  一般地，$n$ 次多项式的 $n$ 阶导数为常数，$n+1$ 阶及以上的导数全为零。
\end{theorem}


\begin{myexample}
  求 $y = \sin x$ 的高阶导数。

  $\sin x$ 的导数呈周期为 $4$ 的循环：
  \[
  (\sin x)' = \cos x,\quad (\sin x)'' = -\sin x,\quad (\sin x)''' = -\cos x,\quad (\sin x)^{(4)} = \sin x
  \]
\end{myexample}

由此，我们还能推出$\cos x$的规律：
\begin{theorem}{三角函数的高阶导数}{}
  一般地，$(\sin x)^{(n)} = \sin (x + \frac{n\pi}{2})$，$(\cos x)^{(n)} = \cos (x + \frac{n\pi}{2})$。
\end{theorem}

\subsection{复合函数求导}

前面讨论的都是直接对 $x$ 求导，但是有时候函数往往是一层套一层的，例如 $y = \sin x^2$，它不是单纯的 $\sin$，而是 $\sin$ 里面又包了一个 $x^2$。这种函数的求导法则称为链式法则。

\begin{theorem}{链式法则}{}
设 $y = f(u)$ 在 $u_0$ 处可导，$u = g(x)$ 在 $x_0$ 处可导，且 $u_0 = g(x_0)$，则复合函数 $y = f(g(x))$ 在 $x_0$ 处可导，且
\[
\frac{dy}{dx} = f'(u_0) \cdot g'(x_0) = f'(g(x_0)) \cdot g'(x_0)
\]
即：
\[
\frac{dy}{dx} = \frac{dy}{du} \cdot \frac{du}{dx}
\]
\end{theorem}

$\frac{dy}{dx} = \frac{dy}{du} \cdot \frac{du}{dx}$的形式上就像是分数约分，$du$ 上下抵消。不过这只是一个记忆技巧，并非严格的数学推导，但它极好地捕捉了链式法则的直觉。

链式法则可以推广到任意多层的复合。三层复合时：
\[
\frac{dy}{dx} = \frac{dy}{du} \cdot \frac{du}{dv} \cdot \frac{dv}{dx}
\]

形象地说，链式法则就像剥洋葱，从最外层一层层往内求导，每一步的导数依次相乘。

我们来看最开始的例子：
\begin{myexample}
  求 $y = \sin x^2$ 的导数。

  令 $u = x^2$，则$y = \sin u$，则有：
  \[
  \frac{dy}{du} = \cos u, \qquad \frac{du}{dx} = 2x
  \]

  由链式法则：
  \[
  \frac{dy}{dx} = \cos u \cdot 2x = 2x \cos x^2
  \]
\end{myexample}

\begin{myexample}
  求 $y = e^{3x}$ 的导数。

  令 $u = 3x$，则$y = e^u$，则有：
  \[
  \frac{dy}{du} = e^u, \qquad \frac{du}{dx} = 3
  \]

  由链式法则：
  \[
  \frac{dy}{dx} = e^u \cdot 3 = 3e^{3x}
  \]
\end{myexample}

熟练之后可以省略设 $u$ 的步骤，直接用“外导乘内导”的口诀：先对整体求导，里面不变，再乘以内部的导数。例如：

\begin{myexample}
  求 $y = \ln(\cos x^2)$ 的导数。

  这是一个三层复合：最外层是 $\ln$，中间层是 $\cos$，最内层是 $x^2$。按照剥洋葱的步骤，从最外层一层层往里求导，每一步的导数依次相乘：

  最外层 $y = \ln u$，对 $u$ 求导得 $\frac{1}{u}$，里面 $u = \cos x^2$ 不变。

  中间层 $u = \cos v$，对 $v$ 求导得 $-\sin v$，里面 $v = x^2$ 不变。

  最内层 $v = x^2$，对 $x$ 求导得 $2x$。

  三步连乘：
  \[
  \frac{dy}{dx} = \frac{1}{\cos x^2} \cdot (-\sin x^2) \cdot 2x = -2x \cdot \frac{\sin x^2}{\cos x^2} = -2x \tan x^2
  \]

\end{myexample}

常犯的一个错误是只求了外层忘了内层。例如把 $\sin x^2$ 的导数写成 $\cos x^2$ 就停了。记住一个简单的自检方法，如果内层不是单纯的 $x$，链式法则就一定要出手，每一步链都不能掉。

\subsection{参数方程求导}

有时候，$x$ 和 $y$ 并不以 $y = f(x)$ 的形式直接关联，而是各自通过第三个变量 $t$ 来表达：
\(\begin{cases}
x = \varphi (t) \\
y = \phi (t)
\end{cases}\)，那么问题来了：怎么求 $y$ 对 $x$ 的导数？

微分给出了答案：

\begin{theorem}{参数方程求导法则}{}
  若由参数方程\(\begin{cases}
x = \varphi (t) \\
y = \phi (t)
\end{cases}\)确定的$y=f(x)$在对应区间上可导，则其导数为：

\[
\frac{dy}{dx} = \frac{\phi' (t)}{\varphi' (t)}
\]

\end{theorem}

由微分定义，我们有：

\[
dx=\varphi' (t)dt, \qquad dy=\phi' (t)dt
\]

两者相除，我们就得到了上述公式。

也就是说，由参数方程确定的$y=f(x)$，求 $y$ 对 $x$ 的导数，等于 $y$ 和 $x$ 各自对参数 $t$ 的导数之商。

\begin{myexample}
  已知圆的参数方程 \(\begin{cases}
x = \cos t \\
y = \sin t
\end{cases}\)，$t \neq k\pi$，求 $\frac{dy}{dx}$。

  先求 $x$ 和 $y$ 对 $t$ 的导数：
  \[
  \frac{dx}{dt} = -\sin t, \qquad \frac{dy}{dt} = \cos t
  \]

  代入公式：
  \[
  \frac{dy}{dx} = \frac{\cos t}{-\sin t} = -\cot t
  \]

\end{myexample}

如果需要求参数方程的二阶导数，按定义推导，由于\(
\frac{d^2y}{dx^2} = \frac{d}{dx} (\frac{dy}{dx}) = \frac{d(\frac{dy}{dx})}{dx}\)，所以先把$\frac{dy}{dx}$看成新的$y=\psi (t)$，得到$d\psi =\psi' (t) dt$，再求$dx$。此时$dx$还是不变，即$dx=\varphi' (t)dt$，所以此时的二阶导数应该是$\frac{d^2y}{dx^2} = \frac{d\psi }{dx} = \frac{\psi' (t)}{\varphi' (t)}$。

全部写开，我们得到：

\begin{theorem}{参数方程的二阶导数}{}
    若由参数方程\(\begin{cases}
x = \varphi (t) \\
y = \phi (t)
\end{cases}\)确定的$y=f(x)$在对应区间上有二阶导数，则其二阶导数为：
\[
\frac{d^2y}{dx^2} = \frac{\phi''(t)\,\varphi'(t) - \phi'(t)\,\varphi''(t)}{[\varphi'(t)]^3}
\]
\end{theorem}

注意二阶导不能直接写成$\frac{\phi ''(t)}{\varphi ''(t)}$，这是一个常见的错误，必须按定义先求一阶导再对 $t$ 求一次。

\begin{myexample}
  已知圆的参数方程 \(\begin{cases}
x = \cos t \\
y = \sin t
\end{cases}\)，$t \neq k\pi$，求$\frac{d^2 y}{dx^2}$。

刚刚我们已经求了$dx$和一阶导：
  \[
  dx = -\sin t dt, \qquad \frac{dy}{dx} = -\cot t
  \]

  下面用两种方法求 $\frac{d^2y}{dx^2}$：

  方法一：二阶导按定义为：
  \[
\frac{d^2y}{dx^2} = \frac{d}{dx} (\frac{dy}{dx}) = \frac{d(\frac{dy}{dx})}{dx}
  \]
  
  已求得 $\frac{dy}{dx} = -\cot t$，令$\psi (t)= \frac{dy}{dx} = -\cot t$，对 $t$ 求导：
  \[
  \frac{d\psi }{dt} = \csc^2 t = \frac{1}{\sin^2 t}
  \]

  代入得：
  \[
  \frac{d^2y}{dx^2} = \frac{d\psi }{dx} = \frac{d\psi }{dt} \cdot \frac{dt}{dx} = \frac{1}{\sin^2 t} \cdot \frac{1}{-\sin t} = -\frac{1}{\sin^3 t} = -\csc^3 t
  \]

  方法二：参数方程的二阶导公式为：
\[
\frac{d^2y}{dx^2} = \frac{\phi''(t)\,\varphi'(t) - \phi'(t)\,\varphi''(t)}{[\varphi'(t)]^3}
\]

  先求二阶导：
  \[\varphi ''(t) = (-\sin t)' = -\cos t, \qquad \phi ''(t) = (\cos t)' = -\sin t\]
  
  代入公式：
  \[
  \frac{d^2y}{dx^2}
  = \frac{(-\sin t)(-\sin t) - (\cos t)(-\cos t)}{(-\sin t)^3}
  = \frac{\sin^2 t + \cos^2 t}{-\sin^3 t}
  = -\frac{1}{\sin^3 t}
  = -\csc^3 t
  \]
  
\end{myexample}

两种方法结果一致。本质上没有区别。

\subsection{隐函数求导}

之前我们求导的函数都是 $y = f(x)$ 的形式，$y$ 被“显式”地写成 $x$ 的表达式。但实际中常有这样的方程：
\[
F(x, y) = 0
\]

它并不直接写出 $y = f(x)$，而是将 $x$ 和 $y$ 的关系统一放在一个式子里。这种形式确定的函数叫隐函数：

\begin{definition}{隐函数}{}
  若函数关系$y = f(x)$或$x = g(y)$由方程
  \[
  F(x, y) = 0
  \]
  所确定，则称该函数为\textbf{隐函数}。
\end{definition}

隐函数求导的思路非常直接：把方程 $F(x,y)=0$ 两边同时对 $x$ 求导，其中 $y$ 视为 $x$ 的函数 $y(x)$，碰到 $y$ 的项就用链式法则。用莱布尼茨记号写出来尤为清晰。

为什么遇到$y$要用链式法则呢？因为，任何含有$y$的项，本质上都是$x$的复合函数。例如，$\sin y$其实是$\sin y(x)$，$\ln y$其实是$\ln y(x)$。按照复合函数求导法则，先对外层函数求导，里面不变，再乘以内部函数$y(x)$的导数，也就是$\frac{dy}{dx}$。

\begin{theorem}{隐函数求导法则}{}
  若函数关系$y = f(x)$由方程
  \[
  F(x, y) = 0
  \]
  所确定，则求 $y$ 对 $x$ 的导数时，方程两边同时对$x$求导，其中 $y$ 视为 $x$ 的函数 $y=y(x)$，最后解出$\frac{dy}{dx}$。
\end{theorem}

\begin{myexample}
  设 $x^2 + y^2 = 1$（$y \geq 0$），求 $\frac{dy}{dx}$。

  两边对 $x$ 求导：
  \[
  \frac{d}{dx}(x^2) + \frac{d}{dx}(y^2) = \frac{d}{dx}(1)
  \]

  左边第一项是 $2x$，右边是 $0$。左边第二项中 $y^2$ 是 $y$ 的函数，$y$ 又是 $x$ 的函数，由链式法则：
  \[
  \frac{d}{dx}(y^2) = 2y \cdot \frac{dy}{dx}
  \]
  
  因此：
  \[
  2x + 2y \frac{dy}{dx} = 0
  \]

  解得：
  \[
  \frac{dy}{dx} = -\frac{x}{y}
  \]

\end{myexample}

结果中同时包含 $x$ 和 $y$，这在隐函数求导中是很正常的，我们不一定要把 $y$ 用 $x$ 表达出来。

  我们可以验证，由 $x^2+y^2=1$ 解得 $y = \sqrt{1-x^2}$，直接求导得 $\frac{dy}{dx} = \frac{-x}{\sqrt{1-x^2}} = -\frac{x}{y}$，这与我们求出来的结果一致。

\begin{myexample}
  设方程 $e^y + xy = x$ 确定了隐函数 $y = y(x)$，求 $\frac{d^2 y}{dx^2}$。
  
  左边 $e^y$ 对 $x$ 求导：
  \[\frac{d}{dx}(e^y) = e^y \cdot \frac{dy}{dx}\]
  
  $xy$ 是乘法，求导要注意：
  \[\frac{d}{dx}(xy) = 1 \cdot y + x \cdot \frac{dy}{dx}\]
  
  右边 $x$ 求导得 $1$。于是：
  \[
  e^y \frac{dy}{dx} + y + x \frac{dy}{dx} = 1
  \]

  把含 $\frac{dy}{dx}$ 的项提到一边：
  \[
  (e^y + x) \frac{dy}{dx} = 1 - y
  \]

  解得：
  \[
  \frac{dy}{dx} = \frac{1 - y}{e^y + x}
  \]

  则二阶导数为：
  \[
  \frac{d^2 y}{dx^2} =\frac{d}{dx} (\frac{1 - y}{e^y + x})
  \]

  所以：
  \[
\begin{aligned}
\frac{d^2y}{dx^2} 
&= \frac{d}{dx}(\frac{1-y}{e^y+x}) \\[6pt]
&= \frac{ (1-y)'(e^y+x) - (1-y)(e^y+x)' }{ (e^y+x)^2 } \\[6pt]
&= \frac{ (-y')(e^y+x) - (1-y)(e^y\cdot y' + 1) }{ (e^y+x)^2 } \\[6pt]
&= \frac{ -y'(e^y+x) - (1-y)e^y y' - (1-y) }{ (e^y+x)^2 } \\[6pt]
&= \frac{ -y'[(e^y+x) + (1-y)e^y] - (1-y) }{ (e^y+x)^2 } \\[6pt]
&= \frac{ -y'(2e^y + x - ye^y) - (1-y) }{ (e^y+x)^2 }
\end{aligned}
\]

代入\(y' = \frac{dy}{dx} = \frac{1-y}{e^y+x}\)：
\[
\begin{aligned}
\frac{d^2y}{dx^2} 
&= \frac{ -\frac{1-y}{e^y+x}( 2e^y + x - ye^y ) - (1-y) }{ (e^y+x)^2 } \\[8pt]
&= \frac{ -(1-y)( 2e^y + x - ye^y ) - (1-y)(e^y+x) }{ (e^y+x)^3 } \\[6pt]
&= \frac{ -(1-y)[ 2e^y + x - ye^y + e^y + x ] }{ (e^y+x)^3 } \\[6pt]
&= -\frac{ (1-y)( 3e^y + 2x - ye^y ) }{ (e^y+x)^3 }
\end{aligned}
\]
\end{myexample}

隐函数求导的核心就一条：把 $y$ 当作 $x$ 的函数，每碰到 $y$ 的项就使用链式法则，乘以 $\frac{dy}{dx}$，然后从方程中把 $\frac{dy}{dx}$ 解出来。

实际上，方程\(F(x, y) = 0\)可看作关于$x,y$的二元函数，也就是说，在多元函数微分学里，我们有更强大的工具可以求隐函数的导数。


\subsection{一阶微分形式不变性}

学完各种各样的导数之后，我们可以推出微分的一个极其优雅的性质。设 $y = f(u)$，如果 $u$ 是自变量，微分自然是 $dy = f'(u)\,du$。但如果 $u$ 本身是另一个变量 $x$ 的函数 $u = g(x)$，那么复合函数 $y = f(g(x))$ 的微分会是什么形式？

用微分的定义来算。先写出 $y$ 对 $x$ 的微分，由链式法则，先对$f$求导，再乘以内部$g$的导数，也就是：
\[
dy = [f(g(x))]'\, dx = f'(g(x))\, g'(x)\, dx
\]

注意到 $g'(x)\, dx$ 正是 $u$ 作为 $x$ 的函数的微分，即 $du = g'(x)\, dx$。代入上式：
\[
dy = f'(u)\, du
\]

也就是说，无论 $u$ 是自变量还是中间变量，微分永远写成 $dy = f'(u)\, du$ 这个形式。

\begin{theorem}{一阶微分形式不变性}{}
设 $y = f(u)$ 可微。则不论 $u$ 是自变量还是中间变量，总有：
\[
dy = f'(u)\, du
\]
\end{theorem}

这个性质是微分的灵魂。在后面的换元积分法中你会反复看到这个性质的身影。它是连接微分学与积分学的第一座桥梁。

一阶微分形式不变性由莱布尼茨发现。不仅如此，莱布尼茨还发现了复合函数求导的链式法则，不过当时没有严谨证明，证明的工作最终由柯西等人完成。而隐函数求导法则的系统化，则在多元函数微分学建立之后，由柯西、雅可比等人发展为著名的隐函数定理，成为现代数学分析的基石。

关于导数，柯西在《分析教程》（1821）中将导数明确表述为：
\[
f'(x) = \lim_{\Delta x \to 0} \frac{f(x+\Delta x)-f(x)}{\Delta x}
\]

柯西还系统讨论了连续性与可导性的关系，证明了可导必连续。更进一步，德国数学家魏尔斯特拉斯构造了一个函数：
\[
f(x) = \sum_{n=0}^{\infty} a^n \cos(b^n \pi x)
\]

其中 \(0 < a < 1\)，\(b\) 是一个正奇数，且 \(ab > 1 + \frac{3\pi}{2}\)。

这个函数在每一处都连续，但在每一处都不可导。画出来大概是一团永远在抖动的锯齿，无论放多大，都看不到任何一段光滑的曲线。

魏尔斯特拉斯早年在波恩大学读法律与财政，成绩平平，随后转入数学，却在毕业考试中因病失败。此后长达十余年，他在偏远的普鲁士小城当中学教师，教授书法、地理和体育。所以，当你说别人的数学是体育老师教的时候，别忘了魏尔斯特拉斯曾经也是体育老师。
% !TEX root = ../main.tex
\chapter{一元函数积分学}

学完微分之后，一个自然的问题是：如果知道了一个函数的导数，能不能把原函数找回来？这就是求导的逆运算。

本章先从不定积分入手，逐步引入换元法和分部积分法这两大核心技巧，最后建立定积分与不定积分之间的联系。

\section{不定积分}
微分学告诉我们，对一个函数求导可以得到它的导数。反过来，已知一个函数的导数，求原来的函数，这个逆过程叫做\textbf{求不定积分}。

\subsection{求不定积分}

\begin{definition}{原函数}{}
设函数 $F(x)$ 与 $f(x)$ 都定义在区间 $I$ 上。若对于 $I$ 上每一点 $x$，都有
\[
F'(x) = f(x)
\]
则称 $F(x)$ 是 $f(x)$ 在区间 $I$ 上的一个\textbf{原函数}。
\end{definition}

例如，$F(x)=x^2$ 的导数是 $f(x)=2x$，所以 $x^2$ 是 $2x$ 的一个原函数。但注意，$(x^2+1)'=2x$，$(x^2-5)'=2x$，甚至 $(x^2+C)'=2x$（$C$ 为任意常数）。这说明原函数不是唯一的。因此我们把所有原函数的集合称为不定积分。

\begin{definition}{不定积分}{}
函数 $f(x)$ 在区间 $I$ 上的全体原函数称为 $f(x)$ 的\textbf{不定积分}，记作
\[
\int f(x)\, dx = F(x) + C
\]
其中 $F(x)$ 是 $f(x)$ 的任意一个原函数，$C$ 称为\textbf{积分常数}，$f(x)$ 称为\textbf{被积函数}，$x$ 称为\textbf{积分变量}，符号 $\int$ 称为\textbf{积分号}。
\end{definition}

简单来说，求不定积分，就是找一个函数，使得它的导数等于被积函数，然后在后面加一个常数 $C$。这个 $C$ 绝不能丢，因为导数会抹掉常数信息，积分必须把它补回来。

你可能会问，为什么积分号后面总要拖一个微分 $dx$？那是因为 $dx$ 有几个作用：

首先，$dx$指明积分变量是$x$，$dy$指明积分变量是$y$。例如 $\int x^2 y\, dx$ 和 $\int x^2 y\, dy$ 是完全不同的积分，前者把 $y$ 当常数，对 $x$ 积分，得到 $\frac{x^3}{3}y+C$。后者把 $x^2$ 当常数，对 $y$ 积分，得到 $\frac{x^2 y^2}{2}+C$。意思就是，看到后面的$dx$，就知道是对$x$积分。

让我们先回忆一下微分的定义：$dy = f'(x)\,dx$。把两边同时积分，左边是 $\int dy = y + C$，右边是 $\int f'(x)\,dx=f(x)+C$。也就是说，积分号 $\int$ 和微分号 $d$ 是互逆的运算，$dx$ 正是被积回来的那个微小变化量。

那么$\int f(x)\,dx$ 的意思就是“把 $f(x)$ 和无穷小量 $dx$ 乘起来，得到原函数”。这也是约定俗成的写法。而积分号$\int$正是莱布尼茨发明的。

由于积分是微分的逆运算，我们可以直接把求导公式反过来，得到最基本的积分公式：

\begin{theorem}{基本积分公式}{}
\[
\begin{aligned}
&\text{(1)}\ \int k\, dx = kx + C \quad (k\text{ 为常数})
&&\text{(2)}\ \int x^{\mu}\, dx = \frac{x^{\mu+1}}{\mu+1} + C \quad (\mu \neq -1) \\[4pt]
&\text{(3)}\ \int \frac{1}{x}\, dx = \ln|x| + C
&&\text{(4)}\ \int e^x\, dx = e^x + C \\[4pt]
&\text{(5)}\ \int a^x\, dx = \frac{a^x}{\ln a} + C \quad (a>0,\ a\neq 1)
&&\text{(6)}\ \int \sin x\, dx = -\cos x + C \\[4pt]
&\text{(7)}\ \int \cos x\, dx = \sin x + C
&&\text{(8)}\ \int \sec^2 x\, dx = \tan x + C \\[4pt]
&\text{(9)}\ \int \csc^2 x\, dx = -\cot x + C
&&\text{(10)}\ \int \sec x \tan x\, dx = \sec x + C \\[4pt]
&\text{(11)}\ \int \csc x \cot x\, dx = -\csc x + C
&&\text{(12)}\ \int \frac{1}{1+x^2}\, dx = \arctan x + C \\[4pt]
&\text{(13)}\ \int \frac{1}{\sqrt{1-x^2}}\, dx = \arcsin x + C
\end{aligned}
\]
\end{theorem}

这些公式不需要死记，每一条都可以通过右边求导来验证。

积分还有两条最基本的运算性质，它们直接继承自求导的线性性质：

\begin{theorem}{不定积分的线性性质}{}
设 $f(x)$ 和 $g(x)$ 的原函数都存在，$k$ 为常数，则
\[
\begin{aligned}
&\text{(1)}\quad \int k f(x)\, dx = k \int f(x)\, dx \\[4pt]
&\text{(2)}\quad \int \bigl[f(x) \pm g(x)\bigr]\, dx = \int f(x)\, dx \pm \int g(x)\, dx
\end{aligned}
\]
\end{theorem}

通俗地说，常数因子可以提到积分号外面去，和的积分等于积分的和，差的积分等于积分的差。这两条性质配合基本积分公式，就能处理大多数简单函数的积分了。

\begin{myexample}
  举例说明：求 $ \int (3x^2 + 2x - 1)\, dx$。

  怎么做呢？我们先利用线性性质把被积函数拆成积分的和，也就是：
  $$\int 3x^2 dx + \int 2x dx - \int 1dx$$

  然后把常数提到外面去，也就是：
  $$3\int x^2\, dx + 2\int x\, dx - \int dx$$

  接着，利用基本积分公式，算出每一项的原函数。因为大家都带一个常数，所以可以统一写成一个$C$，也就是：
  $$3 \cdot \frac{x^3}{3} + 2 \cdot \frac{x^2}{2} - x + C$$

  最后化简得到最终结果。

  把过程写出来，就是：
  \[
  \begin{aligned}
  \int (3x^2 + 2x - 1)\, dx
  &= 3\int x^2\, dx + 2\int x\, dx - \int 1\, dx \\
  &= 3 \cdot \frac{x^3}{3} + 2 \cdot \frac{x^2}{2} - x + C \\
  &= x^3 + x^2 - x + C
  \end{aligned}
  \]

\end{myexample}

  我们可以验证一下：

  对结果求导，$(x^3+x^2-x+C)' = 3x^2+2x-1$，与被积函数一致。那就是算对了。

\begin{myexample}
  举例说明：求 $ \int \frac{x^3 - 3x + 2\sqrt{x}}{x}\, dx$。

  先化简被积函数：$\dfrac{x^3 - 3x + 2\sqrt{x}}{x} = x^2 - 3 + 2x^{-\frac{1}{2}}$。

  \[
  \begin{aligned}
  \int (x^2 - 3 + 2x^{-\frac{1}{2}})\, dx
  &= \int x^2\, dx - 3\int 1\, dx + 2\int x^{-\frac{1}{2}}\, dx \\
  &= \frac{x^3}{3} - 3x + 2 \cdot 2\sqrt{x} + C \\
  &= \frac{x^3}{3} - 3x + 4\sqrt{x} + C
  \end{aligned}
  \]
\end{myexample}

这里有一个重要的经验：遇到分式或根式，先化简再积分。很多看起来复杂的积分，化简之后就变成了基本公式的组合。

\begin{myexample}
  举例说明：求 $ \int (\sin x + 2e^x)\, dx$。

  \[
  \int (\sin x + 2e^x)\, dx = \int \sin x\, dx + 2\int e^x\, dx = -\cos x + 2e^x + C
  \]
\end{myexample}

由上述例子可以看出，我们要牢记积分基本公式，才能把积分算出来。

刚刚说到，不定积分和求导互为逆运算，这个关系可以用两句话记住：

先积后导，回到原函数：$ \frac{d}{dx}(\int f(x)\, dx) = f(x)$。

先导后积，回到原函数加常数：$\displaystyle \int F'(x)\, dx = F(x) + C$。

以及，一定一定别忘记加$C$。


\subsection{第一类换元法}

基本积分公式只能处理最简单的被积函数，一旦被积函数是复合函数的形式，比如 $\int \cos 2x \,dx$，基本公式就用不了了。第一类换元法（也叫凑微分法）就是专门解决这类问题的。

第一类换元法直接来源于链式法则的逆用。回忆链式法则：
\[
\frac{d}{dx} F(u(x)) = F'(u(x)) \cdot u'(x) = f(u(x)) \cdot u'(x)
\]

这个式子的意思是对$F(u(x))$求导，等于$F'(u) \cdot u'(x)$，其中 $ F' = f$，$u$是$x$的函数。然后我们两边积分：
\[
\int f(u(x)) \cdot u'(x)\, dx = F(u(x)) + C
\]

因为$u=u(x)$，那么就有$du = u'(x)\, dx$，上式简洁地写成：

\[
\int f(u(x)) \cdot u'(x)\, dx = \int f(u)\, du = F(u) + C
\]

有没有发现有点似曾相识？是的，这正是我们上一章讲到的一阶微分形式不变性。

由此，我们有第一类换元法：
\begin{theorem}{第一类换元法}{}
  设 $F'(u) = f(u)$，且 $u = u(x)$ 可导，则
\[
\int f(u(x)) \cdot u'(x)\, dx = \int f(u)\, du = F(u) + C
\]
\end{theorem}

看上去很复杂，其实就是找一个函数$u(x)$，再把$u'(x)dx$凑出来，然后将被积函数换成关于$u$的函数，$dx$换成$du$，这样整个积分就换成关于$u$的积分，这个积分是要可以用基本公式积出来的，积出来之后再把$x$代回去。

我们来看最开始的例子：

\begin{myexample}
  求 $\int \cos 2x\, dx$。

  令 $u = 2x$，则 $du = 2\, dx$，即 $dx = \frac{1}{2}\, du$。则有：
  \[
  \int \cos 2x\, dx = \int \cos u \cdot \frac{1}{2}\, du
  = \frac{1}{2} \int \cos u\, du
  = \frac{1}{2} \sin u + C
  = \frac{1}{2} \sin 2x + C
  \]
\end{myexample}

仔细看这个例子可以发现，我们令$u = 2x$，由微分公式，$du = 2\, dx$，这时候就能用$du$表示$dx$了。于是整个积分就可以用$u$来表示。这也说明了，凑微分的难点之一在于怎么用我们设的$du$表示$dx$。

再来看一个例子：

\begin{myexample}
  求 $\int 2x \cos x^2\, dx$。

  令 $u = x^2$，则 $du = 2x\, dx$。则有：
  \[
  \int \cos x^2 \cdot 2x\, dx = \int \cos u\, du = \sin u + C = \sin x^2  + C
  \]

\end{myexample}

  这个例子展示了凑微分法的精髓，我们有时候不需要把 $dx$ 单独解出来再代入，而是把 $2x\,dx$ 整体识别为 $du$。我们令$u = x^2$，由微分公式，$du = 2x\, dx$，这时候我们发现被积函数里刚好就有$du = 2x\, dx$，而$\cos x^2$刚好可以表示成$\cos u$，完美。

\begin{myexample}
  求$\int \sin^3 x \cos x\, dx$。

  令$u = \sin x$，则 $du = \cos x\, dx$，则有：

\[
\int \sin^3 x \cdot \cos x\, dx = \int u^3 \, du = \frac{u^4}{4} + C = \frac{\sin^4 x}{4} + C
\]
\end{myexample}

这就是“凑微分”三个字的真正含义，如何在被积表达式里凑出一个完整的 $du$。我们令$u = \sin x$，由微分公式， $du = \cos x\, dx$，这时候我们发现被积函数里刚好就有$du = \cos x\, dx$，和上题类似，直接一套带走。

\begin{myexample}
  求 $\int \frac{1}{2x+1}\, dx$。

  令 $u = 2x+1$，则 $du = 2\, dx$，$dx = \frac{1}{2}\, du$。
  \[
  \int \frac{1}{2x+1}\, dx = \int \frac{1}{u} \cdot \frac{1}{2}\, du
  = \frac{1}{2} \ln|u| + C
  = \frac{1}{2} \ln|2x+1| + C
  \]
\end{myexample}

\begin{myexample}
  求 $\int \tan x\, dx$。

  先把 $\tan x$ 写成 $\frac{\sin x}{\cos x}$。
  
  令 $u = \cos x$，则 $du = -\sin x\, dx$，即 $\sin x\, dx = -du$。
  \[
  \int \tan x\, dx = \int \frac{\sin x}{\cos x}\, dx
  = \int \frac{1}{u} \cdot (-du)
  = -\ln|u| + C
  = -\ln|\cos x| + C
  = \ln|\sec x| + C
  \]

  这个积分非常有用，建议记住它：$\int \tan x\, dx = -\ln|\cos x| + C = \ln|\sec x| + C$。
\end{myexample}

总结一下第一类换元法的核心：

首先通过敏锐的数学直觉，设一个$u=u(x)$，使得被积函数能用$u$来表示，$dx$能用$du$来表示。

然后把原积分换成能用$u$和$du$来表示的基本积分，最后一通操作，直接带走。

以及，一定一定别忘记加$C$。

你可能会说，主播主播这种方法还是太吃技巧了，有没有不吃技巧的方法？有的有的，大部分第一类换元法都围绕这几个形式变形，把他们记住，遇到类似的直接带走就可以了。如果遇到没有的呢？那就借助万能的AI。

\begin{theorem}{线性函数凑微分公式}{}
  含有$ax+b$的积分，可化为以下形式：
\[
\begin{aligned}
&\text{(1)}\ \int (ax+b)^n \,dx = \frac{1}{a}\cdot\frac{(ax+b)^{n+1}}{n+1}+C, \quad n\neq -1 \\[4pt]
&\text{(2)}\ \int \frac{1}{ax+b} \,dx = \frac{1}{a}\ln|ax+b|+C \\[4pt]
&\text{(3)}\ \int e^{ax+b} \,dx = \frac{1}{a}e^{ax+b}+C \\[4pt]
&\text{(4)}\ \int a^{kx+b} \,dx = \frac{1}{k\ln a}a^{kx+b}+C, \quad a>0,a\neq1 \\[4pt]
&\text{(5)}\ \int \sin(ax+b) \,dx = -\frac{1}{a}\cos(ax+b)+C \\[4pt]
&\text{(6)}\ \int \cos(ax+b) \,dx = \frac{1}{a}\sin(ax+b)+C \\[4pt]
&\text{(7)}\ \int \sec^2(ax+b) \,dx = \frac{1}{a}\tan(ax+b)+C \\[4pt]
&\text{(8)}\ \int \csc^2(ax+b) \,dx = -\frac{1}{a}\cot(ax+b)+C \\[4pt]
&\text{(9)}\ \int \sec(ax+b)\tan(ax+b) \,dx = \frac{1}{a}\sec(ax+b)+C \\[4pt]
&\text{(10)}\ \int \csc(ax+b)\cot(ax+b) \,dx = -\frac{1}{a}\csc(ax+b)+C
\end{aligned}
\]
\end{theorem}

我们来看个例子直观感受公式的使用：

\begin{myexample}
  求 $\int ( \frac{4}{2x-1} - \sin(3x-4) + 3e^{5x+2} ) dx$。

  先逐项计算：

令 $u = 2x-1$，则 $du = 2\,dx$，$dx = \frac{1}{2}\,du$。
    \[
    \int \frac{4}{2x-1}\,dx = 4\int \frac{1}{u}\cdot\frac{1}{2}\,du
    = 2\ln|u| + C_1 = 2\ln|2x-1| + C_1
    \]

令 $v = 3x-4$，则 $dv = 3\,dx$，$dx = \frac{1}{3}\,dv$。
    \[
    \int -\sin(3x-4)\,dx = -\int \sin v \cdot \frac{1}{3}\,dv
    = -\frac{1}{3}(-\cos v) + C_2 = \frac{1}{3}\cos(3x-4) + C_2
    \]

令 $w = 5x+2$，则 $dw = 5\,dx$，$dx = \frac{1}{5}\,dw$。
    \[
    \int 3e^{5x+2}\,dx = 3\int e^{w}\cdot\frac{1}{5}\,dw
    = \frac{3}{5}e^{w} + C_3 = \frac{3}{5}e^{5x+2} + C_3
    \]

  相加并合并常数 $C = C_1 + C_2 + C_3$，得：
  \[
  \int ( \frac{4}{2x-1} - \sin(3x-4) + 3e^{5x+2} ) dx
  = 2\ln|2x-1| + \frac{3}{5}e^{5x+2} + \frac{1}{3}\cos(3x-4) + C
  \]

  但是，我们注意到，该积分每一项都是关于线性函数 $ax+b$ 的复合，可直接使用线性函数凑微分公式：

第一项，由公式：$\int \frac{1}{ax+b}\,dx = \frac{1}{a}\ln|ax+b|+C$，其中 $a=2$，$b=-1$，系数 $4$ 提出，得：
    \[
    \int \frac{4}{2x-1}\,dx = 4\cdot\frac{1}{2}\ln|2x-1| + C_1 = 2\ln|2x-1| + C_1
    \]

第二项，由公式：$\int \sin(ax+b)\,dx = -\frac{1}{a}\cos(ax+b)+C$，其中 $a=3$，$b=-4$，注意被积函数有负号，得：
    \[
    \int -\sin(3x-4)\,dx = -(-\frac{1}{3}\cos(3x-4)) + C_2 = \frac{1}{3}\cos(3x-4) + C_2
    \]

第三项，由公式：$\int e^{ax+b}\,dx = \frac{1}{a}e^{ax+b}+C$，其中 $a=5$，$b=2$，得：
    \[
    \int 3e^{5x+2}\,dx = 3\cdot\frac{1}{5}e^{5x+2} + C_3 = \frac{3}{5}e^{5x+2} + C_3
    \]

  相加并合并积分常数 $C = C_1 + C_2 + C_3$，得最终结果：
  \[
  \int ( \frac{4}{2x-1} - \sin(3x-4) + 3e^{5x+2} ) dx
  = 2\ln|2x-1| + \frac{3}{5}e^{5x+2} + \frac{1}{3}\cos(3x-4) + C
  \]
\end{myexample}


\begin{theorem}{幂函数凑微分公式}{}
  含有幂函数的积分，可化为以下形式：
\[
\begin{aligned}
&\text{(1)}\ \int x^n (x^{n+1})^k \,dx = \frac{1}{n+1}\cdot\frac{x^{(n+1)(k+1)}}{k+1}+C, \quad k\neq -1 \\[4pt]
&\text{(2)}\ \int x e^{x^2} \,dx = \frac{1}{2}e^{x^2}+C \\[4pt]
&\text{(3)}\ \int x^n e^{x^{n+1}} \,dx = \frac{1}{n+1}e^{x^{n+1}}+C \\[4pt]
&\text{(4)}\ \int x\sqrt{x^2\pm a^2} \,dx = \frac{1}{3}(x^2\pm a^2)^{\frac{3}{2}}+C \\[4pt]
&\text{(5)}\ \int \frac{x}{\sqrt{x^2\pm a^2}} \,dx = \sqrt{x^2\pm a^2}+C \\[4pt]
&\text{(6)}\ \int x\sin x^2 \,dx = -\frac{1}{2}\cos x^2+C \\[4pt]
&\text{(7)}\ \int x\cos x^2 \,dx = \frac{1}{2}\sin x^2+C \\[4pt]
&\text{(8)}\ \int x\sec^2 x^2 \,dx = \frac{1}{2}\tan x^2+C
\end{aligned}
\]
\end{theorem}

我们再看个例子直观感受公式的使用：

\begin{myexample}
  求 $\int ( 2x e^{x^2} + 3x^2 e^{x^3} - 4x \sin x^2 ) dx$。

    逐项换元计算：

令 $u = x^2$，则 $du = 2x\,dx$，$x\,dx = \frac{1}{2}\,du$。
    \[
    \int 2x e^{x^2}\,dx = \int e^{u}\,du = e^{u} + C_1 = e^{x^2} + C_1
    \]

令 $v = x^3$，则 $dv = 3x^2\,dx$，$x^2\,dx = \frac{1}{3}\,dv$。
    \[
    \int 3x^2 e^{x^3}\,dx = \int e^{v}\,dv = e^{v} + C_2 = e^{x^3} + C_2
    \]

令 $w = x^2$，则 $dw = 2x\,dx$，$x\,dx = \frac{1}{2}\,dw$。
    \[
    \int -4x \sin x^2\,dx = -4\cdot\frac{1}{2}\int \sin w\,dw = -2(-\cos w) + C_3 = 2\cos x^2 + C_3
    \]

  相加并合并常数 $C = C_1+C_2+C_3$，得最终结果：
  \[
  \int ( 2x e^{x^2} + 3x^2 e^{x^3} - 4x \sin x^2 ) dx
  = e^{x^2} + e^{x^3} + 2\cos x^2 + C
  \]

  我们注意到，该积分每一项都是“幂函数乘复合函数”的形式，可直接使用幂函数凑微分公式：

第一项：由公式：$\int x e^{x^2}\,dx = \frac{1}{2}e^{x^2}+C$，系数 $2$ 提出，得：
    \[
    \int 2x e^{x^2}\,dx = 2\cdot\frac{1}{2}e^{x^2}+C_1 = e^{x^2}+C_1
    \]

第二项：由公式：$\int x^n e^{x^{n+1}}\,dx = \frac{1}{n+1}e^{x^{n+1}}+C$，其中$n=2$，系数 $3$ 提出，得：
    \[
    \int 3x^2 e^{x^3}\,dx = 3\cdot\frac{1}{3}e^{x^3}+C_2 = e^{x^3}+C_2
    \]

第三项：由公式：$\int x \sin x^2\,dx = -\frac{1}{2}\cos x^2+C$，系数 $-4$ 提出，得：
    \[
    \int -4x \sin x^2\,dx = -4\cdot(-\frac{1}{2}\cos x^2)+C_3 = 2\cos x^2+C_3
    \]

  相加并合并常数 $C = C_1+C_2+C_3$，得最终结果：
  \[
  \int ( 2x e^{x^2} + 3x^2 e^{x^3} - 4x \sin x^2 ) dx
  = e^{x^2} + e^{x^3} + 2\cos x^2 + C
  \]
\end{myexample}

\begin{theorem}{对数函数凑微分公式}{}
  含有$\ln x$的积分，可化为以下形式：
  \[
\begin{aligned}
&\text{(1)}\ \int \frac{(\ln x)^k}{x} \,dx = \frac{(\ln x)^{k+1}}{k+1}+C, \quad k\neq -1 \\[4pt]
&\text{(2)}\ \int \frac{1}{x\ln x} \,dx = \ln|\ln x|+C \\[4pt]
&\text{(3)}\ \int \frac{1}{x(\ln x)^k} \,dx = \frac{(\ln x)^{1-k}}{1-k}+C, \quad k\neq 1 \\[4pt]
&\text{(4)}\ \int \frac{1}{x\ln x\ln(\ln x)} \,dx = \ln|\ln(\ln x)|+C \\[4pt]
&\text{(5)}\ \int \frac{f'(\ln x)}{x} \,dx = f(\ln x)+C
\end{aligned}
\]
\end{theorem}

\begin{myexample}
  求 $\int ( \frac{(\ln x)^2}{x} + \frac{1}{x\ln x} + \frac{\sin(\ln x)}{x} ) dx$。

  注意到被积函数中反复出现 $\ln x$ 及其导数 $\frac{1}{x}$，可令 $u = \ln x$ 统一换元。

令 $u = \ln x$，则 $du = \frac{1}{x}\,dx$。
    \[
    \int \frac{(\ln x)^2}{x}\,dx = \int u^2\,du = \frac{u^3}{3} + C_1 = \frac{1}{3}(\ln x)^3 + C_1
    \]

同样代换，第二项化为：
    \[
    \int \frac{1}{x\ln x}\,dx = \int \frac{1}{u}\,du = \ln|u| + C_2 = \ln|\ln x| + C_2
    \]

第三项：
    \[
    \int \frac{\sin(\ln x)}{x}\,dx = \int \sin u\,du = -\cos u + C_3 = -\cos(\ln x) + C_3
    \]

  相加并合并常数 $C = C_1+C_2+C_3$，得：
  \[
  \int ( \frac{(\ln x)^2}{x} + \frac{1}{x\ln x} + \frac{\sin(\ln x)}{x} ) dx
  = \frac{1}{3}(\ln x)^3 + \ln|\ln x| - \cos(\ln x) + C
  \]

  也可直接使用对数函数凑微分公式。

第一项：由公式：$\int \frac{(\ln x)^k}{x} \,dx = \frac{(\ln x)^{k+1}}{k+1}+C$，其中 $k=2$，得：
    \[
    \int \frac{(\ln x)^2}{x}\,dx = \frac{(\ln x)^{3}}{3} + C_1
    \]

第二项：由公式：$\int \frac{1}{x\ln x} \,dx = \ln|\ln x|+C$，得：
    \[
    \int \frac{1}{x\ln x}\,dx = \ln|\ln x| + C_2
    \]

第三项：由公式：$\int \frac{f'(\ln x)}{x} \,dx = f(\ln x)+C$，令 $f'(u)=\sin u$，则 $f(u)=-\cos u$，得：
    \[
    \int \frac{\sin(\ln x)}{x}\,dx = -\cos(\ln x) + C_3
    \]

  相加并合并常数 $C = C_1+C_2+C_3$，得：
  \[
  \int ( \frac{(\ln x)^2}{x} + \frac{1}{x\ln x} + \frac{\sin(\ln x)}{x} ) dx
  = \frac{1}{3}(\ln x)^3 + \ln|\ln x| - \cos(\ln x) + C
  \]
\end{myexample}


\begin{theorem}{指数函数凑微分公式}{}
  含有指数函数的积分，可化为以下形式：
  \[
\begin{aligned}
&\text{(1)}\ \int e^{x} f(e^{x}) \,dx = \int f(u)\,du,\; u=e^x \\[4pt]
&\text{(2)}\ \int \frac{e^x}{1+e^{2x}} \,dx = \arctan(e^x)+C \\[4pt]
&\text{(3)}\ \int \frac{e^x}{\sqrt{1-e^{2x}}} \,dx = \arcsin(e^x)+C \\[4pt]
&\text{(4)}\ \int e^{ax}\sin(e^{ax}) \,dx = -\frac{1}{a}\cos(e^{ax})+C \\[4pt]
&\text{(5)}\ \int e^{ax}\cos(e^{ax}) \,dx = \frac{1}{a}\sin(e^{ax})+C \\[4pt]
&\text{(6)}\ \int a^x f(a^x)\,dx = \frac{1}{\ln a}\int f(u)\,du,\; u=a^x \\[4pt]
&\text{(7)}\ \int f'(x)\,e^{f(x)} \,dx = e^{f(x)}+C
\end{aligned}
\]
\end{theorem}

\begin{myexample}
  求 $\int ( \frac{e^{x}}{1+e^{2x}} + 2e^{2x}\sin(e^{2x}) + 2x e^{x^2} ) dx$。

  逐项换元计算：

令 $u = e^{x}$，则 $du = e^{x}\,dx$。
    \[
    \int \frac{e^{x}}{1+e^{2x}}\,dx = \int \frac{1}{1+u^{2}}\,du = \arctan u + C_1 = \arctan(e^{x}) + C_1
    \]

令 $v = e^{2x}$，则 $dv = 2e^{2x}\,dx$，即 $e^{2x}\,dx = \frac{1}{2}\,dv$。
    \[
    \int 2e^{2x}\sin(e^{2x})\,dx = \int \sin v\,dv = -\cos v + C_2 = -\cos(e^{2x}) + C_2
    \]

令 $w = x^{2}$，则 $dw = 2x\,dx$。
    \[
    \int 2x e^{x^{2}}\,dx = \int e^{w}\,dw = e^{w} + C_3 = e^{x^{2}} + C_3
    \]

  相加并合并常数 $C = C_1+C_2+C_3$，得最终结果：
  \[
  \int ( \frac{e^{x}}{1+e^{2x}} + 2e^{2x}\sin(e^{2x}) + 2x e^{x^2} ) dx
  = \arctan(e^x) - \cos(e^{2x}) + e^{x^2} + C
  \]

  也可以直接使用指数函数凑微分公式：

第一项：由公式：$\int \frac{e^x}{1+e^{2x}}\,dx = \arctan(e^x)+C$，得：
    \[
    \int \frac{e^{x}}{1+e^{2x}}\,dx = \arctan(e^x) + C_1
    \]

第二项：由公式：$\int e^{ax}\sin(e^{ax})\,dx = -\frac{1}{a}\cos(e^{ax})+C$，其中 $a=2$，系数 $2$ 提出，得：
    \[
    \int 2e^{2x}\sin(e^{2x})\,dx = 2\cdot(-\frac{1}{2}\cos(e^{2x}))+C_2 = -\cos(e^{2x}) + C_2
    \]

第三项：由公式：$\int f'(x)e^{f(x)}\,dx = e^{f(x)}+C$，令 $f(x)=x^2$，则 $f'(x)=2x$，得：
    \[
    \int 2x e^{x^2}\,dx = e^{x^2} + C_3
    \]

  相加并合并常数 $C = C_1+C_2+C_3$，得最终结果：
  \[
  \int ( \frac{e^{x}}{1+e^{2x}} + 2e^{2x}\sin(e^{2x}) + 2x e^{x^2} ) dx
  = \arctan(e^x) - \cos(e^{2x}) + e^{x^2} + C
  \]
\end{myexample}

\begin{theorem}{三角函数凑微分公式}{}
  含有三角函数的积分，可化为以下形式：
  \[
\begin{aligned}
&\text{(1)}\ \int \tan x \,dx = -\ln|\cos x|+C \\[4pt]
&\text{(2)}\ \int \cot x \,dx = \ln|\sin x|+C \\[4pt]
&\text{(3)}\ \int \sec x \,dx = \ln|\sec x + \tan x|+C \\[4pt]
&\text{(4)}\ \int \csc x \,dx = \ln|\csc x - \cot x|+C \\[4pt]
&\text{(5)}\ \int \tan^n x \sec^2 x \,dx = \frac{\tan^{n+1}x}{n+1}+C, \quad n\neq -1 \\[4pt]
&\text{(6)}\ \int \cot^n x \csc^2 x \,dx = -\frac{\cot^{n+1}x}{n+1}+C, \quad n\neq -1 \\[4pt]
&\text{(7)}\ \int \sin ax \sin bx\,dx = \frac{\sin(a-b)x}{2(a-b)} - \frac{\sin(a+b)x}{2(a+b)} + C,\quad a^2\neq b^2 \\[4pt]
&\text{(8)}\ \int \cos ax \cos bx\,dx = \frac{\sin(a-b)x}{2(a-b)} + \frac{\sin(a+b)x}{2(a+b)} + C,\quad a^2\neq b^2 \\[4pt]
&\text{(9)}\ \int \sin ax \cos bx\,dx = -\frac{\cos(a-b)x}{2(a-b)} - \frac{\cos(a+b)x}{2(a+b)} + C,\quad a^2\neq b^2 \\[4pt]
&\text{(10)}\ \int \sin^{2k+1}x \cos^n x \,dx = -\int (1-\cos^2 x)^k \cos^n x \,d(\cos x) \\[4pt]
&\text{(11)}\ \int \sin^m x \cos^{2k+1}x \,dx = \int \sin^m x (1-\sin^2 x)^k \,d(\sin x) \\[4pt]
&\text{(12)}\ \int \sin^n x\,dx = -\frac{1}{n}\sin^{n-1}x\cos x + \frac{n-1}{n}\int \sin^{n-2}x\,dx \quad (n\ge 2) \\[4pt]
&\text{(13)}\ \int \cos^n x\,dx = \frac{1}{n}\cos^{n-1}x\sin x + \frac{n-1}{n}\int \cos^{n-2}x\,dx \quad (n\ge 2) \\[4pt]
&\text{(14)}\ \int \tan^n x\,dx = \frac{\tan^{n-1}x}{n-1} - \int \tan^{n-2}x\,dx \quad (n\ge 2) \\[4pt]
&\text{(15)}\ \int \cot^n x\,dx = -\frac{\cot^{n-1}x}{n-1} - \int \cot^{n-2}x\,dx \quad (n\ge 2) \\[4pt]
&\text{(16)}\ \int \sec^n x\,dx = \frac{\sec^{n-2}x\tan x}{n-1} + \frac{n-2}{n-1}\int \sec^{n-2}x\,dx \quad (n\ge 2) \\[4pt]
&\text{(17)}\ \int \csc^n x\,dx = -\frac{\csc^{n-2}x\cot x}{n-1} + \frac{n-2}{n-1}\int \csc^{n-2}x\,dx \quad (n\ge 2)
\end{aligned}
\]
\end{theorem}
三角函数没什么好说的，这些都是必背公式。记住就行。


\begin{myexample}
  求 $\int ( 3\sin^{3}(2x-1) + 4\tan^{2}(3x)\sec^{2}(3x) - 2\sec(4x)\tan(4x) ) dx$。

  注意到，每一项都是“线性函数嵌套三角函数”的形式，所以可先对线性函数凑微分，再使用三角函数凑微分公式：

第一项：先线性凑微分，$\int \sin^{3}(2x-1)\,dx = \frac{1}{2}\int \sin^{3}(2x-1)\,d(2x-1)$。
          再由公式：$\int \sin^{n} u\,du = -\frac{1}{n}\sin^{n-1}u\cos u + \frac{n-1}{n}\int \sin^{n-2}u\,du$，其中 $u=2x-1$，$n=3$，系数 $3$ 提出，得：
          \[
          \begin{aligned}
            \int 3\sin^{3}(2x-1)\,dx
            &= \frac{3}{2}\int \sin^{3}u\,du \\
            &= \frac{3}{2}( -\frac{1}{3}\sin^{2}u\cos u + \frac{2}{3}\int \sin u\,du ) \\
            &= \frac{3}{2}( -\frac{1}{3}\sin^{2}u\cos u - \frac{2}{3}\cos u ) + C_{1} \\
            &= \frac{1}{2}\cos^{3}u - \frac{3}{2}\cos u + C_{1} \\
            &= \frac{1}{2}\cos^{3}(2x-1) - \frac{3}{2}\cos(2x-1) + C_{1}
          \end{aligned}
          \]

第二项：线性凑微分 $\int \tan^{2}(3x)\sec^{2}(3x)\,dx = \frac{1}{3}\int \tan^{2}(3x)\sec^{2}(3x)\,d(3x)$。
          由公式：$\int \tan^{n} u \sec^{2} u\,du = \frac{\tan^{n+1}u}{n+1} + C$，其中 $u=3x$，$n=2$，系数 $4$ 提出，得：
          \[
          \int 4\tan^{2}(3x)\sec^{2}(3x)\,dx = 4\cdot\frac{1}{3}\cdot\frac{1}{3}\tan^{3}(3x) + C_{2}
          = \frac{4}{9}\tan^{3}(3x) + C_{2}
          \]

第三项：线性凑微分 $\int \sec(4x)\tan(4x)\,dx = \frac{1}{4}\int \sec(4x)\tan(4x)\,d(4x)$。
          由公式：$\int \sec u \tan u\,du = \sec u + C$，其中 $u=4x$，系数 $-2$ 提出，得：
          \[
          \int -2\sec(4x)\tan(4x)\,dx = -2\cdot\frac{1}{4}\sec(4x) + C_{3} = -\frac{1}{2}\sec(4x) + C_{3}
          \]

合并各项并令 $C = C_{1}+C_{2}+C_{3}$，得最终结果：
\[
\begin{aligned}
\int \bigl( 3\sin^{3}(2x-1) + 4\tan^{2}(3x)\sec^{2}(3x) - 2\sec(4x)\tan(4x) \bigr) dx \\
= \frac{1}{2}\cos^{3}(2x-1) - \frac{3}{2}\cos(2x-1) + \frac{4}{9}\tan^{3}(3x) - \frac{1}{2}\sec(4x) + C
\end{aligned}
\]
\end{myexample}

这说明，实际操作中，往往会有多种函数的复合，这就需要我们火眼金睛，精准地设$u$换元。有时找不出正确的$u$，这是正常现象。只要再多观察一下被积函数，多尝试一下，最后总能发现其中奥秘。如果观察很久之后还是发现不了，那就放弃它，交给AI。

以及，一定一定别忘记加$C$。

最后还有反三角函数。反三角函数出现的机会不多，我们简单带过：


\begin{theorem}{反三角函数凑微分公式}{}
  含有反三角函数的积分，可化为以下形式：
  \[
\begin{aligned}
&\text{(1)}\ \int \frac{1}{\sqrt{a^2-x^2}} \,dx = \arcsin\frac{x}{a}+C, \quad a>0 \\[4pt]
&\text{(2)}\ \int \frac{1}{a^2+x^2} \,dx = \frac{1}{a}\arctan\frac{x}{a}+C \\[4pt]
&\text{(3)}\ \int \frac{1}{x^2-a^2} \,dx = \frac{1}{2a}\ln|\frac{x-a}{x+a}|+C \\[4pt]
&\text{(4)}\ \int \frac{1}{\sqrt{x^2\pm a^2}} \,dx = \ln|x+\sqrt{x^2\pm a^2}|+C \\[4pt]
&\text{(5)}\ \int \frac{1}{x\sqrt{x^2-1}} \,dx = \operatorname{arcsec}|x|+C \\[4pt]
&\text{(6)}\ \int \frac{1}{x\sqrt{1-x^2}} \,dx = -\ln|\frac{1+\sqrt{1-x^2}}{x}|+C \\[4pt]
&\text{(7)}\ \int \frac{\arcsin x}{\sqrt{1-x^2}} \,dx = \frac{1}{2}\arcsin^2 x + C \\[4pt]
&\text{(8)}\ \int \frac{\arctan x}{1+x^2} \,dx = \frac{1}{2}\arctan^2 x + C
\end{aligned}
\]
\end{theorem}

其中，公式$\int \frac{1}{a^2+x^2} \,dx = \frac{1}{a}\arctan\frac{x}{a}+C$和$\int \frac{1}{x^2-a^2} \,dx = \frac{1}{2a}\ln|\frac{x-a}{x+a}|+C$极为重要，应当牢记。

以及，一定一定别忘记加$C$。

\subsection{第二类换元法}

第一类换元法可以看成“从里往外”凑：先设内层函数 $u(x)$，再凑出 $u'(x)\,dx$。但有时候被积函数里有根号，例如 $\int \sqrt{1-x^2}\,dx$，我们根本看不出该凑什么。令 $u = \sqrt{x}$？显然不对。令 $u = 1-x^2$？那 $dx$ 和 $\sqrt{u}$ 搞不到一起去。好像怎么凑都不对。也许你要说，赶紧用AI救一救啊，别急，我们还有一招。

这时候就轮到第二类换元法出场了。第二类换元法可以看成“从外往里”换：我们主动令 $x = \phi(t)$，把整个积分换成关于 $t$ 的形式，希望换完之后根号消失或者变得好处理。由此我们有：

\begin{theorem}{第二类换元法}{}
设 $x = \phi(t)$ 单调可微，且 $\phi'(t) \neq 0$，则
\[
\int f(x)\, dx = \int f(\phi(t)) \cdot \phi'(t)\, dt
\]
\end{theorem}

等我们把右边的积分积出来之后，再把 $x$ 解出来代回去。

有没有发现和第一类换元法很像？第一类换元法是令 $u = u(x)$，然后用$du$表示$dx$。第二类换元法是令 $x = \phi(t)$，然后用$dx$表示 $\phi'(t)\,dt$。方向相反，但本质都是链式法则的逆用。只不过第一类是顺着已有的函数关系凑，第二类是主动创造一个新的函数关系来换。

那么问题来了，令 $x$ 等于什么好呢？第二类换元法最常见的应用场景是被积函数含有 $\sqrt{a^2 - x^2}$、$\sqrt{x^2 + a^2}$、$\sqrt{x^2 - a^2}$ 这三种根式。怎么处理它们呢？不知道你还记不记得，高中数学里有一招叫做三角换元，你还能想起来吗？

想不起来没关系，我们现在带你回忆一下：

\begin{theorem}{三角换元积分公式}{}
\[
\begin{aligned}
&\text{(1)}\ \int \sqrt{a^2 - x^2}\, dx = \frac{a^2}{2}\arcsin\frac{x}{a} + \frac{x}{2}\sqrt{a^2-x^2} + C, \quad a>0 \\[4pt]
&\text{(2)}\ \int \sqrt{x^2 + a^2}\, dx = \frac{x}{2}\sqrt{x^2+a^2} + \frac{a^2}{2}\ln|x + \sqrt{x^2+a^2}| + C \\[4pt]
&\text{(3)}\ \int \sqrt{x^2 - a^2}\, dx = \frac{x}{2}\sqrt{x^2-a^2} - \frac{a^2}{2}\ln|x + \sqrt{x^2-a^2}| + C, \quad a \leq |x|
\end{aligned}
\]
\end{theorem}

我们以第一个公式为例，看看具体的操作：

\begin{myexample}
  求 $\int \sqrt{a^2 - x^2}\, dx$（$a > 0$）。

  看到 $\sqrt{a^2 - x^2}$，怎么处理呢？刚刚说到，用三角换元。怎么换呢？让我们回忆一下，三角换元里利用的是$\sin^2 t + \cos^2 t = 1$，即 $1 - \sin^2 t = \cos^2 t$。那么我们只需要令 $x = a\sin t$，其中 $-\frac{\pi}{2} \leq t \leq \frac{\pi}{2}$，则 $dx = a\cos t\, dt$，$\sqrt{a^2 - x^2} = a\cos t$。根号消失了。

  代入原积分：
  \[
  \int \sqrt{a^2 - x^2}\, dx = \int a\cos t \cdot a\cos t\, dt = a^2 \int \cos^2 t\, dt
  \]

  用半角公式 $\cos^2 t = \frac{1 + \cos 2t}{2}$：
  \[
  a^2 \int \cos^2 t\, dt = a^2\!(\frac{t}{2} + \frac{\sin 2t}{4}) + C = \frac{a^2}{2}t + \frac{a^2}{2}\sin t\cos t + C
  \]

  解出$x$并代回：

  由 $\sin t = \frac{x}{a}$，$\cos t = \frac{\sqrt{a^2-x^2}}{a}$，$t = \arcsin\frac{x}{a}$，得：

  \[
  \int \sqrt{a^2 - x^2}\, dx = \frac{a^2}{2}\arcsin\frac{x}{a} + \frac{x}{2}\sqrt{a^2-x^2} + C
  \]
\end{myexample}

注意代回这一步。令 $x = a\sin t$ 时，$t = \arcsin \frac{x}{a}$，但 $\cos t$ 是多少？硬算的话， $\cos t = \frac{\sqrt{a^2-x^2}}{a}$也可以。不过我们更推荐的是画直角三角形：因为 $\sin t = \frac{x}{a}$，所以我们设角 $t$ 的对边为 $x$、斜边为 $a$，邻边自然就是 $\sqrt{a^2-x^2}$。然后 $\cos t$、$\tan t$ 等都可以直接从这个三角形里读出来，比硬算三角函数快得多。

\begin{center}
\begin{tikzpicture}[scale=1.5, font=\small]
  % 顶点定义：角 t 在左下 (0,0)，直角在右下 (3,0)，对边顶点在右上 (3,1.8)
  \coordinate (T) at (0,0);   % 锐角 t 的顶点
  \coordinate (R) at (3,0);   % 直角顶点
  \coordinate (P) at (3,1.8); % 对边顶点

  % 三角形
  \draw[thick] (T) -- (R) -- (P) -- cycle;

  % 直角标记（放在 R 处）
  \draw (2.8,0) -- (2.8,0.2) -- (3,0.2);

  % 角 t 的弧线（圆心在 T，从 0° 到 atan(1.8/3)≈31°）
  \draw (0.7,0) arc (0:31:0.7);
  \node at (0.85,0.28) {$t$};

  % 边标注
  \node[below] at (1.5,0) {$\sqrt{a^2-x^2}$};
  \node[right] at (3,0.9) {$x$};
  \node[sloped, above] at (1.5,0.95) {$a$};

  % 右侧公式
  \node[align=left, anchor=west] at (4.2,0.7) {
    \(\displaystyle\sin t = \frac{x}{a}\) \\[8pt]
    \(\displaystyle\cos t = \frac{\sqrt{a^2-x^2}}{a}\) \\[8pt]
    \(\displaystyle\tan t = \frac{x}{\sqrt{a^2-x^2}}\)
  };
\end{tikzpicture}
\end{center}

当然我们不是每次都已知$\sin t = \frac{x}{a}$，可能已知的是其他三角函数值，不过原理是一样的。根据已知三角函数值设边长，求出所有边长之后再算三角函数值即可。后面所有三角换元的代回都可以用这个技巧。

我们再来看第二条公式：

\begin{myexample}
  求 $\int \sqrt{x^2 + a^2}\, dx$（$a > 0$）。

  看到 $\sqrt{x^2 + a^2}$，怎么处理呢？这时候就要用到我们新学的三角恒等式了。因为 $1 + \tan^2 t = \sec^2 t$。设 $x = a\tan t$，其中$-\frac{\pi}{2} < t < \frac{\pi}{2}$，则 $dx = a\sec^2 t\, dt$，$\sqrt{x^2 + a^2} = a\sec t$。根号消失了。

  代入：
  \[
  \int \sqrt{x^2 + a^2}\, dx = \int a\sec t \cdot a\sec^2 t\, dt = a^2 \int \sec^3 t\, dt
  \]

  这里出现了 $\int \sec^3 t\, dt$，由公式：$\int \sec^n t\,dt = \frac{\sec^{n-2}t\tan t}{n-1} + \frac{n-2}{n-1}\int \sec^{n-2}t\,dt$，其中$n=3$，得：
\[
a^2 \int \sec^3 t\,dt =a^2(\frac{\sec t\tan t}{2} + \frac{1}{2}\int \sec t\,dt)
\]

其中 $\int \sec t\,dt$ 由公式：\(\int \sec t\,dt = \ln|\sec t + \tan t| + C\)，得：

\[
a^2 \int \sec^3 t\,dt = a^2 (\frac{1}{2}\sec t \tan t + \frac{1}{2}\ln|\sec t + \tan t|) + C
\]

画三角形：
\begin{center}
\begin{tikzpicture}[scale=1.5, font=\small]
  % 顶点定义：角 t 在左下 (0,0)，直角在右下 (3,0)，对边顶点在右上 (3,1.8)
  \coordinate (T) at (0,0);   % 锐角 t 的顶点
  \coordinate (R) at (3,0);   % 直角顶点
  \coordinate (P) at (3,1.8); % 对边顶点

  % 三角形
  \draw[thick] (T) -- (R) -- (P) -- cycle;

  % 直角标记（放在 R 处）
  \draw (2.8,0) -- (2.8,0.2) -- (3,0.2);

  % 角 t 的弧线（圆心在 T，从 0° 到 atan(1.8/3)≈31°）
  \draw (0.7,0) arc (0:31:0.7);
  \node at (0.85,0.28) {$t$};

  % 边标注
  \node[below] at (1.5,0) {$a$};
  \node[right] at (3,0.9) {$x$};
  \node[sloped, above] at (1.2,0.95) {$\sqrt{a^2+x^2}$};

  % 右侧公式
  \node[align=left, anchor=west] at (4.2,0.7) {
    \(\displaystyle\tan t = \frac{x}{a}\) \\[8pt]
    \(\displaystyle\sec t = \frac{\sqrt{a^2+x^2}}{a}\)
  };
\end{tikzpicture}
\end{center}

代入原积分，得：
\[
\begin{aligned}
a^2 ( \frac{1}{2}\sec t \tan t + \frac{1}{2}\ln|\sec t + \tan t| ) + C
&= \frac{a^2}{2} \cdot \frac{\sqrt{x^2+a^2}}{a} \cdot \frac{x}{a} + \frac{a^2}{2} \ln| \frac{\sqrt{x^2+a^2}}{a} + \frac{x}{a} | + C \\[4pt]
= \frac{x}{2}\sqrt{x^2+a^2} + \frac{a^2}{2} \ln| \frac{x+\sqrt{x^2+a^2}}{a} | + C 
&= \frac{x}{2}\sqrt{x^2+a^2} + \frac{a^2}{2} \ln| x+\sqrt{x^2+a^2} | - \frac{a^2}{2}\ln a + C
\end{aligned}
\]

将常数 $-\frac{a^2}{2}\ln a$ 并入 $C$，最终得到：
\[
\int \sqrt{x^2+a^2}\,dx = \frac{x}{2}\sqrt{x^2+a^2} + \frac{a^2}{2} \ln| x+\sqrt{x^2+a^2} | + C
\]

\end{myexample}

好了，那么剩下一个$\int \sqrt{x^2 - a^2}\, dx$该怎么做呢？我相信有了前两个例子，聪明的你已经想到怎么做了。不过，实际操作中，我们遇到的往往是它们与其他函数的复合：

\begin{myexample}
  求 $\int \frac{1}{\sqrt{x^2 + 4}}\, dx$。

  设 $x = 2\tan t$，其中 $-\frac{\pi}{2} < t < \frac{\pi}{2}$，则 $dx = 2\sec^2 t\, dt$，$\sqrt{x^2 + 4} = 2\sec t$。

代入：
\[
\int \frac{1}{\sqrt{x^2+4}}\, dx = \int \frac{1}{2\sec t} \cdot 2\sec^2 t\, dt = \int \sec t\, dt
\]

由公式 $\displaystyle \int \sec t\,dt = \ln|\sec t + \tan t| + C$，得：
\[
\int \sec t\,dt = \ln|\sec t + \tan t| + C
\]

画三角形：
\begin{center}
\begin{tikzpicture}[scale=1.5, font=\small]
  % 角 t 在左下 (0,0)，直角在右下 (2,0)，对边顶点在右上 (2,1.5)
  \coordinate (T) at (0,0);   % 锐角 t 的顶点
  \coordinate (R) at (3,0);   % 直角顶点
  \coordinate (P) at (3,1.8); % 对边顶点

  \draw[thick] (T) -- (R) -- (P) -- cycle;

  % 直角标记
  \draw (2.8,0) -- (2.8,0.2) -- (3,0.2);

  % 角 t 弧线
  \draw (0.7,0) arc (0:31:0.7);
  \node at (0.85,0.28) {$t$};

  % 边标注
  \node[below] at (1.5,0) {$2$};
  \node[right] at (3,0.75) {$x$};
  \node[sloped, above] at (1,0.85) {$\sqrt{x^2+4}$};

  % 右侧三角函数值
  \node[align=left, anchor=west] at (4.2,0.7) {
    $\displaystyle\tan t = \frac{x}{2}$ \\[8pt]
    $\displaystyle\sec t = \frac{\sqrt{x^2+4}}{2}$
  };
\end{tikzpicture}
\end{center}

代入原积分，得：
\[
\begin{aligned}
\ln|\sec t + \tan t| + C
&= \ln| \frac{\sqrt{x^2+4}}{2} + \frac{x}{2} | + C \\[4pt]
&= \ln| \frac{x + \sqrt{x^2+4}}{2} | + C \\[4pt]
&= \ln| x + \sqrt{x^2+4} | - \ln 2 + C
\end{aligned}
\]

将常数 $-\ln 2$ 并入 $C$，最终得到：
\[
\int \frac{1}{\sqrt{x^2+4}}\,dx = \ln| x + \sqrt{x^2+4} | + C
\]
\end{myexample}

除了三角换元，还有一种常用的换元是\textbf{根式换元}。当被积函数含有 $\sqrt[n]{ax+b}$ 时，我们令 $t = \sqrt[n]{ax+b}$ 就可以直接把根号消掉。

\begin{theorem}{根式换元积分公式}{}
  \[
\int \sqrt[n]{ax+b}\,dx = \frac{n}{a(n+1)}(ax+b)^{\frac{n+1}{n}} + C \qquad (a\neq 0,\, n\neq -1)
\]
\end{theorem}

背后的原理其实是线性函数凑微分公式：$\int (ax+b)^n \,dx = \frac{1}{a}\cdot\frac{(ax+b)^{n+1}}{n+1}+C$，其中$n\neq -1 $。

\begin{myexample}
  求 $\int \frac{1}{\sqrt{x} + 1}\, dx$。

  被积函数里有 $\sqrt{x}$，我们可以尝试根式换元。令 $t = \sqrt{x}$，则 $x = t^2$，$dx = 2t\, dt$。

  代入原积分：
  \[
  \int \frac{1}{\sqrt{x} + 1}\, dx = \int \frac{1}{t + 1} \cdot 2t\, dt = \int \frac{2t}{t+1}\, dt
  \]

  处理分式：$\frac{2t}{t+1} = 2 - \frac{2}{t+1}$，于是：

  \[
  \int (2 - \frac{2}{t+1}) dt = 2t - 2\ln|t+1| + C = 2\sqrt{x} - 2\ln(\sqrt{x}+1) + C
  \]

\end{myexample}

总结一下第二类换元法的使用场景：

看到 $\sqrt{a^2 - x^2}$，$\sqrt{x^2 + a^2}$，$\sqrt{x^2 - a^2}$，用三角换元。


看到 $\sqrt[n]{ax+b}$，用根式换元。


以及，一定一定别忘记加$C$。


\subsection{分部积分法}

换元法解决的是复合函数型的积分。但还有一类常见的情况：被积函数是两个不同类型函数的乘积，比如 $\int x e^x\,dx$、$\int x \sin x\,dx$。换元法怎么做呢？读者可以先想一想。

想不出来是正常的，因为我们需要另一种工具。

让我们先回忆一下乘法求导法则：
\[
(uv)' = u'v + uv'
\]

两边积分：
\[
\int (uv)'\, dx = \int u'v\, dx + \int uv'\, dx
\]

左边就是 $uv$（省略常数）：
\[
uv = \int u'v\, dx + \int uv'\, dx
\]

移项得：
\[
\int uv'\, dx = uv - \int u'v\, dx
\]

写成微分形式就是：

\begin{theorem}{分部积分公式}{}
\[
\int u\, dv = uv - \int v\, du
\]
\end{theorem}

这个公式的精髓在于，右边的 $\int v\, du$ 可能比左边的 $\int u\, dv$ 更好积分。关键在于怎么选 $u$ 和 $dv$。选对了，直接一通带走。选错了，可能越积越复杂。

一般的经验法则是\textbf{把容易求导的函数选为 $u$，把容易积分的函数选为 $dv$}。一般来说，按“对、反、幂、三、指”的优先级选择 $u$，也就是：
\begin{center}
  对数函数$>$反三角函数$>$幂函数$>$三角函数$>$指数函数
\end{center}



来看几个例子。

\begin{myexample}
  求 $\int x e^x\, dx$。

  分析被积函数，这是一个幂函数和一个指数函数相乘。根据“对数函数$>$反三角函数$>$幂函数$>$三角函数$>$指数函数”，那么我们应该选$x$作为$u$。

  令 $u = x$，$dv = e^x\, dx$。则 $du = dx$，$v = e^x$。
  
  代入公式：
  \[
  \int x e^x\, dx = x e^x - \int e^x\, dx = x e^x - e^x + C = (x-1)e^x + C
  \]

\end{myexample}

我们可以验证一下：$((x-1)e^x + C)' = e^x + (x-1)e^x = x e^x$，没问题。

这个例子展示了分部积分法的核心思想：$x$ 求导后变成 $1$，$e^x$ 积分后还是 $e^x$，所以右边的 $\int e^x\,dx$ 变得非常简单。

\begin{myexample}
  求 $\int x \sin x\, dx$。

  分析被积函数，这是一个幂函数和一个三角函数相乘。根据“对数函数$>$反三角函数$>$幂函数$>$三角函数$>$指数函数”，那么我们应该选$x$作为$u$。

  令 $u = x$，$dv = \sin x\, dx$。则 $du = dx$，$v = -\cos x$。
  
  代入公式：
  \[
  \int x \sin x\, dx = -x\cos x - \int (-\cos x)\, dx = -x\cos x + \int \cos x\, dx = -x\cos x + \sin x + C
  \]
\end{myexample}

\begin{myexample}
  求 $\int \ln x\, dx$。

  这个积分没有乘积怎么办？其实，这个积分可以看成 $\int 1 \cdot \ln x\, dx$。分析被积函数，这是一个常数函数和一个对数函数相乘。根据“对数函数$>$反三角函数$>$幂函数$>$三角函数$>$指数函数”，那么我们应该选$\ln x$作为$u$。
  
  令 $u = \ln x$，$dv = dx$。则 $du = \frac{1}{x}\, dx$，$v = x$。
  
  代入公式：
  \[
  \int \ln x\, dx = x\ln x - \int x \cdot \frac{1}{x}\, dx = x\ln x - \int 1\, dx = x\ln x - x + C
  \]

\end{myexample}

  这个结果非常有用，建议记住：$\int \ln x\, dx = x\ln x - x + C$。
  
  有些积分需要连续使用两次分部积分，比如：

\begin{myexample}
  求 $\int e^x \sin x\, dx$。

  分析被积函数，这是一个指数函数和一个三角函数相乘。根据“对数函数$>$反三角函数$>$幂函数$>$三角函数$>$指数函数”，那么我们应该选$\sin x$作为$u$。

  令 $u = \sin x$，$dv = e^x\, dx$。则 $du = \cos x\, dx$，$v = e^x$。

  代入公式：
  \[
  \int e^x \sin x\, dx = e^x \sin x - \int e^x \cos x\, dx
  \]

  右边出现了 $\int e^x \cos x\, dx$，这时候我们可以再用一次分部积分。
  
  令 $u = \cos x$，$dv = e^x\, dx$，则 $du = -\sin x\, dx$，$v = e^x$。
  
  代入公式：
  \[
  \int e^x \cos x\, dx = e^x \cos x - \int e^x (-\sin x)\, dx = e^x \cos x + \int e^x \sin x\, dx
  \]

  代入原积分：
  \[
  \int e^x \sin x\, dx = e^x \sin x - (e^x \cos x + \int e^x \sin x\, dx)
  \]

  这时候我们发现，右边又出现了 $\int e^x \sin x\, dx$，我们把它移到左边：
  \[
  2\int e^x \sin x\, dx = e^x \sin x - e^x \cos x + C
  \]

  两边除以 $2$：
  \[
  \int e^x \sin x\, dx = \frac{e^x(\sin x - \cos x)}{2} + C
  \]
\end{myexample}

这个例子展示了分部积分的一个神奇现象，当某些积分连续用了两次分部积分法的时候，原来的积分“绕回来了”。这时候把它当成一个未知数，解方程就能得到结果。这种技巧叫做“循环分部积分”，遇到 $e^x$ 和三角函数的乘积时经常用到。

总结一下分部积分法的核心，就是选择关键的$u$和$dv$，大部分时候分部积分都是两个函数的乘积，按照优先级选择$u$基本不会选错。

以及，一定一定别忘记加$C$。

虽然说我们用上述方法能解出大部分的不定积分，但是有些积分是我们很难解或者不能用初等函数表示的。

例如我们口算都能算出来\(\int \frac{1}{x^5} dx=-\frac{1}{4x^4}+C\)，但是\(\int \frac{1}{1+x^5}dx\)等于多少呢？

实际上，这是一个非常难算的积分：

\[
\begin{aligned}
\int \frac{dx}{1+x^5} 
&= \frac{1}{5}\ln|x+1| 
- \frac{1}{20}\ln\!(x^4 - x^3 + x^2 - x + 1) \\[4pt]
&\quad + \frac{1}{5}\sqrt{\frac{5+\sqrt{5}}{2}}\;
\arctan\!(\frac{4x - 1 + \sqrt{5}}{\sqrt{10 - 2\sqrt{5}}}) \\[4pt]
&\quad + \frac{1}{5}\sqrt{\frac{5-\sqrt{5}}{2}}\;
\arctan\!(\frac{4x - 1 - \sqrt{5}}{\sqrt{10 + 2\sqrt{5}}}) + C
\end{aligned}
\]

更进一步的，\(\int \frac{1}{1+x+x^5}dx\)等于多少呢？

还有一些积分根本无法用初等函数表达。例如：

\[
\begin{aligned}
&\int e^{-x^2}\,dx,\qquad \int \sin x^2\,dx,\qquad \int \cos x^2\,dx,\qquad \int \frac{\sin x}{x}\,dx,\qquad \int \frac{\cos x}{x}\,dx,\\
&\int \frac{e^x}{x}\,dx,\qquad \int \frac{1}{\ln x}\,dx,\qquad \int \frac{dx}{\sqrt{1-k^2\sin^2 x}},\qquad \int \sqrt{1-k^2\sin^2 x}\,dx,\\
&\int \frac{dx}{\sqrt{x^3+1}},\qquad \int \ln \ln x\,dx,\qquad \int \sin\!(\frac{1}{x})dx,\qquad \int x^x\,dx
\end{aligned}
\]

不相信？自己动手算算呗。

\section{定积分}
我们已经学会了导数，学会了不定积分。现在我们要把无数个瞬间的微小结果重新累积起来，求得一个确定的总量，这就是定积分。至此，我们算是踏进了微积分的大门。

\subsection{定积分的几何意义}

想象一下，你在平面上画了一条曲线 $y = f(x)$（假设 $f(x) \geq 0$），然后在 $x$ 轴上标出区间 $[a, b]$。这条曲线、$x$ 轴以及两条竖直线 $x = a$ 和 $x = b$ 围成的区域，就是所谓的\textbf{曲边梯形}。定积分 $\int_a^b f(x)\,dx$ 的几何意义就是这个曲边梯形的面积。

\begin{center}
  \begin{tikzpicture}[scale=1.3]
    % 定义函数和区间
    \def\a{0.5}     % 左端点
    \def\b{4.5}     % 右端点
    
    % 先填充曲边梯形区域（曲线下方，从 a 到 b）
    \fill[cyan!20] 
        plot[domain=\a:\b, smooth] (\x, {0.1*\x*\x + 0.5}) 
        -- (\b,0) -- (\a,0) -- cycle;
    
    % 绘制区间端点竖虚线
    \draw[dashed, gray] (\a, 0) -- (\a, {0.1*(\a)^2 + 0.5});
    \draw[dashed, gray] (\b, 0) -- (\b, {0.1*(\b)^2 + 0.5});
    
    % 绘制曲线（完整显示，从 0.2 到 5.0 保留一点延伸）
    \draw[domain=0.2:5.0, smooth, very thick, blue] 
        plot (\x, {0.1*\x*\x + 0.5}) 
        node[above left] {$y=f(x)$};
    
    % 坐标轴（最后绘制，使箭头和线条位于最上层）
    \draw[->] (0,0) -- (5.2,0) node[below] {$x$};
    \draw[->] (0,0) -- (0,3.2) node[left] {$y$};
    
    % 原点标签
    \node[below left] at (0,0) {$O$};
    
    % 区间标记
    \node[below] at (\a, 0) {$a$};
    \node[below] at (\b, 0) {$b$};
    
\end{tikzpicture}
\begin{tikzpicture}[scale=1.3]
    % 定义函数和区间
    \def\a{0.5}     % 左端点
    \def\b{4.5}     % 右端点
    \def\n{7}       % 分割份数
    
    % 坐标轴（y轴上限调整为3.2）
    \draw[->] (0,0) -- (5.2,0) node[below] {$x$};
    \draw[->] (0,0) -- (0,3.2) node[left] {$y$};
    
    % 原点标签
    \node[below left] at (0,0) {$O$};
    
    % 画曲线 y = f(x) = 0.1x^2 + 0.5
    \draw[domain=0.2:5.0, smooth, thick, blue] 
        plot (\x, {0.1*\x*\x + 0.5}) 
        node[above left] {$y=f(x)$};
    
    % 计算并绘制矩形（取左端点高度）
    \foreach \i in {0,1,...,6} {
        \pgfmathsetmacro\xleft{\a + (\b-\a)*\i/\n}
        \pgfmathsetmacro\xright{\a + (\b-\a)*(\i+1)/\n}
        \pgfmathsetmacro\fheight{0.1*(\xleft)^2 + 0.5}
        
        \fill[cyan!30, draw=cyan!60, thick] 
            (\xleft, 0) rectangle (\xright, \fheight);
    }
    
    % 再画一次曲线，使其盖在矩形上面
    \draw[domain=0.2:5.0, smooth, very thick, blue] 
        plot (\x, {0.1*\x*\x + 0.5});
    
    % 区间端点竖线
    \draw[dashed, gray] (\a, 0) -- (\a, {0.1*(\a)^2 + 0.5});
    \draw[dashed, gray] (\b, 0) -- (\b, {0.1*(\b)^2 + 0.5});
    
    % 区间标记
    \node[below] at (\a, 0) {$a$};
    \node[below] at (\b, 0) {$b$};
    
    % 突出显示第5个矩形 (索引4)
    \pgfmathsetmacro\xleftfive{\a + (\b-\a)*3/\n}
    \pgfmathsetmacro\xrightfive{\a + (\b-\a)*4/\n}
    \pgfmathsetmacro\fheightfive{0.1*(\xleftfive)^2 + 0.5}
    
    \fill[orange!50, draw=orange!80, very thick, opacity=0.7] 
        (\xleftfive, 0) rectangle (\xrightfive, \fheightfive);
    
    % 标注 xi_i 和 f(xi_i)
    \pgfmathsetmacro\xmid{\xleftfive + (\xrightfive-\xleftfive)/2}
    \node[above left, orange!90] at (\xmid, \fheightfive) {$f(\xi_i)$};
    \node[below left, orange!90] at (\xmid, 0.05) {$\xi_i$};
    
    % 宽度标注
    \draw[|<->|, red!70, thick] 
        (\xleftfive, 0.15) -- (\xrightfive, 0.15);
    \node[above, red!70] at (\xmid, 0.25) {$\Delta x_i$};
    
\end{tikzpicture}
\end{center}

那要怎么求这个曲边梯形的面积呢？我们可以把区间 $[a, b]$ 分成 $n$ 个小段，每段宽度为 $\Delta x_i$，在每段上取一个点 $\xi_i$，用 $f(\xi_i) \cdot \Delta x_i$ 作为这一小条矩形的面积。那么所有小矩形面积之和就是：
\[
\sum_{i=1}^{n} f(\xi_i)\, \Delta x_i
\]

当$n$足够大的时候，曲线就可以近似看成直线，面积就可以用小矩形近似，也就是说，我们把区间$[a, b]$分得越细，近似就越精确。用数学语言表示就是：当$n \to \infty$，$\Delta x_i \to 0$时：
\[
\int_a^b f(x)\, dx = \lim_{\Delta x_i \to 0} \sum_{i=1}^{n} f(\xi_i)\, \Delta x_i
\]

由此，我们得出了定积分的定义：

\begin{definition}{定积分}{}
设函数 \(f(x)\) 在 \([a,b]\) 上有界，在 \([a,b]\) 中任意插入若干个分点
\[
a=x_0<x_1<x_2<\cdots<x_n=b
\]
把区间 \([a,b]\) 分成 \(n\) 个小区间
\[
[x_0,x_1], \ [x_1,x_2], \ \cdots, \ [x_{n-1},x_n]
\]
各个小区间的长度依次为
\[
\Delta x_1=x_1-x_0, \ \Delta x_2=x_2-x_1, \ \cdots, \ \Delta x_n=x_n-x_{n-1}
\]
在每个小区间 \([x_{i-1},x_i]\) 上任取一点 \(\xi_i\) (\(x_{i-1} \leqslant \xi_i \leqslant x_i\))，作函数值 \(f(\xi_i)\) 与小区间长度 \(\Delta x_i\) 的乘积 \(f(\xi_i) \cdot \Delta x_i\)，并作出和

\[
S=\sum_{i=1}^{n} f(\xi_i) \Delta x_i
\]

记 \(\lambda = \max\{\Delta x_1, \Delta x_2, \cdots, \Delta x_n\}\)，如果当 \(\lambda \to 0\) 时，\(S\) 的极限总存在，且与闭区间 \([a,b]\) 的分法及点 \(\xi_i\) 的取法无关，那么我们就说\(f(x)\)在\([a,b]\)上\textbf{可积}，并且称这个极限 \(I\) 为函数 \(f(x)\) 在区间 \([a,b]\) 上的定积分，记作 \(\int_{a}^{b} f(x) dx\)，即：

\[
\int_{a}^{b} f(x) dx = \lim_{\lambda \to 0} \sum_{i=1}^{n} f(\xi_i) \Delta x_i = I
\]

其中 \(f(x)\) 称为\textbf{被积函数}，\(f(x) dx\) 称为\textbf{被积表达式}，\(x\) 称为\textbf{积分变量}，\(a\) 称为\textbf{积分下限}，\(b\) 称为\textbf{积分上限}，\([a,b]\) 称为\textbf{积分区间}。

\end{definition}

定积分定义的背后蕴含着“以直代曲”和“微元法”朴素的数学思想。这也说明了积分号 $\int$ 的来源：$\int$是一个拉长的字母 S（Sum），代表求和。

这个定义看上去很复杂，其实说的就是一件事：分割，近似，求和，取极限。把区间分割成一个个小区间，在小区间里用小矩形近似小曲边梯形，把所有小矩形的面积加起来，再让小区间长度趋近于$0$，取和的极限，定积分本质上说的就是这么一回事。

而几何意义就更直观了。$\int_{a}^{b} f(x) dx$表示由曲线$f(x)$、$x$ 轴以及两条竖直线 $x = a$ 和 $x = b$ 围成的区域的面积。

但要注意，如果 $f(x)$ 在某些地方为负，那“面积”怎么算？定积分对负的部分会给出负值，因为 $\int_a^b f(x)\, dx$ 计算的是\textbf{代数面积}，也就是说，$x$ 轴上方的面积为正，下方的面积为负，最后再加在一起。

根据定义，我们能得到几个性质：

\begin{theorem}{初等函数的可积性}{}
  任何初等函数在其定义域内的任意闭区间上都是可积的。
\end{theorem}

\begin{theorem}{零区间积分}{}
  \[\int_{a}^{a} f(x) dx=0\]
\end{theorem}

\begin{theorem}{线性性质}{} 
  若$f,g$可积，$\alpha,\beta \in \mathbb{R}$，则$\alpha f+\beta g$可积，且：

  \[\int_{a}^{b} [\alpha f(x)+\beta g(x)] dx=\alpha \int_{a}^{b} f(x) dx+\beta \int_{a}^{b} g(x) dx\]
\end{theorem}

\begin{theorem}{上下限互换}{}
  \[\int_{a}^{b} f(x) dx=-\int_{b}^{a} f(x) dx\]
\end{theorem}

\begin{theorem}{区间长度性质}{}
  \[\int_{a}^{b} dx=b-a\]
\end{theorem}


由此得到：

\begin{theorem}{区间可加性}{} 
  对任意$a,b,c\in \mathbb{R} $，有：

  \[\int_{a}^{b} f(x) dx=\int_{a}^{c} f(x) dx+\int_{c}^{b} f(x) dx\]
\end{theorem}

\begin{theorem}{奇偶性}{} 
  设$f(x)$在区间$[-a,a]$上可积。

  若$f(x)$为奇函数，则：\(\int_{-a}^{a} f(x) dx=0\)

  若$f(x)$为偶函数，则：\(\int_{-a}^{a} f(x) dx=2\int_{0}^{a} f(x) dx\)
\end{theorem}

\begin{theorem}{保序性}{} 
  若$f(x)\leq g(x)$，则：
  \[\int_{a}^{b} f(x) dx \leq \int_{a}^{b} g(x) dx\]
\end{theorem}

我们来看几个例子：

\begin{myexample}
  用定积分的几何意义，求 $\int_0^1 x\, dx$。

  $\int_0^1 x\, dx$表示由直线$y=x$、$x$ 轴以及两条竖直线 $x = 0$ 和 $x = 1$ 围成的区域的面积。

  直观上，$f(x) = x$ 在 $[0, 1]$ 上是一条从原点出发到 $(1, 1)$ 的直线，曲边梯形退化为一个直角三角形，底为 $1$，高为 $1$，面积为 $\frac{1}{2}$。

  因此 $\int_0^1 x\, dx = \frac{1}{2}$。
\end{myexample}

\begin{myexample}
  用定积分的几何意义，求 $\int_0^{2\pi} \sin x\, dx$。

  $\int_0^{2\pi} \sin x\, dx$表示由曲线$y=\sin x$、$x$ 轴以及两条竖直线 $x = 0$ 和 $x = 2\pi$ 围成的区域的面积。

  直观上，$\sin x$ 在 $[0, \pi]$ 上为正，在 $[\pi, 2\pi]$ 上为负。两个区域的几何形状完全对称，面积相等但符号相反，所以代数面积为 $0$。
\end{myexample}

某种程度上，这也算是定积分奇偶性的体现。

如果要求实际的面积，那就需要把负的部分取绝对值，分段积分，也就是：\(\int_a^b |f(x)|\, dx\)

在上面的例子中，$\sin x$ 在 $[0, 2\pi]$ 上的实际面积为 $\int_0^\pi \sin x\, dx + \int_\pi^{2\pi} (-\sin x)\, dx$，即一上一下的面积之和。

总结一下：

定积分 $\int_a^b f(x)\, dx$ 的几何意义是代数面积。如果 $f(x) \geq 0$，那它就是普通的面积。如果要求实际面积，需要取绝对值分段积分。

定积分的定义是分割，近似，求和，取极限。


\subsection{积分上限函数}

我们已经知道定积分 $\int_a^b f(x)\,dx$ 表示的是曲边梯形的面积。但如果我们把上限 $b$ 换成一个变量 $x$，情况就变了。$\int_a^x f(t)\,dt$ 不再是一个固定的数，而是随着 $x$ 的变化而变化。当 $x$ 变大时，积分区间变长，面积就变大。$x$ 变小时，面积就变小。也就是说，这个时候 $\int_a^x f(t)\,dt$ 是一个关于 $x$ 的函数。

这里要注意我们的积分变量用 $t$ 而不用 $x$，这是为了避免和积分上限 $x$ 混淆。

由此我们有：

\begin{definition}{积分上限函数}{}
设 $f(x)$ 在 $[a, b]$ 上可积，则对任意 $x \in [a, b]$，定义
\[
\Phi(x) = \int_a^x f(t)\, dt
\]
$\Phi(x)$ 称为 $f(x)$ 在 $[a, b]$ 上的\textbf{积分上限函数}。
\end{definition}

这个函数有一个极其重要的性质：它对 $x$ 的导数恰好就是被积函数 $f(x)$。

\begin{theorem}{积分上限函数的性质}{}
设 $f(x)$ 在 $[a, b]$ 上连续，则 $\Phi(x) = \int_a^x f(t)\, dt$ 在 $[a, b]$ 上可导，且
\[
\Phi'(x) = \frac{d}{dx} (\int_a^x f(t)\, dt) = f(x)
\]
\end{theorem}

我们简单证明一下：

证明它，我们需要分为三步。第一步，先证明估值定理：

根据连续函数的最值定理：若 $f(x)$ 在闭区间 $[a,b]$ 上连续，则 $f(x)$ 在 $[a,b]$ 上一定能取到最大值和最小值。记最大值为$M$，最小值为$m$，则有：\[m\leq f(x)\leq M\]

由区间长度性质得到：\(\int_{a}^{b} m dx=m(b-a)\)，\(\int_{a}^{b} M dx=M(b-a)\)

由保序性：\[m(b-a)\leq \int_{a}^{b} f(x) dx\leq M(b-a)\]

于是我们证明了估值定理。它给出了曲边梯形面积的上下界估计：

\begin{theorem}{估值定理}{}
设 $f(x)$ 在 $[a, b]$ 上连续且可积，且$\exists m,M\in \mathbb{R} $，使得$m\leq f(x)\leq M$，则有：
\[
m(b-a)\leq \int_{a}^{b} f(x) \,dx \leq M(b-a)
\]
\end{theorem}

第二步，证明积分中值定理：

由估值定理：\(m(b-a)\leq \int_{a}^{b} f(x) dx\leq M(b-a)\)，两边除以$b-a$，得到：\[
m\leq \frac{1}{b-a} \int_{a}^{b} f(x) \,dx \leq M
\]

记$\mu =\frac{1}{b-a} \int_{a}^{b} f(x) \,dx$，则$\mu$的大小介于$m$与$M$之间。

由于$f(x)$连续，不妨设$f(c)=m$，$f(d)=M$，其中$c,d\in[a,b]$，且$c<d$。

由连续函数的介值定理：若 $f(x)$ 在闭区间 $[c,d]$ 上连续，且 $f(c) \neq f(d)$，则对于 $f(c)$ 与 $f(d)$ 之间的任意一个数 $\mu$，都存在至少一个 $\xi \in (c,d)$，使得 $f(\xi) = \mu$。

则有：\[\int_{a}^{b} f(x) \,dx =f(\xi )(b-a)\]

于是我们证明了积分中值定理。它在微积分中有重要地位：

\begin{theorem}{积分中值定理}{}
设 $f(x)$ 在 $[a, b]$ 上连续且可积，则$\exists \xi \in [a,b]$，使得：
\[\int_{a}^{b} f(x) \,dx =f(\xi )(b-a)\]
\end{theorem}

第三步，证明积分上限函数的性质：

由区间可加性与上下限互换：\[\Phi (x+h)-\Phi (x)=\int_a^{x+h} f(t)\, dt-\int_a^x f(t)\, dt=\int_x^{x+h} f(t)\, dt\]

其中$h>0$且$x+h\in [a,b]$。

由积分中值定理：\(\int_x^{x+h} f(t)\, dt =h \cdot f(\xi )\)，两边除以$h$，即：\(\frac{1}{h} \int_x^{x+h} f(t)\, dt =f(\xi )\)

其中$\xi \in [x,x+h]$，所以$0\leq \xi -x\leq h$。

令$h\rightarrow 0$，由夹逼定理：$\lim_{h \to 0} (\xi -x) = 0 $，则$\lim_{h \to 0} \xi = x $。由于$f(x)$连续，因此$\lim_{h \to 0} f(\xi) = f(x) $。

综上，我们得到：

\[\Phi '(x)=\lim_{h \to 0}  \frac{\Phi (x+h)-\Phi (x)}{h} =\lim_{h \to 0} f(\xi)=f(x)\]

证毕。

这个性质还可以进一步推广：如果积分上限不是 $x$ ，而是一个关于 $x$ 的函数 $u(x)$，则需要加上链式法则：

\[
\Phi '(x)=\frac{d}{dx} (\int_a^{u(x)} f(t)\, dt) = f(u(x)) \cdot u'(x)
\]


我们举例说明一下。

\begin{myexample}
  设 $\Phi(x) = \int_0^x t^2\, dt$，求 $\Phi'(x)$。

  被积函数 $f(t) = t^2$，积分上限是 $x$，所以：
  \[
  \Phi'(x) = \frac{d}{dx} (\int_0^x t^2\, dt) = x^2
  \]
\end{myexample}

\begin{myexample}
  设 $\Phi(x) = \int_1^x \frac{\sin t}{t}\, dt$，求 $\Phi'(x)$。

  被积函数 $f(t) = \frac{\sin t}{t}$，积分上限是 $x$，所以：
  \[
  \Phi'(x) = \frac{\sin x}{x}
  \]

\end{myexample}

  $\int_{}^{} \frac{\sin t}{t} \,dt $ 是无法用初等函数表达的，但积分上限函数的导数仍然可以直接写出。这就是这个性质的强大之处，即使你不会算积分，也能求导。

\begin{myexample}
  求 $\frac{d}{dx} (\int_0^{x^2} \sin t\, dt)$。

  令 $u(x) = x^2$，则被积函数在 $t = x^2$ 处取值，再乘以 $u'(x) = 2x$：
  \[
  \frac{d}{dx} (\int_0^{x^2} \sin t\, dt) = \sin x^2 \cdot 2x = 2x\sin x^2
  \]
\end{myexample}

积分上限函数的导数公式是连接不定积分和定积分的桥梁，也是接下来微积分基本定理的基石。没有这个公式，我们就没有办法把“求面积”和“求原函数”联系起来。


\subsection{微积分基本定理}

前面我们学了不定积分，接下来我们学习著名的微积分基本定理，这个定理将揭示不定积分与定积分的内在联系，也是整个微积分学中最核心、最深刻的定理。

\begin{theorem}{微积分基本定理}{}
设 $f(x)$ 在 $[a, b]$ 上连续，$F(x)$ 是 $f(x)$ 在 $[a, b]$ 上的任意一个原函数，则：
\[
\int_a^b f(x)\, dx = F(b) - F(a)
\]
\end{theorem}

有时为了书写方便，也记作\(\int_a^b f(x)\, dx = [F(x)]_a^b\)。

这个定理也被称为牛顿-莱布尼茨公式，因为牛顿和莱布尼茨各自独立地发明了微积分，并分别从不同角度揭示了微分与积分之间的互逆关系。他们的贡献共同奠定了微积分的基础。

为什么这个定理成立？我们简要证明一下：

我们已经证明了积分上限函数 $\Phi(x) = \int_a^x f(t)\,dt$ 满足 $\Phi'(x) = f(x)$，所以 $\Phi(x)$ 本身就是 $f(x)$ 的一个原函数。那么 $\Phi(x)$ 和 $F(x)$ 只差一个常数 $C$，也就是：

\[F(x) = \Phi(x) + C\]

令 $x = a$，得： \[F(a)=\Phi(a) +C = \int_a^a f(t)\,dt + C \]

由零区间积分，得\(\int_a^a f(t)\,dt=0\)，即$C = F(a)$。

同理，令 $x = b$，得：
\[
F(b)=\Phi(b) +C = \int_a^b f(t)\,dt + C
\]

代入$C = F(a)$，移项得：

\[
\int_a^b f(t)\,dt=F(b)-F(a)
\]

证毕。

这个定理告诉我们，\textbf{定积分等于原函数在上限和下限处的值之差}。它把“求面积”转化成了“求原函数再代值”，也就是把一个几何问题转化为了代数问题。之前我们求 $\int_0^1 x\,dx$ 靠的是几何直觉，现在有了这个定理，我们可以用代数方法算出结果：

\begin{myexample}
  求 $\int_0^1 x\, dx$。

  $x$ 的一个原函数是 $\frac{x^2}{2}$，所以：
  
  \[
  \int_0^1 x\, dx= \frac{1^2}{2} - \frac{0^2}{2} = \frac{1}{2}
  \]
\end{myexample}

\begin{myexample}
  求 $\int_0^\pi \sin x\, dx$。

  $\sin x$ 的一个原函数是 $-\cos x$，所以：
  \[
  \int_0^\pi \sin x\, dx = (-\cos\pi) - (-\cos 0) = 1 - (-1) = 2
  \]
\end{myexample}

几何上，这正是 $\sin x$ 在 $[0, \pi]$ 上的面积。

\begin{myexample}
  求 $\int_1^e \frac{1}{x}\, dx$。

  $\frac{1}{x}$ 的一个原函数是 $\ln|x|$，所以：
  \[
  \int_1^e \frac{1}{x}\, dx = \ln e - \ln 1 = 1 - 0 = 1
  \]
\end{myexample}

\begin{myexample}
  求 $\int_0^1 e^x\, dx$。

  $e^x$ 的一个原函数是 $e^x$，所以：
  \[
  \int_0^1 e^x\, dx = e^1 - e^0 = e - 1
  \]
\end{myexample}

需要注意的是，使用牛顿-莱布尼茨公式时，$f(x)$ 必须在 $[a, b]$ 上连续。如果 $f(x)$ 在区间内有间断点，直接套公式可能出错。不连续的函数，我们可以分段求定积分，再将结果相加。

形象地说，积分是在总结过去，微分是在展望未来。积分从堆积的过往还原完整的面积，微分从当下的斜率预见未来的趋势。积分把一个个音符变成一首完整的歌，微分把一首歌分解为一个个音符。

微积分基本定理是整个微积分的枢纽。正如牛顿和莱布尼茨各自独立发现的那样，它把微分学和积分学统一了起来。

它是一座连接微分和积分的桥梁。牛顿用它描述了天体的运行，莱布尼茨用它书写了自然的规律。

这正是牛顿-莱布尼茨公式所告诉我们的。

那座桥梁，就是枢纽本身。

\subsection{定积分的计算方法}

有了微积分基本定理，定积分的计算就变成了“找原函数、代上限、减下限”的公式化操作。因此我们前面学过的所有不定积分技巧都可以直接用在定积分上。只不过有一些细节上的区别。

不定积分的换元法用 $u$ 替换后要把 $x$ 代回去。但定积分不需要代回去。换元的同时把积分上下限也换掉，最后直接用新变量的上下限代值计算。

\begin{theorem}{定积分的换元法}{}
设 $f(x)$ 在 $[a, b]$ 上连续，$x = \phi(t)$ 在 $[\alpha, \beta]$ 上单调可微，且 $\phi(\alpha) = a$，$\phi(\beta) = b$，则：
\[
\int_a^b f(x)\, dx = \int_\alpha^\beta f(\phi(t)) \cdot \phi'(t)\, dt
\]
\end{theorem}

也就是说，换元之后的上下限不一样了，因此，\textbf{换元时上下限要跟着换}。$\phi(\alpha) = a$ 意味着新下限 $\alpha$ 对应旧下限 $a$，$\phi(\beta) = b$意味着新上限 $\beta$ 对应旧上限 $b$。如果用 $u$ 代换，也是完全一样的，$u$ 的上下限由 $x$ 的上下限决定。

我们来看几个例子：

\begin{myexample}
  求 $\int_0^4 \frac{1}{1+\sqrt{x}}\, dx$。

  令 $t = \sqrt{x}$，则 $x = t^2$，$dx = 2t\, dt$。换上下限：$x = 0$ 时 $t = 0$，$x = 4$ 时 $t = 2$。
  \[
  \int_0^4 \frac{1}{1+\sqrt{x}}\, dx = \int_0^2 \frac{1}{1+t} \cdot 2t\, dt = \int_0^2 \frac{2t}{1+t}\, dt
  \]

  所以：
  \[
  \int_0^2 \frac{2t}{1+t}\, dt = [2t - 2\ln|1+t|]_0^2 = (4 - 2\ln 3) - (0 - 0) = 4 - 2\ln 3
  \]
\end{myexample}

这个例子说明，最后代入的是新变量 $t$ 的上下限 $0$ 和 $2$，不需要代回 $x$。

\begin{myexample}
  求 $\int_0^1 \sqrt{1-x^2}\, dx$。

  令 $x = \sin t$，则 $dx = \cos t\, dt$，$\sqrt{1-x^2} = \cos t$。换上下限：$x = 0$ 时 $t = 0$，$x = 1$ 时 $t = \frac{\pi}{2}$。
  \[
  \int_0^1 \sqrt{1-x^2}\, dx = \int_0^{\frac{\pi}{2}} \cos t \cdot \cos t\, dt = \int_0^{\frac{\pi}{2}} \cos^2 t\, dt
  \]

  用降幂公式 $\cos^2 t = \frac{1+\cos 2t}{2}$，得：
  \[
  \int_0^{\frac{\pi}{2}} \frac{1+\cos 2t}{2}\, dt = \frac{1}{2}[t + \frac{\sin 2t}{2}]_0^{\frac{\pi}{2}} = \frac{1}{2}(\frac{\pi}{2} + 0) - 0 = \frac{\pi}{4}
  \]
\end{myexample}
  
  几何上，这正是单位圆在第一象限的面积，大小正是$\frac{\pi}{4}$，完全合理。

和换元法一样，分部积分法也可以直接用在定积分上，只是公式里多了上下限：

\begin{theorem}{定积分的分部积分法}{}
\[
\int_a^b u\, dv = [uv]_a^b - \int_a^b v\, du
\]
\end{theorem}

\begin{myexample}
  求 $\int_0^1 x e^x\, dx$。

  令 $u = x$，$dv = e^x\, dx$，则 $du = dx$，$v = e^x$。代入公式：
  \[
  \int_0^1 x e^x\, dx = [xe^x]_0^1 - \int_0^1 e^x\, dx = (1 \cdot e - 0) - [e^x]_0^1 = e - (e - 1) = 1
  \]
\end{myexample}

\begin{myexample}
  求 $\int_0^{\frac{\pi}{2}} x \cos x\, dx$。

  令 $u = x$，$dv = \cos x\, dx$，则 $du = dx$，$v = \sin x$。代入公式：
  \[
  \int_0^{\frac{\pi}{2}} x \cos x\, dx = [x\sin x]_0^{\frac{\pi}{2}} - \int_0^{\frac{\pi}{2}} \sin x\, dx = \frac{\pi}{2} \cdot 1 - 0 - [-\cos x]_0^{\frac{\pi}{2}} = \frac{\pi}{2} - (0 - (-1)) = \frac{\pi}{2} - 1
  \]
\end{myexample}

总结一下：

定积分的计算方法和不定积分完全一致，区别只有两点，换元时上下限跟着变，计算结果直接用新变量的上下限代值计算。

\subsection{反常积分}

到目前为止，我们讨论的定积分都有两个隐含的前提：积分区间是有限的，被积函数是有界的。但实际中常常遇到不满足这两个条件的情形。比如积分区间是无穷的，例如$\int_1^{+\infty} \frac{1}{x^2}\,dx$，或者被积函数在某点无界，例如$\int_0^1 \frac{1}{\sqrt{x}}\,dx$。这类积分我们统一叫做\textbf{反常积分}。

反常积分的计算思路很简单：先用有限区间上的定积分来近似，再取极限。

\begin{definition}{无穷区间上的反常积分}{}
设 $f(x)$ 在 $[a, +\infty)$ 上连续，若极限
\[
\lim_{b \to +\infty} \int_a^b f(x)\, dx
\]
存在，则称反常积分 $\int_a^{+\infty} f(x)\, dx$ \textbf{收敛}，并定义$\int_a^{+\infty} f(x)\, dx$的值为该极限。若极限不存在，则称该反常积分\textbf{发散}。
\end{definition}

类似地可以定义 $\int_{-\infty}^b f(x)\,dx$ 和 $\int_{-\infty}^{+\infty} f(x)\,dx$。后者的计算要拆成两段：
\[
\int_{-\infty}^{+\infty} f(x)\, dx = \int_{-\infty}^{0} f(x)\, dx + \int_0^{+\infty} f(x)\, dx
\]

当两段都收敛时才说整个积分收敛。

\begin{myexample}
  判断 $\int_1^{+\infty} \frac{1}{x^2}\, dx$ 是否收敛，若收敛则求值。

  先算有限区间上的定积分：
  \[
  \int_1^b \frac{1}{x^2}\, dx = [-\frac{1}{x}]_1^b = -\frac{1}{b} + 1 = 1 - \frac{1}{b}
  \]

  再取极限：
  \[
  \int_1^{+\infty} \frac{1}{x^2}\, dx = \lim_{b \to +\infty} (1 - \frac{1}{b}) = 1
  \]

  因此该积分收敛，值为 $1$。
\end{myexample}

\begin{myexample}
  判断 $\int_1^{+\infty} \frac{1}{x}\, dx$ 是否收敛。

  \[
  \int_1^b \frac{1}{x}\, dx = [\ln x]_1^b = \ln b
  \]

  再取极限：
  \[\lim_{b \to +\infty} \ln b = +\infty\]
  
  极限不存在，所以积分发散。
\end{myexample}

接下来我们讨论被积函数在某点无界的情况。我们定义：

\begin{definition}{瑕点}{}
  若函数$f(x)$在$a$点的任一邻域内都无界，则称$a$为函数$f(x)$的\textbf{瑕点}。
\end{definition}

瑕点一定是间断点或函数未定义点，但间断点不一定是瑕点。

由此引出无界函数的反常积分：

\begin{definition}{无界函数的反常积分}{}
设 $f(x)$ 在 $(a, b]$ 上连续，点$a$ 为$f(x)$的瑕点，若极限
\[
\lim_{t \to a^+} \int_t^b f(x)\, dx
\]
存在，则称反常积分 $\int_a^b f(x)\, dx$ \textbf{收敛}，并定义$\int_a^b f(x)\, dx$的值为该极限。若极限不存在，则称该反常积分\textbf{发散}。
\end{definition}

类似地可以定义点$b$ 为$f(x)$的瑕点时函数 \(f(x)\) 在 \([a, b)\) 上的反常积分和点$c$ 为$f(x)$的瑕点时函数 \(f(x)\) 在 \([a, c)\) 及 \((c, b]\) 上的反常积分，这两个积分仍然记为 \(\int_a^b f(x) \, dx\)，后者的计算要拆成两段：
\[
\int_a^b f(x) \, dx = \int_a^c f(x) \, dx + \int_c^b f(x) \, dx
\]

无界函数的反常积分又称为\textbf{瑕积分}。

\begin{myexample}
  判断 $\int_0^1 \frac{1}{\sqrt{x}}\, dx$ 是否收敛。

  $x = 0$ 是瑕点。先算：
  \[
  \int_t^1 \frac{1}{\sqrt{x}}\, dx = [2\sqrt{x}]_t^1 = 2 - 2\sqrt{t}
  \]

  再取极限：
  \[
  \int_0^1 \frac{1}{\sqrt{x}}\, dx = \lim_{t \to 0^+} (2 - 2\sqrt{t}) = 2
  \]

  极限存在，收敛，值为 $2$。
\end{myexample}

\begin{myexample}
  判断 $\int_0^1 \frac{1}{x}\, dx$ 是否收敛。

  $x = 0$ 是瑕点。先算：
  \[
  \int_t^1 \frac{1}{x}\, dx = [\ln x]_t^1 = 0 - \ln t = -\ln t
  \]
  
  再取极限：
  \[
  \lim_{t \to 0^+} (-\ln t) = +\infty
  \]
  
  极限不存在，发散。
\end{myexample}

上面两个例子的瑕点都在左端点 $a$。如果瑕点在右端点 $b$，方法类似，只是极限方向不同。

\begin{myexample}
  判断 $\int_0^1 \frac{1}{\sqrt{1-x}}\, dx$ 是否收敛。

  $x = 1$ 是瑕点。按定义，先算：
  \[
  \int_0^t \frac{1}{\sqrt{1-x}}\, dx = [-2\sqrt{1-x}]_0^t = -2\sqrt{1-t} + 2
  \]

  再取极限：
  \[
  \int_0^1 \frac{1}{\sqrt{1-x}}\, dx = \lim_{t \to 1^-} (-2\sqrt{1-t} + 2) = 2
  \]

  极限存在，收敛，值为 $2$。
\end{myexample}

如果瑕点在区间内部，则需要拆成两段分别处理：

\begin{myexample}
  判断 $\int_{-1}^1 \frac{1}{x^2}\, dx$ 是否收敛。

  $x = 0$ 是瑕点，它在区间内部。拆成两段：
  \[
  \int_{-1}^1 \frac{1}{x^2}\, dx = \int_{-1}^0 \frac{1}{x^2}\, dx + \int_0^1 \frac{1}{x^2}\, dx
  \]

  分别计算。第一段，$x = 0$ 是右端瑕点：
  \[
  \int_{-1}^t \frac{1}{x^2}\, dx = [-\frac{1}{x}]_{-1}^t = -\frac{1}{t} - 1
  \]

  再取极限：
  \[\lim_{t \to 0^-} (-\frac{1}{t} - 1) = +\infty\]
  
  这个瑕积分第一段就发散了，既然有一段发散，那么整个积分就发散，不需要再算第二段。
\end{myexample}

总结一下反常积分的计算步骤：

如果是无穷区间上的积分，就把无穷上限换成有限上限 $b$，先算定积分，再令 $b \to +\infty$ 或 $b \to -\infty$ 取极限。极限存在就收敛，不存在就发散。

如果是瑕积分，瑕点在端点$a$，就把瑕点换成$t$算定积分，再令$t \to a^+$ 取极限。如果瑕点在中间，就拆成两段，分别处理。极限存在就收敛，不存在就发散。

本质上，反常积分就是“先积分，后取极限”。它把定积分的概念从有限推广到了无穷，因此所有不定积分技巧和定积分技巧都可以直接用在反常积分上。

\subsection{$\Gamma$函数}

最后介绍一个在数学中极其重要的特殊函数，$\Gamma$ 函数。它的定义是一个反常积分：

\begin{definition}{$\Gamma$函数}{}
\[
\Gamma(s) = \int_0^{+\infty} x^{s-1} e^{-x}\, dx, \quad s > 0
\]
\end{definition}

这个积分把 $s$ 从离散的自然数推广到了连续的正实数。它有一个非常漂亮的性质：

\begin{theorem}{$\Gamma$函数的递推公式}{}
\[
\Gamma(s+1) = s \cdot \Gamma(s)
\]
\end{theorem}

我们用分部积分法简要证明一下：

令 $u = x^s$，$dv = e^{-x}\,dx$，则 $du = sx^{s-1}\,dx$，$v = -e^{-x}$，则：
\[
\Gamma(s+1) = \int_0^{+\infty} x^s e^{-x}\, dx = [-x^s e^{-x}]_0^{+\infty} + s\int_0^{+\infty} x^{s-1} e^{-x}\, dx = 0 + s\,\Gamma(s)
\]

这个递推公式让我们可以算出所有正整数处的 $\Gamma$ 值。首先计算 $\Gamma(1)$：
\[
\Gamma(1) = \int_0^{+\infty} e^{-x}\, dx = [-e^{-x}]_0^{+\infty} = 0 - (-1) = 1
\]

然后由递推公式：
\[
\Gamma(2) = 1 \cdot \Gamma(1) = 1,\quad
\Gamma(3) = 2 \cdot \Gamma(2) = 2,\quad
\Gamma(4) = 3 \cdot \Gamma(3) = 6
\]

你发现它们有什么规律了吗？

一般地，对正整数 $n$有：
\[
\Gamma(n) = (n-1)!
\]

也就是说，$\Gamma$ 函数其实就是阶乘的连续版本。不过$\Gamma(n) = (n-1)!$差了一个 $1$ 的偏移。这个偏移是历史遗留问题，当初欧拉定义时就是这么写的，后来大家也就习惯了。

还有一个特殊值值得记住：
\[
\Gamma(\frac{1}{2}) = \int_0^{+\infty} x^{-\frac{1}{2}} e^{-x}\, dx = \sqrt{\pi}
\]

这个结果的证明需要用到换元 $x = t^2$ 和高斯积分 $\int_{-\infty}^{+\infty} e^{-t^2}\, dt = \sqrt{\pi}$，后者是概率论中最核心的积分之一。我们这里省略不证。

1729年，年轻的欧拉与哥德巴赫讨论一个问题：既然$n!$对所有正整数都有定义，那么对于更一般的正实数 $s$，是否也能定义一个类似的$s!$，将阶乘推广到所有正实数？后来，他在给哥德巴赫的信中，构造了 $\Gamma$ 函数的前身：

\[
\int_{0}^{1} (-\ln t)^n  \,dt 
\]

令$t = e^{-x}$，则 $dt = -e^{-x} dx$，当 $t = 0$ 时，$x \to +\infty$，当 $t = 1$ 时，$x = 0$。于是：

\[\int_{0}^{1} (-\ln t)^n  \,dt = \int_{+\infty}^{0} -x^n e^{-x} dx=\int_{0}^{+\infty} x^n e^{-x} dx\]

随后，欧拉用这个积分作为“广义阶乘”的定义。法国数学家勒让德在 1811 年首次引入了希腊字母 $\Gamma$ 来表示这个函数，并给出了现代标准定义。之后 $\Gamma$ 函数就沿用至今。

$\Gamma$ 函数在概率论、数论、物理学中都有广泛应用。它把离散的阶乘推广到了连续的实数域，将许多复杂的积分和无穷乘积转换为已知量，是连接离散数学和连续数学的一座桥梁。
% !TEX root = ../main.tex
\chapter{多元函数微分学}

一元函数微分学只处理 $y=f(x)$ 这种一个因变量对应一个自变量的情况。我们把视野扩展到多个自变量，进入多元函数微分学。多元函数微分学主要研究“一对多”和“多对多”关系下的求导问题，核心思想与一元情形一脉相承，只是多了一些方向和维度的考量。

\section{多元函数初步}
和一元函数一样，我们需要先定义多元函数和它的相关概念，才能更好地研究多元函数本身和它的相关性质。不过本书主要讨论二元函数，对更多元的函数不作展开。

\subsection{空间中的点集与多元函数}

我们先回忆一下第一章中邻域的概念：$x_0$ 的 $\delta$ 邻域是数轴上以 $x_0$ 为中心、左右各延伸 $\delta$ 的一个开区间。也就是说，邻域在一元函数的时候，是一个一维的区域，那么，推广到二元函数，就是二维区域，更进一步的，推广到我们的三维空间，就能得到：

\begin{definition}{三维空间的邻域}{}
设 $P_0(x_0, y_0, z_0)$ 是空间中的一点，$\delta > 0$。称集合
\[
U(P_0, \delta) = \{ P \mid |PP_0| < \delta \} = \{ (x,y,z) \mid (x-x_0)^2 + (y-y_0)^2 + (z-z_0)^2 < \delta^2 \}
\]
为点 $P_0$ 的 $\delta$ \textbf{邻域}。它是一个以 $P_0$ 为球心、$\delta$ 为半径的\textbf{开球}。
\end{definition}

开球的意思是这个球不包含球面：

\begin{center}
  \begin{tikzpicture}
    % 设置三维视角
    \begin{axis}[
      scale=1.5,
      scale only axis=true,
      axis lines=center,       % 坐标轴过原点
      axis equal image,        % 三轴等比例
      xlabel={$x$}, ylabel={$y$}, zlabel={$z$},
      xmin=-3.5, xmax=3.5,
      ymin=-2.5, ymax=2.5,
      zmin=-2.5, zmax=2.5,
      view={60}{30},           % 视角，可调整
      ticks=none,              % 不显示刻度
      enlargelimits=0.1,
      domain=0:360,
      y domain=0:180,
      samples=40,              % 网格密度
      samples y=20,
      surf,                    % 曲面着色
      faceted color=cyan!30,   % 网格线颜色
      fill opacity=0.6,        % 透明度
      colormap={white}{color(0cm)=(cyan!30); color(1cm)=(cyan!60)} % 颜色映射
      ]
      
      % 绘制球面：半径为2，用球坐标 (x=r sinθ cosφ, y=r sinθ sinφ, z=r cosθ)
      \addplot3[surf, shader=interp, draw=blue!30, dashed, very thin]
               ({2*sin(y)*cos(x)},
                {2*sin(y)*sin(x)},
                {2*cos(y)});
                
      % 标记球心
      \addplot3[mark=*, mark size=2pt, color=black] coordinates {(0,0,0)}
          node[below left] {$P_0$};
      
      % 绘制半径线段（从球心到(2,0,0)）
      \addplot3[red!70, thick, no markers] coordinates {(0,0,0) (2,0,0)}
          node[midway, above, sloped] {$\delta$};
          
    \end{axis}
  \end{tikzpicture}
\end{center}

接下来我们要定义点集。类似的，一元函数的定义域是数轴上的一个区间，二元函数的定义域是平面上的一个区域，三元函数的定义域就是空间中的一个点集：

\begin{definition}{点集}{}
  由点构成的集合称为点集，记为 $E$。
\end{definition}

有了邻域和点集，我们可以描述点和点集之间的关系：

\begin{definition}{点和点集之间的关系}{}
设 $E$ 是空间中的一个点集，$P$ 是空间中的一点，定义：

\textbf{内点}：若存在 $\delta > 0$，使得 $U(P, \delta) \subseteq E$，则称 $P$ 是 $E$ 的\textbf{内点}。

\textbf{外点}：若存在 $\delta > 0$，使得 $U(P, \delta) \cap E = \varnothing$，则称 $P$ 是 $E$ 的\textbf{外点}。

\textbf{边界点}：若在 $P$ 的任意邻域内，既有 $E$ 中的点，又有不在 $E$ 中的点，则称 $P$ 是 $E$ 的\textbf{边界点}。
\end{definition}

简单来说，内点是“完全在里面的点”，外点是“完全在外面的点”，边界点是“踩在边界上的点”。比如开圆盘 $x^2+y^2 < 1$ 的内点是所有严格在圆内的点，边界点是圆周 $x^2+y^2=1$ 上的点。

有了这些概念，我们可以定义开集和闭集：

\begin{definition}{开集和闭集}{}
\textbf{开集}：若点集 $E$ 中的每一个点都是内点，则称 $E$ 为\textbf{开集}。

\textbf{闭集}：若点集 $E$ 的补集是开集，则称 $E$ 为\textbf{闭集}。
\end{definition}

例如，开球 $U(P_0, \delta)$ 是开集，包含球面的球是闭集。区间 $(a,b)$ 是开集，$[a,b]$ 是闭集，$[a,b)$ 既不开也不闭。这些和一元情形完全类似。

最后，我们还需要“连通”的概念。多元函数的定义域通常要求是连通的。直观地说，就是“一整块”，通常不能分成互不相交的几块。

\begin{definition}{连通}{}
  若\textbf{不存在}两个不相交的开集 $A$ 和 $B$，使得点集 $E \subseteq A \cup B$ 且 $E \cap A \neq \varnothing$，$E \cap B \neq \varnothing$，则称点集 $E$ 是\textbf{连通的}。
\end{definition}

\begin{center}
  \begin{tikzpicture}[scale=1.5]
  
    % ========== 左：连通集 ==========
    \begin{scope}[local bounding box=left]
  
      % 画一个不规则但连通的区域（用光滑闭合曲线）
      \path[fill=green!20, draw=green!60!black, thick]
           plot[smooth cycle] coordinates {
               (0,0) (1.2,-0.6) (2.2,0.2) (2.8,1.0)
               (2.2,2.0) (1.0,2.6) (0.0,2.2) (-0.6,1.2)
           };
  
      % 区域内的两点 A 和 B
      \coordinate (A) at (0.5, 0.5);
      \coordinate (B) at (1.8, 1.8);
  
      \fill (A) circle[radius=1.5pt] node[below] {$A$};
      \fill (B) circle[radius=1.5pt] node[above] {$B$};
  
      % 红色实线表示能在区域内连接 A、B
      \draw[red, very thick, ->, >=stealth]
           (A) to[out=70, in=210] (B);
    \end{scope}
  
    % ========== 右：不连通集 ==========
    \begin{scope}[xshift=4.5cm, local bounding box=right]
  
      % 左侧圆盘
      \filldraw[fill=green!20, draw=green!60!black, thick]
           (0, 1) circle (1cm);
      % 右侧圆盘
      \filldraw[fill=green!20, draw=green!60!black, thick]
           (3, 1) circle (1cm);
  
      % A 在左边，B 在右边
      \coordinate (A2) at (0, 1);
      \coordinate (B2) at (3, 1);
  
      \fill (A2) circle[radius=1.5pt] node[below] {$A$};
      \fill (B2) circle[radius=1.5pt] node[below] {$B$};
  
      % 红色虚线表示无法连通，中间画叉
      \draw[red, very thick, dashed]
           (A2) -- (B2) node[midway, fill=white, inner sep=2pt, font=\bfseries] {$\times$};
    \end{scope}
  
  \end{tikzpicture}
\end{center}

也就是说，连通的点集可以用一条完全在点集内的连续曲线连接点集内的两点。不连通的点集可以分成两块或更多块，两点的连线可能不在点集内。例如圆环虽然有一个洞，但它是连通的，因为可以绕过洞连接两点。而两个不相交的圆盘就是不连通的。

由此我们可以定义我们之前一直在说的“区域”：

\begin{definition}{区域}{}
若点集 $D$ 是\textbf{连通的开集}，则称 $D$ 为一个\textbf{开区域}。开区域连同它的边界一起，称为\textbf{闭区域}。开区域和闭区域统称为\textbf{区域}。
\end{definition}

准备工作做好了，我们就可以定义多元函数了。和一元函数类似，多元函数就是一个“多对一”的规则：

\begin{definition}{二元函数}{}
设 $D$ 是平面上的一个点集。若对于 $D$ 中的每一个点 $P(x,y)$，按照某种确定的对应法则 $f$，都有唯一确定的实数 $z$ 与之对应，则称 $f$ 是定义在 $D$ 上的一个\textbf{二元函数}，记作
\[
z = f(x, y), \quad (x,y) \in D
\]
$D$ 称为函数的\textbf{定义域}，$x$ 和 $y$ 称为\textbf{自变量}，$z$ 称为\textbf{因变量}。
\end{definition}

推广到一般情形，$n$ 元函数就是 $z = f(x_1, x_2, \dots, x_n)$，但本书主要讨论二元函数。

二元函数的几何意义非常直观，$z = f(x,y)$ 在空间直角坐标系中表示一张\textbf{曲面}。例如 $z = x^2 + y^2$ 是一个旋转抛物面，$z = \sqrt{1-x^2-y^2}$ 是上半球面。每给定一个 $(x,y)$，函数值 $z$ 就是曲面在该点的高度。由无数个 $(x,y)$ 对应的 $z$ 值，形成无数个三维点$(x,y,z)$，无数多个三维点就在空间中组成一张曲面。这点和一元函数也是一样的，由无数个 $x$ 对应的 $y$ 值，形成无数个二维点$(x,y)$，无数多个二维点就在平面中组成一条曲线。

我们以一个具体的函数作为例子：
\begin{myexample}
  求二元函数 $z = \ln(x+y) + \sqrt{4-x^2-y^2}$ 的定义域。

  要使函数有意义，也就是说：

  $\ln(x+y)$ 要求 $x + y > 0$，$\sqrt{4-x^2-y^2}$ 要求 $x^2 + y^2 \leq 4$。

  因此定义域为：
  \[
  D = \{ (x,y) \mid x + y > 0,\; x^2 + y^2 \leq 4 \}
  \]
\end{myexample}
  
  画出来就是半径为 $2$ 的闭圆盘，被直线 $x+y=0$ 切去一半，且只保留 $x+y>0$ 的那一半。这是一个有界区域。

\begin{center}
\begin{tikzpicture}[scale=1.2]
  
  % 坐标轴
  \draw[->] (-2.5,0) -- (2.5,0) node[below] {$x$};
  \draw[->] (0,-2.5) -- (0,2.5) node[left] {$y$};
  \node[below left] at (0,0) {$O$};

  % 填充区域：只填充不描边
  \path[fill=blue!20, opacity=0.6]
       (-45:2) arc[start angle=-45, end angle=135, radius=2]
       -- cycle;

  % 绘制圆弧边界（实线，表示包含边界）
  \draw[blue!60, thick]
       (-45:2) arc[start angle=-45, end angle=135, radius=2];

  % 绘制直线边界（虚线，表示 x+y>0 不包含边界）
  \draw[blue!60, thick, dashed]
       (-45:2) -- (135:2);

  % 整条直线 x+y=0（灰色虚线，作参照）
  \draw[gray, dashed, thin] (-2,2) -- (2,-2) node[below right] {$x+y=0$};

  % 标记两个交点
  \filldraw (135:2) circle[radius=0.03] node[left] {$(-\sqrt2,\sqrt2)$};
  \filldraw (-45:2) circle[radius=0.03] node[right] {$(\sqrt2,-\sqrt2)$};

\end{tikzpicture}
\end{center}

二元函数和一元函数的一个重要区别在于，一元函数的定义域是数轴上的区间，只有左右两个方向。二元函数的定义域是平面上的区域，可以从无穷多个方向趋近于同一点。正是这个区别，使得多元函数的极限比一元情形复杂得多。


\subsection{多元函数的极限}

在一元函数中，极限 $\lim_{x \to x_0} f(x) = L$ 的意思是，无论 $x$ 从左边还是右边趋近于 $x_0$，$f(x)$ 都趋近于 $L$。只有两条路可走，要么左要么右。

到了二元函数，情况就变了。$(x,y)$ 趋近于 $(x_0, y_0)$ 时，可以有无数条路径。因为平面上有无数条曲线可以从不同方向趋近于同一点。为了保证极限存在，必须保证沿着所有路径趋近时，函数值都趋近于同一个数。

\begin{definition}{二元函数的极限}{}
设函数 $z = f(x, y)$ 在 $(x_0, y_0)$ 的某个去心邻域内有定义，$L$ 为实数。若对于任意给定的 $\epsilon > 0$，总存在 $\delta > 0$，使得当
\[
0 < \sqrt{(x - x_0)^2 + (y - y_0)^2} < \delta
\]
时，就有
\[
|f(x, y) - L| < \epsilon
\]
则称当 $(x, y)$ 趋近于 $(x_0, y_0)$ 时，$f(x, y)$ 的\textbf{极限}为 $L$，记作
\[
\lim_{(x,y) \to (x_0, y_0)} f(x, y) = L
\]
\end{definition}

有没有发现这个定义和第二章的 $\epsilon-\delta$ 定义几乎一模一样？唯一的区别是 $|x - x_0|$ 被换成了两点间的距离 $\sqrt{(x - x_0)^2 + (y - y_0)^2}$。所以核心思想完全一致：\textbf{无论我把误差 $\epsilon$ 定得多小，你都能找到一个以 $(x_0, y_0)$ 为中心，半径为 $\delta$的小圆盘，只要 $(x, y)$ 进入这个圆盘，$f(x, y)$ 就一定落在 $L$ 的 $\epsilon$ 邻域内}。

但这个定义在实际计算中有一层隐藏的严格性，它要求极限值与趋近路径无关。也就是说，无论你沿着哪条曲线趋近于 $(x_0, y_0)$，得到的结果必须一样。

这就给了我们一个判断极限不存在的工具，\textbf{找两条不同的路径，如果极限值不同，就说明极限不存在}。

\begin{myexample}
  讨论 $f(x, y) = \dfrac{xy}{x^2 + y^2}$ 在 $(0, 0)$ 处的极限。

  沿着 $y = kx$趋近于原点，那么代入$y = kx$，我们有：
  \[
  f(x, kx) = \frac{x \cdot kx}{x^2 + (kx)^2} = \frac{kx^2}{x^2(1 + k^2)} = \frac{k}{1 + k^2}
  \]

  当 $(x, y) \to (0, 0)$ 时，这个值等于 $\frac{k}{1+k^2}$，它依赖 $k$ 的取值。因此极限不存在。
\end{myexample}

怎么理解呢？我们沿着 $y = x$得到极限值为 $\frac{1}{2}$，沿着 $y = 0$得到极限值为$0$。不同路径得到不同结果，就像一元函数里左极限不等于右极限一样，所以极限不存在。

再看一个例子：

\begin{myexample}
  讨论 $f(x, y) = \dfrac{x^2 y}{x^4 + y^2}$ 在 $(0, 0)$ 处的极限。

  沿着 $y = kx$ 趋近原点，那么代入$y = kx$，得：
  \[
  f(x, kx) = \frac{x^2 \cdot kx}{x^4 + (kx)^2} = \frac{kx^3}{x^4 + k^2 x^2} = \frac{kx}{x^2 + k^2}
  \]

  \[\lim_{x \to 0} \frac{kx}{x^2 + k^2} = \frac{0}{k^2} = 0\]

  沿着任意一条直线，极限都是 $0$。但这能说明极限就是 $0$ 吗？让我们试试抛物线 $y = x^2$，那么代入$y = x^2$，得：
  \[
  f(x, x^2) = \frac{x^2 \cdot x^2}{x^4 + (x^2)^2} = \frac{x^4}{x^4 + x^4} = \frac{1}{2}
  \]

  沿着抛物线取到的极限是 $\frac{1}{2}$，而直线是 $0$。所以极限还是不存在。
\end{myexample}

这个例子说明，沿着所有直线极限都相同，不意味着极限存在，因为还有曲线路径。所以在二元函数中，要证明极限存在，必须用定义严格证明。要证明极限不存在，找到两条不同路径得到不同结果即可。

和一元函数一样，二元函数的极限也有类似的运算律。若两个函数的极限都存在，则函数四则运算的极限等于极限的四则运算结果。证明方法也和一元完全类似，此处不再重复。

\subsection{多元函数的连续性}

和一元情形一样，连续性的定义就是把极限中的 $L$ 换成函数值 $f(x_0, y_0)$，并去掉“去心”条件。

\begin{definition}{二元函数的连续性}{}
设函数 $z = f(x, y)$ 在 $(x_0, y_0)$ 的某个邻域内有定义。若
\[
\lim_{(x,y) \to (x_0, y_0)} f(x, y) = f(x_0, y_0)
\]
则称 $f(x, y)$ 在点 $(x_0, y_0)$ 处\textbf{连续}。如果函数在区域 $D$ 上的每一点都连续，就称 $f(x, y)$ 在 $D$ 上连续。
\end{definition}

这个定义拆开来看就是三个条件，$f(x_0, y_0)$ 有定义，极限 $\lim_{(x,y) \to (x_0, y_0)} f(x, y)$ 存在，极限等于函数值。三者缺一不可。这也和一元函数极限与连续的关系完全一样。因此之前在一元函数成立的定理在二元函数仍然适用，包括：

\begin{theorem}{多元初等函数的连续性}{}
  多元初等函数在其定义域内都是连续的。
\end{theorem}

这里的多元初等函数和一元初等函数类似，由一个解析表达式定义，包含多项式、有理函数、根式、指数对数、三角函数、反三角函数等，以及它们经有限次四则运算和复合得出的函数。

\begin{theorem}{连续函数的运算性质}{}
若 $f(x, y)$ 和 $g(x, y)$ 都在 $(x_0, y_0)$ 处连续，则
\[
f \pm g,\quad f \cdot g,\quad \frac{f}{g}\ (g \neq 0)
\]
都在 $(x_0, y_0)$ 处连续。
\end{theorem}


\begin{theorem}{有界性定理}{}
若 $f(x, y)$ 在有界闭区域 $D$ 上连续，则$f(x, y)$ 在 $D$ 上有界。
\end{theorem}

\begin{theorem}{最值定理}{}
  若 $f(x, y)$ 在有界闭区域 $D$ 上连续，则$f(x, y)$ 在 $D$ 上能取到最大值和最小值。
\end{theorem}

\begin{theorem}{介值定理}{}
若 $f(x, y)$ 在区域 $D$ 上连续，$P_1, P_2 \in D$ 且 $f(P_1) < f(P_2)$，则对于 $f(P_1)$ 与 $f(P_2)$ 之间的任意一个数 $\mu$，都存在至少一个点 $P_\xi \in D$，使得 $f(P_\xi) = \mu$。
\end{theorem}

这些性质和一元闭区间上的连续性性质如出一辙，核心条件都是“有界闭区域”加“连续”。反过来，如果区域不是闭的或者函数不连续，这些结论就不保证成立。这些定理本质上和一元函数完全一致，只是从一维推广到二维。从闭区间推广到有界闭区域。

\section{偏导数}
在一元函数中，导数 $f'(x)$ 描述了函数值随自变量变化的快慢。对于二元函数 $z = f(x,y)$，我们自然也想考察函数值随自变量变化的情况。但是自变量有两个，怎么处理呢？

\subsection{偏导数及其几何意义}
面对两个自变量，核心思想是先固定一个，再看另一个。因为二元函数 $z = f(x,y)$的自变量在平面上可以沿无数个方向变化，但我们没法一次性把所有这些方向都研究完。所以，我们只能分别看沿坐标轴方向的变化。于是，我们就引出了偏导数：

\begin{definition}{偏导数}{}
设函数 $z = f(x,y)$ 在点 $(x_0, y_0)$ 的某邻域内有定义。

若极限
\[
\lim_{\Delta x \to 0} \frac{f(x_0 + \Delta x, y_0) - f(x_0, y_0)}{\Delta x}
\]
存在，则称此极限为 $f(x,y)$ 在点 $(x_0, y_0)$ 处\textbf{对 $x$ 的偏导数}，记作
\[
\frac{\partial z}{\partial x}\Big|_{(x_0,y_0)},\quad \frac{\partial f}{\partial x}\Big|_{(x_0,y_0)},\quad f_x(x_0,y_0),\quad \text{或}\quad z_x(x_0,y_0)
\]

若极限
\[
\lim_{\Delta y \to 0} \frac{f(x_0, y_0 + \Delta y) - f(x_0, y_0)}{\Delta y}
\]
存在，则称此极限为 $f(x,y)$ 在点 $(x_0, y_0)$ 处\textbf{对 $y$ 的偏导数}，记作
\[
\frac{\partial z}{\partial y}\Big|_{(x_0,y_0)},\quad \frac{\partial f}{\partial y}\Big|_{(x_0,y_0)},\quad f_y(x_0,y_0),\quad \text{或}\quad z_y(x_0,y_0)
\]
\end{definition}

这个定义和一元导数的定义长得几乎一模一样，唯一的区别是求对 $x$ 的偏导数时，把 $y$ 固定为 $y_0$，只让 $x$ 变化。求对 $y$ 的偏导数时，把 $x$ 固定为 $x_0$，只让 $y$ 变化。换句话说，\textbf{偏导数就是把其他变量视作常数后，对该变量求一元导数}。

因此，偏导数的计算法则其实非常直接。对哪个变量求导，就把其余变量当作常数，然后利用一元函数的所有求导公式和法则。

来看一个例子：

\begin{myexample}
设 $z = x^3y + 2x^2 + e^{xy}$，求 $\frac{\partial z}{\partial x}$ 和 $\frac{\partial z}{\partial y}$。

对 $x$ 求偏导时，把 $y$ 当作常数：
\[
\frac{\partial z}{\partial x} = 3x^2 \cdot y + 4x + e^{xy} \cdot y = 3x^2y + 4x + y e^{xy}
\]

对 $y$ 求偏导时，把 $x$ 当作常数：
\[
\frac{\partial z}{\partial y} = x^3 \cdot 1 + 0 + e^{xy} \cdot x = x^3 + x e^{xy}
\]
\end{myexample}

是不是很简单？所谓偏导数，本质上就是“假装”其他变量是常数，然后用一元导数的方法去算。

接下来看几何意义。在一元函数中，导数 $f'(x_0)$ 是曲线 $y = f(x)$ 在点 $(x_0, f(x_0))$ 处切线的斜率。二元函数 $z = f(x,y)$ 的图像是空间中的一张曲面，偏导数 $f_x(x_0,y_0)$ 和 $f_y(x_0,y_0)$ 也对应着切线的斜率，只不过需要用平面去截曲面，再退化成二维曲线来看。

固定 $y = y_0$，平面 $y = y_0$ 与曲面 $z = f(x,y)$ 相交，得到一条曲线
\[
C_x:\; z = f(x, y_0), \quad y = y_0
\]

这条曲线可以看作关于 $x$ 的一元函数。它在点 $P_0(x_0, y_0, f(x_0, y_0))$ 处的切线对 $x$ 轴的斜率，正好就是偏导数 $f_x(x_0, y_0)$。同样，固定 $x = x_0$，平面 $x = x_0$ 与曲面相交得到曲线
\[
C_y:\; z = f(x_0, y), \quad x = x_0
\]

该曲线在 $P_0$ 处的切线对 $y$ 轴的斜率就是偏导数 $f_y(x_0, y_0)$。

简而言之，\textbf{偏导数 $f_x$ 衡量的是曲面沿 $x$ 轴方向的变化趋势，$f_y$ 衡量的是曲面沿 $y$ 轴方向的变化趋势。} 曲面被坐标面“切开”后，我们看到的正是两个一元函数，因此偏导数的几何意义与一元导数其实是一脉相承的。

\begin{center}
  \begin{tikzpicture}
  \begin{axis}[
      scale=1.3,
      view={65}{35},
      xmin=-2.5, xmax=2.5,
      ymin=-2.5, ymax=2.5,
      zmin=0, zmax=3.5,
      xlabel=$x$, ylabel=$y$, zlabel=$z$,
      zlabel style={rotate=90},
      axis lines=left,
      colormap/cool,
      clip=false
  ]
  
  % 绘制曲面 z = f(x,y) = 2 - 0.2(x^2 + y^2)
  \addplot3[
      surf,
      shader=interp, 
      domain=-2:2, 
      y domain=-2:2, 
      opacity=0.3, 
      faceted color=black!30
  ] {2 - 0.2*(x^2 + y^2)};
  
  % 用 x = 1 的平面“切开”曲面（对应 y 方向的偏导）
  \addplot3[
      surf, 
      mesh/rows=2, 
      shader=flat, 
      draw=blue!60, 
      fill=blue!10, 
      opacity=0.4
  ] coordinates {(1,-2.5,0) (1,2.5,0) (1,-2.5,2.5) (1,2.5,2.5)};
  
  % 绘制“切痕”曲线 C_y (平面 x=1 与曲面的交线)
  \addplot3[
      domain=-2:2, 
      samples=50, 
      thick, 
      blue!80!black
  ] ({1}, {y}, {2 - 0.2*(1^2 + y^2)});
  
  % 绘制切线 T_y (斜率 = f_y)
  \addplot3[
      domain=-1.5:1.5, 
      samples=2, 
      dashed, 
      thick, 
      red!80!black
  ] ({1}, {1.5+x}, {1.35 - 0.6*x});
  
  % 标记切点 P(x0, y0)
  \addplot3[only marks, mark=*, mark size=2, black] coordinates {(1,1.5,1.35)};
  
  \node[label={[label distance=0.2cm]135:{$P(x_0,y_0)$}}] at (1,2.1,1.15) {};
  
  \node[label={$x=1$}] at (1,-3,1.35) {};
  
  \end{axis}
  \end{tikzpicture}
\end{center}

上图的例子就很好展现了偏导数的几何意义。曲面被平面 $x = 1$ 截出一条蓝色曲线，在点 $P(x_0,y_0)$ 处，该曲线沿 $y$ 轴的红色切线斜率就是 $\frac{\partial z}{\partial y}$。

有了偏导数的几何直观，我们就可以研究二元函数沿坐标轴方向的变化率了。不过要注意，这只是两个特殊方向，沿着$x$和沿着$y$的方向，并不能代表所有方向的变化率。偏导数只反映沿坐标轴方向的信息，无法控制其他路径的行为。这暗示了多元函数的偏导存在和连续性之间的关系比一元函数复杂得多。

回想我们之前讨论过的函数：

\begin{myexample}
  函数\(f(x,y) = \frac{xy}{x^2 + y^2}\)补充定义 $f(0,0)=0$，讨论$f(x,y)$的连续性。

补充定义 $f(0,0)=0$后函数表述为：
\[
f(x,y) = \begin{cases}
\frac{xy}{x^2 + y^2}, & (x,y) \neq (0,0) \\
0, & (x,y) = (0,0)
\end{cases}
\]

根据偏导数定义：
\[
f_x(0,0) = \lim_{\Delta x\to0}\frac{f(\Delta x,0)-f(0,0)}{\Delta x} = 0
\]
\[
f_y(0,0) = \lim_{\Delta y\to0}\frac{f(0,\Delta y)-f(0,0)}{\Delta y} = 0
\]

两个偏导数都存在且为 $0$。但我们已经证明过，$(x,y)\to(0,0)$ 时该函数的极限根本不存在，自然也就不连续。
\end{myexample}

这个反例说明，在多元函数中，\textbf{偏导数存在并不能保证函数连续}。这一点和一元函数截然不同，需要特别记忆。

\subsection{高阶偏导数}

偏导数 $f_x(x,y)$ 和 $f_y(x,y)$ 本身仍然是 $x,y$ 的二元函数。如果它们对 $x$ 或 $y$ 再次可求偏导，我们就得到\textbf{二阶偏导数}。依此类推，还可以定义三阶、四阶乃至更高阶的偏导数。

\begin{definition}{高阶偏导数的记号}{}
设函数 $z = f(x, y)$ 具有任意阶偏导数，其高阶偏导数用以下记号表示：

拉格朗日记号：
\[
f_x,\ f_y,\quad 
f_{xx},\ f_{xy},\ f_{yx},\ f_{yy},\quad 
f_{xxx},\ f_{xxy},\ f_{xyx},\ \dots,\ 
f_{\underbrace{x\cdots x}_{p}\,\underbrace{y\cdots y}_{q}}
\]

莱布尼茨记号：
\[
\frac{\partial z}{\partial x},\ \frac{\partial z}{\partial y},\quad 
\frac{\partial^2 z}{\partial x^2},\ \frac{\partial^2 z}{\partial x\partial y},\ \frac{\partial^2 z}{\partial y\partial x},\ \frac{\partial^2 z}{\partial y^2},\quad \frac{\partial^3 z}{\partial x^3},\ \frac{\partial^3 z}{\partial x^2\partial y},\ \frac{\partial^3 z}{\partial x\partial y\partial x},\ \dots,\ \frac{\partial^{p+q} z}{\partial x^p\,\partial y^q}
\]
其中 $f_{xy}$ 、 $f_{yx}$ 、 $f_{xyx}$ 等称为\textbf{混合偏导数}。
\end{definition}

混合偏导数有求导次序的区别，$f_{xy}$ 表示先对 $x$ 再对 $y$ 求偏导，$f_{yx}$ 表示先对 $y$ 再对 $x$ 求偏导。也就是谁放前面就对谁先求偏导。

有时我们也用$f_1$，$f_2$ 等来表示一阶偏导数，用 $f_{11}$，$f_{12}$ 等来表示二阶偏导数。$f_1$表示对第一个变量求偏导，$f_{12}$表示先对第一个变量求偏导再对第二个变量求偏导。$f_{11}$，$f_{12}$以此类推。这种记法在抽象函数里很常见。

计算高阶偏导数并没有新技巧，就是老老实实一次又一次地用一元求导法则，记住每次只对一个变量求导，其他变量视为常数。看一个例子：

\begin{myexample}
设 $z = x^3y^2 + \sin xy$，求所有二阶偏导数。

先算一阶偏导：
\[
\frac{\partial z}{\partial x} = 3x^2y^2 + y\cos xy,\qquad
\frac{\partial z}{\partial y} = 2x^3y + x\cos xy
\]

再对每个一阶偏导继续求偏导：
\begin{align*}
\frac{\partial^2 z}{\partial x^2}
&= \frac{\partial}{\partial x}(3x^2y^2 + y\cos xy)
= 6xy^2 - y^2\sin xy \\[4pt]
\frac{\partial^2 z}{\partial x\partial y}
&= \frac{\partial}{\partial y}(3x^2y^2 + y\cos xy)
= 6x^2y + \cos xy - xy\sin xy \\[4pt]
\frac{\partial^2 z}{\partial y\partial x}
&= \frac{\partial}{\partial x}(2x^3y + x\cos xy)
= 6x^2y + \cos xy - xy\sin xy \\[4pt]
\frac{\partial^2 z}{\partial y^2}
&= \frac{\partial}{\partial y}(2x^3y + x\cos xy)
= 2x^3 - x^2\sin xy
\end{align*}
\end{myexample}

仔细观察结果，我们发现 $f_{xy}$ 和 $f_{yx}$ 结果一样。这说明，在某种条件下，混合偏导数的值与求导次序无关：

\begin{theorem}{克莱罗定理}{}
若函数 $z = f(x,y)$ 的两个混合偏导数 $f_{xy}$ 和 $f_{yx}$ 在区域 $D$ 内连续，则它们在 $D$ 内处处相等：
\[
f_{xy}(x,y) = f_{yx}(x,y)
\]
\end{theorem}

这个定理的意义在于，当我们看到对称出现的混合偏导时，只要它们连续，就可以放心地交换求导次序，不必担心结果不同。上面的例题正是由于 $\sin xy$ 无穷次可微，混合偏导数连续，自然满足这个定理。

对于更高阶的偏导数，比如 $f_{xxy}$，只要涉及的混合偏导数都连续，求导次序就可以任意排列。比如 $f_{xxy}$、$f_{xyx}$、$f_{yxx}$ 在连续条件下都相等。这使得我们在计算时可以灵活选择最方便的次序，往往能大幅简化运算。

高阶偏导数的几何意义不如一阶那样直观。一阶偏导数衡量曲面沿坐标轴的倾斜程度，二阶偏导数 $f_{xx}$ 和 $f_{yy}$ 则衡量这种倾斜本身变化的快慢，即曲面的\textbf{弯曲程度}。混合偏导数 $f_{xy}$ 反映了曲面在 $x$ 方向变化时，$y$ 方向的斜率如何变化，这可以理解为一种“扭转”。



\subsection{多元复合函数求偏导数}

在一元函数中，我们有链式法则。到了多元函数，复合的方式更加多样，自变量和中间变量都可以有多个，链式法则依然存在，只是要把所有“影响路径”都考虑进来。

我们从一个最典型的情形入手。设 $z = f(u,v)$ 是 $u,v$ 的函数，而 $u = u(x,y)$，$v = v(x,y)$ 又都是 $x,y$ 的函数。把它们复合起来，就得到 $z$ 关于 $x,y$ 的二元函数：
\[
z = f(u(x,y),\, v(x,y))
\]

现在要怎么求 $\frac{\partial z}{\partial x}$ 和 $\frac{\partial z}{\partial y}$呢？

直观上，$x$ 发生微小变化时，会同时引起 $u$ 和 $v$ 的变化，而这两个变化又会“传递”到 $z$。因此 $z$ 对 $x$ 的变化率应该是两条路径上变化率的叠加。把这种关系画出来，会得到一张“变量依赖关系图”：

\begin{center}
\begin{tikzpicture}[scale=1.1,
    >=stealth,
    var/.style={draw, circle, minimum size=6mm, inner sep=1pt},
    edge/.style={->, thick}]

  \node[var] (z) at (0, 0)    {$z$};
  \node[var] (u) at (2, 0.8)  {$u$};
  \node[var] (v) at (2,-0.8)  {$v$};
  \node[var] (x) at (4, 0.8)  {$x$};
  \node[var] (y) at (4,-0.8)  {$y$};

  \draw[edge] (z) -- (u);
  \draw[edge] (z) -- (v);
  \draw[edge] (u) -- (x);
  \draw[edge] (u) -- (y);
  \draw[edge] (v) -- (x);
  \draw[edge] (v) -- (y);
\end{tikzpicture}
\end{center}

也就是说，函数$z$ 依赖于中间变量$u$ 和 $v$，而$u$ 和 $v$又分别依赖于自变量 $x$和$y$。所以每一条从 $z$ 到 $x$ 的路径都对应着一种影响，例如$z \to u \to x$ 的影响是 $\frac{\partial z}{\partial u}\frac{\partial u}{\partial x}$，$z \to v \to x$ 的影响是 $\frac{\partial z}{\partial v}\frac{\partial v}{\partial x}$。所有这些影响加起来，就是 $z$ 对 $x$ 的总变化率。这个规律总结起来，就是多元复合函数求导的链式法则：

\begin{theorem}{多元函数链式法则}{}
设 $z = f(u,v)$ 在对应点 $(u,v)$ 处具有连续偏导数，$u = u(x,y)$，$v = v(x,y)$ 在点 $(x,y)$ 处具有连续偏导数，则复合函数 $z = f(u(x,y), v(x,y))$ 在 $(x,y)$ 处的偏导数存在，且：
\[
\frac{\partial z}{\partial x} = \frac{\partial z}{\partial u}\frac{\partial u}{\partial x} + \frac{\partial z}{\partial v}\frac{\partial v}{\partial x}
\]
\[
\frac{\partial z}{\partial y} = \frac{\partial z}{\partial u}\frac{\partial u}{\partial y} + \frac{\partial z}{\partial v}\frac{\partial v}{\partial y}
\]
\end{theorem}

这个公式说的是，\textbf{因变量对某个自变量的偏导数，等于所有从因变量通往该自变量的路径上偏导数乘积的总和。} 有多少条路径，就有多少项相加，每条路径经过几层，就乘上几个偏导数。

有了它，我们就能处理各种复杂的复合情形。

下面看几个例子：

\begin{myexample}
  设函数$z=u^2+v^2$，其中 $u = \sin t$，$v = \cos t$，求 $\frac{dz}{dt}$。

  先画出路径关系图：

\begin{center}
\begin{tikzpicture}[scale=1.2,
    >=stealth,
    var/.style={draw, circle, minimum size=7mm, inner sep=0pt},
    edge/.style={->, thick}]

  \node[var] (z) at (0,0)   {$z$};
  \node[var] (u) at (2,0.7) {$u$};
  \node[var] (v) at (2,-0.7){$v$};
  \node[var] (t) at (4,0)   {$t$};

  \draw[edge] (z) -- (u);
  \draw[edge] (z) -- (v);
  \draw[edge] (u) -- (t);
  \draw[edge] (v) -- (t);
\end{tikzpicture}
\end{center}

我们可以清晰地看到，$z$ 对 $t$ 的变化率有两条路径：$z \to u \to t$ 和 $z \to v \to t$。因此：
\[
\frac{dz}{dt} = \frac{\partial z}{\partial u}\frac{du}{dt} + \frac{\partial z}{\partial v}\frac{dv}{dt}
\]

所以：
\[
\frac{dz}{dt} = \frac{\partial z}{\partial u}\cos t - \frac{\partial z}{\partial v}\sin t
=2u\cos t-2v\sin t=2\sin t \cos t-2\cos t \sin t=0
\]
\end{myexample}

由这个例子我们得到，若 $z = f(u,v)$，而 $u = u(t)$，$v = v(t)$ 都只是一元函数，那么复合后 $z$ 就是 $t$ 的一元函数。此时对 $t$ 求导就是“全导数”：
\[
\frac{dz}{dt} = \frac{\partial z}{\partial u}\frac{du}{dt} + \frac{\partial z}{\partial v}\frac{dv}{dt}
\]

这是一元函数链式法则在多元中的直接推广。

不过，不是所有复合函数都要这样算，我们可以把 $u = \sin t$和$v = \cos t$代入到$z=u^2+v^2$里，化简后我们会发现$z=1$，因此导数自然是$0$了。


\begin{myexample}
  设函数$z=u^2+v^2$，其中 $u = x-y$，$v = x+y$，求 $\frac{\partial z}{\partial x}$ 和 $\frac{\partial z}{\partial y}$。

分析函数，$u$和$v$都是$x$和$y$的函数，所以$z$是$x$和$y$的函数。画出路径关系图：

\begin{center}
\begin{tikzpicture}[scale=1.2,
    >=stealth,
    var/.style={draw, circle, minimum size=7mm, inner sep=0pt},
    edge/.style={->, thick}]
    
  \node[var] (z) at (0,0)   {$z$};
  \node[var] (u) at (2,0.7) {$u$};
  \node[var] (v) at (2,-0.7){$v$};
  \node[var] (x) at (4,0.7) {$x$};
  \node[var] (y) at (4,-0.7){$y$};

  \draw[edge] (z) -- (u);
  \draw[edge] (z) -- (v);
  \draw[edge] (u) -- (x);
  \draw[edge] (u) -- (y);
  \draw[edge] (v) -- (x);
  \draw[edge] (v) -- (y);
\end{tikzpicture}
\end{center}

从图中可以看到，$z$ 对 $x$ 有两条路径：$z \to u \to x$ 和 $z \to v \to x$，对 $y$ 同样有两条路径：$z \to u \to y$ 和 $z \to v \to y$。因此：
  \[
  \frac{\partial z}{\partial x} = \frac{\partial z}{\partial u}\frac{\partial u}{\partial x}
    + \frac{\partial z}{\partial v}\frac{\partial v}{\partial x}
  \]
  \[
  \frac{\partial z}{\partial y} = \frac{\partial z}{\partial u}\frac{\partial u}{\partial y}
    + \frac{\partial z}{\partial v}\frac{\partial v}{\partial y}
  \]

  计算偏导数：
  \[
  \frac{\partial z}{\partial u} = 2u \quad \frac{\partial z}{\partial v} = 2v
  \quad
  \frac{\partial u}{\partial x} = 1 \quad \frac{\partial u}{\partial y} = -1
  \quad
  \frac{\partial v}{\partial x} = 1 \quad \frac{\partial v}{\partial y} = 1
  \]

  代入得：
  \begin{align*}
    \frac{\partial z}{\partial x}
    &= 2u \cdot 1 + 2v \cdot 1 = 2(u+v) \\
    &= 2[(x-y)+(x+y)] = 4x\\[4pt]
    \frac{\partial z}{\partial y}
    &= 2u \cdot (-1) + 2v \cdot 1 = 2(v-u) \\
    &= 2[(x+y)-(x-y)] = 4y
  \end{align*}

\end{myexample}

由这个例子我们发现，若 $z = f(u,v)$，而 $u = u(x,y)$，$v = v(x,y)$ 都是二元函数，那么复合后 $z$ 就是 $x,y$ 的二元函数。这时候就要用多元复合函数链式法则，画出路径关系图，找出所有从 $z$ 到自变量的路径，然后把每条路径上偏导数的乘积加起来。

\begin{myexample}
  设函数 $z = f(u+v, u-v)$，其中 $u = x$，$v = xy$，求 $\frac{\partial z}{\partial x}$ 和 $\frac{\partial z}{\partial y}$。

  分析函数，$z$ 通过 $u$ 和 $v$ 依赖于 $x,y$，其中 $u$ 仅依赖于 $x$，是$x$的一元函数，$v$ 依赖于 $x$ 和 $y$，是$x,y$的二元函数。画出路径关系图：
  \begin{center}
  \begin{tikzpicture}[scale=1.2,
      >=stealth,
      var/.style={draw, circle, minimum size=7mm, inner sep=0pt},
      edge/.style={->, thick}]

    \node[var] (z) at (0,0)   {$z$};
    \node[var] (u) at (2,0.7) {$u$};
    \node[var] (v) at (2,-0.7){$v$};
    \node[var] (x) at (4,0.7) {$x$};
    \node[var] (y) at (4,-0.7){$y$};

    \draw[edge] (z) -- (u);
    \draw[edge] (z) -- (v);
    \draw[edge] (u) -- (x);
    \draw[edge] (v) -- (x);
    \draw[edge] (v) -- (y);
  \end{tikzpicture}
  \end{center}

  从图中可以看出，$z$ 对 $x$ 有两条路径：$z \to u \to x$ 和 $z \to v \to x$，对 $y$ 有一条路径：$z \to v \to y$，因此，由链式法则：
  \[
  \frac{\partial z}{\partial x} = \frac{\partial z}{\partial u}\frac{\partial u}{\partial x}
    + \frac{\partial z}{\partial v}\frac{\partial v}{\partial x}
  \]
  \[
  \frac{\partial z}{\partial y} = \frac{\partial z}{\partial v}\frac{\partial v}{\partial y}
  \]

  令 $p = u+v$, $q = u-v$，则 $z = f(p,q)$。由复合函数求导法则：
  \[
  \frac{\partial z}{\partial u} = \frac{\partial f}{\partial p}\frac{\partial p}{\partial u}
    + \frac{\partial f}{\partial q}\frac{\partial q}{\partial u}
    = f_1 \cdot 1 + f_2 \cdot 1 = f_1 + f_2
  \]
  \[
  \frac{\partial z}{\partial v} = \frac{\partial f}{\partial p}\frac{\partial p}{\partial v}
    + \frac{\partial f}{\partial q}\frac{\partial q}{\partial v}
    = f_1 \cdot 1 + f_2 \cdot (-1) = f_1 - f_2
  \]

  又由 $u = x$, $v = xy$ 得：
  \[
  \frac{\partial u}{\partial x} = 1,\quad \frac{\partial u}{\partial y} = 0,\quad
  \frac{\partial v}{\partial x} = y,\quad \frac{\partial v}{\partial y} = x
  \]

  代入链式法则公式：
  \begin{align*}
    \frac{\partial z}{\partial x}
    &= (f_1 + f_2) \cdot 1 + (f_1 - f_2) \cdot y \\
    &= f_1 + f_2 + y f_1 - y f_2 \\
    &= (1+y)f_1 + (1-y)f_2 \\[6pt]
    \frac{\partial z}{\partial y}
    &= (f_1 - f_2) \cdot x = x(f_1 - f_2)
  \end{align*}

  因为我们没有 $f$ 的显式，最终结果只能保留$f_1$和$f_2$，这也是常见的处理方式。
\end{myexample}

这个例子其实是三层复合抽象函数求偏导数，外层函数 $f$ 并没有给出具体表达式，只告诉我们 $z = f(u,v)$ 的某种结构。其中 \(f_1, f_2\) 分别是 \(f\) 对第一个中间变量 \(p\) 和第二个中间变量 \(q\) 的偏导数，它们仍是 \(u, v\) 的函数，即 \(f_1 = f_1(p,q),\; f_2 = f_2(p,q)\)。


\begin{myexample}
  设\(z = f(x^2 + y^2,\; e^{xy})\)，求 \(\frac{\partial^2 z}{\partial x^2}\) 和 \(\frac{\partial^2 z}{\partial y^2}\)。

  分析函数，令 \(u = x^2 + y^2\)，\(v = e^{xy}\)，则 \(z = f(u, v)\)。\(u\) 和 \(v\) 都是 \(x, y\) 的二元函数，\(z\) 通过 \(u, v\) 依赖于 \(x, y\)。画出路径关系图：

\begin{center}
\begin{tikzpicture}[scale=1.2,
    >=stealth,
    var/.style={draw, circle, minimum size=7mm, inner sep=0pt},
    edge/.style={->, thick}]
    
  \node[var] (z) at (0,0)   {$z$};
  \node[var] (u) at (2,0.7) {$u$};
  \node[var] (v) at (2,-0.7){$v$};
  \node[var] (x) at (4,0.7) {$x$};
  \node[var] (y) at (4,-0.7){$y$};

  \draw[edge] (z) -- (u);
  \draw[edge] (z) -- (v);
  \draw[edge] (u) -- (x);
  \draw[edge] (u) -- (y);
  \draw[edge] (v) -- (x);
  \draw[edge] (v) -- (y);
\end{tikzpicture}
\end{center}

从图中可以看出，\(z\) 对 \(x\) 有两条路径 \(z \to u \to x\) 和 \(z \to v \to x\)，对 \(y\) 同样有两条路径 \(z \to u \to y\) 和 \(z \to v \to y\)。因此，由链式法则先求一阶偏导数：
\[
\frac{\partial z}{\partial x} = \frac{\partial f}{\partial u}\frac{\partial u}{\partial x} + \frac{\partial f}{\partial v}\frac{\partial v}{\partial x} = f_1 \cdot 2x + f_2 \cdot y e^{xy}
\]
\[
\frac{\partial z}{\partial y} = \frac{\partial f}{\partial u}\frac{\partial u}{\partial y} + \frac{\partial f}{\partial v}\frac{\partial v}{\partial y} = f_1 \cdot 2y + f_2 \cdot x e^{xy}
\]

计算二阶偏导数。先求 \(\dfrac{\partial^2 z}{\partial x^2}\)，对 \(\dfrac{\partial z}{\partial x}\) 再求关于 \(x\) 偏导：
\[
\frac{\partial^2 z}{\partial x^2} = \frac{\partial}{\partial x}(2x f_1 + y e^{xy} f_2)
\]

注意 \(f_1, f_2\) 对 \(x\) 求导时仍需用链式法则：
\[
\frac{\partial f_1}{\partial x} = \frac{\partial f_1}{\partial u}\frac{\partial u}{\partial x} + \frac{\partial f_1}{\partial v}\frac{\partial v}{\partial x} = f_{11}\cdot 2x + f_{12}\cdot y e^{xy}
\]
\[
\frac{\partial f_2}{\partial x} = \frac{\partial f_2}{\partial u}\frac{\partial u}{\partial x} + \frac{\partial f_2}{\partial v}\frac{\partial v}{\partial x} = f_{21}\cdot 2x + f_{22}\cdot y e^{xy}
\]

其中 \(f_{11}, f_{12}, f_{21}, f_{22}\) 为 \(f\) 的二阶偏导数。代入计算：
\begin{align*}
\frac{\partial^2 z}{\partial x^2}
&= 2f_1 + 2x\frac{\partial f_1}{\partial x} + \frac{\partial}{\partial x}(y e^{xy})\cdot f_2 + y e^{xy}\frac{\partial f_2}{\partial x} \\
&= 2f_1 + 2x(2x f_{11} + y e^{xy} f_{12}) + y^2 e^{xy} f_2 + y e^{xy}(2x f_{21} + y e^{xy} f_{22}) \\
&= 2f_1 + y^2 e^{xy} f_2 + 4x^2 f_{11} + 2xy e^{xy} f_{12} + 2xy e^{xy} f_{21} + y^2 e^{2xy} f_{22}
\end{align*}

由于混合偏导数连续，有 \(f_{12} = f_{21}\)，合并得：
\[
\frac{\partial^2 z}{\partial x^2} = 2f_1 + y^2 e^{xy} f_2 + 4x^2 f_{11} + 4xy e^{xy} f_{12} + y^2 e^{2xy} f_{22}
\]

再求 \(\dfrac{\partial^2 z}{\partial y^2}\)，对 \(\dfrac{\partial z}{\partial y}\) 求关于 \(y\) 偏导：
\[
\frac{\partial^2 z}{\partial y^2} = \frac{\partial}{\partial y}(2y f_1 + x e^{xy} f_2)
\]

\(f_1, f_2\) 对 \(y\) 的偏导数为：
\[
\frac{\partial f_1}{\partial y} = f_{11}\cdot 2y + f_{12}\cdot x e^{xy},\quad
\frac{\partial f_2}{\partial y} = f_{21}\cdot 2y + f_{22}\cdot x e^{xy}
\]

其中 \(\frac{\partial}{\partial y}(x e^{xy}) = x^2 e^{xy}\)。代入：
\begin{align*}
\frac{\partial^2 z}{\partial y^2}
&= 2f_1 + 2y\frac{\partial f_1}{\partial y} + x^2 e^{xy} f_2 + x e^{xy}\frac{\partial f_2}{\partial y} \\
&= 2f_1 + 2y(2y f_{11} + x e^{xy} f_{12}) + x^2 e^{xy} f_2 + x e^{xy}(2y f_{21} + x e^{xy} f_{22}) \\
&= 2f_1 + x^2 e^{xy} f_2 + 4y^2 f_{11} + 2xy e^{xy} f_{12} + 2xy e^{xy} f_{21} + x^2 e^{2xy} f_{22}
\end{align*}

利用 \(f_{12} = f_{21}\) 合并，得：
\[
\frac{\partial^2 z}{\partial y^2} = 2f_1 + x^2 e^{xy} f_2 + 4y^2 f_{11} + 4xy e^{xy} f_{12} + x^2 e^{2xy} f_{22}
\]

\end{myexample}

这个例子其实是复合函数求二阶偏导数，虽然很复杂，但是思路是清晰的，就是不断使用链式法则求导，找出所有从因变量到自变量的路径，然后把每条路径上偏导数的乘积加起来。

\begin{myexample}
  设函数 $z = f(u,v,w)$，其中 $u = x^2 + y^2$，$v = e^x$，$w = \ln y$，求 $\frac{\partial^2 z}{\partial x \partial y}$。

  分析函数，\(u\) 依赖于 \(x,y\)，\(v\) 仅依赖于 \(x\)，\(w\) 仅依赖于 \(y\)。画出路径关系图：

\begin{center}
\begin{tikzpicture}[scale=1.3,
    >=stealth,
    var/.style={draw, circle, minimum size=7mm, inner sep=0pt},
    edge/.style={->, thick}]
    
  \node[var] (z) at (0,0)    {$z$};
  \node[var] (u) at (2,1.2)  {$u$};
  \node[var] (v) at (2,0)    {$v$};
  \node[var] (w) at (2,-1.2) {$w$};
  \node[var] (x) at (4,1.2)  {$x$};
  \node[var] (y) at (4,-1.2) {$y$};

  \draw[edge] (z) -- (u);
  \draw[edge] (z) -- (v);
  \draw[edge] (z) -- (w);
  \draw[edge] (u) -- (x);
  \draw[edge] (u) -- (y);
  \draw[edge] (v) -- (x);
  \draw[edge] (w) -- (y);
\end{tikzpicture}
\end{center}

从图中可以看出：\(z\) 对 \(x\) 有两条路径：\(z \to u \to x\)，\(z \to v \to x\)。\(z\) 对 \(y\) 有两条路径：\(z \to u \to y\)，\(z \to w \to y\)。

由链式法则，先计算一阶偏导：
\[
\frac{\partial z}{\partial x} = \frac{\partial z}{\partial u} \frac{\partial u}{\partial x} + \frac{\partial z}{\partial v} \frac{d v}{d x}
= f_1 \cdot 2x + f_2 \cdot e^x = 2x f_1 + e^x f_2
\]
\[
\frac{\partial z}{\partial y} = \frac{\partial z}{\partial u} \frac{\partial u}{\partial y} + \frac{\partial z}{\partial w} \frac{d w}{d y}
= f_1 \cdot 2y + f_3 \cdot \frac{1}{y} = 2y f_1 + \frac{1}{y} f_3
\]

下面求混合偏导数 \(\dfrac{\partial^2 z}{\partial x \partial y} = \dfrac{\partial}{\partial y}\!\left(\dfrac{\partial z}{\partial x}\right)\)。对 \(2x f_1 + e^x f_2\) 关于 \(y\) 求偏导时，\(2x\) 与 \(e^x\) 与 \(y\) 无关，视为系数：
\[
\frac{\partial^2 z}{\partial x \partial y} = 2x \frac{\partial f_1}{\partial y} + e^x \frac{\partial f_2}{\partial y}
\]

现在求 \(f_1\) 和 \(f_2\) 对 \(y\) 的偏导数。它们通过中间变量 \(u,v,w\) 依赖于 \(y\)，但 \(v\) 与 \(y\) 无关，因此链式法则中只出现 \(u\) 和 \(w\) 的路径：
\[
\frac{\partial f_1}{\partial y} = \frac{\partial f_1}{\partial u}\frac{\partial u}{\partial y} + \frac{\partial f_1}{\partial w}\frac{\partial w}{\partial y},
\qquad
\frac{\partial f_2}{\partial y} = \frac{\partial f_2}{\partial u}\frac{\partial u}{\partial y} + \frac{\partial f_2}{\partial w}\frac{\partial w}{\partial y}
\]

代入 \(\frac{\partial u}{\partial y} = 2y\)，\(\frac{\partial w}{\partial y} = \frac{1}{y}\)，得：
\[
\frac{\partial f_1}{\partial y} = f_{11} \cdot 2y + f_{13} \cdot \frac{1}{y},\qquad
\frac{\partial f_2}{\partial y} = f_{21} \cdot 2y + f_{23} \cdot \frac{1}{y}
\]

将上述结果代回，整理得：
\begin{align*}
\frac{\partial^2 z}{\partial x \partial y}
&= 2x( 2y f_{11} + \frac{1}{y} f_{13} ) + e^x( 2y f_{21} + \frac{1}{y} f_{23} ) \\
&= 4xy f_{11} + \frac{2x}{y} f_{13} + 2y e^x f_{21} + \frac{e^x}{y} f_{23}
\end{align*}

因为没有 \(f\) 的显式，最终结果保留二阶偏导记号，它们仍然都是 \(u = x^2 + y^2\)，\(v = e^x\)，\(w = \ln y\) 的函数。
\end{myexample}

这是一个三元的函数，但是只要我们抓住路径关系图，找出所有从因变量到自变量的路径，即使是无限元函数我们也能求偏导数。

最后要说明的是，对 $u$ 求偏导，把$v$看作常数，与 $u$ 最终依赖什么自变量毫无关系。最终代入 $u,v$ 关于 $x,y$ 的表达式时，才真正显含 $x,y$。每次求偏导时，都把其他变量视为常数，就不会出错。

多元复合函数求偏导数的链式法则，是整个多元微分学运算的骨架。它来源于一元函数，但比一元函数复杂。只要掌握了它，就能处理各种复杂的复合函数求偏导数问题。

\subsection{隐函数求偏导数}

有时变量之间的关系常常不是以 $y = f(x)$ 或 $z = f(x,y)$ 这种显式给出的，而是以隐函数的形式给出。例如我们之前学过的一元隐函数$F(x, y) = 0$，虽然不能简单地解出 $y$ 关于 $x$ 的初等函数表达式，但在一定条件下，它仍然确定了一个 $y$ 对 $x$ 的函数。我们现在要研究如何直接从方程本身求出这个“隐藏”函数的偏导数。

我们先从我们学过的一元隐函数入手。设方程$F(x, y) = 0$在某个区域内确定了一个函数 $y = y(x)$。也就是说，把 $y(x)$ 代入后，$F(x, y(x)) = 0$ 是一个恒等式。我们的目标是求 $\dfrac{dy}{dx}$，而不必先解出 $y$ 的显式。

现在我们有了链式法则和偏导数，这个想法实施起来就很简单了。把 $F(x,y)$ 看作 $x$ 和 $y$ 的二元函数，而 $y$ 又是 $x$ 的函数。对恒等式 $F(x, y(x)) = 0$ 两边关于 $x$ 求偏导，根据链式法则，有：
\[
\frac{\partial F}{\partial x} \cdot 1 + \frac{\partial F}{\partial y} \cdot \frac{dy}{dx} = 0
\]

也就是：
\[
F_x + F_y \cdot y' = 0
\]

只要 $F_y \neq 0$，就可以解出：
\[
\frac{dy}{dx} = -\frac{F_x}{F_y}
\]

这就是一元隐函数的求导公式。

\begin{theorem}{一元隐函数求导公式}{}
  设方程 $F(x, y) = 0$ 在某个区间内确定了一个隐函数 $y = y(x)$，则
  \[
  \frac{dy}{dx} = -\frac{F_x}{F_y}
  \]
\end{theorem}

接下来看一个例子：

\begin{myexample}
设方程 $e^y + xy = 0$ 确定了函数 $y = y(x)$，求 $\dfrac{dy}{dx}$。

令 $F(x, y) = e^y + xy$，则：
\[
F_x = y,\qquad F_y = e^y + x
\]

代入公式，得：
\[
\frac{dy}{dx} = -\frac{F_x}{F_y} = -\frac{y}{e^y + x}
\]
\end{myexample}

这就是我们在第二章里说的，除了方程两边对$x$求导外，求隐函数导数的更强大的工具。

现在推广到二元隐函数，设方程$F(x, y, z) = 0$在某个区域内确定了一个函数 $z = z(x, y)$，即 $F(x, y, z(x,y)) = 0$。我们现在要求 $\dfrac{\partial z}{\partial x}$ 和 $\dfrac{\partial z}{\partial y}$，怎么做呢？

和一元隐函数的操作一样，我们把 $F(x,y,z)$ 看作三元函数，而 $z$ 是 $x$ 和 $y$ 的函数。固定 $y$，对恒等式两边关于 $x$ 求偏导，链式法则出手：
\[
\frac{\partial F}{\partial x} \cdot 1 + \frac{\partial F}{\partial y} \cdot 0 + \frac{\partial F}{\partial z} \cdot \frac{\partial z}{\partial x} = 0
\]

也就是：
\[
F_x + F_z \frac{\partial z}{\partial x} = 0
\]

类似地，固定 $x$ 对 $y$ 求偏导，得到：
\[
F_y + F_z \frac{\partial z}{\partial y} = 0
\]

因此，只要 $F_z \neq 0$，就有：
\[
\frac{\partial z}{\partial x} = -\frac{F_x}{F_z},\qquad
\frac{\partial z}{\partial y} = -\frac{F_y}{F_z}
\]

这就是二元隐函数的求偏导公式。

\begin{theorem}{二元隐函数求导公式}{}
  设方程 $F(x,y,z)$ 在某个区域内确定了一个隐函数 $z = z(x,y)$，则
\[
\frac{\partial z}{\partial x} = -\frac{F_x}{F_z},\qquad
\frac{\partial z}{\partial y} = -\frac{F_y}{F_z}
\]
\end{theorem}

直观上说，求 $\dfrac{\partial z}{\partial x}$ 时，把 $y$ 当成常数，然后按照一元隐函数的方式处理即可。

接下来看一个例子：

\begin{myexample}
设方程 $x^2 + y^2 + z^2 = 4$ 在 $z > 0$ 时确定了隐函数 $z = z(x,y)$，求 $\dfrac{\partial z}{\partial x}$ 和 $\dfrac{\partial z}{\partial y}$。

令 $F(x, y, z) = x^2 + y^2 + z^2 - 4$，则有：
\[
F_x = 2x,\quad F_y = 2y,\quad F_z = 2z
\]

因此：
\[
\frac{\partial z}{\partial x} = -\frac{2x}{2z} = -\frac{x}{z},\qquad
\frac{\partial z}{\partial y} = -\frac{2y}{2z} = -\frac{y}{z}
\]

这个结果和从显式 $z = \sqrt{4 - x^2 - y^2}$ 直接求偏导的结果完全一样。读者可自行验证。
\end{myexample}

这种方法的最大优点是，即使 $F(x,y,z)=0$ 无法解出 $z$ 的显式，也能求出偏导数。例如方程 $z + e^z = xy$，根本无法写出 $z$ 的初等表达式，但我们可以轻易求出偏导。

\begin{myexample}
设方程 $z + e^z = xy$ 确定了隐函数 $z = z(x,y)$，求 $\dfrac{\partial z}{\partial x}$ 和 $\dfrac{\partial z}{\partial y}$。

令 $F(x, y, z) = z + e^z - xy$，则有：
\[
F_x = -y,\quad F_y = -x,\quad F_z = 1 + e^z
\]

于是：
\[
\frac{\partial z}{\partial x} = -\frac{-y}{1 + e^z} = \frac{y}{1 + e^z},\qquad
\frac{\partial z}{\partial y} = -\frac{-x}{1 + e^z} = \frac{x}{1 + e^z}
\]
\end{myexample}

如果还需要求二阶偏导数，就在一阶偏导数的基础上再求一次偏导，此时要注意一阶偏导的表达式里可能仍含有 $z$，而 $z$ 本身又是 $x,y$ 的函数，需要再次使用链式法则。

\begin{myexample}
设方程 $z + e^z = xy$ 确定了隐函数 $z = z(x,y)$，求 $\dfrac{\partial^2 z}{\partial x^2}$。

已知 $\dfrac{\partial z}{\partial x} = \dfrac{y}{1 + e^z}$，再对 $x$ 求偏导：

\begin{align*}
\frac{\partial^2 z}{\partial x^2}
&= \frac{\partial}{\partial x}( \frac{y}{1 + e^z} )
 = y \cdot \frac{-\; e^z \frac{\partial z}{\partial x}}{(1 + e^z)^2} \\
&= -\frac{y e^z}{(1 + e^z)^2} \cdot \frac{y}{1 + e^z}
 = -\frac{y^2 e^z}{(1 + e^z)^3}
\end{align*}

这样我们就求得了二阶偏导数。
\end{myexample}

隐函数求偏导的方法与前面多元复合函数链式法则是一脉相承的。我们可以这样总结：

面对一个方程，认定哪些是自变量，哪些是因变量，然后对等式两边关于某个自变量求偏导，因变量求导时按照链式法则展开，最后解出所需的偏导数。


\section{全微分}
与偏导数对应的，我们有全微分。顾名思义，它是函数在一点处因所有自变量同时微小变化所引起的全部微分。偏导数，就是偏向某个变量的导数。全微分，就是全部变量的微分。

\subsection{函数可微的条件}

在一元函数中，可导和可微是等价的，导数存在也就意味着微分存在。但在多元函数中，情况要微妙得多。

让我们先回忆一下一元可微的定义：$f(x)$ 在 $x_0$ 处可微，是指存在常数 $A$ 使得当$\Delta x \to 0$时，
$\Delta y = A \cdot \Delta x + o(\Delta x)$，其中这个 $A$ 就是 $f'(x_0)$。推广到二元函数，我们自然希望：
\[
\Delta z = A \cdot \Delta x + B \cdot \Delta y + o\!(\sqrt{(\Delta x)^2 + (\Delta y)^2})
\]

其中 $A$ 和 $B$ 是两个偏导数 $f_x$ 和 $f_y$。那么我们仿照一元函数，可以定义：

\begin{definition}{全微分的定义}{}
设函数 $z = f(x, y)$ 在 $(x_0, y_0)$ 的某个邻域内有定义。记自变量的增量为 $\Delta x$和$\Delta y$，函数值的增量为$\Delta z$，自变量的增量$(\Delta x,\Delta y)$的模长为$\rho$，则：
  \[
  \Delta z = f(x_0+\Delta x, y_0+\Delta y)
  \]
  \[
  \rho = \sqrt{(\Delta x)^2 + (\Delta y)^2}
  \]
  如果存在两个与 $\Delta x$ 和 $\Delta y$ 无关的常数 $A$ 和 $B$，使得：
\begin{center}
    当$\rho \to 0$时，$\Delta z = A \Delta x + B \Delta y + o(\rho)$
\end{center}

我们就说 $f(x, y)$ 在 $(x_0, y_0)$ 处\textbf{可微}，并称 $A \Delta x + B \Delta y$ 为 $f$ 在 $(x_0, y_0)$ 处的\textbf{全微分}，记作
\[
dz = A\, dx + B\, dy
\]
\end{definition}

有没有发现这和一元可微的定义几乎一样？唯一的区别是一元的误差是 $o(\Delta x)$，二元的误差是 $o(\rho)$，其中 $\rho$ 的几何意义是两点间的距离。

和一元情形类似，如果 $f$ 在 $(x_0, y_0)$ 处可微，则 $A = f_x(x_0, y_0)$，$B = f_y(x_0, y_0)$。也就是说，全微分可以写成：
\[
dz = f_x(x_0, y_0)\, dx + f_y(x_0, y_0)\, dy
\]

这就给出了可微的一个必要条件：

\begin{theorem}{可微的必要条件}{}
若 $f(x, y)$ 在 $(x_0, y_0)$ 处可微，则 $f$ 在该点的两个偏导数 $f_x$ 和 $f_y$ 都存在，且全微分为：
\[
dz = f_x(x_0, y_0)\, dx + f_y(x_0, y_0)\, dy
\]
\end{theorem}

但关键是，反过来并不成立。\textbf{偏导数都存在，不能保证函数可微}。看一下这个例子：

\begin{myexample}
  设 $f(x, y) = \begin{cases}
    \dfrac{xy}{\sqrt{x^2 + y^2}}, & (x, y) \neq (0, 0) \\
    0, & (x, y) = (0, 0)
  \end{cases}$，判断函数在原点是否可微。

  先算 $f$ 在原点处的两个偏导数：
  \[
  f_x(0, 0) = \lim_{\Delta x \to 0} \frac{f(\Delta x, 0) - f(0, 0)}{\Delta x}
  = \lim_{\Delta x \to 0} \frac{0 - 0}{\Delta x} = 0
  \]

  同理求得 $f_y(0, 0) = 0$。两个偏导数都存在。

  再看可微性。如果它可微，则全微分应该为 $dz = 0 \cdot dx + 0 \cdot dy = 0$，且 $\Delta z = 0 \cdot \Delta x + 0 \cdot \Delta y + o(\rho)$，也就是 $\frac{\Delta z}{\rho} \to 0$。

  沿 $y = x$ 方向计算：
  \[
  \frac{\Delta z}{\rho} = \frac{f(x, x) - 0}{\sqrt{x^2 + x^2}}
  = \frac{\frac{x^2}{\sqrt{2x^2}}}{\sqrt{2x^2}}
  = \frac{\frac{x^2}{\sqrt{2}\,|x|}}{\sqrt{2}\,|x|}
  = \frac{1}{2}
  \]

  显然$\frac{1}{2}$并不趋于$0$，所以函数在原点不可微，尽管偏导数都存在。
\end{myexample}

那么偏导数满足什么条件才能保证可微呢？我们给出：

\begin{theorem}{可微的充分条件}{}
若 $f(x, y)$ 在 $(x_0, y_0)$ 的某个邻域内偏导数 $f_x$ 和 $f_y$ 都存在，且 $f_x$ 和 $f_y$ 在 $(x_0, y_0)$ 处连续，则 $f(x, y)$ 在 $(x_0, y_0)$ 处可微。
\end{theorem}

这个定理是最常用的判断可微的方法。我们用这个定理再判断一下刚刚的例子：

\begin{myexample}
  设 $f(x, y) = \begin{cases}
    \dfrac{xy}{\sqrt{x^2 + y^2}}, & (x, y) \neq (0, 0) \\
    0, & (x, y) = (0, 0)
  \end{cases}$，判断函数在原点是否可微。

已求得原点处的两个偏导数：$f_x(0, 0) = 0$，$f_y(0, 0) = 0$，再求其他情况的偏导数：

\[f_x=\frac{\partial f}{\partial x}=\frac{y^3}{\sqrt{(x^2+y^2)^3}}\]

\[f_y=\frac{\partial f}{\partial y}=\frac{x^3}{\sqrt{(x^2+y^2)^3}}\]

沿路径$y=kx$取$f_x$在$(0,0)$处的极限：

\[\lim_{x \to 0} f_x =\lim_{x \to 0} \frac{k^3x^3}{\sqrt{(x^2+k^2x^2)^3}}=\frac{k^3}{\sqrt{(1+k^2)^3}}\]

也就是说，$f_x$在$(0,0)$处的极限与$k$的取值有关，因此，$f_x$在$(0,0)$处不连续，所以函数在原点不可微。
\end{myexample}

把偏导数的存在性、连续性、函数的可微性之间的关系总结一下：

\begin{center}
\begin{tikzpicture}[scale=1.3,
    >=stealth,
    var/.style={draw, rectangle, minimum size=7mm, inner sep=3pt},
    edge/.style={->, thick}]
    
  \node[var] (pc) at (0,0)    {偏导连续};
  \node[var] (diff) at (2.5,0)  {函数可微};
  \node[var] (cont) at (5,0)    {函数连续};
  \node[var] (pe) at (2.5,-1.6) {偏导存在};

  \draw[edge] (pc) -- (diff);
  \draw[edge] (diff) -- (cont);
  \draw[edge] (diff) -- (pe);
\end{tikzpicture}
\end{center}

注意箭头都是单向的。多元函数可导和可微的关系比一元函数更复杂，需要特别记忆。

\subsection{一阶全微分形式不变性}

在一元函数中，我们有一阶微分形式不变性。到了多元函数，全微分是否也具备这种“形式不依赖变量身份”的性质呢？答案是肯定的。我们以二元函数为例：

\begin{theorem}{一阶全微分形式不变性}{}
设函数 $z = f(u, v)$ 可微。则不论 $u, v$ 是自变量还是中间变量，总有：
\[
dz = \frac{\partial z}{\partial u}\,du + \frac{\partial z}{\partial v}\,dv
\]
\end{theorem}

证明方法也和一元的情形是类似的：

当 $u, v$ 为中间变量时，由复合函数链式法则，可得：
\[
\frac{\partial z}{\partial x} = \frac{\partial z}{\partial u}\frac{\partial u}{\partial x} + \frac{\partial z}{\partial v}\frac{\partial v}{\partial x},\qquad
\frac{\partial z}{\partial y} = \frac{\partial z}{\partial u}\frac{\partial u}{\partial y} + \frac{\partial z}{\partial v}\frac{\partial v}{\partial y}
\]

于是 $z$ 作为 $x, y$ 函数的全微分为：
\begin{align*}
dz &= \frac{\partial z}{\partial x}\,dx + \frac{\partial z}{\partial y}\,dy \\
&= \Bigl(\frac{\partial z}{\partial u}\frac{\partial u}{\partial x} + \frac{\partial z}{\partial v}\frac{\partial v}{\partial x}\Bigr)dx
   + \Bigl(\frac{\partial z}{\partial u}\frac{\partial u}{\partial y} + \frac{\partial z}{\partial v}\frac{\partial v}{\partial y}\Bigr)dy \\
&= \frac{\partial z}{\partial u}\Bigl(\frac{\partial u}{\partial x}dx + \frac{\partial u}{\partial y}dy\Bigr)
   + \frac{\partial z}{\partial v}\Bigl(\frac{\partial v}{\partial x}dx + \frac{\partial v}{\partial y}dy\Bigr) \\
&= \frac{\partial z}{\partial u}\,du + \frac{\partial z}{\partial v}\,dv
\end{align*}

这就说明，无论从哪一层变量去写，一阶全微分的形式都不改变。

这个性质在实际计算中非常方便。求偏导数时，有时不必分清“谁是自变量谁是中间变量”，只要对最外层的函数直接写出全微分，再代入内层变量的微分，整理后 $dx, dy$ 的系数自然就是要求的偏导数。

我们看个例子说明一下：

\begin{myexample}
设 $z = f(u, v)$，其中 $u = x^2 + y^2$，$v = e^{xy}$，求 $\dfrac{\partial z}{\partial x}$ 和 $\dfrac{\partial z}{\partial y}$。

对 $z$ 写出全微分：
\[
dz = f_u\,du + f_v\,dv
\]

分别计算 $du, dv$：
\[
du = 2x\,dx + 2y\,dy,\qquad
dv = y e^{xy}\,dx + x e^{xy}\,dy
\]

代入 $dz$ 并整理：
\begin{align*}
dz &= f_u(2x\,dx + 2y\,dy) + f_v(y e^{xy}\,dx + x e^{xy}\,dy) \\
&= (2x f_u + y e^{xy} f_v)\,dx + (2y f_u + x e^{xy} f_v)\,dy
\end{align*}

因此同时得到两个偏导数：
\[
\frac{\partial z}{\partial x} = 2x f_u + y e^{xy} f_v,\qquad
\frac{\partial z}{\partial y} = 2y f_u + x e^{xy} f_v
\]
\end{myexample}

可以看到，利用全微分形式不变性，我们可以在一步运算中同时算出所有偏导数。之前学习的隐函数求偏导数也可以从这个角度重新理解：

对于方程 $F(x, y, z) = 0$ 所确定的隐函数 $z = z(x, y)$，两边求全微分：
\[
F_x\,dx + F_y\,dy + F_z\,dz = 0
\]

当 $F_z \neq 0$ 时，解出：
\[
dz = -\frac{F_x}{F_z}\,dx - \frac{F_y}{F_z}\,dy
\]

由此立即得到 $\frac{\partial z}{\partial x} = -\frac{F_x}{F_z}$，$\frac{\partial z}{\partial y} = -\frac{F_y}{F_z}$，这与前面的隐函数求导公式完全一致。

\subsection{参数方程求偏导数}

前面我们处理了由一个方程 $F(x,y,z)=0$ 确定的隐函数。现在有了全微分，我们可以进一步处理由参数方程确定的函数。我们依旧以二元函数为例，那么方程组$\begin{cases}
F(x,y,u,v)=0 \\
G(x,y,u,v)=0
\end{cases}$在一定的条件下可以把 $u,v$ 同时确定为 $x,y$ 的函数，即：$u=u(x,y)$，$v=v(x,y)$。这可以看成是一种以 $x,y$ 为“参数”的参数方程形式。

解决的思路还是我们熟悉的链式法则，不过这回要同时处理两个函数，得到的是一个关于偏导数的方程组。

把 $u=u(x,y)$，$v=v(x,y)$ 代入 $F$ 和 $G$，得到：
\[
F(x,\,y,\,u(x,y),\,v(x,y))=0,\qquad
G(x,\,y,\,u(x,y),\,v(x,y))=0
\]

两边对 $x$ 求偏导。注意这时候的 $F$ 是四元函数，变量为 $x,y,u,v$，其中 $u,v$ 又分别是 $x,y$ 的函数。画出路径关系：

\begin{center}
\begin{tikzpicture}[scale=1.3,
    >=stealth,
    var/.style={draw, circle, minimum size=7mm, inner sep=0pt},
    edge/.style={->, thick}]
    
  % 函数 (左)
  \node[var] (F) at (0,1.2)  {$F$};
  \node[var] (G) at (0,-1.2) {$G$};

  % 中间变量 (中，向中间靠拢)
  \node[var] (u) at (2.2,0.8) {$u$};
  \node[var] (v) at (2.2,-0.8){$v$};

  % 自变量 (右)
  \node[var] (x) at (4.5,1.2)  {$x$};
  \node[var] (y) at (4.5,-1.2) {$y$};

  % 函数到中间变量 (F->u, F->v, G->u, G->v)
  \draw[edge] (F) -- (u);
  \draw[edge] (F) -- (v);
  \draw[edge] (G) -- (u);
  \draw[edge] (G) -- (v);

  % 函数直接到自变量 (F->x, F->y, G->x, G->y)
  \draw[edge] (F) -- (x);
  \draw[edge] (F) -- (y);
  \draw[edge] (G) -- (x);
  \draw[edge] (G) -- (y);

  % 中间变量到自变量 (u->x, u->y, v->x, v->y)
  \draw[edge] (u) -- (x);
  \draw[edge] (u) -- (y);
  \draw[edge] (v) -- (x);
  \draw[edge] (v) -- (y);

\end{tikzpicture}
\end{center}

从图中可以看到，$F$ 对 $x$ 的总变化率有三条路径：直接到 $x$，以及经过 $u$ 或 $v$ 再到 $x$。对 $G$ 也是如此。因此链式法则给出：
\[
\begin{cases}
F_x \cdot 1 + F_u \cdot u_x + F_v \cdot v_x = 0\\[4pt]
G_x \cdot 1 + G_u \cdot u_x + G_v \cdot v_x = 0
\end{cases}
\]

把已知的偏导数 $F_x,F_u,F_v$ 和 $G_x,G_u,G_v$ 当作系数，这就是关于未知量 $u_x,v_x$ 的二元一次方程组。解出它，就得到了 $\frac{\partial u}{\partial x}$ 和 $\frac{\partial v}{\partial x}$。

同理，对 $y$ 求偏导可以得到关于 $u_y,v_y$ 的方程组：
\[
\begin{cases}
F_y \cdot 1 + F_u \cdot u_y + F_v \cdot v_y = 0\\[4pt]
G_y \cdot 1 + G_u \cdot u_y + G_v \cdot v_y = 0
\end{cases}
\]

我们以关于未知量 $u_x,v_x$ 的二元一次方程组为例，把$\frac{\partial u}{\partial x}$解出来：

将第一个方程两边同乘 \(G_v\)，第二个方程两边同乘 \(F_v\)：
$$
\begin{cases}
F_u G_v u_x + F_v G_v v_x = -F_x G_v \\
G_u F_v u_x + G_v F_v v_x = -G_x F_v
\end{cases}
$$

两式相减，\(v_x\) 的项消掉，得到：
$$
(F_u G_v - F_v G_u) u_x = -F_x G_v + G_x F_v = F_v G_x - F_x G_v
$$

移项得：
$$
\frac{\partial u}{\partial x} = u_x = \frac{F_v G_x - F_x G_v}{F_u G_v - F_v G_u}
$$

其他偏导数同理。

仔细观察，我们会发现分子分母可以用行列式表示，于是我们有：

\begin{definition}{雅可比行列式}{}
  定义\textbf{雅可比行列式}：
  \[
J = \frac{\partial(F,G)}{\partial(u,v)} =
\begin{vmatrix}
\frac{\partial F}{\partial u} & \frac{\partial F}{\partial v} \\[6pt]
\frac{\partial G}{\partial u} & \frac{\partial G}{\partial v}
\end{vmatrix}
= 
\begin{vmatrix}
F_u & F_v \\
G_u & G_v
\end{vmatrix}
= F_u G_v - F_v G_u
\]

\end{definition}

也就是说，只要 $J \neq 0$，我们就可以写出偏导数的解：
\begin{theorem}{参数方程求偏导公式}{}
  若方程组$\begin{cases}
F(x,y,u,v)=0 \\
G(x,y,u,v)=0
\end{cases}$确定了隐函数$u = u(x,y)$，$v = v(x,y)$，则偏导数为：
\[
\frac{\partial u}{\partial x} = -\frac{\begin{vmatrix}
F_x & F_v \\
G_x & G_v
\end{vmatrix}}{J}
,\qquad
\frac{\partial v}{\partial x} = -\frac{\begin{vmatrix}
F_u & F_x \\
G_u & G_x
\end{vmatrix}}{J}
\]
\[
\frac{\partial u}{\partial y} = -\frac{\begin{vmatrix}
F_y & F_v \\
G_y & G_v
\end{vmatrix}}{J}
,\qquad
\frac{\partial v}{\partial y} = -\frac{\begin{vmatrix}
F_u & F_y \\
G_u & G_y
\end{vmatrix}}{J}
\]
\end{theorem}

观察这些公式，规律其实很清楚，偏导数的分母总是雅可比行列式，分子则是把雅可比行列式的对应列换成 $x$ 或 $y$ 的偏导数列。例如要求$\frac{\partial u}{\partial x}$，$\partial u$在雅可比行列式的第一列，就把第一列所有$\partial u$换成 $\partial x$，其余不变。要求$\frac{\partial v}{\partial y}$，$\partial v$在雅可比行列式的第二列，就把第二列所有$\partial v$换成 $\partial y$，其余不变。

我们来看个例子实践一下：

\begin{myexample}
  设方程组$\begin{cases}
u +\cos v=e^x+\ln y\\
\sin u +v=\ln x+e^y
\end{cases}$确定了隐函数 $u = u(x,y)$，$v = v(x,y)$，求 $\dfrac{\partial u}{\partial x}$ 与 $\dfrac{\partial v}{\partial y}$。

将方程组写成隐函数形式：
\[
\begin{aligned}
F(x,y,u,v) &= u + \cos v - e^x - \ln y = 0 \\
G(x,y,u,v) &= \sin u + v - \ln x - e^y = 0
\end{aligned}
\]

计算一阶偏导数：
\[
\begin{aligned}
F_u &= 1, & F_v &= -\sin v, & F_x &= -e^x, & F_y &= -\frac{1}{y}\\
G_u &= \cos u, & G_v &= 1, & G_x &= -\frac{1}{x}, & G_y &= -e^y
\end{aligned}
\]

关于 $(u,v)$ 的雅可比行列式为：
\[
J = \frac{\partial(F,G)}{\partial(u,v)} =
\begin{vmatrix}
\frac{\partial F}{\partial u} & \frac{\partial F}{\partial v} \\[6pt]
\frac{\partial G}{\partial u} & \frac{\partial G}{\partial v}
\end{vmatrix}
=
\begin{vmatrix}
F_u & F_v \\
G_u & G_v
\end{vmatrix}
= F_u G_v - F_v G_u
= 1\cdot 1 - (-\sin v)\cos u = 1 + \sin v \cos u
\]

要求$\frac{\partial u}{\partial x}$，$\partial u$都在雅可比行列式的第一列，将第一列$\partial u$换为$\partial x$，其余不变，得到：
\[
\frac{\partial u}{\partial x} = -\frac{\begin{vmatrix}
\frac{\partial F}{\partial x} & \frac{\partial F}{\partial v} \\[6pt]
\frac{\partial G}{\partial x} & \frac{\partial G}{\partial v}
\end{vmatrix}}{J} = -\frac{\begin{vmatrix}
F_x & F_v \\
G_x & G_v
\end{vmatrix}}{J}
\]

其中：
\[
\begin{vmatrix}
F_x & F_v \\
G_x & G_v
\end{vmatrix}
= F_xG_v - F_vG_x = (-e^x)\cdot 1 - (-\sin v)(-\frac{1}{x})
= -e^x - \frac{\sin v}{x}
\]

因此：
\[
\frac{\partial u}{\partial x} = -\frac{-e^x - \dfrac{\sin v}{x}}{1 + \sin v \cos u}
= \frac{x e^x + \sin v}{x(1 + \sin v \cos u)}
\]

要求$\frac{\partial v}{\partial y}$，$\partial u$都在雅可比行列式的第二列，将第二列$\partial v$换为$\partial y$，其余不变，得到：
\[
\frac{\partial v}{\partial y} = -\frac{\begin{vmatrix}
\frac{\partial F}{\partial u} & \frac{\partial F}{\partial y} \\[6pt]
\frac{\partial G}{\partial u} & \frac{\partial G}{\partial y}
\end{vmatrix}}{J} = -\frac{\begin{vmatrix}
F_u & F_y \\
G_u & G_y
\end{vmatrix}}{J}
\]

其中：
\[
\begin{vmatrix}
F_u & F_y \\
G_u & G_y
\end{vmatrix}
= F_uG_y - F_yG_u = 1\cdot(-e^y) - (-\frac{1}{y})\cos u
= -e^y + \frac{\cos u}{y}
\]

于是：
\[
\frac{\partial v}{\partial y} = -\frac{-e^y + \dfrac{\cos u}{y}}{1 + \sin v \cos u}
= \frac{y e^y - \cos u}{y(1 + \sin v \cos u)}
\]

综上，
\[
\frac{\partial u}{\partial x} = \frac{x e^x + \sin v}{x(1+\sin v\cos u)},\qquad
\frac{\partial v}{\partial y} = \frac{y e^y - \cos u}{y(1+\sin v\cos u)}
\]

\end{myexample}

看完这个例子，你可能会说这计算量也太大了吧。确实，所以我们还有一招，利用全微分形式不变性，在某种程度上可以更省力地处理这类问题。

我们对两个方程两边同时求全微分：
\[
\begin{cases}
F_x\,dx + F_y\,dy + F_u\,du + F_v\,dv = 0 \\
G_x\,dx + G_y\,dy + G_u\,du + G_v\,dv = 0
\end{cases}
\]

把 $dx,dy$ 看作已知，$du,dv$ 看作未知，这又是一个二元一次方程组。解出 $du,dv$：
\[
du = \frac{\partial u}{\partial x}\,dx + \frac{\partial u}{\partial y}\,dy,\qquad
dv = \frac{\partial v}{\partial x}\,dx + \frac{\partial v}{\partial y}\,dy
\]

也就是说，结果中 $dx$ 的系数就是 $\frac{\partial u}{\partial x}$ 和 $\frac{\partial v}{\partial x}$，$dy$ 的系数就是 $\frac{\partial u}{\partial y}$ 和 $\frac{\partial v}{\partial y}$。这样就求出了所有偏导数。

\begin{theorem}{全微分求偏导数法则}{}
  若方程组$\begin{cases}
F(x,y,u,v)=0 \\
G(x,y,u,v)=0
\end{cases}$确定了隐函数$u = u(x,y)$，$v = v(x,y)$，则：
\[
du = \frac{\partial u}{\partial x}\,dx + \frac{\partial u}{\partial y}\,dy,\qquad
dv = \frac{\partial v}{\partial x}\,dx + \frac{\partial v}{\partial y}\,dy
\]
然后比较系数得到对应偏导数。
\end{theorem}

我们还是以刚刚的例子作示范：

\begin{myexample}
  设方程组$\begin{cases}
u +\cos v=e^x+\ln y\\
\sin u +v=\ln x+e^y
\end{cases}$确定了隐函数 $u = u(x,y)$，$v = v(x,y)$，求 $\dfrac{\partial u}{\partial x}$ 与 $\dfrac{\partial v}{\partial y}$。

对两式两边同时求全微分：
\[
\begin{cases}
du - \sin v\,dv = e^x\,dx + \dfrac{1}{y}\,dy\\[6pt]
\cos u\,du + dv = \dfrac{1}{x}\,dx + e^y\,dy
\end{cases}
\]

先求 $\dfrac{\partial u}{\partial x}$。将第一个方程加上 $\sin v$ 乘以第二个方程，消去 $dv$：

\[
(du - \sin v\,dv) + \sin v(\cos u\,du + dv)
= (e^x\,dx + \frac{1}{y}\,dy) + \sin v(\frac{1}{x}\,dx + e^y\,dy)
\]

左端化简为：
\[
du + \sin v \cos u\,du = (1 + \sin v\cos u)\,du
\]

右端按 $dx,dy$ 合并：
\[
(e^x + \frac{\sin v}{x})\,dx + (\frac{1}{y} + e^y\sin v)\,dy
\]

因此：
\[
du = \frac{e^x + \dfrac{\sin v}{x}}{1 + \sin v\cos u}\,dx
    + \frac{\dfrac{1}{y} + e^y\sin v}{1 + \sin v\cos u}\,dy
\]

由全微分公式 $du = \frac{\partial u}{\partial x}\,dx + \frac{\partial u}{\partial y}\,dy$，即得：
\[
\frac{\partial u}{\partial x} = \frac{e^x + \dfrac{\sin v}{x}}{1 + \sin v\cos u}
      = \frac{x e^x + \sin v}{x(1 + \sin v\cos u)}
\]

再求 $\dfrac{\partial v}{\partial y}$。将第二个方程减去 $\cos u$ 乘以第一个方程，消去 $du$：
\[
(\cos u\,du + dv) - \cos u(du - \sin v\,dv)
= (\frac{1}{x}\,dx + e^y\,dy) - \cos u(e^x\,dx + \frac{1}{y}\,dy)
\]

左端化简为：
\[
dv + \cos u\sin v\,dv = (1 + \sin v\cos u)\,dv
\]

右端按 $dx,dy$ 合并：
\[
(\frac{1}{x} - e^x\cos u)\,dx + (e^y - \frac{\cos u}{y})\,dy
\]

因此：
\[
dv = \frac{\dfrac{1}{x} - e^x\cos u}{1 + \sin v\cos u}\,dx
    + \frac{e^y - \dfrac{\cos u}{y}}{1 + \sin v\cos u}\,dy
\]

由全微分公式 $dv = \frac{\partial v}{\partial x}\,dx + \frac{\partial v}{\partial y}\,dy$，即得：
\[
\frac{\partial v}{\partial y} = \frac{e^y - \dfrac{\cos u}{y}}{1 + \sin v\cos u}
      = \frac{y e^y - \cos u}{y(1 + \sin v\cos u)}
\]

\end{myexample}

对于不同的参数方程，我们可以选取不同的计算方法以简便运算。不过一般而言，这两种方法都差不多，掌握一种就可以求出绝大部分由参数方程确定的隐函数的偏导数了。

\section{方向导数与梯度}
偏导数 $f_x$ 和 $f_y$ 只告诉我们函数沿 $x$ 轴和 $y$ 轴方向的变化率。但实际中我们可能想知道，沿着任意一个方向，函数变化得有多快？比如山坡上某一点，沿正北方向的坡度是多少？沿西北方向呢？这就是我们接下来要讨论的问题。

\subsection{方向导数}
我们想知道，从函数上一点 $P_0$ 出发沿某个方向走一小步 $t$，到达点 $P'_0$，我们关心的是，当这一步越来越小时，函数值的变化是多少？换句话说，如果这一步足够小，是不是可以求得函数在此方向的变化率？由此引出方向导数：

\begin{definition}{方向导数}{}
设 $z = f(x, y)$ 在 $P_0(x_0, y_0)$ 的某个邻域内有定义，$\vec{l} = (\cos\alpha, \cos\beta)$ 是一个单位方向向量。若极限
\[
\lim_{t \to 0^+} \frac{f(x_0 + t\cos\alpha,\; y_0 + t\cos\beta) - f(x_0, y_0)}{t}
\]
存在，则称此极限为 $f$ 在 $P_0$ 处沿方向 $\vec{l}$ 的\textbf{方向导数}，记作 $\dfrac{\partial f}{\partial \vec{l}}\Big|_{P_0}$。
\end{definition}

方向导数就是函数沿某个特定方向的变化率。当 $\vec{l} = (1, 0)$（沿 $x$ 轴正方向）时，方向导数就是 $f_x$，当 $\vec{l} = (0, 1)$（沿 $y$ 轴正方向）时，方向导数就是 $f_y$。所以，偏导数其实是方向导数的特例。

如果函数可微，方向导数有一个非常简洁的计算公式：

\begin{theorem}{方向导数的计算公式}{}
若 $f(x, y)$ 在 $P_0(x_0, y_0)$ 处可微，则沿任意方向 $\vec{l} = (\cos\alpha, \cos\beta)$ 的方向导数为：
\[
\frac{\partial f}{\partial \vec{l}} = f_x(x_0, y_0)\cos\alpha + f_y(x_0, y_0)\cos\beta
\]
\end{theorem}

这个公式说明，只要算出两个偏导数 $f_x$ 和 $f_y$，再乘以方向的分量 $\cos\alpha$ 和 $\cos\beta$，再相加就得到方向导数了。你可能会说，哎呀每次都要算$\cos\alpha$ 和 $\cos\beta$实在是太麻烦了，有没有更简便一点的方法？

其实是有的。刚刚定义里提到，$ (\cos\alpha, \cos\beta)$ 是一个单位方向向量，也就是说，对任意的$\vec{l}$，我们只需要求$\vec{l}$的单位向量就可以了，而$\vec{l}$的单位向量是好求的，我们用$\vec{l}^0$表示$\vec{l}$的单位向量，则$\vec{l}^0=\frac{\vec{l}}{|\vec{l}|}$，这个实际上是高中的数学知识了。

我们来看个例子：

\begin{myexample}
  求 $f(x, y) = x^2 + y^2$ 在点 $(1, 2)$ 处沿方向 $\vec{l} = (1, \sqrt{3})$ 的方向导数。

  先求偏导：$f_x = 2x$，$f_y = 2y$。在 $(1, 2)$ 处，$f_x = 2$，$f_y = 4$。

  方向 $\vec{l} = (1, \sqrt{3})$，则$\vec{l}^0=\frac{\vec{l}}{|\vec{l}|}=(\frac{1}{2},\frac{\sqrt{3}}{2})$。

  代入公式：
  \[
  \frac{\partial f}{\partial \vec{l}} = 2 \cdot \frac{1}{2} + 4 \cdot \frac{\sqrt{3}}{2} = 1 + 2\sqrt{3}
  \]

  结果为正，说明沿这个方向函数值在增加。
\end{myexample}

这就是说，实际计算中，我们并不是一定要求$(\cos\alpha,\cos\beta)$这个向量，绝大多数时候都是求出$\vec{l}^0$，用$\vec{l}^0$代替$(\cos\alpha,\cos\beta)$，因为它们本质上是等价的。

方向导数的几何意义也是好理解的。类比偏导数，方向导数实际上就是函数$z=f(x,y)$沿方向$\vec{l}$的切线斜率。方向导数为正，就说明函数值沿方向$\vec{l}$增加，方向导数为负，就说明函数值沿方向$\vec{l}$减少。

\subsection{梯度及其几何意义}

我们先来仔细观察一下方向导数的公式。我们发现，它可以写成向量的点积形式：
\[
\frac{\partial f}{\partial \vec{l}} = (f_x,\; f_y) \cdot \vec{l}^0=f_x \cos\alpha + f_y \cos\beta
\]

这个公式里面的向量 $(f_x, f_y)$ 表示什么呢？直观上看，这个向量由两个偏导数组成。我们知道，偏导数表示函数对某个自变量的变化率，那么向量 $(f_x, f_y)$表示的是变化率的方向吗？

在此之前，我们先定义Nabla算子：

\begin{definition}{Nabla算子}{}
  定义向量微分算子
\[
\nabla = (\frac{\partial}{\partial x},\; \frac{\partial}{\partial y},\; \frac{\partial}{\partial z})
\]
\end{definition}

这是三维形式的Nabla算子。二维形式也类似，就是去掉$\frac{\partial}{\partial z}$一项。要注意的是，Nabla算子本身不是一个数，你可以把它理解为“$+$”“$-$”这种符号，它是运算符，单独出现时没有意义，一般作用在函数身上，尤其是多元函数。

我们还知道，方向导数是函数$z=f(x,y)$沿方向$\vec{l}$的切线斜率，反映的是函数沿方向$\vec{l}$的变化率，那么，函数沿着哪一个方向变化率最大呢？沿着哪一个方向变化率最小呢？

我们再来观察一下改写之后的方向导数的公式。我们发现，要让方向导数取到最大，$\vec{l}^0$就必须和向量$(f_x, f_y)$同向，因为偏导数由函数决定，而$\vec{l}^0$的方向是可以任取的。当$\vec{l}^0$和向量$(f_x, f_y)$同向时，方向导数正好就等于向量 $(f_x, f_y)$ 的模长，也就是说，向量 $(f_x, f_y)$ 表示的正是方向导数最大的方向，沿着向量 $(f_x, f_y)$ 的函数值变化率最大。同理，当$\vec{l}^0$和向量$(f_x, f_y)$垂直时，点积结果为$0$，此时向量 $\vec{l}^0$ 表示的是方向导数为$0$的方向，沿着向量 $\vec{l}^0$ 的函数值变化率为$0$。那么我们就能够定义：

\begin{definition}{梯度}{}
设 $z = f(x, y)$ 在 $P_0(x_0, y_0)$ 处偏导数存在，则称向量
\[
\nabla f\big|_{P_0} = (f_x(x_0, y_0),\; f_y(x_0, y_0))
\]
为 $f$ 在 $P_0$ 处的\textbf{梯度}。记作 $\nabla f$ 或“grad $f$”。
\end{definition}

梯度的方向决定了函数值变化最快的方向，它的模长决定了函数值变化最快时的速率。\textbf{梯度是个向量}。

引入了Nabla算子，我们可以直接定义梯度的运算。梯度实际上等于Nabla算子与$f$的数乘。
\[
\text{grad} f=\nabla f= (\frac{\partial f}{\partial x},\; \frac{\partial f}{\partial y},\; \frac{\partial f}{\partial z})
\]

当然，这是三维的。二维的写成$(\frac{\partial f}{\partial x},\; \frac{\partial f}{\partial y})$。

有了梯度，我们再来重写方向导数公式：
\[
\frac{\partial f}{\partial \vec{l}}=\nabla f \cdot \vec{l}^0
\]


从这个式子可以立刻看出几件事：

当 $\theta = 0$，即 $\vec{l}$ 与梯度同方向时，方向导数取最大值 $|\nabla f|$，梯度方向是函数值增加最快的方向。

当 $\theta = \pi$，即 $\vec{l}$ 与梯度反方向时，方向导数取最小值 $-|\nabla f|$，沿梯度反方向，函数值减少最快。

当 $\theta = \frac{\pi}{2}$，即 $\vec{l}$ 与梯度垂直时，方向导数为 $0$，沿这个方向，函数值不增不减。

这和我们刚刚讨论的结果是一致的。下图的蓝色箭头表示的是梯度的方向，红色箭头表示的是这个曲面最陡的方向。

\begin{center}
  \begin{tikzpicture}
  \begin{axis}[
      view={55}{10},
      axis lines=none,           % 去掉坐标轴
      domain=-2:2,
      y domain=-2:2,
      samples=35,
      samples y=35,
      colormap/viridis,
  ]
  
  % 曲面 z = 0.2x^2 + 0.1y^2
  \addplot3[
      surf,
      shader=interp,             % 平滑着色，去掉网格线
      opacity=0.65,
  ]
  {0.2*x^2 + 0.1*y^2};
  
  % 点 P = (1,1,f(1,1))
  \addplot3[
      only marks,
      mark=*,
      mark size=1.8pt,
      color=red,
  ]
  coordinates {(1,1,0.3)};
  
  % 梯度向量 (曲面上的上升方向)
  \addplot3[
      ->,
      very thick,
      red,
  ]
  coordinates {
      (1,1,0.3)
      (1.8,1.4,0.7)
  };
  
  \node[red] at (axis cs:1.05,1.85,0.75) {最陡方向};

  \addplot3[
    dashed,
    gray,
    thin,
]
coordinates {
    (1,1,0.3)
    (1,1,0)
};

  \addplot3[
    dashed,
    gray,
    thin,
]
coordinates {
    (1.8,1.4,0.7)
    (1.8,1.4,0)
};

\addplot3[
      only marks,
      mark=*,
      mark size=1.8pt,
      color=blue,
  ]
  coordinates {(1,1,0)};

% 三维向量的水平投影（灰色箭头，在 z=0 平面）
\addplot3[
    ->,
    very thick,
    blue,
]
coordinates {
    (1,1,0)
    (1.8,1.4,0)
};
\node[blue, above] at (axis cs:1.45,1.2,0) {$\nabla f$};
  
  \end{axis}
  \end{tikzpicture}
\end{center}

我们来看个例子：

\begin{myexample}
  求 $f(x, y) = x^2 + y^2$ 在点 $(1, 2)$ 处的梯度，以及函数值增加最快的方向和速率。

  $f_x = 2x$，$f_y = 2y$。在 $(1, 2)$ 处，$\nabla f = (2, 4)$。

  梯度方向 $(2, 4)$ 就是函数值增加最快的方向，增加速率为 $|\nabla f| = \sqrt{4 + 16} = 2\sqrt{5}$。
\end{myexample}

梯度还有一个非常重要的几何意义。设函数$f(x, y)$，平面$z = C$， 那么$f(x, y) = C$就是一条等高线，在等高线上的每一点的梯度 $\nabla f$ 都垂直于等高线。因为沿着等高线走，函数值不变，方向导数为 $0$，即 $\nabla f \cdot \vec{l} = 0$，所以 $\nabla f$ 与等高线的切线方向 $\vec{l}$ 垂直。这和我们之前讲“法向量垂直于平面”的思想完全一致。

想象你在爬山。等高线就是地图上那些闭合的圈。梯度方向就是最陡的上坡方向，它总是垂直于等高线，指向山顶。沿着等高线走，既不上升也不下降。

推广到三元函数 $w = f(x, y, z)$，梯度的定义完全类似，几何上，$\nabla f$ 垂直于等值面 $f(x, y, z) = C$。

\subsection{多元函数微分学的几何应用}

和一元函数类似，导数用来求切线斜率，推广到二元函数，偏导数和全微分就会在切线和切平面上大显身手。

我们先以空间曲线为例。设空间曲线以参数方程$\Gamma: 
\begin{cases}
x = x(t)\\
y = y(t)\\
z = z(t)
\end{cases}$给出，其中 $x(t), y(t), z(t)$ 可导且导数不同时为零。仿照一元函数，我们有：

\begin{theorem}{空间曲线的切向量}{}
若曲线 $\Gamma$ 的参数方程中 $x(t), y(t), z(t)$ 在 $t_0$ 处可导，且导数不全为零，则曲线 $\Gamma$在 $t_0$ 处的切向量为：
\[
\vec{T} = ( x'(t_0),\, y'(t_0),\, z'(t_0) )
\]
\end{theorem}

切向量的几何意义是显然的，当三个方向都取导数时，那么这个向量指向的就是曲线在该点的瞬时方向。也就是说，切向量是曲线在该点的切线的方向向量。已知一点和方向向量，我们就能写出直线：

\begin{theorem}{空间曲线的切线}{}
  若曲线 $\Gamma$ 的参数方程中 $x(t), y(t), z(t)$ 在 $t_0$ 处可导，且导数不全为零，则曲线在 $P_0$ 处的切线方程为：
\[
\frac{x-x_0}{x'(t_0)} = \frac{y-y_0}{y'(t_0)} = \frac{z-z_0}{z'(t_0)}
\]
切向量是切线的一个方向向量。
\end{theorem}

现在有了一点和切向量，我们根据平面的点法式，就能写出一张平面的方程：

\begin{theorem}{法平面}{}
若曲线 $\Gamma$在 $t_0$ 处存在切向量与切线，则曲线 $\Gamma$在 $t_0$ 处的法平面方程为：
\[
x'(t_0)(x-x_0) + y'(t_0)(y-y_0) + z'(t_0)(z-z_0) = 0
\]
法平面过曲线 $\Gamma$ 上一点 $P_0$ 且与该点切线垂直。
\end{theorem}

由于我们把切向量当作平面的法向量，用点法式写出了平面方程，所以切线自然垂直于这个平面。

有了切向量，切线，法平面，我们就能更具体地研究空间曲线了。

\begin{myexample}
求螺旋线 $\Gamma: 
\begin{cases}
x = \cos t\\
y = \sin t\\
z = \frac{t}{10}
\end{cases}$ 在 $t = \frac{\pi}{2}$ 处的切线与法平面方程。

直接对$t$求导，得切向量为 $\vec{T} = (-\sin t,\; \cos t,\; \frac{1}{10})$。在 $t=\frac{\pi}{2}$ 时，$\vec{T} = (-1,0,\frac{1}{10})$。

切线方程：
\[
\frac{x-0}{-1} = \frac{y-1}{0} = \frac{z-\frac{\pi}{20}}{\frac{1}{10}}
\]

这里 $y-1=0$ 一项单独写出，即： 
\[
\begin{cases} x + 10z = \frac{\pi}{2} \\ y = 1 \end{cases}
\]

法平面方程：
\[
-1\cdot(x-0) + 0\cdot(y-1) + \frac{1}{10} \cdot(z-\frac{\pi}{20}) = 0
\]

化简得： $10x - z + \frac{\pi}{20} = 0$。

从图象来看，就是：
\begin{center}
\begin{tikzpicture}[scale=2]
  \tdplotsetmaincoords{55}{135}
  \begin{scope}[tdplot_main_coords]

    % 坐标轴
    \draw[->,thick] (0,0,0) -- (1.5,0,0) node[below left]{$x$};
    \draw[->,thick] (0,0,0) -- (0,1.5,0) node[below right]{$y$};
    \draw[->,thick] (0,0,0) -- (0,0,2.5)   node[above]{$z$};

    % 螺旋线: x=cos(t), y=sin(t), z=t/10
    \draw[
      domain=0:6*pi,
      samples=400,
      smooth,
      variable=\t,
      red,
      thick
    ] plot ({cos(\t r)}, {sin(\t r)}, {\t/10}) node[black, right] {$\Gamma$};

    % 切点 P0 对应 t=π/2
    \coordinate (P0) at (0,1,{pi/20});
    \filldraw[black] (P0) circle[radius=0.8pt] node[above left] {$P_0$};

    % 切向量 T = (-1,0,1/10)，缩放以便观察
    \draw[->, blue, thick] (P0) -- ++({-0.6},0,{0.06}) node[below] {$\vec{T}$};

    % 法平面（半透明矩形）
    % 法向量为 T，两个垂直方向：u=(0,1,0), v=(1,0,10)
    % 验证垂直：u·T=0, v·T=-1+1=0
    % 矩形参数：P0 + a*u + b*v, a∈[-0.5,0.5], b∈[-0.08,0.08]
    \fill[blue!20, opacity=0.5, draw=blue!40, thick]
      ({-0.08}, {1-0.5}, {pi/20 - 0.8}) --
      ({0.08}, {1-0.5}, {pi/20 + 0.8}) --
      ({0.08}, {1+0.5}, {pi/20 + 0.8}) --
      ({-0.08}, {1+0.5}, {pi/20 - 0.8}) -- cycle;

  \end{scope}
\end{tikzpicture}
\end{center}

\end{myexample}

我们知道，空间曲线并不总是由参数方程确定，当空间曲线由两张空间曲面的交线的形式给出时，我们可以借助梯度来求切向量：

\begin{theorem}{曲面交线的切向量}{}
设空间曲线 $\Gamma$ 由方程组
\[
\begin{cases}
F(x,y,z) = 0 \\
G(x,y,z) = 0
\end{cases}
\]
确定，其中 $F,G$ 具有连续偏导数，且在点 $P_0$ 处 $\nabla F(P_0)$ 与 $\nabla G(P_0)$ 不平行，则 $\Gamma$ 在 $P_0$ 处的切向量为：
\[
\vec{T} = \nabla F(P_0) \times \nabla G(P_0)
\]
\end{theorem}

这是因为曲线同时落在两张曲面上，其切向量必须与两张曲面的法向量都垂直，而叉积正好给出这样一个向量。有了切向量，切线方程和法平面方程就能用刚刚的方法写出，此处不再赘述。

我们来看个例子：

\begin{myexample}
求曲线 $\begin{cases} x^2+y^2+z^2=6 \\ x+y+z=0 \end{cases}$ 在点 $(1,1,-2)$ 处的切线。

令 $F=x^2+y^2+z^2-6$，$G=x+y+z$。则：
\[
\nabla F = (2x,2y,2z),\quad \nabla G = (1,1,1)
\]

在 $(1,1,-2)$ 处，$\nabla F = (2,2,-4)$，$\nabla G = (1,1,1)$。因此：
\[
\vec{T} = \nabla F \times \nabla G =
\begin{vmatrix}
\vec{i} & \vec{j} & \vec{k} \\
2 & 2 & -4 \\
1 & 1 & 1
\end{vmatrix}
= (6, -6, 0)
\]

所以切线方程为 $\frac{x-1}{6} = \frac{y-1}{-6} = \frac{z+2}{0}$，或写为 $\begin{cases} x+y=2 \\ z = -2 \end{cases}$的形式。
\end{myexample}

研究完曲线，现在我们更进一步，从曲线转向曲面。先回忆一下梯度的几何意义，其中有一条性质：$\nabla F$ 垂直于等值面。那么，如果现在要求曲面法向量的话，你有没有什么思路呢？是的，我们可以直接用梯度求空间曲面的法向量：

\begin{theorem}{空间曲面的法向量}{}
设曲面 $\Sigma$ 由方程 $F(x,y,z)=0$ 确定，$F$ 在点 $P_0(x_0,y_0,z_0)$ 处可微且 $\nabla F(P_0)\neq \vec{0}$，则曲面在 $P_0$ 处的法向量为
\[
\vec{n} = \nabla F(P_0) = (F_x(P_0), F_y(P_0), F_z(P_0))
\]

特别地，当曲面以函数 $z=f(x,y)$ 形式给出时，令 $F(x,y,z)=f(x,y)-z$，则法向量为：
\[
\vec{n} = (f_x(x_0,y_0),\; f_y(x_0,y_0),\; -1)
\]
\end{theorem}

有了法向量，就说明我们能确定一条直线。显然，这条直线就是空间曲面的法线，直观上看，法向量的方向就是曲面在该点最“直”地离开曲面的方向，也就是法线的方向。

\begin{theorem}{空间曲面的法线}{}
  设曲面 $\Sigma$ 由方程 $F(x,y,z)=0$ 确定，$F$在 $P_0$ 处的法向量为$\vec{n} = \nabla F(P_0) = (F_x(P_0), F_y(P_0), F_z(P_0))$，则法线方程为：
\[
\frac{x-x_0}{F_x(P_0)} = \frac{y-y_0}{F_y(P_0)} = \frac{z-z_0}{F_z(P_0)}
\]

特别地，当曲面以函数 $z=f(x,y)$ 形式给出时，法线方程为：
\[
\frac{x-x_0}{f_x(x_0,y_0)} = \frac{y-y_0}{f_y(x_0,y_0)} = \frac{z-z_0}{-1}
\]
法向量是法线的一个方向向量。
\end{theorem}

现在有了一点和法向量，我们根据平面的点法式，就能写出一张平面的方程：

\begin{theorem}{切平面}{}
  设曲面 $\Sigma$ 由方程 $F(x,y,z)=0$ 确定，$F$在 $P_0$ 处的法向量为$\vec{n} = \nabla F(P_0) = (F_x(P_0), F_y(P_0), F_z(P_0))$，则切平面方程为：
\[
F_x(P_0) \cdot(x-x_0) + F_y(P_0) \cdot(y-y_0) + F_z(P_0) \cdot(z-z_0) = 0
\]

特别地，当曲面以函数 $z=f(x,y)$ 形式给出时，切平面方程为：
\[
z - z_0 = f_x(x_0,y_0)(x-x_0) + f_y(x_0,y_0)(y-y_0)
\]
切平面过曲面$\Sigma$一点$P_0$且与该点法线垂直。
\end{theorem}

要注意的是，切平面虽说是切平面，但是和曲面可能不止一个交点或一条交线。类比一元函数里的切线，有时与曲线的交点可能不止切点一个。

我们比较一下切平面方程与全微分$dz = f_x(x_0,y_0)\,dx + f_y(x_0,y_0)\,dy$就会发现，全微分 $dz$ 恰好是切平面上当自变量从 $(x_0,y_0)$ 变到 $(x_0+dx,y_0+dy)$ 时 $z$ 坐标的增量。这就说明，全微分是曲面在一点处近似的具体表现。也就是说，切平面在局部很好地逼近了曲面，就像切线逼近曲线一样。

\begin{myexample}
求旋转抛物面 $z = x^2 + y^2$ 在点 $(1,1,2)$ 处的切平面与法线方程。

由 $f(x,y) = x^2+y^2$ 可得 $f_x = 2x$，$f_y = 2y$。所以在 $(1,1)$ 处，$f_x=2, f_y=2$。

切平面方程为：
\[
z-2 = 2(x-1) + 2(y-1)
\]

化简得：$2x + 2y - z = 2$。

法线方程为：
\[
\frac{x-1}{2} = \frac{y-1}{2} = \frac{z-2}{-1}
\]
\end{myexample}

本质上，求法线和切平面的思想与一元函数求切线的思想是一脉相承的。

\section{多元函数的极值}
多元函数和一元函数类似，同样可以求极值和最值，不过情况有点不同。我们这里主要以二元函数为例，更多元的函数的研究方法是类似的。

\subsection{无条件极值}

在一元函数中，极值点就是曲线由上升转为下降的“山峰”，或由下降转为上升的“山谷”，在极值点处，切线是水平的，即导数等于零。现在我们要把这个思路搬到二元函数的曲面上。

我们先给出局部极值的精确定义。和一元情形完全类似，“局部”意味着只关心一点附近的小范围。

\begin{definition}{二元函数的极值}{}
设函数 $z = f(x, y)$ 在点 $P_0(x_0, y_0)$ 的某邻域内有定义。

若对该邻域内所有异于 $P_0$ 的点 $(x, y)$，都有：
\[
f(x, y) < f(x_0, y_0)
\]
则称 $f(x_0, y_0)$ 为 $f$ 的一个\textbf{极大值}，$P_0$ 称为\textbf{极大值点}。

若对该邻域内所有异于 $P_0$ 的点 $(x, y)$，都有：
\[
f(x, y) > f(x_0, y_0)
\]
则称 $f(x_0, y_0)$ 为 $f$ 的一个\textbf{极小值}，$P_0$ 称为\textbf{极小值点}。

极大值与极小值统称\textbf{极值}。
\end{definition}

想象一下，如果你站在曲面的一个“山峰”或“谷底”，那么该点的切平面必定是水平的。切平面水平，意味着两个偏导数都是零。这就引出了极值的必要条件。

\begin{theorem}{极值存在的必要条件}{}
若函数 $z = f(x, y)$ 在点 $(x_0, y_0)$ 处具有偏导数，且在该点取得极值，则必有
\[
f_x(x_0, y_0) = 0,\quad f_y(x_0, y_0) = 0
\]
\end{theorem}

我们把满足 $f_x(x_0, y_0)=0$ 且 $f_y(x_0, y_0)=0$ 的点称为\textbf{驻点}。和一元函数一样，驻点不一定是极值点。偏导数不存在的点也可能有极值，不过我们主要讨论可微的函数，这类异常点相对少见。

一个最典型的例子是双曲抛物面，也就是马鞍面。

\begin{myexample}
  求 $f(x, y) = xy$ 的极值。

  先求偏导：$f_x = y$，$f_y = x$。令$f_x=f_y=0$，解得$x=y=0$，所以原点是驻点。

  但它是极值点吗？我们取 $y=x$ 方向，$f(x, x) = x^2 \ge 0$，原点处取极小值 $0$。我们再取 $y=-x$ 方向，$f(x, -x) = -x^2 \le 0$，原点处又取极大值 $0$。
  
  原点一会是极小值，一会是极大值，那么它究竟是不是极值点呢？很遗憾，它并不是。$f(x, y) = xy$ 实际上是没有极值的。
\end{myexample}

这种既非极大也非极小的驻点，或者既是极大也是极小的驻点，我们称为\textbf{鞍点}。直观上看就是这样：

\begin{center}
  \begin{tikzpicture}
  \begin{axis}[
      view={20}{45},
      axis lines=center,           % 显示穿过原点的坐标轴
      axis line style={->},        % 轴末端带箭头
      xlabel={$x$}, ylabel={$y$}, zlabel={$z$},
      xtick=\empty,
      ytick=\empty,
      ztick=\empty,
      domain=-2.5:2.5,
      y domain=-2.5:2.5,
      samples=35,
      samples y=35,
      colormap/viridis,
  ]
  
  % 马鞍面
  \addplot3[
      surf,
      shader=interp,               % 平滑着色，无网格线
      opacity=0.65,
  ]
  {x^2 - y^2};

  \addplot3[
      only marks,
      mark=*,
      mark size=1.8pt,
      color=red,
  ]
  coordinates {(0,0,0)};
  
  \end{axis}
  \end{tikzpicture}
\end{center}


顾名思义，鞍点就是在马鞍中心的点。那么，如何区分驻点到底是极值点还是鞍点？我们先回忆一下，在一元函数里我们是怎么判断极值点是极大值点还是极小值点的？没错，看二阶导数的正负。二元函数也同样，只不过稍微复杂一点：

\begin{theorem}{极值存在的充分条件}{}
设 $z = f(x, y)$ 在驻点 $P_0(x_0, y_0)$ 的某邻域内具有连续的二阶偏导数。令
\[
A = f_{xx}(P_0),\quad B = f_{xy}(P_0),\quad C = f_{yy}(P_0)
\]
并记判别式 $\Delta = AC - B^2$，则有：

若 $\Delta > 0$，则 $f$ 在 $P_0$ 处取极值，且当 $A > 0$ 时为极小值，当 $A < 0$ 时为极大值。

若 $\Delta < 0$，则 $f$ 在 $P_0$ 处不取极值。

若 $\Delta = 0$，则该方法失效，需用其他办法判断。

\end{theorem}

这个判别法与一元函数的二阶导数判别法本质是一样的。

我们来看一个例子：

\begin{myexample}
  求函数 $f(x, y) = x^3 + y^3 - 3xy$ 的极值。

  首先求偏导，令偏导等于零，解方程组：
  \[
  \begin{cases}
    f_x = 3x^2 - 3y = 0\\
    f_y = 3y^2 - 3x = 0
  \end{cases}
  \]

  解得 $\begin{cases}
    y = x^2\\
    x = y^2
  \end{cases}$，代入得 $x = x^4$，即 $x(x^3 - 1) = 0$。解得两个驻点：$(0, 0)$ 和 $(1, 1)$。

计算二阶偏导数：
  \[
  f_{xx} = 6x,\quad f_{xy} = -3,\quad f_{yy} = 6y
  \]

逐点判别：

在 $(0, 0)$ 处：$A = 0$，$B = -3$，$C = 0$，$\Delta = 0\cdot0 - (-3)^2 = -9 < 0$。判别式小于零，所以 $(0, 0)$ 不是极值点。

在 $(1, 1)$ 处：$A = 6$，$B = -3$，$C = 6$，$\Delta = 6\cdot6 - (-3)^2 = 36-9=27 > 0$。又 $A=6>0$，所以 $f$ 在 $(1, 1)$ 处取得极小值，极小值为 $f(1, 1) = 1+1-3 = -1$。

\end{myexample}

总结一下找无条件极值的基本流程：

先求两个偏导数，找出所有驻点，然后用二阶判别法逐一检验。

\subsection{条件极值}

无条件极值是在定义域内部自由地找最大值最小值。但实际中常常有附加的约束条件。例如在周长固定的矩形中，什么时候面积最大？这时候自变量 $(x, y)$ 不能在整个定义域上随便取，而必须满足某个方程 $g(x, y) = 0$。这种带约束的极值问题，叫做\textbf{条件极值}。

最直接的想法是，从约束方程 $g(x, y) = 0$ 中解出一个变量，代入 $f(x, y)$，化为一元函数的无条件极值。比如约束是 $y = 2x$，直接代入 $f(x, 2x)$ 变成 $x$ 的一元函数就行。但如果约束方程很复杂，解不出 $y$ 关于 $x$ 的显式表达式。这时候要怎么办呢？

拉格朗日在他1788年的巨著《分析力学》中介绍了一个方法：

\begin{definition}{拉格朗日乘数法}{}
  设目标函数$f(x, y)$和$m$个约束函数$g_k(x, y)=0 \ (k=1,2,\dots,m)$均连续可微，构造\textbf{拉格朗日函数}：
  \[
  \mathcal{L}(x, y, \lambda_1, \lambda_2, \dots, \lambda_m)=f(x, y) +\sum_{k=1}^{m} \lambda_k g_k(x, y)
  \]
  其中$\lambda_1, \lambda_2, \dots, \lambda_m$称为\textbf{拉格朗日乘数}。

  若点$P(x_0,y_0)$是目标函数$f(x, y)$在约束$g_k(x, y)=0$下的极值点，则必然存在一组拉格朗日乘数$\lambda_1, \lambda_2, \dots, \lambda_m$，
  使得：
  \[
    \frac{\partial \mathcal{L}}{\partial x}=0, \quad
    \frac{\partial \mathcal{L}}{\partial y}=0, \quad
    \frac{\partial \mathcal{L}}{\partial \lambda_1}=0, \quad
    \frac{\partial \mathcal{L}}{\partial \lambda_2}=0, \quad
    \dots, \quad
    \frac{\partial \mathcal{L}}{\partial \lambda_m}=0
  \]
\end{definition}

拉格朗日乘数法的核心思想是构造一个新的辅助函数，然后对这个多变量函数求无条件极值。令拉格朗日函数对所有变量的偏导数全为零，解这个方程组，得到的解就是可能的条件极值点。也就是说，拉格朗日乘数法的巧妙之处在于把约束条件“吞”进了辅助函数的偏导数为零的条件里，把条件极值转换为了无条件极值。

为什么这样做是对的？我们从几何上看，约束 $g(x, y) = 0$ 是平面上的一条曲线。$f(x, y)$ 在这条曲线上取极值时，$f$ 的梯度 $\nabla f$ 和约束曲线的法向量 $\nabla g$ 必须是平行的，因为如果它们不平行，沿着曲线走一小步，$f$ 的值就会变化。因此极值点处 $\nabla f = -\lambda\, \nabla g$，而这正是$\frac{\partial \mathcal{L}}{\partial x}=0$和
    $\frac{\partial \mathcal{L}}{\partial y}=0$的含义。

我们来看开头的那个例子：

\begin{myexample}
  求周长为定值 $L$ 的矩形中面积的最大值。

  设矩形的长为 $x$、宽为 $y$，则面积 $S = xy$。周长约束为 $2x + 2y = L$，即 $x + y = \frac{L}{2}$。

  构造拉格朗日函数：
  \[
  \mathcal{L}(x, y, \lambda) = xy + \lambda (x + y - \frac{L}{2})
  \]

  求偏导：
  \[
  \frac{\partial \mathcal{L}}{\partial x} = y + \lambda = 0,\quad
  \frac{\partial \mathcal{L}}{\partial y} = x + \lambda = 0,\quad
  \frac{\partial \mathcal{L}}{\partial \lambda} = x + y - \frac{L}{2} = 0
  \]

  前两个方程得 $y = -\lambda$，$x = -\lambda$，即 $x = y$。

  代入约束，$2x = \frac{L}{2}$，得 $x = y = \frac{L}{4}$。此时面积最大，为 $S = \frac{L^2}{16}$。

\end{myexample}

  因此，在周长固定的矩形中，正方形面积最大。

\begin{myexample}
  求 $f(x, y) = x^2 + y^2$ 在约束 $x + y = 1$ 下的极值。

  约束方程写成 $g(x, y) = x + y - 1 = 0$。构造拉格朗日函数：
  \[
  \mathcal{L}(x, y, \lambda) = x^2 + y^2 + \lambda(x + y - 1)
  \]

  求偏导并令其为零：
  \[
  \frac{\partial \mathcal{L}}{\partial x} = 2x + \lambda = 0,\quad
  \frac{\partial \mathcal{L}}{\partial y} = 2y + \lambda = 0,\quad
  \frac{\partial \mathcal{L}}{\partial \lambda} = x + y - 1 = 0
  \]

  前两个方程得 $x = -\frac{\lambda}{2}$，$y = -\frac{\lambda}{2}$，即 $x = y$。
  
  代入第三个方程：$x + x = 1$，得 $x = y = \frac{1}{2}$，$\lambda = -1$。

  因此极值为 $f(\frac{1}{2}, \frac{1}{2}) = \frac{1}{2}$。
  
  几何上，这表示原点 $(0, 0)$ 到直线 $x + y = 1$ 的最短距离的平方。
\end{myexample}

拉格朗日乘数法可以推广到多个变量和多个约束的情形。不过我们主要讨论二元函数带一个约束。

我们用拉格朗日乘数法找到的极值点通常要进一步判断是极大值还是极小值。但在很多实际应用中，问题的背景已经暗示了是极大值还是极小值，即使不做严格判定，依据实际背景也能得到正确答案。

最后我们给一道不等式的题目。也许你高中的时候可能会觉得这种题目很难，但现在有了拉格朗日乘数法这个工具，你不妨用它来算一算。

\begin{myexample}
  已知非负实数 \(a, b, c\) 满足\(a + b + c = 1\)，求 \(a^3 + b^3 + c^3 + 6abc\) 的最大值。

我们在这里只给出不等式的解法。拉格朗日乘数法留给读者。

由$(a+b+c)^3 = a^3+b^3+c^3 + 3(a+b)(b+c)(c+a)$，代入 \(a+b+c=1\)，原式化为：

\[
a^3+b^3+c^3+6abc = 1 - 3(a+b)(b+c)(c+a) + 6abc
\]

由均值不等式，对非负实数有：$a+b \ge 2\sqrt{ab}$，$b+c \ge 2\sqrt{bc}$，$c+a \ge 2\sqrt{ca}$，相乘得：
\[
(a+b)(b+c)(c+a) \ge 8abc
\]

代入得：
\[
1 - 3(a+b)(b+c)(c+a) + 6abc \le 1 - 3 \cdot 8abc + 6abc = 1 - 18abc \le 1
\]

当且仅当\(a,b,c\) 中有两个为 \(0\)，另一个为\(1\)时两个等号同时成立。
\end{myexample}



% !TEX root = ../main.tex
\chapter{多重黎曼积分}
前面我们学了定积分 $\int_a^b f(x)\,dx$，它的核心思想可以概括成“分割、近似、求和、取极限”，表示的是曲边梯形的面积。然后我们又研究了多元函数，把视野从平面直角坐标系拓展到了空间直角坐标系。自然地，我们想到，一元函数的定积分是否也能拓展呢？
\section{二重积分}
和研究多元函数一样，我们先以二元函数为例。回忆一元函数定积分，积分区域是一段区间长度，那么，二元函数的积分区域就应该是一块面积了。由此我们就从一维到了二维，得到了二重积分。
\subsection{二重积分}
二重积分就是把定积分的思想从一维推广到二维，只不过具体的操作有些区别，积分区域不再是切一条线段，而是切一块平面区域，求和不再是用矩形近似，而是用小长方体近似。

设 $z = f(x, y)$ 是定义在平面区域 $D$ 上的非负函数，$z = f(x, y)$ 的图形是一张曲面，我们称曲面和区域 $D$ 之间的空间为“曲顶柱体”，意思就是这个空间是一个柱体，但是顶面是曲面。

仿照定积分的几何意义说明，我们自然地将它“升维”，给出二重积分的几何意义。

想象一下，在空间直角坐标系中，\(xOy\) 平面上有一个有界闭区域 \(D\)，上方盖着一张曲面 \(z = f(x,y)\)，以 \(D\) 的边界为准线，作母线平行于 \(z\) 轴的柱面，这个柱面与曲面以及 \(xOy\) 平面围成的立体，就是“曲顶柱体”。二重积分 \(\iint_D f(x,y)\,d\sigma\) 的几何意义就是这个曲顶柱体的体积。

\begin{center}
% 曲顶柱体：二重积分的几何意义
\pgfmathdeclarefunction{F}{2}{%
  \pgfmathparse{1.05 + 0.22*#1 + 0.12*#2 + 0.32*sin(70*#1)}%
}

\begin{tikzpicture}[
    scale=1.05,
    x={(1cm,0cm)},
    y={(0.55cm,0.32cm)},
    z={(0cm,1cm)},
    line join=round,
    line cap=round
]
    % 椭圆形底域 D 的参数
    \def\cx{2.15}
    \def\cy{1.15}
    \def\ra{1.55}
    \def\rb{0.78}

    % 坐标轴
    \draw[->] (-0.25,0,0) -- (4.45,0,0) node[below] {$x$};
    \draw[->] (0,-0.2,0) -- (0,2.85,0) node[right] {$y$};
    \draw[->] (0,0,0) -- (0,0,2.65) node[left] {$z$};
    \node[below left] at (0,0,0) {$O$};

    % 侧面 —— 用不透明的浅色填充，不加渐变
    \fill[cyan!20]   % 不透明，颜色自行调整
        plot[domain=0:360, samples=120, smooth]
        ({\cx + \ra*cos(\x)}, {\cy + \rb*sin(\x)},
         {F(\cx + \ra*cos(\x), \cy + \rb*sin(\x))})
        -- plot[domain=360:0, samples=120, smooth]
        ({\cx + \ra*cos(\x)}, {\cy + \rb*sin(\x)}, 0)
        -- cycle;

    % 母线 —— 用虚线表达柱面的弯曲，无光照感
    \foreach \ang in {0,1,...,359} {
        \pgfmathsetmacro\xx{\cx + \ra*cos(\ang)}
        \pgfmathsetmacro\yy{\cy + \rb*sin(\ang)}
        \pgfmathsetmacro\zz{F(\xx,\yy)}
        \draw[cyan!40, line width=0.4mm] (\xx,\yy,0) -- (\xx,\yy,\zz);
    }

    % 曲顶 z = f(x,y)
    \fill[blue!18, opacity=0.72]
        plot[domain=0:360, samples=120, smooth]
        ({\cx + \ra*cos(\x)}, {\cy + \rb*sin(\x)},
         {F(\cx + \ra*cos(\x), \cy + \rb*sin(\x))})
        -- cycle;

    % 底域 D（浅色填充，不透明）
    \draw[cyan!70!black, thick]
        plot[domain=195:360, samples=80, smooth]
        ({\cx + \ra*cos(\x)}, {\cy + \rb*sin(\x)}, 0);

    \draw[cyan!70!black, thick, dashed]
        plot[domain=0:195, samples=80, smooth]
    ({\cx + \ra*cos(\x)}, {\cy + \rb*sin(\x)}, 0);

    % 底面与顶面的轮廓线
    \draw[blue!75!black, very thick]
        plot[domain=0:360, samples=120, smooth]
        ({\cx + \ra*cos(\x)}, {\cy + \rb*sin(\x)},
         {F(\cx + \ra*cos(\x), \cy + \rb*sin(\x))});

         \node[blue!70!black, above] at (3.35,1.8,{F(3.35,1.8)}) {$z=f(x,y)$};
         
         \node[blue!70!black, above] at (3.2,1.4,0) {$D$};

\end{tikzpicture}
\begin{tikzpicture}[
    scale=1.05,
    x={(1cm,0cm)},
    y={(0.55cm,0.32cm)},
    z={(0cm,1cm)},
    line join=round,
    line cap=round
]
    % 椭圆形底域 D 的参数
    \def\cx{2.15}
    \def\cy{1.15}
    \def\ra{1.55}
    \def\rb{0.78}

    % 坐标轴
    \draw[->] (-0.25,0,0) -- (4.45,0,0) node[below] {$x$};
    \draw[->] (0,-0.2,0) -- (0,2.85,0) node[right] {$y$};
    \draw[->] (0,0,0) -- (0,0,2.65) node[left] {$z$};
    \node[below left] at (0,0,0) {$O$};

    % 侧面 —— 不透明浅色填充（无渐变）
    \fill[cyan!20]
        plot[domain=0:360, samples=120, smooth]
        ({\cx + \ra*cos(\x)}, {\cy + \rb*sin(\x)},
         {F(\cx + \ra*cos(\x), \cy + \rb*sin(\x))})
        -- plot[domain=360:0, samples=120, smooth]
        ({\cx + \ra*cos(\x)}, {\cy + \rb*sin(\x)}, 0)
        -- cycle;

    % 母线（密集细线，模拟柱面纹理）
    \foreach \ang in {0,1,...,359} {
        \pgfmathsetmacro\xx{\cx + \ra*cos(\ang)}
        \pgfmathsetmacro\yy{\cy + \rb*sin(\ang)}
        \pgfmathsetmacro\zz{F(\xx,\yy)}
        \draw[cyan!40, line width=0.4mm] (\xx,\yy,0) -- (\xx,\yy,\zz);
    }

    % 曲顶 z = f(x,y)
    \fill[blue!18, opacity=0.72]
        plot[domain=0:360, samples=120, smooth]
        ({\cx + \ra*cos(\x)}, {\cy + \rb*sin(\x)},
         {F(\cx + \ra*cos(\x), \cy + \rb*sin(\x))})
        -- cycle;

    % 底域 D 的轮廓（部分实线，部分虚线）
    \draw[cyan!70!black, thick]
        plot[domain=195:360, samples=80, smooth]
        ({\cx + \ra*cos(\x)}, {\cy + \rb*sin(\x)}, 0);
    \draw[cyan!70!black, thick, dashed]
        plot[domain=0:195, samples=80, smooth]
        ({\cx + \ra*cos(\x)}, {\cy + \rb*sin(\x)}, 0);

    % 顶面轮廓线（粗蓝）
    \draw[blue!75!black, very thick]
        plot[domain=0:360, samples=120, smooth]
        ({\cx + \ra*cos(\x)}, {\cy + \rb*sin(\x)},
         {F(\cx + \ra*cos(\x), \cy + \rb*sin(\x))});

    % ---------- 平行截平面 x = c 产生的分割线 ----------
    \foreach \xx in {0.9, 1.3, 1.7, 2.15, 2.6, 3.0, 3.4} {
        % 判别式
        \pgfmathsetmacro\disc{1 - (\xx-\cx)*(\xx-\cx)/(\ra*\ra)}
        \pgfmathparse{\disc > 0 ? sqrt(\disc) : 0}
        \pgfmathsetmacro\dyy{\rb * \pgfmathresult}   % 半长
        % 两个 y 值
        \pgfmathsetmacro\yyA{\cy - \dyy}             % 较小的 y（可见）
        \pgfmathsetmacro\yyB{\cy + \dyy}             % 较大的 y（不可见）
        \pgfmathsetmacro\zzA{F(\xx,\yyA)}
        \pgfmathsetmacro\zzB{F(\xx,\yyB)}
        % 实线：y 较小的竖线（可见侧面）
        \draw[thick, green!40!black] (\xx,\yyA,0) -- (\xx,\yyA,\zzA);
        % 虚线：y 较大的竖线（不可见侧面）
        \draw[thick, green!40!black, dashed] (\xx,\yyB,0) -- (\xx,\yyB,\zzB);
        % 底面弦（全虚线，被遮挡）
        \draw[thick, green!40!black, dashed] (\xx,\yyA,0) -- (\xx,\yyB,0);
        % 曲顶截线（全实线，视为可见）
        \draw[thick, green!40!black]
            plot[domain=\yyA:\yyB, samples=90, smooth]
            (\xx, \x, {F(\xx, \x)});
    }

    % ---------- 微元标注（典型小柱体，橙色） ----------
    \def\xs{1.60}
    \def\ys{1.01}
    \def\dx{0.58}
    \def\dy{0.46}
    \pgfmathsetmacro\zA{F(\xs,\ys)}
    \pgfmathsetmacro\zB{F(\xs+\dx,\ys)}
    \pgfmathsetmacro\zC{F(\xs+\dx,\ys+\dy)}
    \pgfmathsetmacro\zD{F(\xs,\ys+\dy)}
    % 四根竖线
    \draw[orange!90!black, ultra thick, dashed]
        (\xs,\ys,0) -- (\xs,\ys,\zA)
        ({\xs+\dx},\ys,0) -- ({\xs+\dx},\ys,\zB)
        ({\xs+\dx},{\ys+\dy},0) -- ({\xs+\dx},{\ys+\dy},\zC)
        (\xs,{\ys+\dy},0) -- (\xs,{\ys+\dy},\zD);
    % 顶面四边形（实线）
    \draw[orange!90!black, ultra thick]
        (\xs,\ys,\zA) -- ({\xs+\dx},\ys,\zB) --
        ({\xs+\dx},{\ys+\dy},\zC) -- (\xs,{\ys+\dy},\zD) -- cycle;
    % 底面矩形（虚线）
    \draw[orange!90!black, ultra thick, dashed]
        (\xs,\ys,0) -- ({\xs+\dx},\ys,0) --
        ({\xs+\dx},{\ys+\dy},0) -- (\xs,{\ys+\dy},0) -- cycle;
    % 样本点
    \pgfmathsetmacro\xmid{\xs+\dx/2}
    \pgfmathsetmacro\ymid{\ys+\dy/2}
    \pgfmathsetmacro\zmid{F(\xmid,\ymid)}
    \fill[orange!90!black] (\xmid,\ymid,\zmid) circle (1.4pt);
    \fill[orange!90!black] (\xmid,\ymid,0) circle (1.4pt);
    \draw[orange!80!black, dashed] (\xmid,\ymid,0) -- (\xmid,\ymid,\zmid);
    \node[orange!85!black, above left] at (\xmid,\ymid,\zmid+0.15) {$f(\xi_i,\eta_i)$};
    \node[orange!85!black, below] at (\xmid,\ymid,-0.2) {$(\xi_i,\eta_i)$};
    \node[orange!85!black, above] at (\xmid+0.1,\ymid,\zmid-0.05) {$\Delta\sigma_i$};

    % 曲顶标注
    \node[blue!70!black, above] at (3.35,1.8,{F(3.35,1.8)}) {$z=f(x,y)$};

    \node[blue!70!black, above] at (3.2,1.4,0) {$D$};
\end{tikzpicture}
\end{center}

\begin{definition}{二重积分}{}
设函数 \(f(x,y)\) 在有界闭区域 \(D\) 上有界。将 \(D\) 任意分成 \(n\) 个小闭区域
\[
\Delta\sigma_1,\ \Delta\sigma_2,\ \cdots,\ \Delta\sigma_n
\]
每个小区域既表示该区域本身，也表示该区域的面积。在每个小区域 \(\Delta\sigma_i\) 上任取一点 \((\xi_i,\eta_i)\)，作函数值 \(f(\xi_i,\eta_i)\) 与小区域面积 \(\Delta\sigma_i\) 的乘积 \(f(\xi_i,\eta_i) \cdot \Delta\sigma_i\)，并作出和
\[
S = \sum_{i=1}^{n} f(\xi_i,\eta_i) \Delta\sigma_i
\]

记 \(d_i\) 为小区域 \(\Delta\sigma_i\) 的直径，即该小区域内任意两点间距离的最大值，令 \(\lambda = \max\{d_1,d_2,\cdots,d_n\}\)。如果当 \(\lambda \to 0\) 时，\(S\) 的极限总存在，且与闭区域 \(D\) 的分法及点 \((\xi_i,\eta_i)\) 的取法无关，那么称 \(f(x,y)\) 在 \(D\) 上\textbf{可积}，并称这个极限 \(I\) 为函数 \(f(x,y)\) 在闭区域 \(D\) 上的\textbf{二重积分}，记作
\[
\iint_{D} f(x,y)\,d\sigma
\quad\text{或}\quad
\iint_{D} f(x,y)\,dxdy
\]
即：
\[
\iint_{D} f(x,y)\,d\sigma = \lim_{\lambda \to 0} \sum_{i=1}^{n} f(\xi_i,\eta_i) \Delta\sigma_i = I
\]

其中 \(f(x,y)\) 称为\textbf{被积函数}，\(f(x,y)\,d\sigma\) 称为\textbf{被积表达式}，\(x,y\) 称为\textbf{积分变量}，\(D\) 称为\textbf{积分区域}，$\iint$称为\textbf{二重积分号}。
\end{definition}

说人话就是，先把区域 $D$ 分成 $n$ 个小块 $\Delta\sigma_1, \Delta\sigma_2, \dots, \Delta\sigma_n$，每小块的面积记为 $\Delta\sigma_i$。在每小块上取一点 $(\xi_i, \eta_i)$，用 $f(\xi_i, \eta_i) \cdot \Delta\sigma_i$ 作为这一小条“细柱体”的体积。然后求和，记为$\sum_{i=1}^n f(\xi_i, \eta_i)\, \Delta\sigma_i$。再令所有小块的直径的最大值 $\lambda \to 0$，取极限$\lim_{\lambda \to 0} \sum_{i=1}^{n} f(\xi_i,\eta_i) \Delta\sigma_i$，最后得到的极限值就是二重积分的值。

既然二重积分的定义与定积分类似，那么二重积分也应该有一些和定积分类似的性质，由二重积分的定义可以得到：

\begin{theorem}{初等函数的可积性}{}
任何初等函数在其定义域内的任意有界闭区域上都是可积的。
\end{theorem}

\begin{theorem}{零面积区域积分}{}
若积分区域 \(D\) 的面积为 \(0\)，则
\[\iint_{D} f(x,y)\,d\sigma = 0\]
\end{theorem}

\begin{theorem}{线性性质}{}
若 \(f,g\) 在 \(D\) 上可积，\(\alpha,\beta \in \mathbb{R}\)，则 \(\alpha f+\beta g\) 在 \(D\) 上可积，且
\[\iint_{D} [\alpha f(x,y)+\beta g(x,y)]\,d\sigma = \alpha \iint_{D} f(x,y)\,d\sigma + \beta \iint_{D} g(x,y)\,d\sigma\]
\end{theorem}

\begin{theorem}{区域面积性质}{}
  若积分区域 \(D\) 的面积为 \(S\)，则
\[\iint_{D} d\sigma = S\]
\end{theorem}

\begin{theorem}{区域可加性}{}
若将 \(D\) 分为两个没有公共内点的子区域 \(D_1\) 与 \(D_2\)，则
\[\iint_{D} f(x,y)\,d\sigma = \iint_{D_1} f(x,y)\,d\sigma + \iint_{D_2} f(x,y)\,d\sigma\]
\end{theorem}

\begin{theorem}{对称性}{}
设 \(D\) 关于 \(y\) 轴对称，记 \(D_1\) 为 \(x \geq 0\) 的部分。

若 \(f(x,y)\) 关于 \(x\) 为奇函数，则
\[\iint_{D} f(x,y)\,d\sigma = 0\]

若 \(f(x,y)\) 关于 \(x\) 为偶函数，则
\[\iint_{D} f(x,y)\,d\sigma = 2 \iint\limits_{D_1} f(x,y)\,d\sigma\]

设 \(D\) 关于 \(x\) 轴对称，记 \(D_2\) 为 \(y \geq 0\) 的部分。

若 \(f(x,y)\) 关于 \(y\) 为奇函数，则
\[\iint_{D} f(x,y)\,d\sigma = 0\]

若 \(f(x,y)\) 关于 \(y\) 为偶函数，则
\[\iint_{D} f(x,y)\,d\sigma = 2 \iint_{D_2} f(x,y)\,d\sigma\]
\end{theorem}

\begin{theorem}{保序性}{}
若在 \(D\) 上 \(f(x,y) \leq g(x,y)\)，则
\[\iint_{D} f(x,y)\,d\sigma \leq \iint_{D} g(x,y)\,d\sigma\]
\end{theorem}

这些性质都和定积分是类似的，只不过是把区间换成了区域，面积换成了体积。

\begin{myexample}
  求 $\iint_D -3\, d\sigma$，其中 $D$ 是以原点为圆心、半径为 $R$ 的圆。

  被积函数为$z(x,y)=-3$，也就是说，这个二重积分求的曲顶柱体的顶面是平面$z(x,y)=-3$，底面区域是圆，面积为 $\pi R^2$。这时曲顶柱体就退化成圆柱体，体积等于底面积乘高。

  所以：

  \[
  \iint_D -3\, d\sigma=-3 \cdot \pi R^2=-3\pi R^2
  \]
\end{myexample}

二重积分和定积分一样，计算结果取的是代数体积。也就是在区域 $D$ 上方为正，在区域 $D$ 下方为负。

直接用定义来计算二重积分显然不太可能，和定积分一样，我们也有二重积分的计算法。

\subsection{二重积分的直角坐标形式}

面对二重积分，我们要怎么计算曲顶柱体的体积呢？让我们想一想，要积出来体积，可以先积一个面积，然后再对面积积分就是体积了。那要怎么积面积呢？没错，就是定积分。所以，二重积分完全可以转化为两次定积分来算。这种“先积一个，再积另一个”的做法，就叫累次积分。这种把未知转化为已知的思想，就叫“化归”思想。

如果区域 \(D\) 可以写成夹在两条竖直线 \(x=a\) 与 \(x=b\) 之间，且上下两条边分别是曲线 \(y=\varphi_1(x)\) 和 \(y=\varphi_2(x)\)的形式，即：
\[
D=\{(x,y)\mid a\le x\le b,\ \varphi_1(x)\le y\le\varphi_2(x)\}
\]

那么我们习惯上把这种区域称为 $X$ 型区域。这时二重积分可以化成先对 \(y\) 后对 \(x\) 的二次积分：

\begin{theorem}{$X$ 型区域的二重积分}{}
    若区域 \(D\) 是 $X$ 型区域，则在区域 \(D\) 上的二重积分可化为：
    \[
\iint_D f(x,y)\,dxdy = \int_a^b\left[\int_{\varphi_1(x)}^{\varphi_2(x)} f(x,y)\,dy\right]dx = \int_a^b\ \int_{\varphi_1(x)}^{\varphi_2(x)} f(x,y)\,dydx
\]  
\end{theorem}


\begin{center}
    \begin{tikzpicture}
    \begin{scope}[shift={(0,0)},
        declare function={
            a = 1.2; b = 3.8;
            y1(\x) = 1.5 + 0.5*sin((\x-2.5)*160);
            y2(\x) = 3.5 + 0.3*sin((\x-2.5)*180);
        }]
        % 坐标轴
        \draw[->, thick] (-0.3,0) -- (4.5,0) node[right] {\(x\)};
        \draw[->, thick] (0,-0.3) -- (0,4.5) node[above] {\(y\)};
        \node[below left] at (0,0) {\(O\)};
    
        % 填充区域（包含左右竖直线段）
        \fill[gray!30]
            (a, {y1(a)}) --
            plot[domain=a:b, smooth] (\x, {y1(\x)}) --
            (b, {y2(b)}) --
            plot[domain=b:a, smooth] (\x, {y2(\x)}) --
            cycle;
    
        % 下边界曲线 + 右端点标签
        \draw[thick]
            plot[domain=a:b, smooth] (\x, {y1(\x)})
            node[below right, xshift=-2pt, yshift=-4pt] {\(y=\varphi_1(x)\)};
        % 上边界曲线 + 右端点标签
        \draw[thick]
            plot[domain=a:b, smooth] (\x, {y2(\x)})
            node[above right, xshift=-2pt, yshift=4pt] {\(y=\varphi_2(x)\)};
    
        % 左右竖直实线：连接上下端点
        \draw[thick] (a, {y1(a)}) -- (a, {y2(a)});
        \draw[thick] (b, {y1(b)}) -- (b, {y2(b)});
    
        % 虚线：x轴 → 上方曲线端点
        \draw[dashed] (a,0) -- (a, {y2(a)});
        \draw[dashed] (b,0) -- (b, {y2(b)});
        \node[below] at (a,0) {\(a\)};
        \node[below] at (b,0) {\(b\)};
    
        \node at (2.5, 2.8) {\(D\)};
    \end{scope}
    \begin{scope}[shift={(7,0)},
        declare function={
            a = 1.2; b = 3.8; k = 0.7;
            yLo(\x) = 2.0 - k*sin((\x-a)*180/(b-a));
            yHi(\x) = 2.0 + k*sin((\x-a)*180/(b-a));
        }]
        % 坐标轴
        \draw[->, thick] (-0.3,0) -- (4.5,0) node[right] {\(x\)};
        \draw[->, thick] (0,-0.3) -- (0,4.5) node[above] {\(y\)};
        \node[below left] at (0,0) {\(O\)};
    
        % 填充区域
        \fill[gray!30]
            (a, {yLo(a)}) --
            plot[domain=a:b, smooth] (\x, {yLo(\x)}) --
            (b, {yHi(b)}) --
            plot[domain=b:a, smooth] (\x, {yHi(\x)}) --
            cycle;
    
        % 下边界曲线 + 右端点标签
        \draw[thick]
            plot[domain=a:b, smooth] (\x, {yLo(\x)})
            node[below right, xshift=-2pt, yshift=-6pt] {\(y=\varphi_1(x)\)};
        % 上边界曲线 + 右端点标签
        \draw[thick]
            plot[domain=a:b, smooth] (\x, {yHi(\x)});
            
            \node[above, xshift=2pt, yshift=12pt] at (a, {yHi(a)}) {\(y=\varphi_2(x)\)};
    
        % 左右竖直实线：连接上下端点（此处上下端点重合，实线长度为零，可省略，但为保证一致性仍然画出）
        \draw[thick] (a, {yLo(a)}) -- (a, {yHi(a)});
        \draw[thick] (b, {yLo(b)}) -- (b, {yHi(b)});
    
        % 虚线：x轴 → 上方曲线端点
        \draw[dashed] (a,0) -- (a, {yHi(a)});
        \draw[dashed] (b,0) -- (b, {yHi(b)});
        \node[below] at (a,0) {\(a\)};
        \node[below] at (b,0) {\(b\)};
    
        \node at (2.5, 2) {\(D\)};
    \end{scope}
    \end{tikzpicture}
\end{center}

在区间 \([a,b]\) 上任意取定一点 \(x_0\)，作平行于 \(yOz\) 面的平面\(x=x_0\)，这平面截曲顶柱体所得的截面是一个以区间 \([\varphi_1(x_0),\varphi_2(x_0)]\) 为底、曲线\(z=f(x_0,y)\) 为曲边的曲边梯形，所以这截面的面积为$\int_{\varphi_1(x_0)}^{\varphi_2(x_0)} f(x_0,y)\,dy$，那么任意$x$处的截面面积就是$\int_{\varphi_1(x)}^{\varphi_2(x)} f(x,y)\,dy$，再对$x$积分，把面积累加起来，就得到曲顶柱体体积，也就是$\int_a^b\left[\int_{\varphi_1(x)}^{\varphi_2(x)} f(x,y)\,dy\right]dx$。

\begin{center}
    \begin{tikzpicture}[
      % 三维坐标在纸面上的斜投影方向
      x={(1.10cm,0cm)},
      y={(-0.68cm,0.48cm)},
      z={(0cm,1.08cm)},
      >=stealth,
      line cap=round,
      line join=round,
    ]
    
    % 关键点：下标 1、2 分别对应 y=phi_1(x)、y=phi_2(x)
    \coordinate (O)   at (0,0,0);
    \coordinate (A1)  at (1.80,0.35,0);
    \coordinate (A2)  at (1.80,2.25,0);
    \coordinate (M1)  at (3.55,0.25,0);
    \coordinate (M2)  at (3.55,2.45,0);
    \coordinate (B1)  at (5.25,0.50,0);
    \coordinate (B2)  at (5.25,2.10,0);
    
    % 上方曲面边界上的对应点
    \coordinate (A1t) at (1.80,0.35,2.35);
    \coordinate (A2t) at (1.80,2.25,2.70);
    \coordinate (M1t) at (3.55,0.25,2.55);
    \coordinate (M2t) at (3.55,2.45,2.95);
    \coordinate (B1t) at (5.25,0.50,2.30);
    \coordinate (B2t) at (5.25,2.10,2.65);
    
    % 坐标轴
    \draw[->,thick] (O) -- (6.35,0,0) node[below right=-1pt] {$x$};
    \draw[->,thick] (O) -- (0,3.15,0) node[above left=-1pt] {$y$};
    \draw[->,thick] (O) -- (0,0,4.25) node[left] {$z$};
    \node[below left] at (O) {$O$};
    
    % x 轴上的标记
    \node[below=5pt] at (1.80,0,0) {$a$};
    \node[below=5pt] at (3.55,0,0) {$x_0$};
    \node[below=5pt] at (5.25,0,0) {$b$};
    
    % 底面区域的边界
    \draw[thick]
      (A1) .. controls (2.45,0.20,0) and (2.95,0.16,0) .. (M1)
           .. controls (4.20,0.25,0) and (4.75,0.36,0) .. (B1);
    \draw[thick,dashed]
      (A2) .. controls (2.45,2.48,0) and (2.95,2.52,0) .. (M2)
           .. controls (4.15,2.42,0) and (4.72,2.26,0) .. (B2);
    \draw[thick] (A1) -- (A2);
    \draw[thick,dashed] (B1) -- (B2);
    
    % 背后的竖直母线
    \draw[thick] (A2) -- (A2t);
    \draw[thick,dashed] (B2) -- (B2t);
    
    
    
    
    % 两侧的竖直边界
    \draw[thick] (A1) -- (A1t);
    \draw[thick] (B1) -- (B1t);
    
    % 文字标注及引线
    \draw[thick]
      (4.15,1.32,3.10) -- (5.35,0.72,3.52)
      node[right] {$z=f(x,y)$};
    \draw[thick]
      (4.15,2.33,0) -- (5.55,1.55,1.00)
      node[right] {$y=\varphi_2(x)$};
    \draw[thick]
      (4.70,0.42,0) -- (5.72,0.20,0.35)
      node[right] {$y=\varphi_1(x)$};
    
    % 上方曲面的轮廓
    % 第一条光滑曲线 A1t → M1t → B1t
    \draw[thick]
      (A1t) .. controls (2.40,0.24,2.62) and (2.95,0.20,2.68) .. (M1t)
            .. controls (4.15,0.30,2.42) and (4.75,0.40,2.42) .. (B1t);
    
    \draw[thick,dashed]
      (A2t) .. controls (2.98,2.54,2.90) and (3.55,2.45,2.78) .. (B2t);
    
    % 第二条光滑曲线 A2t → M2t → B2t
    \draw[thick]
      (A2t) .. controls (2.42,2.48,2.95) and (2.98,2.54,3.08) .. (M2t)
            .. controls (4.12,2.36,2.82) and (4.73,2.25,2.82) .. (B2t);
    
    \draw[thick]
      (A1t) .. controls (1.80,0.82,3.05) and (1.80,1.32,3.13) .. (A2t);
    \draw[thick]
      (B1t) .. controls (5.25,0.90,2.76) and (5.25,1.65,2.78) .. (B2t);
    
    % x=x_0 处的截面 A(x_0)
    % 先填充，再描边；透明度可以改为 1 得到实心灰色
    \fill[gray!55,opacity=.62]
      (M1) -- (M2) -- (3.52,2.40,2.79) .. controls (3.53,2.42,2.91) .. (3.7,2.45,2.91) -- (3.7,2.45,2.91)
      .. controls (3.55,1.85,3.12) and (3.55,0.82,3.35) .. (M1t)
      -- cycle;
    \draw[thick,dashed] (M1) -- (M2);
    \draw[thick] (M1) -- (M1t);
    \draw[thick,dashed] (M2) -- (3.52,2.40,2.79);
    \draw[thick]
      (3.7,2.45,2.91) .. controls (3.55,1.85,3.12) and (3.55,0.82,3.35) .. (M1t);
    \draw[thick,dashed] (3.7,2.45,2.91) .. controls (3.53,2.42,2.91) .. (3.52,2.40,2.79);
    \node at (3.55,1.22,1.48) {$x=x_0$};
    
    
    \draw[thick,dashed] (1.80,0.35,0) -- (1.80,0,0);
    
    \draw[thick,dashed] (5.25,0.50,0) -- (5.25,0,0);
    
    \draw[thick,dashed] (3.55,0.25,0) -- (3.55,0,0);
    
    \end{tikzpicture}
\end{center}

如果区域 \(D\) 可以写成夹在两条水平线 \(y=c\) 与 \(y=d\) 之间，且左右两边分别是曲线 \(x=\psi_1(y)\) 和 \(x=\psi_2(y)\)的形式，即：
\[
D=\{(x,y)\mid c\le y\le d,\ \psi_1(y)\le x\le\psi_2(y)\}
\]

那么我们习惯上把这种区域称为 $Y$ 型区域。这时二重积分可以化成先对 \(x\) 后对 \(y\) 的二次积分：

\begin{theorem}{$Y$ 型区域的二重积分}{}
    若区域 \(D\) 是 $Y$ 型区域，则在区域 \(D\) 上的二重积分可化为：
\[
\iint_D f(x,y)\,dxdy = \int_c^d\left[ \int_{\psi_1(y)}^{\psi_2(y)} f(x,y)\,dx\right]dy = \int_c^d\ \int_{\psi_1(y)}^{\psi_2(y)} f(x,y)\,dxdy
\]
\end{theorem}

\begin{center}
\begin{tikzpicture}
% ========= 左图：一般 y 型区域 =========
\begin{scope}[shift={(0,0)},
    declare function={
        c = 1.2; d = 3.8;
        psi1(\y) = 1.5 + 0.35*sin((\y-2.5)*160);
        psi2(\y) = 3.5 + 0.35*cos((\y-2.5)*160);
    }]
    % 坐标轴
    \draw[->, thick] (-0.3,0) -- (4.5,0) node[right] {\(x\)};
    \draw[->, thick] (0,-0.3) -- (0,4.5) node[above] {\(y\)};
    \node[below left] at (0,0) {\(O\)};

    % 填充区域
    \fill[gray!30]
        ({psi1(c)}, c) --
        plot[domain=c:d, smooth] ({psi1(\x)}, \x) --
        ({psi2(d)}, d) --
        plot[domain=d:c, smooth] ({psi2(\x)}, \x) --
        cycle;

    % 左边界曲线 + 标签（放在上端点）
    \draw[thick]
        plot[domain=c:d, smooth] ({psi1(\x)}, \x)
        node[above left, xshift=16pt, yshift=-32pt] {\(x=\psi_1(y)\)};
    % 右边界曲线 + 标签（放在上端点）
    \draw[thick]
        plot[domain=c:d, smooth] ({psi2(\x)}, \x)
        node[above right, xshift=-2pt, yshift=-12pt] {\(x=\psi_2(y)\)};

    % 上下水平实线（连接左右端点）
    \draw[thick] ({psi1(c)}, c) -- ({psi2(c)}, c);
    \draw[thick] ({psi1(d)}, d) -- ({psi2(d)}, d);

    % 虚线：从 y 轴到右曲线端点，标出 c 和 d
    \draw[dashed] (0,c) -- ({psi2(c)}, c);
    \draw[dashed] (0,d) -- ({psi2(d)}, d);
    \node[left] at (0,c) {\(c\)};
    \node[left] at (0,d) {\(d\)};

    % 区域标签
    \node at (2.5, 2.5) {\(D\)};
\end{scope}

% ========= 右图：上下端点重合的 y 型区域 =========
\begin{scope}[shift={(7,0)},
    declare function={
        c = 1.2; d = 3.8; k = 0.8;
        psiLo(\y) = 2.5 - k*sin((\y-c)*180/(d-c));
        psiHi(\y) = 2.5 + k*sin((\y-c)*180/(d-c));
    }]
    % 坐标轴
    \draw[->, thick] (-0.3,0) -- (4.5,0) node[right] {\(x\)};
    \draw[->, thick] (0,-0.3) -- (0,4.5) node[above] {\(y\)};
    \node[below left] at (0,0) {\(O\)};

    % 填充区域
    \fill[gray!30]
        ({psiLo(c)}, c) --
        plot[domain=c:d, smooth] ({psiLo(\x)}, \x) --
        ({psiHi(d)}, d) --
        plot[domain=d:c, smooth] ({psiHi(\x)}, \x) --
        cycle;

    % 左边界曲线 + 标签（放在上端点）
    \draw[thick]
        plot[domain=c:d, smooth] ({psiLo(\x)}, \x)
        node[above left, xshift=-16pt, yshift=-24pt] {\(x=\psi_1(y)\)};
    % 右边界曲线 + 标签（放在下端点，仿照 x 型对应图的标注方式）
    \draw[thick]
        plot[domain=c:d, smooth] ({psiHi(\x)}, \x);
    \node[above, xshift=48pt, yshift=12pt] at ({psiHi(c)}, c) {\(x=\psi_2(y)\)};

    % 上下水平实线：连接左右端点（此处端点重合，长度为零，可省略，保留以保持一致性）
    \draw[thick] ({psiLo(c)}, c) -- ({psiHi(c)}, c);
    \draw[thick] ({psiLo(d)}, d) -- ({psiHi(d)}, d);

    % 虚线：从 y 轴到重合点，标出 c 和 d
    \draw[dashed] (0,c) -- ({psiHi(c)}, c);
    \draw[dashed] (0,d) -- ({psiHi(d)}, d);
    \node[left] at (0,c) {\(c\)};
    \node[left] at (0,d) {\(d\)};

    % 区域标签
    \node at (2.5, 2.5) {\(D\)};
\end{scope}
\end{tikzpicture}
\end{center}

具体的操作和$X$ 型区域的二重积分类似，不再赘述。

先对哪个自变量积分，关键是积分区域。如果是$ X $型区域，就最后对$x$积分，先对$y$积分。如果是$ Y $型区域，就最后对$y$积分，先对$x$积分。也就是说，是什么区域就最后积哪个。

如果区域本身既不是$ X $型也不是$ Y $型，就用平行于坐标轴的直线把它切成几块，每块单独处理，再利用积分的区域可加性加起来就行了。

一般地，当我们写$\int_a^b\ \int_{\varphi_1(x)}^{\varphi_2(x)} f(x,y)\,dxdy$或者$\int_c^d\ \int_{\psi_1(y)}^{\psi_2(y)} f(x,y)\,dxdy$的时候，表示的其实是二重积分。形式上看着像是先对$\int_{\varphi_1(x)}^{\varphi_2(x)} f(x,y)\,dxdy$或者$\int_{\psi_1(y)}^{\psi_2(y)} f(x,y)\,dxdy$积分，再从$a$积到$b$或者从$c$积到$d$。然而这样其实是不合理的，因为第二个积分没有积分变量。

二重积分转换为两次定积分之后，就要处理这两个定积分的上下限。在实际应用中，我们通常用“画图穿线法”来确定它们的上下限。比如区域是$ X $型，就用一根平行于 \(y\) 轴的直线从下往上穿过区域，穿进去的边界就是下边界 \(y=\varphi_1(x)\)，穿出来的边界就是上边界 \(y=\varphi_2(x)\)，然后把这条线从左扫到右，就得到 \(x\) 的范围 \([a,b]\)。$Y $型区域也是类似的。

这么说肯定很难理解，我们看一个例子就明白了：

\begin{myexample}
    计算二重积分$\iint_D x \, d\sigma$，其中区域$D$是由$y$ 轴，直线$y=1$以及直线$x=y$围成的区域。

    我们首先要判断区域$D$是什么类型的区域。我们把直线$y=1$和直线$x=y$画出来，会发现它们与$y$ 轴围成了一个三角形。这个区域看成$X$ 型区域或者$Y$ 型区域都可以，这里作为示例，我们把它看成$Y$ 型区域计算。

    接下来就是“画图穿线法”的操作。判断区域为$Y$ 型，我们就要用一根平行于 \(x\) 轴的直线从左往右穿过这个区域，那么这根直线和区域$D$一定会有两个交点，我们分别记为$A$和$B$。然后比较大小，在左面的点小，作为积分下限，这里是\(x = 0\)。在右面的点大，作为积分上限，这里是\(x = y\)。

 \begin{center}
    \begin{tikzpicture}[x=4cm, y=4cm]
        % 坐标轴
        \draw[->] (-0.1,0) -- (1.2,0) node[right] {$x$};
        \draw[->] (0,-0.1) -- (0,1.2) node[above] {$y$};
    
        % 填充区域 D: 0 <= x <= 1, x <= y <= 1
        \fill[gray!30] (0,0) 
            -- plot[domain=0:1, smooth] (\x, \x) 
            -- (1,1) 
            -- plot[domain=1:0, smooth] (\x, 1) 
            -- cycle;
    
        % 下边界曲线 y = x
        \draw[thick, domain=0:1, smooth] plot (\x, \x);
        \node[below right] at (0.5,0.5) {$x = y$};
        % 上边界直线 y = 1
        \draw[thick] (0,1) -- (1,1);
        \node[above left] at (0.5,1) {$y=1$};
        \draw[thick] (0,0) -- (0,1);
    
        % 右边界虚线及积分限 a, b
        \draw[dashed] (1,0) -- (1,1);
        \node[below] at (1,0) {$b$};
        \node[below right] at (0,0) {$a$};
    
        % 区域标签
        \node at (0.5,0.75) {$D$};
    \end{tikzpicture}
    \begin{tikzpicture}[x=4cm, y=4cm]
        % 坐标轴
        \draw[->] (-0.1,0) -- (1.2,0) node[right] {$x$};
        \draw[->] (0,-0.1) -- (0,1.2) node[above] {$y$};
    
        % 填充区域 D: 0 <= x <= 1, x <= y <= 1
        \fill[gray!30] (0,0) 
            -- plot[domain=0:1, smooth] (\x, \x) 
            -- (1,1) 
            -- plot[domain=1:0, smooth] (\x, 1) 
            -- cycle;
    
        % 下边界曲线 y = x
        \draw[thick, domain=0:1, smooth] plot (\x, \x);
        \node[below right] at (0.5,0.5) {$x = y$};
        % 上边界直线 y = 1
        \draw[thick] (0,1) -- (1,1);
        \node[above left] at (0.5,1) {$y=1$};
        \draw[thick] (0,0) -- (0,1);
    
        % 右边界虚线及积分限 a, b
        \draw[dashed] (1,0) -- (1,1);
        \node[below] at (1,0) {$b$};
        \node[below right] at (0,0) {$a$};
        
    \draw[thick] (-0.1, 0.6) -- (1.1, 0.6);
    \node[above left] at (0,0.6) {$A$};
    \fill (0,0.6) circle (1.5pt);
    \node[above] at (0.6,0.6) {$B$};
    \fill (0.6,0.6) circle (1.5pt);
    
        % 区域标签
        \node at (0.3,0.5) {$D$};
    \end{tikzpicture}
\end{center}
   
那么我们就确定了上下限，得到第一重积分$\int_{0}^y x \, dx$，现在判断区间长度。这个区间从左扫到右，也就是$[0,1]$。由此我们就确定了区间，也得到了第二重积分$\int_0^1\ \int_{0}^y x \, dxdy$。那么我们现在就可以计算了。

先算第一重积分，这是一个对$x$积分的关于$y$的积分上限函数，积出来就是：

$$\int_{0}^y x \, dx=\left[\frac{1}{2}x^2\right]_0^y=\frac{1}{2}y^2$$

再代入第二重积分：

$$\int_0^1 \frac{1}{2}y^2 dy =\left[ \frac{1}{6}y^3\right]_0^1=\frac{1}{6}$$


所以二重积分$\iint_D x \, d\sigma=\frac{1}{6}$。
\end{myexample}


\begin{myexample}
    计算二重积分 $\iint_D (x+y) \, d\sigma$，其中区域$D$为抛物线\(y = x^2\) 和直线 \(y = x\) 围成的区域。

    我们首先先判断区域$D$是什么类型的区域。我们把抛物线\(y = x^2\) 和直线 \(y = x\) 画出来，就会发现它们围成的图形类似于一个弓形。这个区域看成$X$ 型区域或者$Y$ 型区域都可以，但是这里既然已经给出了\(y = x^2\) 和 \(y = x\)，为了方便，我们就把区域$D$看成$X$ 型区域。
    
    接下来就是“画图穿线法”的操作。判断区域为$X$ 型，我们就要用一根平行于 \(y\) 轴的直线从下往上穿过这个区域，那么这根直线和区域$D$一定会有两个交点，我们分别记为$A$和$B$。然后比较大小，在下面的点小，作为积分下限，这里是\(y = x^2\)。在上面的点大，作为积分上限，这里是\(y = x\)。

\begin{center}
\begin{tikzpicture}[x=4cm, y=4cm]
    % 坐标轴
    \draw[->] (-0.1,0) -- (1.2,0) node[right] {$x$};
    \draw[->] (0,-0.1) -- (0,1.2) node[above] {$y$};

    % 填充 X 型区域 D: 0 <= x <= 1, x^2 <= y <= x
    \fill[gray!30] (0,0) 
        -- plot[domain=0:1, smooth] (\x, {\x*\x}) 
        -- (1,1) 
        -- plot[domain=1:0, smooth] (\x, {\x}) 
        -- cycle;

    % 下边界曲线 y = x^2
    \draw[thick, domain=0:1, smooth] plot (\x, {\x*\x});
        \node[below right] at (0.5,0.25) {$y=x^2$};
    % 上边界曲线 y = x
    \draw[thick, domain=0:1, smooth] plot (\x, {\x});
        \node[above left] at (0.5,0.5) {$y=x$};

    % 右边界虚线及 a, b 标注
    \draw[dashed] (1,0) -- (1,1);
    \node[below] at (1,0) {$b$};
    \node[below right] at (0,0) {$a$};

    % 区域标签
    \node at (0.5,0.4) {$D$};
\end{tikzpicture}
\begin{tikzpicture}[x=4cm, y=4cm]
    % 坐标轴
    \draw[->] (-0.1,0) -- (1.2,0) node[right] {$x$};
    \draw[->] (0,-0.1) -- (0,1.2) node[above] {$y$};

    % 填充 X 型区域 D: 0 <= x <= 1, x^2 <= y <= x
    \fill[gray!30] (0,0) 
        -- plot[domain=0:1, smooth] (\x, {\x*\x}) 
        -- (1,1) 
        -- plot[domain=1:0, smooth] (\x, {\x}) 
        -- cycle;

    % 下边界曲线 y = x^2
    \draw[thick, domain=0:1, smooth] plot (\x, {\x*\x});
        \node[below right] at (0.5,0.25) {$y=x^2$};
    % 上边界曲线 y = x
    \draw[thick, domain=0:1, smooth] plot (\x, {\x});
        \node[above left] at (0.5,0.5) {$y=x$};

    % 右边界虚线及 a, b 标注
    \draw[dashed] (1,0) -- (1,1);
    \node[below] at (1,0) {$b$};
    \node[below right] at (0,0) {$a$};
    
    \draw[thick] (0.6, -0.1) -- (0.6, 1.1);
    \node[right] at (0.6, 0.6) {$A$};
    \fill (0.6,0.6) circle (1.5pt);
    \node[right] at (0.6, 0.36) {$B$};
    \fill (0.6, 0.36) circle (1.5pt);

    % 区域标签
    \node at (0.5,0.4) {$D$};
\end{tikzpicture}
\end{center}

那么我们就确定了上下限，得到第一重积分$\int_{x^2}^x (x+y) \, dy$，现在判断区间长度。这个区间从左扫到右，其实就是\(y = x^2\) 和 \(y = x\)的交点，也就是$[0,1]$。由此我们就确定了区间，也得到了第二重积分$\int_0^1\ \int_{x^2}^x (x+y) \, dydx$。那么我们现在就可以计算了。

第一重积分其实是一个关于$x$的积分上限函数，但是这里是对$y$积分，我们先把它积出来，也就是：

$$\int_{x^2}^x (x+y) \, dy
= \left[ xy + \frac{y^2}{2} \right]_{x^2}^{x} = \frac{3}{2}x^2 - x^3 - \frac{1}{2}x^4$$

再代入到第二重积分里，就是：

$$\int_0^1 (\frac{3}{2}x^2 - x^3 - \frac{1}{2}x^4) dx=\left[ \frac{1}{2}x^3 - \frac{1}{4}x^4 - \frac{1}{10}x^5 \right]_0^1=\frac{3}{20}$$

所以二重积分$\iint_D (x+y) \, d\sigma=\frac{3}{20}$。
\end{myexample}

\begin{myexample}
    计算二重积分$\iint_D (x+y) \, d\sigma$，其中区域$D$ 是由曲线 $y = \sin x$、$y = \cos x$、$x = 0$ 和 $x = \dfrac{\pi}{2}$ 围成的闭区域。

    我们首先判断区域 $D$ 是什么类型的区域。把曲线 $y = \sin x$、$y = \cos x$ 与直线 $x = 0$、$x = \dfrac{\pi}{2}$ 画出来，就会发现它们围成了两个曲边三角形。这个区域其实是 $X$ 型区域，但是在区间 $[0,\frac{\pi}{2}]$ 上，两条曲线相交于 $x = \frac{\pi}{4}$，将区域分成了左右两部分。

    接下来就是“画图穿线法”的操作。判断区域为 $X$ 型，我们用一根平行于 $y$ 轴的直线从下往上穿过这个区域，直线与区域 $D$ 会有两个交点。但是区域分左右两部分，我们就要分别讨论。

    我们记区间 $[0,\frac{\pi}{4}]$ 上的区域为$D_1$。在该区间上，曲线 $y = \cos x$ 在 $y = \sin x$ 的上方，因此直线先交 $y = \sin x$，再交 $y = \cos x$。所以这部分的积分下限为 $y = \sin x$，积分上限为 $y = \cos x$。

    我们记区间 $[\frac{\pi}{4},\frac{\pi}{2}]$ 上的区域为$D_2$。在该区间上，曲线 $y = \sin x$ 在 $y = \cos x$ 的上方，因此直线先交 $y = \cos x$，再交$y = \sin x$。所以这部分的积分下限为 $y = \cos x$，积分上限为 $y = \sin x$。

\begin{center}
\begin{tikzpicture}[x=3cm, y=3cm, scale=1]
    % 坐标轴
    \draw[->] (-0.1,0) -- (1.8,0) node[right] {$x$};
    \draw[->] (0,-0.1) -- (0,1.3) node[above] {$y$};

    % 左半区域填充 (0≤x≤π/4, sin x ≤ y ≤ cos x)
    \fill[gray!30] (0,0) -- plot[domain=0:pi/4, smooth] (\x, {sin(\x r)}) --
        (pi/4, {cos(pi/4 r)}) -- plot[domain=pi/4:0, smooth] (\x, {cos(\x r)}) -- cycle;
    % 右半区域填充 (π/4≤x≤π/2, cos x ≤ y ≤ sin x)
    \fill[gray!30] (pi/4, {cos(pi/4 r)}) -- plot[domain=pi/4:pi/2, smooth] (\x, {cos(\x r)}) --
        (pi/2,1) -- plot[domain=pi/2:pi/4, smooth] (\x, {sin(\x r)}) -- cycle;

        \draw[thick] (pi/2,0) -- (pi/2,1);
        \draw[thick] (0,0) -- (0,1);

    % 曲线
    \draw[thick, domain=0:pi/2, smooth] plot (\x, {sin(\x r)});
    \node[above left] at (1.5, {sin(1.5 r)}) {$y=\sin x$};
    \draw[thick, domain=0:pi/2, smooth] plot (\x, {cos(\x r)});
    \node[above right] at (0.2, {cos(0.2 r)}) {$y=\cos x$};

    % 交点虚线
    \draw[dashed] (pi/4,0) -- (pi/4, {sin(pi/4 r)});
    \node[below] at (pi/4,0) {$\frac{\pi}{4}$};

    % 边界标注
    \node[below left] at (0,0) {$0$};
    \node[below] at (pi/2,0) {$\frac{\pi}{2}$};

    % 区域标签
    \node at (0.4,0.7) {$D$};
\end{tikzpicture}
\begin{tikzpicture}[x=3cm, y=3cm, scale=1]
    % 坐标轴
    \draw[->] (-0.1,0) -- (1.8,0) node[right] {$x$};
    \draw[->] (0,-0.1) -- (0,1.3) node[above] {$y$};

    % 填充区域（简化填充）
    \fill[gray!30] (0,0) -- plot[domain=0:pi/4, smooth] (\x, {sin(\x r)}) --
        (pi/4, {cos(pi/4 r)}) -- plot[domain=pi/4:0, smooth] (\x, {cos(\x r)}) -- cycle;
    \fill[gray!30] (pi/4, {cos(pi/4 r)}) -- plot[domain=pi/4:pi/2, smooth] (\x, {cos(\x r)}) --
        (pi/2,1) -- plot[domain=pi/2:pi/4, smooth] (\x, {sin(\x r)}) -- cycle;
        \draw[thick] (pi/2,0) -- (pi/2,1);
        \draw[thick] (0,0) -- (0,1);

    % 曲线
    \draw[thick, domain=0:pi/2, smooth] plot (\x, {sin(\x r)});
    \draw[thick, domain=0:pi/2, smooth] plot (\x, {cos(\x r)});

    \draw[dashed] (pi/4,0) -- (pi/4, {sin(pi/4 r)});
    \node[below] at (pi/4,0) {$\frac{\pi}{4}$};

    % 边界标注
    \node[below left] at (0,0) {$0$};
    \node[below] at (pi/2,0) {$\frac{\pi}{2}$};

    % 穿线法：左段 x=0.5
    \draw[thick] (0.5, -0.1) -- (0.5, 1.1);
    \fill (0.5, {cos(0.5 r)}) circle (1.5pt) node[above right] {$A$};
    \fill (0.5, {sin(0.5 r)}) circle (1.5pt) node[right] {$B$};

    % 穿线法：右段 x=1.2
    \draw[thick] (1.2, -0.1) -- (1.2, 1.1);
    \fill (1.2, {sin(1.2 r)}) circle (1.5pt) node[above right] {$C$};
    \fill (1.2, {cos(1.2 r)}) circle (1.5pt) node[right] {$D$};

    \node at (0.2,0.5) {$D_1$};
    \node at (1.3,0.7) {$D_2$};
\end{tikzpicture}
\end{center}

    那么原积分就可以写成两部分之和。我们记区域$D_1$上的二重积分为$I_1$，记区域$D_2$上的二重积分为$I_2$，则$\iint_D (x+y) \, d\sigma=I_1+I_2$。

    先算$I_1$，我们刚刚确定了$I_1$的第一重积分的积分下限为 $y = \sin x$，积分上限为 $y = \cos x$，积出来就是：
    \[
    \int_{\sin x}^{\cos x} (x+y) \, dy=\left[ xy + \frac{y^2}{2} \right]_{\sin x}^{\cos x} = x\cos x - x\sin x + \frac{\cos 2x}{2}
    \]

    $I_1$的区间是$[0,\frac{\pi}{4}]$，所以第二重积分就是：
    \[
    \int_0^{\frac{\pi}{4}} (x\cos x - x\sin x + \frac{\cos 2x}{2}) \, dx=[x\sin x + \cos x + x\cos x - \sin x + \frac{1}{4}\sin 2x]_0^{\frac{\pi}{4}}=\frac{\pi\sqrt{2}-3}{4}
    \]

    同理可计算$I_2$。$I_2$的第一重积分的积分下限为$y = \cos x$，积分上限为 $y = \sin x$，所以：
    \[
    \int_{\cos x}^{\sin x} (x+y) \, dy=\left[ xy + \frac{y^2}{2} \right]_{\cos x}^{\sin x}=x\sin x - x\cos x - \frac{\cos 2x}{2}
    \]

    $I_2$的区间是$[\frac{\pi}{4},\frac{\pi}{2}]$，所以：
    \[
    \int_{\frac{\pi}{4}}^{\frac{\pi}{2}} (x\sin x - x\cos x - \frac{\cos 2x}{2}) \, dx=[-x\cos x + \sin x - x\sin x - \cos x - \frac{1}{4}\sin 2x]_{\frac{\pi}{4}}^{\frac{\pi}{2}}=\frac{5+\pi\sqrt{2}-2\pi}{4}
    \]

    最后相加得：
    \[\iint_D (x+y) \, d\sigma = I_1 + I_2 = \frac{\pi(\sqrt{2}-1)+1}{2}\]
\end{myexample}

有时我们会遇到一个写好了的累次积分，但按原顺序很难积或者积不出来。这时候我们可以根据上下限画出积分区域，然后交换积分次序，往往就能化难为易。

所谓交换积分次序，其实就是把$X$ 型区域看成$Y$ 型区域再去算，把$Y$ 型区域看成$X$ 型区域再去算。

\begin{myexample}
交换累次积分 \(\int_0^2\left[\int_{y^2}^{4} f(x,y)\,dx\right]dy\) 的积分次序。

观察原积分，我们发现这是一个先对$x$积分，再对$y$积分的二重积分，那么也就是说区域 \(D\) 是个$Y$ 型区域。根据原积分上下限，区域 \(D\) 可以表示为：
\[
D=\{(x,y)\mid 0\le y\le 2,\ y^2\le x\le 4\}
\]

我们把它画出来，就是：
\begin{center}
    \begin{tikzpicture}[x=2cm, y=2cm]
        % 坐标轴
        \draw[->] (-0.3,0) -- (4.8,0) node[right] {$x$};
        \draw[->] (0,-0.2) -- (0,2.5) node[above] {$y$};
    
        % 填充积分区域 D: 0<=y<=2, y^2<=x<=4
        \fill[gray!30] (0,0) 
            -- (4,0) 
            -- (4,2) 
            -- plot[domain=2:0, smooth] ({\x*\x}, \x) 
            -- cycle;
    
        % 左边界曲线 x = y^2  (即 y = sqrt{x})
        \draw[thick, domain=0:2, smooth] plot ({\x*\x}, \x);
        \node[above left] at (2.25,1.5) {$x=y^2$};
    
        % 右边界直线 x = 4
        \draw[thick] (4,0) -- (4,2);
        \draw[thick,dashed] (0,2) -- (4,2);
        \draw[thick] (0,0) -- (4,0);
    
        % 坐标刻度
        \node[below] at (4,0) {$4$};
        \node[left] at (0,2) {$2$};
        \node[below left] at (0,0) {$0$};
    
        % 区域标签
        \node at (2,0.6) {$D$};
    \end{tikzpicture}
\end{center}

因此把区域 \(D\) 转写成$X$型，得：
\[
D=\{(x,y)\mid 0\le x\le 4,\ 0\le y\le\sqrt{x}\}
\]

交换次序后得到：
\[
\int_0^2\left[\int_{y^2}^{4} f(x,y)\,dx\right]dy = \int_0^4\left[\int_0^{\sqrt{x}} f(x,y)\,dy\right]dx
\]

\end{myexample}


\begin{myexample}
求$\int_{0}^{1} \left[ \int_{x}^{1} e^{y^{2}} \, dy \right] dx$。

我们发现，这个累次积分其中的$e^{y^{2}}$是一个典型的“积不出”积分，也就是说如果直接求，我们根本没办法求出它的值。那么这个时候我们就要考虑一下交换积分次序了。

原积分是先对$y$积分，再对$x$积分，也就是说这个区域是个$X$型区域。我们把这个区域画出来：
    
\begin{center}
        \begin{tikzpicture}[x=4cm, y=4cm]
        % 坐标轴
        \draw[->] (-0.1,0) -- (1.2,0) node[right] {$x$};
        \draw[->] (0,-0.1) -- (0,1.2) node[above] {$y$};
    
        % 填充区域 D: 0<=x<=1, x<=y<=1
        \fill[gray!30] (0,0) -- (0,1) -- (1,1) -- cycle;
    
        % 下边界 y = x
        \draw[thick, domain=0:1, smooth] plot (\x, \x);
        \node[below right] at (0.5,0.5) {$y=x$};
    
        % 上边界 y = 1
        \draw[thick] (0,1) -- (1,1);
        \node[above] at (0.5,1) {$y=1$};
    
        % 左边界 x=0 (与y轴重合)
        \draw[thick] (0,0) -- (0,1);
    
        \draw[thick,dashed] (1,0) -- (1,1);
    
        % 坐标标注
        \node[below] at (1,0) {$1$};
        \node[left] at (0,1) {$1$};
        \node[below left] at (0,0) {$0$};
    
        % 区域标签
        \node at (0.25,0.75) {$D$};
    \end{tikzpicture}
\end{center}

再把它转写成$Y$型区域，就得到：
\[
\int_{0}^{1} \left[ \int_{0}^{y} e^{y^{2}} \, dx \right] dy
\]

这里要注意，括号里的积分是对$x$的积分，也就是说和$y$完全无关，因此我们可以直接把$e^{y^{2}}$提出来，这个积分就变成：
\[
\int_{0}^{1} e^{y^{2}} \left[ \int_{0}^{y} \, dx \right] dy
\]

那么原积分就是：
\[
\int_{0}^{1} \left[ \int_{x}^{1} e^{y^{2}} \, dy \right] dx=\int_{0}^{1} e^{y^{2}} \left[ \int_{0}^{y} \, dx \right] dy=\int_{0}^{1} y \, e^{y^{2}} \, dy=\left[\frac{1}{2}e^{y^2}\right]_{0}^{1}=\frac{1}{2} (e - 1)
\]
\end{myexample}

总结一下，要求二重积分，首先要确定积分区域的类型。然后用“画图穿线法”确定第一重积分的上下限，如果是$X$型区域就从下往上穿，$Y$型区域就从左往右穿，也就是从小往大穿，小的作为积分下限，大的作为积分上限。把第一重积分积出来之后，再根据区间，确定第二重积分的上下限，把第一重积分积出来的结果代入第二重积分，就能求出值了。


\subsection{二重积分的极坐标形式}
有些二重积分的积分区域的边界是圆弧或射线，或者被积函数里出现了 \(x^2+y^2\) 类似的项。这时要是还用直角坐标去算，不仅上下限很难写，积分本身也可能变得十分复杂。这时，就可以考虑利用极坐标来计算二重积分。

我们假定读者已经对极坐标有基本的了解，所以我们直接进入极坐标形式下二重积分的计算。与直角坐标不同的是，在极坐标下，面积微元$d\sigma$不再是简单的$dxdy$，而是变成另一种形式。


\begin{center}
\begin{tikzpicture}

% 4. 绘制图形区域 D (粗线圆)
    \draw[thick] (45:4.2) circle (2.5);
    \fill[gray!50] (45:4.2) circle (2.5);

    % 1. 定义原点和极轴
    \coordinate (O) at (0,0);
    \coordinate (A) at (7.5,0);
    
    % 2. 绘制极网格 (灰度射线和圆弧)
    \foreach \angle in {10,20,30,40,50,60,70,80,90} {
        \draw (O) -- (\angle:7);
    }
    % 圆弧: 半径从 1.5 到 7.5，步长 1
    \foreach \r in {1.5,2.5,3.5,4.5,5.5,6.5} {
        \draw (0:\r) arc (0:90:\r);
    }

    % 3. 绘制极轴 OA (粗线，带箭头)
    \draw[->] (O) -- (A) node[below] {\(A\)};
    \node[below left] at (O) {\(O\)};

    

    

    % 5. 绘制阴影微元 \Delta\sigma_i (填充并描边)
    \fill[red!75] (40:4.5) arc (40:50:4.5) -- (50:5.5) arc (50:40:5.5) -- cycle;
    \draw (40:4.5) arc (40:50:4.5) -- (50:5.5) arc (50:40:5.5) -- cycle;

    % 6. 添加标注和向量
    % 标注 D
    \node at (55:6) {\(D\)};
    
    % 标注微元 \Delta\sigma_i
    \node at (46:5) {\(\Delta\sigma\)};

    % 标注 \Delta\theta_i (使用带箭头的弧线)
    \draw[<->, thin] (40:4.3) arc (40:50:4.3);
    \node at (46:3.9) {\(\rho\Delta\theta\)};

    % 标注 \Delta \r (使用带箭头的直线)
    \draw[<->, thin] (38:4.5) -- (38:5.5);
    \node at (34:5.0) {\(\Delta\rho\)};


\end{tikzpicture}
\end{center}

我们知道，在极坐标下，确定一个点需要确定极径$\rho$和极角$\theta$。与直角坐标类似，我们取一个极小的$\Delta\rho$和$\Delta \theta$，那么极径从$\rho$变化到$\rho+\Delta\rho$，极角从$\theta$变化到$\theta+\Delta \theta$，这时候的面积微元就是一个“曲边矩形”，矩形的边长就是$\Delta\rho$，弧长就是$\rho\Delta \theta$和$(\rho+\Delta\rho)\cdot \Delta \theta$。当$\Delta \theta$和$\Delta\rho$足够小的时候，我们就可以认为$\rho\Delta \theta$与$(\rho+\Delta\rho)\cdot \Delta \theta$相等，“曲边矩形”可以近似为矩形，面积也可以认为是$\Delta\rho \cdot \rho\Delta \theta$。然后我们取微分，那么面积就是$d\rho \cdot \rho d\theta$，也就是$\rho d\rho d\theta$。这就是极坐标下的面积微元。

如果区域 \(D\) 可以写成夹在两条射线 \(\theta=\alpha\) 与 \(\theta=\beta\) 之间，且内外边界分别是曲线 \(\rho=\rho_1(\theta)\) 和 \(\rho=\rho_2(\theta)\) 的形式，即：
\[
D=\{(\rho,\theta)\mid \alpha\le\theta\le\beta,\ \rho_1(\theta)\le \rho\le \rho_2(\theta)\}
\]

那么我们习惯上把这种区域称为极坐标下的标准区域。这时二重积分就可以化成先对 \(\rho\) 后对 \(\theta\) 的二次积分：

\begin{theorem}{极坐标下的二重积分}{}
    若区域 \(D\) 在极坐标下可表示为 \(\alpha\le\theta\le\beta,\ \rho_1(\theta)\le \rho\le \rho_2(\theta)\)，则
    \[
    \iint_D f(\rho,\theta)\,d\sigma = \iint_D f(\rho,\theta)\,\rho d\rho d\theta = \int_\alpha^\beta\left[\int_{\rho_1(\theta)}^{\rho_2(\theta)} f(\rho,\theta)\,\rho\,d\rho\right]d\theta
    \]
\end{theorem}

其中$\rho\ge 0$，$0\le \theta \le2\pi$。一般我们先积$\rho$，再积$\theta$。注意极坐标下的面积微元是$\textcolor{red}{\rho}d\rho d\theta$，不是$d\rho d\theta$。

有时候，我们想把直角坐标的二重积分转换为极坐标来算，那么这要怎么操作呢？我们只需要令$x=\rho\cos\theta$，$y=\rho\sin\theta$，再用极坐标计算就可以了。

\begin{theorem}{二重积分变换公式}{}
    \[
    \iint_D f(x,y)\,d\sigma = \iint_D f(x,y)\,dxdy = \iint_D f(\rho\cos\theta,\rho\sin\theta)\,\rho d\rho d\theta
    \]
\end{theorem}

计算极坐标下的二重积分的操作思路和直角坐标时如出一辙。我们先把 \(\theta\) 固定住，那么从极点出发、极角为 \(\theta\) 的射线穿进区域 \(D\) 时碰到的第一条边界就是 \(\rho=\rho_1(\theta)\)，穿出区域时碰到的边界就是 \(\rho=\rho_2(\theta)\)。于是沿着这条射线把函数值从 \(\rho_1(\theta)\) 到 \(\rho_2(\theta)\) 积起来，就得到了对应于这个 \(\theta\) 的一片“累积量”，然后再让 \(\theta\) 从 \(\alpha\) 扫到 \(\beta\)，把所有这样的片累加起来，就得到了整个二重积分。这和直角坐标下“先积一条线，再扫一个面”的做法完全是一样的，只不过是把直线换成了射线。

那么，怎么确定上下限呢？我们依旧用“画图穿线法”，只是在极坐标下，穿区域的线变成了从极点出发的射线。让射线绕着极点，从角度 \(\alpha\) 开始，沿逆时针方向扫到 \(\beta\)。在每一个角度上，射线进入区域的那一点就是积分下限 \(\rho=\rho_1(\theta)\)，穿出区域的那一点就是积分上限 \(\rho=\rho_2(\theta)\)。扫过的角度范围就是 \(\theta\) 的积分区间。也正因如此，所以极坐标下的二重积分基本都是先对$\rho$积分再对$\theta$积分。

和直角坐标一样，我们来看几个例子：

\begin{myexample}
    计算二重积分 \(\iint_D (x^2+y^2)\,d\sigma\)，其中区域 \(D\) 是由圆 \(x^2+y^2=1\) 围成的闭区域。

    我们将原积分化为极坐标形式。令$x = \rho\cos\theta$，$y = \rho\sin\theta$，则$d\sigma = \rho\,d\rho\,d\theta$，原积分化为$\iint_D \rho^2 \cdot \rho\,d\rho\,d\theta$。区域 \(D\) 化为$\rho^2\le1$。

    我们画出区域 \(D\)，再用“画图穿线法”确定上下限。先从极点引出一条射线，这条射线在从极点 \(\rho=0\) 出发，穿出圆时到达 \(\rho=1\)，所以积分下限是 \(\rho=0\)，积分上限是 \(\rho=1\)。让它逆时针扫过 \(0\) 到 \(2\pi\)，积分区间就是$[0,2\pi]$。

    \begin{center}
    \begin{tikzpicture}[scale=2]
        % 单位圆
        \draw[thick] (0,0) circle (1);
        \fill[gray!30] (0,0) circle (1);

        % 极轴
        \draw[->] (-1.3,0) -- (1.3,0) node[right] {$x$};
        \draw[->] (0,-1.3) -- (0,1.3) node[above] {$y$};

        % 射线示例
        \draw[] (0,0) -- (60:1.2);
        \draw[->] (0.35,0) arc (0:60:0.35) node[midway, right] {$\theta$};

        \node at (0.5,0.5) {$D$};
        \node at (0.7,1.2) {$\rho=1$};

    \end{tikzpicture}
    \end{center}

    于是原积分化为：
    \[
\iint_D \rho^2 \cdot \rho\,d\rho\,d\theta = \int_{0}^{2\pi} \int_{0}^{1} \rho^3\,d\rho\,d\theta
    \]

    先算第一重积分：
    \[
    \int_{0}^{1} \rho^3\,d\rho = \left[ \frac{\rho^4}{4} \right]_{0}^{1} = \frac{1}{4}
    \]

    再代入第二重积分：
    \[
    \int_{0}^{2\pi} \frac{1}{4}\,d\theta = \frac{1}{4} \cdot 2\pi = \frac{\pi}{2}
    \]

    所以二重积分 \(\iint_D (x^2+y^2)\,d\sigma=\frac{\pi}{2}\)。
\end{myexample}

\begin{myexample}
    计算二重积分 \(\iint_D \sqrt{x^2+y^2}\,d\sigma\)，其中区域 \(D\) 是由心脏线 \(\rho=1+\cos\theta\) 围成的闭区域。

    同理，我们令$x = \rho\cos\theta$，$y = \rho\sin\theta$，则$d\sigma = \rho\,d\rho\,d\theta$，原积分化为$\iint_D \rho \cdot \rho\,d\rho\,d\theta$，区域 \(D\) 的边界由极坐标方程直接给出了。

    我们先画出区域 \(D\)，再用“画图穿线法”确定上下限。对于每一个角度 \(\theta\)，射线从极点出发，穿进区域时 \(\rho=0\)，穿出时上边界 \(\rho=1+\cos\theta\)。因此积分下限就是\(\rho=0\)，积分上限就是\(\rho=1+\cos\theta\)。\(\theta\)从\(0\) 到 \(2\pi\)能画出整个边界，所以积分区间就是$[0,2\pi]$。

    \begin{center}
    \begin{tikzpicture}[scale=1.8]

        % 心脏线
        \draw[thick, domain=0:360, smooth,samples=200] plot (\x:{1+cos(\x)});
        \fill[gray!30] plot[domain=0:360, smooth,samples=200] (\x:{1+cos(\x)}) -- cycle;

        \draw[->] (-0.5,0) -- (2.5,0) node[right] {$x$};
        \draw[->] (0,-1.3) -- (0,1.3) node[above] {$y$};

        % 射线示例
        \draw (0,0) -- (40:{1.2+cos(40)});
        \draw[->] (0.4,0) arc (0:40:0.4) node[midway, right] {$\theta$};

        \node at (1.5,0.4) {$D$};
        \node at (2.2,1.1) {$\rho=1+\cos\theta$};
    \end{tikzpicture}
    \end{center}

    于是原积分化为：
    \[
    \iint_D \rho\,d\sigma = \int_0^{2\pi}\left[\int_0^{1+\cos\theta} \rho\cdot \rho\,d\rho\right]d\theta
    = \int_0^{2\pi}\left[\int_0^{1+\cos\theta} \rho^2\,d\rho\right]d\theta
    \]

    先算第一重积分：
    \[
    \int_0^{1+\cos\theta} \rho^2\,d\rho = \left[\frac{\rho^3}{3}\right]_0^{1+\cos\theta} = \frac{(1+\cos\theta)^3}{3}
    \]

    再代入第二重积分：
    \[
    \frac13\int_0^{2\pi} (1+\cos\theta)^3\,d\theta = \frac13\left[\int_0^{2\pi}1\,d\theta + 3\int_0^{2\pi}\cos^2\theta\,d\theta\right]
    = \frac13(2\pi + 3\pi) = \frac{5\pi}{3}
    \]

    所以二重积分 \(\iint_D \sqrt{x^2+y^2}\,d\sigma = \frac{5\pi}{3}\)。
\end{myexample}

\begin{myexample}
    计算二重积分 \(\iint_D \arctan\frac{y}{x}\,d\sigma\)，其中区域 \(D\) 是由圆 \(x^2+y^2=1\)，\(x^2+y^2=4\) 与射线 \(y=x\ (x\ge 0)\)，\(y=0\ (x\ge 0)\) 所围成的在第一象限的部分。

    令$x = \rho\cos\theta$，$y = \rho\sin\theta$，先化简被积函数：
    \[
    \arctan\frac{y}{x}=\arctan\frac{\rho\sin\theta}{\rho\cos\theta}=\arctan\tan\theta=\theta
    \]

    所以原积分就是 \(\iint_D \theta\,d\sigma\)，区域 \(D\)是夹在两个圆和两条射线之间的环形扇区。从极点出发的射线，在角度 \(\theta\) 从 \(0\) 扫到 \(\frac{\pi}{4}\) 的过程中，先穿入小圆边界 \(\rho=1\)，再穿出大圆边界 \(\rho=2\)。所以积分下限是\(\rho=1\)，积分上限是\(\rho=2\)，积分区间是\([0,\frac{\pi}{4}]\)。

\begin{center}
    \begin{tikzpicture}[scale=1.6]

        % 填充区域（仍保持在 0°~45°，r=1~2）
        \fill[gray!30] 
            (0:1) arc (0:45:1) -- 
            (45:2) arc (45:0:2) -- cycle;

        % 坐标轴（仅正半轴）
        \draw[->] (0,0) -- (2.5,0) node[right] {$x$};
        \draw[->] (0,0) -- (0,2) node[above] {$y$};

        % 射线
        \draw[dashed] (0,0) -- (45:2.5);
        \draw[thick] (0,0) -- (2.5,0);
        \draw[thick] (0,0) -- (45:2.5);

        % 第一象限圆
        \draw[thick] (1,0) arc (0:45:1);
        \draw[thick] (2,0) arc (0:45:2);


        % 标注
        \draw[->] (0.4,0) arc (0:45:0.4) node[midway, right] {$\theta$};
        \node[below] at (1,0) {$1$};
        \node[below] at (2,0) {$2$};
        \node at (1.3,0.5) {$D$};
        \node at (1.4,0.15) {$\rho=1$};
        \node at (2.4,0.4) {$\rho=2$};
    \end{tikzpicture}
\end{center}

    于是原积分化为：
    \[
    \iint_D \theta\,d\sigma = \int_0^{\frac{\pi}{4}}\left[\int_1^2 \theta \cdot \rho\,d\rho\right]d\theta
    \]

    先算第一重积分，\(\theta\) 对 \(\rho\) 来说相当于常数，可以直接提出：
    \[
    \int_1^2 \theta\,\rho\,d\rho = \theta\left[\frac{\rho^2}{2}\right]_1^2 = \theta(2-\frac12) = \frac32\,\theta
    \]

    再代入第二重积分：
    \[
    \int_0^{\frac{\pi}{4}} \frac32\,\theta\,d\theta = \frac32\left[\frac{\theta^2}{2}\right]_0^{\frac{\pi}{4}} = \frac32\cdot\frac12\cdot\frac{\pi^2}{16} = \frac{3\pi^2}{64}
    \]

    所以 \(\iint_D \arctan\frac{y}{x}\,d\sigma = \frac{3\pi^2}{64}\)。
\end{myexample}

有时我们也会遇到直角坐标下已经写好的累次积分，按原顺序去算几乎无从下手。这时候可以考虑把积分区域画出来，如果它在极坐标下表达特别简单，就可以通过极坐标换元来化难为易。

\begin{myexample}
    求 \(\int_0^1\left[\int_0^{\sqrt{1-x^2}} (x^2+y^2)\,dy\right]dx\)。

    这个累次积分如果直接按直角坐标去算，被积函数虽然能积，但过程略繁琐。我们把它看作某个二重积分，画出区域 \(D\)：

    \begin{center}
    \begin{tikzpicture}[scale=2.5]
        \draw[->] (-0.2,0) -- (1.3,0) node[right] {$x$};
        \draw[->] (0,-0.2) -- (0,1.3) node[above] {$y$};
        \fill[gray!30] (0,0) -- (1,0) arc (0:90:1) -- cycle;
        \draw[thick] (0,0) -- (1,0) arc (0:90:1) -- cycle;
        \node[below] at (1,0) {$1$};
        \node[left] at (0,1) {$1$};
        \node at (0.5,0.5) {$D$};
    \end{tikzpicture}
    \end{center}

    这个区域其实是单位圆在第一象限的部分。被积函数 \(x^2+y^2=\rho^2\)，面积微元是 \(\rho\,d\rho\,d\theta\)。于是原积分变成极坐标下的累次积分：
    \[
    \int_0^{\frac{\pi}{2}}\left[\int_0^1 \rho^2\cdot \rho\,d\rho\right]d\theta
    = \int_0^{\frac{\pi}{2}}\left[\int_0^1 \rho^3\,d\rho\right]d\theta
    \]

    先对 \(\rho\) 积分：
    \[
    \int_0^1 \rho^3\,d\rho = \left[\frac{\rho^4}{4}\right]_0^1 = \frac14
    \]

    再对 \(\theta\) 积分：
    \[
    \int_0^{\frac{\pi}{2}} \frac14\,d\theta = \frac14\cdot\frac{\pi}{2} = \frac{\pi}{8}
    \]

    所以原积分 \(\int_0^1\left[\int_0^{\sqrt{1-x^2}} (x^2+y^2)\,dy\right]dx = \frac{\pi}{8}\)。
\end{myexample}

\begin{myexample}
    计算二重积分$\iint_D e^{x^2+y^2}\,d\sigma$，其中 $D$ 是由圆心在原点、半径为 $a$ 的圆周所围成的闭区域。

    被积函数有$e^{x^2+y^2}$，这是一个典型的“积不出”函数。但是我们可以通过极坐标把它积出来。

    令$x = \rho\cos\theta$，$y = \rho\sin\theta$，则：
\[
\iint_D e^{x^2+y^2}\,d\sigma = \iint_D e^{\rho^2}\,\rho d\rho d\theta
= \int_0^{2\pi} d\theta \int_0^a \rho e^{\rho^2} \,d\rho
\]

对 $\rho$ 积分，令 $u = \rho^2$，$du = 2\rho\,d\rho$，则：
\[
\int_0^a \rho e^{\rho^2}\,d\rho = \frac{1}{2} \int_0^{a^2} e^u\,du
= \frac{1}{2}(e^{a^2} - 1)
\]

再对 $\theta$ 积分：
\[
\int_0^{2\pi} \frac{1}{2}(e^{a^2} - 1)\, d\theta = 2\pi \cdot \frac{1}{2}(e^{a^2} - 1) =\pi(e^{a^2} - 1)
\]

所以
\(\iint_D e^{x^2+y^2}\,d\sigma
= \pi(e^{a^2} - 1)\)。
\end{myexample}

总结一下，用极坐标计算二重积分，先看被积函数和积分区域是否适合用极坐标，通常含有 \(x^2+y^2\)、\(\frac{y}{x}\) 这类式子，或者区域是圆、环、扇形等形状时，极坐标都是很好的选择。然后画出积分区域，用“画图穿线法”确定 \(\rho\) 的上下限和 \(\theta\) 的范围。射线从极点出发，穿入点作为 \(\rho\) 的下限，穿出点作为\(\rho\) 的上限，所有这样的射线扫过的角度范围就是 \(\theta\) 的积分区间。列式时切记面积微元是 $\textcolor{red}{\rho}d\rho d\theta$，不是$d\rho d\theta$。最后先 \(\rho\) 后 \(\theta\)化为累次积分，一步步算出来就行了。

\section{三重积分}
定积分和二重积分的定义可以很自然地扩展到三重积分。不过在进入三重积分之前，我们需要重新认识定积分和二重积分。
\subsection{三重积分}

在定积分和二重积分的定义里，我们总是将一点的函数值与我们要积分的微元相乘，比如定积分是函数值与线段微元相乘，二重积分是函数值与面积微元相乘。如果我们将函数值看成密度函数，那么我们就能用积分求质量。也就是说，对一个函数积分，就能得到积分区域的质量。

具体到定积分，就是已知细杆的线密度$f(x)$，积分求细杆质量。到了二重积分，就是已知薄板的面密度$f(x,y)$，积分求薄板质量。而我们生活在三维空间里，接触的都是空间体，仿照定积分和二重积分，如果我们已知物体的体密度$f(x,y,z)$，就能求出这个物体的质量了。这就是三重积分的物理意义。

在几何上，三重积分求的是四维空间中，以三维区域为底，三元函数$f(x,y,z)$为顶的四维曲顶柱体的“四维体积”。这很难想象，物理意义相比之下更直观也更容易理解。

仿照二重积分，我们再推广一步就能得到三重积分的定义。设 $w = f(x, y, z)$ 是定义在空间区域 $\Omega$ 上的函数。把 $\Omega$ 分成 $n$ 个小块$\Delta V_1,\ \Delta V_2,\ \cdots,\ \Delta V_n$，在每小块上取一点 $(\xi_i, \eta_i, \zeta_i)$，作和 $\sum_{i=1}^n f(\xi_i, \eta_i, \zeta_i)\, \Delta V_i$，取极限就得到三重积分。这本质上还是“分割、近似、求和、取极限”的思想。

\begin{definition}{三重积分}{}
设函数 \(f(x,y,z)\) 在有界闭区域 \(\Omega\) 上有界。将 \(\Omega\) 任意分成 \(n\) 个小闭区域
\[
\Delta V_1,\ \Delta V_2,\ \cdots,\ \Delta V_n
\]
每个小区域既表示该区域本身，也表示该区域的体积。在每个小区域 \(\Delta V_i\) 上任取一点 \((\xi_i,\eta_i,\zeta_i)\)，作函数值 \(f(\xi_i,\eta_i,\zeta_i)\) 与小区域体积 \(\Delta V_i\) 的乘积 \(f(\xi_i,\eta_i,\zeta_i) \cdot \Delta V_i\)，并作出和
\[
S = \sum_{i=1}^{n} f(\xi_i,\eta_i,\zeta_i) \Delta V_i
\]

记 \(d_i\) 为小区域 \(\Delta V_i\) 的直径，即该小区域内任意两点间距离的最大值，令 \(\lambda = \max\{d_1,d_2,\cdots,d_n\}\)。如果当 \(\lambda \to 0\) 时，\(S\) 的极限总存在，且与闭区域 \(\Omega\) 的分法及点 \((\xi_i,\eta_i,\zeta_i)\) 的取法无关，那么称 \(f(x,y,z)\) 在 \(\Omega\) 上\textbf{可积}，并称这个极限 \(I\) 为函数 \(f(x,y,z)\) 在闭区域 \(\Omega\) 上的\textbf{三重积分}，记作
\[
\iiint\limits_{\Omega} f(x,y,z)\,dV
\quad\text{或}\quad
\iiint\limits_{\Omega} f(x,y,z)\,dxdydz
\]
即：
\[
\iiint\limits_{\Omega} f(x,y,z)\,dV = \lim_{\lambda \to 0} \sum_{i=1}^{n} f(\xi_i,\eta_i,\zeta_i) \Delta V_i = I
\]

其中 \(f(x,y,z)\) 称为\textbf{被积函数}，\(f(x,y,z)\,dV\) 称为\textbf{被积表达式}，\(x,y,z\) 称为\textbf{积分变量}，\(\Omega\) 称为\textbf{积分区域}，$\iiint$称为\textbf{三重积分号}。
\end{definition}

一重积分切线段，二重积分切平面，三重积分切空间。积分号从 $\int$ 到 $\iint$ 到 $\iiint$，代表累加的维度越来越多。这个模式还可以继续推广到更高维。

三重积分的性质和二重积分完全一样。这里我们不再重复。


\begin{myexample}
  求 $\iiint_\Omega 2\, dV$，其中 $\Omega$ 是以原点为球心、半径为 $R$ 的球体。

  根据三重积分的物理意义，$\Omega$的体密度是$2$，也就是说这个球体密度处处为$2$。体积就是整个球体体积，那么质量就是$2\cdot \frac43 \pi R^3=\frac83 \pi R^3$。

所以$\iiint_\Omega 2\, dV=\frac83 \pi R^3$。
\end{myexample}

三重积分的定义直接用定义来计算不现实。我们需要像二重积分一样把它转化为累次积分才好计算。

\subsection{三重积分的直角坐标形式}

和二重积分类似，三重积分也是将积分化为一次次的定积分进行计算。三重积分的体积微元在直角坐标下就是 \(dV = dx\,dy\,dz\)，所以三重积分归根结底就是对 \(x,y,z\) 依次做定积分。不过具体按照什么顺序，就要看积分区域长什么样。一般我们有两种思路，一种是把空间区域投影到平面上，先穿一条线，再扫一个面。另一种是用垂直于某条坐标轴的平面去切，像切面包片，先把每片上的二重积分算出来，再把片累加起来。它们分别叫“投影法”和“截面法”。按照积分的顺序，它们又叫作“先一后二法”和“先二后一法”。

设空间区域 \(\Omega\) 在 \(xOy\) 平面上的投影为 \(D_{xy}\)，并且 \(\Omega\) 恰好夹在两张曲面之间，下方曲面是 \(z=z_1(x,y)\)，上方曲面是 \(z=z_2(x,y)\)，即：
\[
\Omega = \{(x,y,z) \mid (x,y)\in D_{xy},\; z_1(x,y) \le z \le z_2(x,y)\}
\]

我们依旧用“画图穿线法”确定上下限。先固定 \((x,y)\)，沿着$z$轴竖着穿一条线。这条线穿进 \(\Omega\) 里的地方就是下边界 \(z=z_1(x,y)\)，穿出去的地方是上边界 \(z=z_2(x,y)\)。先对 \(z\) 积分，积分下限就是\(z=z_1(x,y)\)，积分上限就是\(z=z_2(x,y)\)，积分得到这条“细杆”的累积量，然后把 \(D_{xy}\) 上所有的“细杆”加起来，也就是以区域\(D_{xy}\)做二重积分。最后就求得了三重积分的值。这就是投影法的思想。

\begin{theorem}{三重积分的投影法}{}
设空间闭区域 \(\Omega\) 可表示为
\[
\Omega = \{(x,y,z) \mid (x,y) \in D_{xy},\; z_1(x,y) \le z \le z_2(x,y)\}
\]
其中 \(D_{xy}\) 是 \(\Omega\) 在 \(xOy\) 面上的投影区域，\(z_1(x,y)\)，\(z_2(x,y)\) 在 \(D_{xy}\) 上连续。若 \(f(x,y,z)\) 在 \(\Omega\) 上可积，则
\[
\iiint_{\Omega} f(x,y,z)\,dxdydz
= \iint_{D_{xy}} \left( \int_{z_1(x,y)}^{z_2(x,y)} f(x,y,z)\,dz \right) dxdy
\]
\end{theorem}

这跟二重积分中“先穿一条线，再扫一个面”的做法一模一样，只不过现在多了一根 \(z\) 轴，线是竖着的，扫的是投影面\(D_{xy}\)。

用物理意义理解，就是根据密度函数求出了一条极细的“细杆”质量，然后做二重积分，求出所有这些“细杆”在“薄板”上的质量，这样就是整个空间体的质量。

\begin{center}
    \pgfmathdeclarefunction{F}{2}{%
  \pgfmathparse{1.05 + 0.22*#1 + 0.12*#2 + 0.32*sin(70*#1)}%
    }
\begin{tikzpicture}[
    scale=1.05,
    x={(1cm,0cm)},
    y={(0.55cm,0.32cm)},
    z={(0cm,1cm)},
    line join=round,
    line cap=round
]
    % 椭圆形底域 D 的参数
    \def\cx{2.15}
    \def\cy{1.15}
    \def\ra{1.55}
    \def\rb{0.78}

    % 坐标轴
    \draw[->] (-0.25,0,0) -- (4.45,0,0) node[below] {$x$};
    \draw[->] (0,-0.2,0) -- (0,2.85,0) node[right] {$y$};
    \draw[->] (0,0,0) -- (0,0,3.8) node[left] {$z$};
    \node[below left] at (0,0,0) {$O$};

    

    % 侧面 —— 用不透明的浅色填充，不加渐变
    \fill[cyan!20]   % 不透明，颜色自行调整
        plot[domain=0:360, samples=120, smooth]
        ({\cx + \ra*cos(\x)}, {\cy + \rb*sin(\x)},
         {F(\cx + \ra*cos(\x), \cy + \rb*sin(\x)) + 1})
        -- plot[domain=360:0, samples=120, smooth]
        ({\cx + \ra*cos(\x)}, {\cy + \rb*sin(\x)}, 1)
        -- cycle;

    % 母线 —— 用虚线表达柱面的弯曲，无光照感
    \foreach \ang in {0,1,...,359} {
        \pgfmathsetmacro\xx{\cx + \ra*cos(\ang)}
        \pgfmathsetmacro\yy{\cy + \rb*sin(\ang)}
        \pgfmathsetmacro\zz{F(\xx,\yy)}
        \draw[cyan!40, line width=0.4mm] (\xx,\yy,1) -- (\xx,\yy,\zz+1);
    }

    % 曲顶 z = f(x,y) + 1
    \fill[blue!18, opacity=0.72]
        plot[domain=0:360, samples=120, smooth]
        ({\cx + \ra*cos(\x)}, {\cy + \rb*sin(\x)},
         {F(\cx + \ra*cos(\x), \cy + \rb*sin(\x)) + 1})
        -- cycle;

    % 底域 D（浅色填充，不透明）
    \draw[cyan!70!black, thick]
        plot[domain=195:360, samples=80, smooth]
        ({\cx + \ra*cos(\x)}, {\cy + \rb*sin(\x)}, 1);

    \draw[cyan!70!black, thick, dashed]
        plot[domain=0:195, samples=80, smooth]
        ({\cx + \ra*cos(\x)}, {\cy + \rb*sin(\x)}, 1);

    % 底面与顶面的轮廓线
    \draw[blue!75!black, very thick]
        plot[domain=0:360, samples=120, smooth]
        ({\cx + \ra*cos(\x)}, {\cy + \rb*sin(\x)},
         {F(\cx + \ra*cos(\x), \cy + \rb*sin(\x)) + 1});

         \draw[red!70!black, thick]
        plot[domain=0:360, samples=80, smooth]
        ({\cx + \ra*cos(\x)}, {\cy + \rb*sin(\x)}, 0);

        \fill [red!24!white,opacity=0.72]
        plot[domain=0:360, samples=120, smooth] ({\cx + \ra*cos(\x)}, {\cy + \rb*sin(\x)}, 0) -- cycle;
        
    \draw[thick] (2.15,1.15,2.82) -- (2.15,1.15,3.62);
    \draw[thick] (2.15,1.15,-0.1) -- (2.15,1.15,0.75);
    \draw[thick,dashed] (2.15,1.15,0.75) -- (2.15,1.15,2.82);

    
    \node[blue!70!black, above] at (3.2,1.4,1) {$\Omega$};
    \node[red!70!black, above] at (3.2,1.4,0) {$D_{xy}$};

    % ----- 引线标注 -----
    % 顶面标注：从右上方引线指向曲面
    \draw[thick]
        (4.0, 2.3, {F(3.4,1.65)+1}) node[right] {$z_2=z_2(x,y)$}
        -- (3.4, 1.65, {F(3.4,1.65)+1});

    % 底面标注：从右侧引线指向底域 D
    \draw[thick]
    (5.2, 0.2, 1.1) node[right] {$z_1=z_1(x,y)$}
    -- (3.692, 1.223, 1);
\end{tikzpicture}
\end{center}

当然，投影不一定非要投到 \(xOy\) 面，根据区域\(\Omega\)形状，投到 \(xOz\) 面或 \(yOz\) 面也一样，道理不变。

如果沿着某条轴看过去，区域 \(\Omega\) 的截面形状特别简单，那就更省事了。比如用垂直于 \(z\) 轴的平面 \(z = k\)（$k$为常数）去截 \(\Omega\)，截出来一个平面区域 \(D_z\)。先固定 \(z\)，在 \(D_z\) 上对 \(x,y\) 做二重积分，得到这一薄片的累积量，然后把所有薄片从 \(z=c\) 积到 \(z=d\)。这就是截面法的思想。

\begin{theorem}{三重积分的截面法}{}
设空间闭区域 \(\Omega\) 夹在平面 \(z = c\) 与 \(z = d\) 之间（\(c < d\)），对于每个 \(z \in [c,d]\)，用平面 \(z = k\)（$k$为常数）截 \(\Omega\) 所得的截面区域为 \(D_z\)。若 \(f(x,y,z)\) 在 \(\Omega\) 上可积，则
\[
\iiint_{\Omega} f(x,y,z)\,dxdydz
= \int_{c}^{d} \left( \iint_{D_z} f(x,y,z)\,dxdy \right) dz
\]
\end{theorem}

\begin{center}
    \begin{tikzpicture}[
        scale=1.05,
        x={(1cm,0cm)},
        y={(0.55cm,0.32cm)},
        z={(0cm,1cm)},
        line join=round,
        line cap=round
    ]
        % 椭圆形底域 D 的参数
        \def\cx{2.15}
        \def\cy{1.15}
        \def\ra{1.55}
        \def\rb{0.78}
        \def\H{1.56}   % 椭圆柱的高度（常数）

        % 坐标轴
        \draw[->] (-0.25,0,0) -- (4.45,0,0) node[below] {$x$};
        \draw[->] (0,-0.2,0) -- (0,2.85,0) node[right] {$y$};
        \draw[->] (0,0,0) -- (0,0,2.65) node[left] {$z$};
        \node[below left] at (0,0,0) {$O$};

        % 侧面 —— 用不透明的浅色填充，顶部高度为常数 \H
        \fill[cyan!20]
            plot[domain=0:360, samples=120, smooth]
            ({\cx + \ra*cos(\x)}, {\cy + \rb*sin(\x)}, \H)
            -- plot[domain=360:0, samples=120, smooth]
            ({\cx + \ra*cos(\x)}, {\cy + \rb*sin(\x)}, 0)
            -- cycle;

        % 母线 —— 用虚线表达柱面的棱边
        \foreach \ang in {0,1,...,359} {
            \pgfmathsetmacro\xx{\cx + \ra*cos(\ang)}
            \pgfmathsetmacro\yy{\cy + \rb*sin(\ang)}
            \draw[cyan!40, line width=0.4mm] (\xx,\yy,0) -- (\xx,\yy,\H);
        }

        % 顶面（平面 z = \H）
        \fill[blue!18, opacity=0.72]
            plot[domain=0:360, samples=120, smooth]
            ({\cx + \ra*cos(\x)}, {\cy + \rb*sin(\x)}, \H)
            -- cycle;

        % 底域 D 的轮廓（可见部分实线，不可见部分虚线）
        \draw[cyan!70!black, thick]
            plot[domain=195:360, samples=80, smooth]
            ({\cx + \ra*cos(\x)}, {\cy + \rb*sin(\x)}, 0);
        \draw[cyan!70!black, thick, dashed]
            plot[domain=0:195, samples=80, smooth]
            ({\cx + \ra*cos(\x)}, {\cy + \rb*sin(\x)}, 0);

        % 顶面轮廓线
        \draw[blue!75!black, very thick]
            plot[domain=0:360, samples=120, smooth]
            ({\cx + \ra*cos(\x)}, {\cy + \rb*sin(\x)}, \H);

        % 标签
        \node[blue!70!black, above] at (3.35,1.8,\H) {$z=d$};
        \node[blue!70!black, above] at (1.6,1.2,0) {$\Omega$};
        \node[cyan!70!black, right] at (3.35,1.8,0) {$z=c$};

        % 红色的截面 D_z（位于 z=0.78，可根据需要调整）
        \draw[red!70!black, thick]
            plot[domain=195:360, samples=80, smooth]
            ({\cx + \ra*cos(\x)}, {\cy + \rb*sin(\x)}, 0.78);
        \draw[red!70!black, thick, dashed]
            plot[domain=0:195, samples=80, smooth]
            ({\cx + \ra*cos(\x)}, {\cy + \rb*sin(\x)}, 0.78);
        \fill[red!24!white, opacity=0.72]
            plot[domain=0:360, samples=120, smooth]
            ({\cx + \ra*cos(\x)}, {\cy + \rb*sin(\x)}, 0.78) -- cycle;
        \node[red!70!black] at (4.76,0,1.28) {$D_z$};
    \end{tikzpicture}
\end{center}

当被积函数只和 \(z\) 有关，或者截面面积很容易写出来的时候，截面法往往比投影法简洁得多。

仔细观察我们发现，投影法是先对一个变量积分，再对两个变量积分，所以投影法又叫“先一后二法”。截面法是先对两个变量积分，再对一个变量积分，所以截面法又叫“先二后一法”。

下面看一个例子具体运用这两种方法。

\begin{myexample}
计算三重积分 \(\ \iiint_{\Omega} z\,dV\)，其中 \(\Omega\) 是由平面 \(x=0,\;y=0,\;z=0\) 以及 \(x+y+z=1\) 所围成的四面体。

区域\(\Omega\)是第一卦限里被坐标面切出来的一个四面体。我们先用投影法，把它投到 \(xOy\) 面上。投影区域 \(D\) 就是直线 \(x+y=1\) 与 \(x\) 轴、\(y\) 轴围成的三角形。

\begin{center}
    \begin{tikzpicture}[
        scale=3,
        line join=round,
        x={(-0.9cm,0cm)},
    y={(0.4cm,-0.2cm)},
    z={(0cm,1.2cm)},
        % 面样式
        bottom/.style={fill=gray!50,fill opacity=0.5},
        side/.style={fill=gray!50,fill opacity=0.4},
        slant/.style={fill=gray!50,fill opacity=0.5},
        edge/.style={thick},
    ]
    
    % 定义顶点
    \coordinate (O) at (0,0,0);
    \coordinate (A) at (1,0,0);
    \coordinate (B) at (0,1,0);
    \coordinate (C) at (0,0,1);
    
    % 填充各个面（半透明，所有棱都可见）
    \fill[bottom] (O) -- (A) -- (B) -- cycle;       % z=0 底面
    \fill[side]   (O) -- (A) -- (C) -- cycle;       % y=0 面
    \fill[side]   (O) -- (B) -- (C) -- cycle;       % x=0 面
    \fill[slant]  (A) -- (B) -- (C) -- cycle;       % x+y+z=1 斜面
    
    % 绘制所有棱
    \draw[edge] (A) -- (B);
    \draw[edge] (A) -- (C);
    \draw[edge] (B) -- (C);
    
    % 顶点标签
    \node [above left] at (O) {$O$};
    \node [above left] at (A) {$A$};
    \node [above right] at (B) {$B$};
    \node [above left] at (C) {$C$};

    \def\px{0.35}                     % 选取 x 坐标
    \def\py{0.2}                     % 选取 y 坐标
    \pgfmathsetmacro\pz{1-\px-\py}   % 计算上端点 z 坐标
    \coordinate (P0) at (\px,\py,0);
    \coordinate (P1) at (\px,\py,\pz);
    % 画穿线（带箭头）
    \draw[red!70!black, line width=1.2pt] (0.35,0.2,-0.1) -- (P0);
    \draw[red!70!black, line width=1.2pt,dashed] (P0) -- (P1);
    \draw[red!70!black, line width=1.2pt] (P1) -- (0.35,0.2,1);

    \fill[red!70!black] (P0) circle (0.5pt);
    \fill[red!70!black] (P1) circle (0.5pt);

    \node[left, red!70!black] at (P0) {$P_0$};
    \node[left, red!70!black] at (P1) {$P_1$};
    
    % 坐标轴（可选）
    \draw[-stealth,thick] (1,0,0) -- (1.3,0,0) node[below right] {$x$};
    \draw[-stealth,thick] (0,1,0) -- (0,1.3,0) node[below left] {$y$};
    \draw[-stealth,thick] (0,0,1) -- (0,0,1.3) node[above] {$z$};
    \draw[dashed,thick] (1,0,0) -- (O);
    \draw[dashed,thick] (0,1,0) -- (O);
    \draw[dashed,thick] (0,0,1) -- (O);
    \end{tikzpicture}
    \qquad
    \begin{tikzpicture}[scale=3]
    % 填充投影区域（灰色半透明）
    \fill[gray!50, opacity=0.5] (0,0) -- (1,0) -- (0,1) -- cycle;
    % 边界
    \draw[thick] (0,0) -- (1,0) -- (0,1) -- cycle;
    % 顶点
    \node at (0,0)[below left] {$O$};
    \node at (1,0)[below right] {$A$};
    \node at (0,1)[above left] {$B$};
    % 坐标轴
    \draw[-stealth, thick] (0,0) -- (1.3,0) node[below] {$x$};
    \draw[-stealth, thick] (0,0) -- (0,1.3) node[left] {$y$};
\end{tikzpicture}
\end{center}

我们用“画图穿线法”沿着 \(z\) 轴方向穿线，从底面 \(z=0\) 进入，从平面 \(z = 1-x-y\) 穿出，交点分别是$P_0$，$P_1$。所以：
\[
\iiint_{\Omega} z\,dV
= \iint_{D} \left( \int_{0}^{1-x-y} z \, dz \right) dxdy
\]

内层对 \(z\) 积分时 \(x,y\) 是常数，积出来就是 \(\frac12(1-x-y)^2\)。于是原积分化成二重积分：
\[
\iiint_{\Omega} z\,dV = \iint_{D} \frac12(1-x-y)^2\,dxdy
= \frac12 \int_{0}^{1} dx \int_{0}^{1-x} (1-x-y)^2\,dy
\]

先对 \(y\) 积分：
\[
\int_{0}^{1-x} (1-x-y)^2\,dy
= \left[ -\frac13 (1-x-y)^3 \right]_{0}^{1-x}
= \frac13 (1-x)^3
\]

再对 \(x\) 积分：
\[
\frac12 \cdot \frac13 \int_{0}^{1} (1-x)^3\,dx
= \frac16 \left[ -\frac14 (1-x)^4 \right]_{0}^{1}
= \frac16 \cdot \frac14 = \frac{1}{24}
\]

所以 \(\iiint_{\Omega} z\,dV = \frac{1}{24}\)。

再用截面法，用平面 \(z = k\)（$k$为常数）截 \(\Omega\)，我们能得到一个等腰直角三角形，这个等腰直角三角形的边长随\(z\)的变化而变化。由于截面上的点在区域\(\Omega\)内部，所以它们满足\(x+y+z\le1\)，也就是\(x+y\le1-z\)，这说明等腰直角三角形的边长就是\(1-z\)。

\begin{center}
    \begin{tikzpicture}[
        scale=3,
        line join=round,
        x={(-0.9cm,0cm)},
        y={(0.4cm,-0.2cm)},
        z={(0cm,1.2cm)},
        % 面样式
        bottom/.style={fill=gray!50,fill opacity=0.5},
        side/.style={fill=gray!50,fill opacity=0.4},
        slant/.style={fill=gray!50,fill opacity=0.5},
        edge/.style={thick},
    ]
    
    % 定义顶点
    \coordinate (O) at (0,0,0);
    \coordinate (A) at (1,0,0);
    \coordinate (B) at (0,1,0);
    \coordinate (C) at (0,0,1);
    
    % 填充各个面（半透明，所有棱都可见）
    \fill[bottom] (O) -- (A) -- (B) -- cycle;       % z=0 底面
    \fill[side]   (O) -- (A) -- (C) -- cycle;       % y=0 面
    \fill[side]   (O) -- (B) -- (C) -- cycle;       % x=0 面
    \fill[slant]  (A) -- (B) -- (C) -- cycle;       % x+y+z=1 斜面
    
    % 绘制所有棱
    \draw[edge] (A) -- (B);
    \draw[edge] (A) -- (C);
    \draw[edge] (B) -- (C);
    
    \def\px{0.35}                     % 选取 x 坐标
    \def\py{0.2}                     % 选取 y 坐标
    \pgfmathsetmacro\pz{1-\px-\py}   % 截面高度：与穿线终点一致 (0.45)
    % 截面三角形顶点 (在高度 z = \pz 处)
    \coordinate (Dz0) at (0,0,\pz);
    \coordinate (Dz1) at ({1-\pz},0,\pz);
    \coordinate (Dz2) at (0,{1-\pz},\pz);
    % 绘制截面（半透明绿色填充，边界线加粗）
    \fill[green!30, opacity=0.6] (Dz0) -- (Dz1) -- (Dz2) -- cycle;
    \draw[green!60!black, thick, dashed] (Dz2) -- (Dz0) -- (Dz1);
    \draw[green!60!black, thick] (Dz2) -- (Dz1);
    % 标签
    \node[green!60!black, above left] at (0,0,\pz) {$D_z$};
    
    \node[green!60!black, above right] at (Dz0) {$P_0$};
    \node[green!60!black, above left] at (Dz1) {$P_1$};
    \node[green!60!black, above right] at (Dz2) {$P_2$};

    % 顶点标签
    \node [above left] at (O) {$O$};
    \node [above left] at (A) {$A$};
    \node [above right] at (B) {$B$};
    \node [above left] at (C) {$C$};


    % 坐标轴
    \draw[-stealth,thick] (1,0,0) -- (1.3,0,0) node[below right] {$x$};
    \draw[-stealth,thick] (0,1,0) -- (0,1.3,0) node[below left] {$y$};
    \draw[-stealth,thick] (0,0,1) -- (0,0,1.3) node[above] {$z$};
    \draw[dashed,thick] (1,0,0) -- (O);
    \draw[dashed,thick] (0,1,0) -- (O);
    \draw[dashed,thick] (0,0,1) -- (O);
    \end{tikzpicture}
    \qquad
    \begin{tikzpicture}[scale=3]
    % 填充投影区域（灰色半透明）
    \fill[gray!50, opacity=0.5] (0,0) -- (1,0) -- (0,1) -- cycle;
    % 边界
    \draw[thick] (0,0) -- (1,0) -- (0,1) -- cycle;
    % 顶点
    \node at (0,0)[below left] {$P_0$};
    \node at (1,0)[below right] {$P_1$};
    \node at (0,1)[above left] {$P_2$};
    \node at (0.25,0.25) {$D_z$};
    % 坐标轴
    \draw[-stealth, thick] (0,0) -- (1.3,0) node[below] {$x$};
    \draw[-stealth, thick] (0,0) -- (0,1.3) node[left] {$y$};
    \end{tikzpicture}
\end{center}

所以截面\(D_z\)的面积为\(\frac12 (1-z)^2\)。先对截面积分，即：
\[
\iiint_{\Omega} z\,dV
= \int_{0}^{1} \left( \iint_{D_z} z\,dxdy \right) dz
\]

二重积分里对\(x,y\)积分，因此\(z\)是常数。把\(z\)提出来，就是：
\[
\int_{0}^{1} \left( \iint_{D_z} z\,dxdy \right) dz=\int_{0}^{1} z \left( \iint_{D_z} \,dxdy \right) dz=\int_{0}^{1} [z\cdot\frac12 (1-z)^2] \,dz
\]

所以：
\[
\int_{0}^{1} [z\cdot\frac12 (1-z)^2] \,dz
= \frac12 \int_{0}^{1} \bigl( z - 2z^2 + z^3 \bigr) \, dz
= \frac12 \left[ \frac{z^2}{2} - \frac{2z^3}{3} + \frac{z^4}{4} \right]_{0}^{1}
= \frac12 (\frac12 - \frac23 + \frac14)
= \frac{1}{24}
\]

所以 \(\iiint_{\Omega} z\,dV = \frac{1}{24}\)。
\end{myexample}

\begin{myexample}
计算三重积分 \(\iiint_{\Omega} z\,dV\)，其中 \(\Omega\) 是由圆锥面 \(z = \sqrt{x^2+y^2}\) 与平面 \(z = 1\) 所围成的区域。

我们用截面法计算这个积分。用垂直于 \(z\) 轴的平面去截这个区域，所得的截面都是圆。圆的半径随\(z\)的变化而变化，对于 \(z \in [0,1]\)，截面 \(D_z\) 是圆 \(x^2+y^2 \le z^2\)，面积就是 \(\pi z^2\)。

\begin{center}
    \begin{tikzpicture}[
        scale=2.5,
        line join=round,
        % 视角调整：使顶点在下，底面在上
        x={(-0.866cm,-0.5cm)},
        y={(0.866cm,-0.5cm)},
        z={(0cm,1cm)},
    ]
    
    % 圆锥参数
    \def\R{1}    % 底面半径
    \def\H{1}    % 高度
    
    % 顶点与底面圆心
    \coordinate (V) at (0,0,0);
    \coordinate (Otop) at (0,0,\H);
    
    
    \fill[cyan!40, opacity=0.65, even odd rule]
        (V) -- ({\R*cos(0)}, {\R*sin(0)}, \H) -- ({\R*cos(90)}, {\R*sin(90)}, \H) -- cycle
        % 弓形：从 A 沿弧到 B，再 -- cycle 闭合（弦 BA）
        plot[domain=0:90, samples=120, variable=\t, smooth]
            ({\R*cos(\t)}, {\R*sin(\t)}, \H)
        -- cycle;
    
    % ---- 顶面（半透明填充） ----
    \fill[blue!20, opacity=0.65]
        plot[domain=0:360, samples=120, smooth]
        ({\R*cos(\x)}, {\R*sin(\x)}, \H) -- cycle;
    
    % ---- 底面圆的轮廓线 ----
    \draw[thick]
        plot[domain=0:360, samples=120, smooth]
        ({\R*cos(\x)}, {\R*sin(\x)}, \H);
    
    
    \draw[very thick, cyan!70!black] (V) -- ({\R*cos(0)}, {\R*sin(0)}, \H);
    
    \draw[very thick, cyan!70!black] (V) -- ({\R*cos(90)}, {\R*sin(90)}, \H);
    
    
    % ---- 坐标轴（从顶点向外延伸） ----
    \draw[-stealth, thick] (V) -- (1.3,0,0) node[below] {$x$};
    \draw[-stealth, thick] (V) -- (0,1.3,0) node[below left] {$y$};
    \draw[-stealth, thick] (0,0,1) -- (0,0,1.3) node[above] {$z$};
    \draw[thick, dashed] (V) -- (0,0,1);
    
    % 标签
    \node[below] at (V) {$O$};
    \node[above right] at (0,0,\H) {$z=1$};
    
    \end{tikzpicture}
    \quad
    \begin{tikzpicture}[
    scale=2.5,
    line join=round,
    % 视角调整：使顶点在下，底面在上
    x={(-0.866cm,-0.5cm)},
    y={(0.866cm,-0.5cm)},
    z={(0cm,1cm)},
]

% 圆锥参数
\def\R{1}    % 底面半径
\def\H{1}    % 高度

% 顶点与底面圆心
\coordinate (V) at (0,0,0);
\coordinate (Otop) at (0,0,\H);


\fill[cyan!40, opacity=0.65, even odd rule]
    (V) -- ({\R*cos(0)}, {\R*sin(0)}, \H) -- ({\R*cos(90)}, {\R*sin(90)}, \H) -- cycle
    % 弓形：从 A 沿弧到 B，再 -- cycle 闭合（弦 BA）
    plot[domain=0:90, samples=120, variable=\t, smooth]
        ({\R*cos(\t)}, {\R*sin(\t)}, \H)
    -- cycle;

% ---- 顶面（半透明填充） ----
\fill[blue!20, opacity=0.65]
    plot[domain=0:360, samples=120, smooth]
    ({\R*cos(\x)}, {\R*sin(\x)}, \H) -- cycle;

    \fill[red!20, opacity=0.72]
    plot[domain=0:360, samples=120, smooth]
    ({0.5*\R*cos(\x)}, {0.5*\R*sin(\x)}, 0.5) -- cycle;
    
    \draw[red!70, thick]
    plot[domain=0:90, samples=120, smooth]
    ({0.5*\R*cos(\x)}, {0.5*\R*sin(\x)}, 0.5);

    \draw[red!70, thick, dashed]
    plot[domain=90:360, samples=120, smooth]
    ({0.5*\R*cos(\x)}, {0.5*\R*sin(\x)}, 0.5);

    \node[red!70, right] at (0,0,0.75) {$D_z$};

% ---- 底面圆的轮廓线 ----
\draw[thick]
    plot[domain=0:360, samples=120, smooth]
    ({\R*cos(\x)}, {\R*sin(\x)}, \H);


\draw[very thick, cyan!70!black] (V) -- ({\R*cos(0)}, {\R*sin(0)}, \H);

\draw[very thick, cyan!70!black] (V) -- ({\R*cos(90)}, {\R*sin(90)}, \H);


% ---- 坐标轴（从顶点向外延伸） ----
\draw[-stealth, thick] (V) -- (1.3,0,0) node[below] {$x$};
\draw[-stealth, thick] (V) -- (0,1.3,0) node[below left] {$y$};
\draw[-stealth, thick] (0,0,1) -- (0,0,1.3) node[above] {$z$};
\draw[thick, dashed] (V) -- (0,0,1);

% 标签
\node[below] at (V) {$O$};
\node[above right] at (0,0,\H) {$z=1$};


\end{tikzpicture}
\end{center}

被积函数只含 \(z\)，所以：
\[
\iiint_{\Omega} z\,dV
= \int_{0}^{1} z \left( \iint_{D_z} dxdy \right) dz
\]

已知截面面积为 \(\pi z^2\)，代入二重积分，得：
\[
\int_{0}^{1} z \left( \iint_{D_z} dxdy \right) dz
= \int_{0}^{1} z \cdot \pi z^2 \, dz
= \pi \int_{0}^{1} z^3\,dz
= \frac{\pi}{4}
\]

所以 \( \iiint_{\Omega} z\,dV = \frac{\pi}{4}\)。
\end{myexample}

总的来说，直角坐标下算三重积分，关键的是先看清区域的形状，选择合适的方法。投影法投到平面再穿线，和算二重积分思路十分相似。截面法像切面包片，用垂直于 \(z\) 轴的平面去截区域，然后计算截面面积，可以理解为升维的“画图穿线法”。

\subsection{三重积分的柱面坐标形式}

在研究二重积分时我们发现，当积分区域是圆、环或扇形，或者被积函数含有 \(x^2+y^2\) 这类项时，用极坐标往往能大大简化计算。在三重积分里，类似的场景也会出现。如果空间区域 \(\Omega\) 在某个平面上的投影是圆形或扇形区域，而被积函数又含有 \(x^2+y^2\) 之类的项，那么就可以考虑用柱面坐标。

柱面坐标其实就是平面极坐标加上一个竖坐标 \(z\)。我们知道，在平面上确定一个点可以用极坐标 \((\rho,\theta)\)，那么在空间中，要确定一个点，只需要在 \((\rho,\theta)\) 的基础上再给出它的竖坐标 \(z\) 就可以了。也就是说，空间中一点 \(M\) 的柱面坐标是 \((\rho,\theta,z)\)，其中 \(0\le \rho<+\infty\)，\(0\le|\theta|\le 2\pi\)，\(-\infty<z<+\infty\)。

\begin{center}
\begin{tikzpicture}[
    scale=1.4,
    x={(1cm,0cm)},
    y={(0.4cm,0.35cm)},
    z={(0cm,1cm)},
    line join=round,
    line cap=round
]
    % 坐标轴
    \draw[->] (-0.2,0,0) -- (3.2,0,0) node[below] {$x$};
    \draw[->] (0,-0.2,0) -- (0,2.8,0) node[right] {$y$};
    \draw[->] (0,0,0) -- (0,0,2.8) node[left] {$z$};
    \node[below left] at (0,0,0) {$O$};

    % 定义点 P 的坐标
    \def\rr{2.2}
    \def\th{55}
    \def\zz{2.2}

    % 计算点 P 的直角坐标
    \pgfmathsetmacro\xx{\rr*cos(\th)}
    \pgfmathsetmacro\yy{\rr*sin(\th)}

    % 点 P
    \coordinate (P) at (\xx,\yy,\zz);
    \fill (P) circle (1.2pt) node[above right] {$M(\rho,\theta,z)$};

    % 投影点 P' 在 xOy 平面上
    \coordinate (Pxy) at (\xx,\yy,0);
    \fill (Pxy) circle (1.2pt) node[below right] {$M'(\rho,\theta,0)$};

    % 连接
    \draw[dashed] (P) -- (Pxy);
    \draw[dashed] (0,0,0) -- (Pxy);

    % 标出 z
    \draw[thick] (\xx,\yy,0) -- (\xx,\yy,\zz);
    \node[right] at (\xx,\yy,\zz/2) {$z$};

    % 标出 rho
    \node[above] at (\xx/2,\yy/2,0) {$\rho$};

    % 标出 theta 角
    \draw[->] (0.5,0,0) arc (0:\th:0.5);
    \node[right] at (0.54,0.34,0) {$\theta$};

\end{tikzpicture}
\end{center}

也就是说，柱面坐标与极坐标之间的关系只不过是多了竖坐标 \(z\) ，那么直角坐标转换为柱面坐标也就和极坐标类似：
\[
x = \rho\cos\theta,\quad
y = \rho\sin\theta,\quad
z = z
\]

那么柱面坐标下的体积微元是什么形式呢？回忆二重极坐标的面积微元，是用无数多个同心圆和无数条从原点出发的射线把平面分割成小网格，每个小网格近似为一个“曲边矩形”，边长分别是 \(\Delta \rho\) 和 \(\rho\Delta\theta\)，面积近似为 \(\rho\,\Delta \rho\,\Delta\theta\)，取极限得到 \(d\sigma = \rho\,d\rho\,d\theta\)。

\begin{center}
    \begin{tikzpicture}[
        x={(-0.866cm, -0.5cm)},
        y={(0.866cm, 0cm)},
        z={(0cm, 1cm)},
        scale=2,
        font=\small
    ]
    
    % ---------- 几何参数 ----------
    \def\myr{2}       % 内径
    \def\mydr{1}      % 径向增量
    \def\dtheta{20}   % 角度增量（度）
    \def\z{1.2}         % 底面高度
    \def\dz{0.5}      % 高度增量
    
    % 外径
    \pgfmathtruncatemacro{\myR}{\myr+\mydr}
    
    % 起止角度（使微元在视角下对称一些）
    \pgfmathsetmacro{\thetaStart}{-0.5*\dtheta+60}
    \pgfmathsetmacro{\thetaEnd}{0.5*\dtheta+60}
    
    % ---------- 关键点坐标 ----------
    % 底面内圆柱面端点
    \coordinate (A) at ({\myr*cos(\thetaStart)}, {\myr*sin(\thetaStart)}, \z);
    \coordinate (B) at ({\myr*cos(\thetaEnd)},   {\myr*sin(\thetaEnd)},   \z);
    % 底面外圆柱面端点
    \coordinate (C) at ({\myR*cos(\thetaStart)}, {\myR*sin(\thetaStart)}, \z);
    \coordinate (D) at ({\myR*cos(\thetaEnd)},   {\myR*sin(\thetaEnd)},   \z);
    % 顶面内圆柱面端点
    \coordinate (E) at ({\myr*cos(\thetaStart)}, {\myr*sin(\thetaStart)}, \z+\dz);
    \coordinate (F) at ({\myr*cos(\thetaEnd)},   {\myr*sin(\thetaEnd)},   \z+\dz);
    % 顶面外圆柱面端点
    \coordinate (G) at ({\myR*cos(\thetaStart)}, {\myR*sin(\thetaStart)}, \z+\dz);
    \coordinate (H) at ({\myR*cos(\thetaEnd)},   {\myR*sin(\thetaEnd)},   \z+\dz);
    
    \fill[gray!50,opacity=0.5] 
      (E) arc[start angle=\thetaStart, end angle=\thetaEnd, radius=\myr] 
      -- (H) -- (D) arc[start angle=\thetaEnd, end angle=\thetaStart, radius=\myR] -- (A)
      -- cycle;
    
    % ---------- 绘制线框 ----------
    % 底面
    \draw[dashed]    (A) arc[start angle=\thetaStart, end angle=\thetaEnd, samples=120, radius=\myr];
    \draw[thick]     (C) arc[start angle=\thetaStart, end angle=\thetaEnd, samples=120, radius=\myR];
    \draw[thick]    (A) -- (C);
    \draw[dashed]     (B) -- (D);
    
    % 顶面
    \draw[thick]    (E) arc[start angle=\thetaStart, end angle=\thetaEnd, samples=120, radius=\myr];
    \draw[thick]     (G) arc[start angle=\thetaStart, end angle=\thetaEnd, samples=120, radius=\myR];
    \draw[thick]     (E) -- (G);
    \draw[thick]     (F) -- (H);
    
    % 侧棱
    \draw[thick] (A) -- (E);
    \draw[dashed] (B) -- (F);
    
    
    
    % ---------- 尺寸标注 ----------
    % dz（沿竖棱）
    \draw[thick] (D) -- (H) 
        node[midway, right=2pt, inner sep=0] {$dz$};
    
    % d\rho（沿径向外表面）
    \draw[thick] (C) -- (G) 
        node[midway, left=12pt, below=2pt, inner sep=2pt] {$d\rho$};
    
    % \rho d\theta（沿顶面外圆弧）
    \draw[thick] 
        (G) arc[start angle=\thetaStart, end angle=\thetaEnd, radius=\myR, samples=120] 
        node[midway, above=22pt, left=10pt, inner sep=0] {$\rho\,d\theta$};
    
        \draw[->] ({0.3*\myr*cos(\thetaStart)},{0.3*\myr*sin(\thetaStart)},0) arc(\thetaStart:\thetaEnd:0.6)
        node[midway, left=10pt, inner sep=0]{$d\theta$};
    
    
    % 可选：标注内半径 \rho（参考线）
    \draw[dashed] (0,0,\z+\dz) -- (E) 
        node[midway, right, inner sep=0] {$\rho$};
    
        \draw[dashed] (0,0,\z+\dz) -- (F);
    
        \draw[dashed] (0,0,\z) -- (A);
    
        \draw[dashed] (0,0,\z) -- (B);
    
        \draw[dashed] ({\myr*cos(\thetaStart)}, {\myr*sin(\thetaStart)}, 0) -- (A);
    
        \draw[dashed] ({\myr*cos(\thetaEnd)}, {\myr*sin(\thetaEnd)}, 0) -- (B);
    
        \draw[dashed] ({\myR*cos(\thetaStart)}, {\myR*sin(\thetaStart)}, 0) -- (C);
    
        \draw[dashed] ({\myR*cos(\thetaEnd)}, {\myR*sin(\thetaEnd)}, 0) -- (D);
    
        \draw[dashed] ({\myR*cos(\thetaStart)}, {\myR*sin(\thetaStart)}, 0) -- (0,0,0);
    
        \draw[dashed] ({\myR*cos(\thetaEnd)}, {\myR*sin(\thetaEnd)}, 0) -- (0,0,0);
    
        \draw[dashed]    ({\myr*cos(\thetaStart)}, {\myr*sin(\thetaStart)}, 0) arc[start angle=\thetaStart, end angle=\thetaEnd, samples=120, radius=\myr];
    
        \draw[dashed]    ({\myR*cos(\thetaStart)}, {\myR*sin(\thetaStart)}, 0) arc[start angle=\thetaStart, end angle=\thetaEnd, samples=120, radius=\myR];
    
    % ---------- 坐标轴（可选） ----------
    \draw[->,thick] (0,0,0) -- (1.6,0,0) node[below left] {$x$};
    \draw[->,thick] (0,0,0) -- (0,2.2,0) node[below right] {$y$};
    \draw[->,thick] (0,0,0) -- (0,0,2.2) node[above] {$z$};
    
    \end{tikzpicture}
\end{center}

和二重积分类似，在三重积分里，我们取一个小区域，它在 \(\rho\) 方向上的长度是 \(\Delta \rho\)，在 \(\theta\) 方向上的“弧长”是 \(\rho\Delta\theta\)，在 \(z\) 方向上的高度是 \(\Delta z\)。当 \(\Delta \rho,\Delta\theta,\Delta z\) 都很小时，这个小区域近似为一个长方体，长宽高分别为 \(d\rho\)、\(\rho\,d\theta\)、\(dz\)。因此体积微元就是这三者的乘积，也就是\(dV = \rho\,d\rho\,d\theta\,dz\)，要注意的是，\(dV \ne d\rho\,d\theta\,dz\)，而是\(dV = \textcolor{red}{\rho}\,d\rho\,d\theta\,dz\)。

既然有了体积微元，三重积分在柱面坐标下的计算就呼之欲出了。我们只需要令\(x = \rho\cos\theta,y = \rho\sin\theta,z = z\)，就得到：
\begin{theorem}{三重积分柱面坐标变换公式}{}
    \[
\iiint_{\Omega} f(x,y,z)\,dxdydz
= \iiint_{\Omega} f(\rho\cos\theta,\, \rho\sin\theta,\, z)\; \rho\,d\rho\,d\theta\,dz
\]
\end{theorem}


在积分次序上，柱面坐标和直角坐标的投影法非常相似。我们通常把 \(\Omega\) 往 \(xOy\) 平面投影，得到一个用极坐标描述的平面区域 \(D_{\rho\theta}\)，然后在 \(D_{\rho\theta}\) 的每一个点从下往上穿线，穿入的是积分下限 \(z=z_1(\rho,\theta)\)，穿出的是积分上限 \(z=z_2(\rho,\theta)\)，先对\(z\)积分，再对 \(\rho,\theta\) 做二重积分就能求出三重积分了：

\begin{theorem}{柱面坐标下的三重积分}{}
设空间闭区域 \(\Omega\) 可表示为
\[
\Omega 
= \{(\rho\cos\theta, \rho\sin\theta, z) \mid(\rho\cos\theta, \rho\sin\theta) \in D_{\rho\theta},\;
z_1(\rho,\theta) \le z \le z_2(\rho,\theta)\}
\]
其中 \(D_{\rho\theta}\) 是 \(\Omega\) 在 \(xOy\) 面上的投影区域，\(z_1(\rho,\theta)\)，\(z_2(\rho,\theta)\) 在 \(D_{\rho\theta}\) 上连续。若 \(f(x,y,z)\) 在 \(\Omega\) 上可积，则

\[
\iiint_{\Omega} f(\rho\cos\theta, \rho\sin\theta, z)\,\rho\,d\rho\,d\theta\,dz
= \iint_{D_{\rho\theta}} \left( \int_{z_1(\rho,\theta)}^{z_2(\rho,\theta)} f(\rho\cos\theta, \rho\sin\theta, z)\,dz \right) \rho\,d\rho\,d\theta
\]
\end{theorem}

在极坐标平面区域 \(D_{\rho\theta}\) 上继续算二重积分时，具体操作就和极坐标的二重积分一样。确定上下限的操作依旧是“画图穿线法”，只不过现在是在柱面坐标系下进行，也就是说要先投影到某个平面上再确定。这样一步步做下来，三重积分就被拆解为三次定积分。

下面看几个例子。

\begin{myexample}
    计算三重积分 \(\iiint_{\Omega} (x^2+y^2)\,dV\)，其中 \(\Omega\) 是由圆柱面 \(x^2+y^2=2x\) 与平面 \(z=0\)、\(z=4\) 围成的区域。

    先令 \(x=\rho\cos\theta\)，\(y=\rho\sin\theta\)，\(z=z\)，则体积微元 \(dV=\rho\,d\rho\,d\theta\,dz\)，被积函数 \(x^2+y^2=\rho^2\)。

    再看积分区域。方程 \(x^2+y^2=2x\) 化为极坐标就是 \(\rho^2=2\rho\cos\theta\)，即 \(\rho=2\cos\theta\)。因为 \(z\) 的范围是 \(0\) 到 \(4\)，所以 \(\Omega\) 是一个底面为\(\rho=2\cos\theta\)、高为 \(4\) 的圆柱体。

    先看 \(z\) 方向。我们依然用一条直线穿入\(\Omega\)，直线从底面 \(z=0\) 穿入，从顶面 \(z=4\) 穿出。所以 \(z\) 的积分下限是 \(0\)，上限是 \(4\)。
    
    再把 \(\Omega\) 投影到 \(xOy\) 平面，得到\(D_{\rho\theta}\)。为保证 \(\rho\ge0\)，所以 \(\cos\theta\ge0\)，因此 \(\theta \in [-\frac{\pi}{2},\frac{\pi}{2}]\)。在每一个 \(\theta\) 处，从极点出发的射线穿入区域时 \(\rho=0\)，碰到圆的边界穿出时 \(\rho=2\cos\theta\)。所以积分下限为\(\rho=0\)，积分上限为\(\rho=2\cos\theta\)。

    \begin{center}
\begin{tikzpicture}[scale=1.8]
    % 极坐标网格与区域填充
    \fill[gray!30] 
        plot[domain=-90:90, smooth] (\x:{2*cos(\x)}) -- cycle;
    
    % 画出圆 \rho=2\cos\theta
    \draw[thick]
        plot[domain=-90:90, smooth] (\x:{2*cos(\x)});
    
    % 坐标轴
    \draw[->] (-0.2,0) -- (2.5,0) node[right] {$x$};
    \draw[->] (0,-1.2) -- (0,1.2) node[above] {$y$};
    
    % 一条示例射线
    \draw[thick] (0,0) -- (40:{2.5*cos(40)});
    \draw[->] (0.35,0) arc (0:40:0.35) node[midway, right] {$\theta$};
    
    % 标注
    \node at (1.5,0.5) {$D_{\rho\theta}$};
    \node at (2.2,0.9) {$\rho=2\cos\theta$};

\end{tikzpicture}
\end{center}

    
    于是三重积分化为：
    \[
    \iiint_{\Omega} (x^2+y^2)\,dV
    = \iint_{D_{\rho\theta}}\left(\int_{0}^{4} \rho^2 \, dz\right) \, \rho d\rho d\theta
    = \int_{-\frac{\pi}{2}}^{\frac{\pi}{2}} d\theta \int_{0}^{2\cos\theta} d\rho \int_{0}^{4} \rho^3 \, dz
    \]

    先对 \(z\) 积分，被积函数 \(\rho^3\) 与 \(z\) 无关，所以：
    \[
    \int_{0}^{4} \rho^3 \, dz = 4\rho^3
    \]

    再对 \(\rho\) 积分：
    \[
    \int_{0}^{2\cos\theta} 4\rho^3\,d\rho = 4\left[ \frac{\rho^4}{4} \right]_{0}^{2\cos\theta} = (2\cos\theta)^4 = 16\cos^4\theta
    \]

    最后对 \(\theta\) 积分：
    \[
    \int_{-\frac{\pi}{2}}^{\frac{\pi}{2}} 16\cos^4\theta\,d\theta
    = 32 \int_{0}^{\frac{\pi}{2}} \cos^4\theta\,d\theta
    = 32 \cdot \frac{3}{4}\cdot\frac{1}{2}\cdot\frac{\pi}{2} = 6\pi
    \]

    所以 \(\iiint_{\Omega} (x^2+y^2)\,dV = 6\pi\)。
\end{myexample}

\begin{myexample}
    计算三重积分 \(\iiint_{\Omega} z\sqrt{x^2+y^2}\,dV\)，其中 \(\Omega\) 是由圆锥面 \(z=\sqrt{x^2+y^2}\) 与平面 \(z=1\) 围成的区域。

    这个区域是一个圆锥体，顶点在原点，底面是 \(z=1\) 平面上的圆 \(x^2+y^2\le 1\)。令 \(x=\rho\cos\theta\)，\(y=\rho\sin\theta\)，\(z=z\)，则被积函数化为 \(z\cdot \rho\)，体积微元是 \(\rho\,d\rho\,d\theta\,dz\)。

    \begin{center}
    \begin{tikzpicture}[
        scale=2.5,
        line join=round,
        % 视角调整：使顶点在下，底面在上
        x={(-0.866cm,-0.5cm)},
        y={(0.866cm,-0.5cm)},
        z={(0cm,1cm)},
    ]
    
    % 圆锥参数
    \def\R{1}    % 底面半径
    \def\H{1}    % 高度
    
    % 顶点与底面圆心
    \coordinate (V) at (0,0,0);
    \coordinate (Otop) at (0,0,\H);
    
    
    \fill[cyan!40, opacity=0.65, even odd rule]
        (V) -- ({\R*cos(0)}, {\R*sin(0)}, \H) -- ({\R*cos(90)}, {\R*sin(90)}, \H) -- cycle
        % 弓形：从 A 沿弧到 B，再 -- cycle 闭合（弦 BA）
        plot[domain=0:90, samples=120, variable=\t, smooth]
            ({\R*cos(\t)}, {\R*sin(\t)}, \H)
        -- cycle;
    
    % ---- 顶面（半透明填充） ----
    \fill[blue!20, opacity=0.65]
        plot[domain=0:360, samples=120, smooth]
        ({\R*cos(\x)}, {\R*sin(\x)}, \H) -- cycle;
    
    % ---- 底面圆的轮廓线 ----
    \draw[thick]
        plot[domain=0:360, samples=120, smooth]
        ({\R*cos(\x)}, {\R*sin(\x)}, \H);
    
    
    \draw[very thick, cyan!70!black] (V) -- ({\R*cos(0)}, {\R*sin(0)}, \H);
    
    \draw[very thick, cyan!70!black] (V) -- ({\R*cos(90)}, {\R*sin(90)}, \H);
    
    
    % ---- 坐标轴（从顶点向外延伸） ----
    \draw[-stealth, thick] (V) -- (1.3,0,0) node[below] {$x$};
    \draw[-stealth, thick] (V) -- (0,1.3,0) node[below left] {$y$};
    \draw[-stealth, thick] (0,0,1) -- (0,0,1.3) node[above] {$z$};
    \draw[thick, dashed] (V) -- (0,0,1);
    
    % 标签
    \node[below] at (V) {$O$};
    \node[above right] at (0,0,\H) {$z=1$};
    
    \end{tikzpicture}
    \quad
    \begin{tikzpicture}[
    scale=2.5,
]
  % 填充圆盘并画出边界
  \filldraw[fill=blue!10, draw=blue!60, very thick] (0,0) circle (1);

  \draw[->] (-1.3,0) -- (1.3,0) node[right] {$x$};
  \draw[->] (0,-1.3) -- (0,1.3) node[above] {$y$}; 

  % 说明文字
  \node at (0.5,0.6) {$D_{\rho\theta}$};

  \draw[thick] (0,0) -- (30:{1.5*cos(30)});

  \draw[->] (0.35,0) arc (0:30:0.35) node[midway, right] {$\theta$};

  \node at (-0.5,0.6) {$\rho=1$};

\end{tikzpicture}
\end{center}

    接下来用“画图穿线法”确定积分限。先看 \(z\) 方向。直线从圆锥面 \(z=\rho\) 穿入，从顶面 \(z=1\) 穿出。所以 \(z\) 的积分下限是 \(\rho\)，上限是 \(1\)。
    
    再把\(\Omega\)投影到\(xOy\) 平面，得到 \(D_{\rho\theta}\)。投影区域是圆 \(\rho\le 1\)，所以 \(\theta\) 从 \(0\) 扫到 \(2\pi\)，对每一个 \(\theta\)，射线从 \(\rho=0\) 穿入，从 \(\rho=1\) 穿出。

    于是：
    \[
    \iiint_{\Omega} z\sqrt{x^2+y^2}\,dV
    =\iint_{D_{\rho\theta}} \left(\int_{\rho}^{1} (z\cdot \rho) \, dz\right) \, \rho d\rho d\theta
    = \int_{0}^{2\pi} d\theta \int_{0}^{1} \rho^2\,d\rho \int_{\rho}^{1} z \, dz
    \]

    先对 \(z\) 积分：
    \[
    \int_{\rho}^{1} z\,dz = \left[ \frac{z^2}{2} \right]_{\rho}^{1} = \frac{1}{2} - \frac{\rho^2}{2} = \frac{1-\rho^2}{2}
    \]

    再对 \(\rho\) 积分：
    \[
    \int_{0}^{1} \rho^2 \cdot \frac{1-\rho^2}{2} \, d\rho
    = \frac{1}{2} \int_{0}^{1} (\rho^2 - \rho^4) \, d\rho
    = \frac{1}{2} \left[ \frac{\rho^3}{3} - \frac{\rho^5}{5} \right]_{0}^{1}
    = \frac{1}{2} \left( \frac{1}{3} - \frac{1}{5} \right)
    = \frac{1}{2} \cdot \frac{2}{15} = \frac{1}{15}
    \]

    最后对 \(\theta\) 积分，被积函数已经是常数：
    \[
    \int_{0}^{2\pi} \frac{1}{15}\,d\theta = \frac{1}{15} \cdot 2\pi = \frac{2\pi}{15}
    \]

    所以 \(\iiint_{\Omega} z\sqrt{x^2+y^2}\,dV = \frac{2\pi}{15}\)。
\end{myexample}

总结一下，用柱面坐标计算三重积分的思路和直角坐标下的投影法一脉相承。当区域在 \(xOy\) 平面上的投影是圆、扇形或圆环，且被积函数出现 \(x^2+y^2\) 时，柱面坐标通常是最好的选择。要注意的就是体积微元的 \(\textcolor{red}{\rho}\) 绝对不能漏掉，还有“画图穿线法”判断上下限时要分别对 \(z\)、\(\theta\)、\(\rho\)依次操作，看清楚穿入和穿出的边界各是什么。

\subsection{三重积分的球面坐标形式}

当积分区域是球体、球壳或者球体的一部分，又或者被积函数里出现了 \(x^2+y^2+z^2\) 这类项，用柱面坐标有时还是会比较麻烦。这时就可以搬出球面坐标，它是专门对付球形区域的有力工具。

球面坐标用三个数来确定空间中一点的位置，分别是向径 \(r\)，方位角 \(\theta\)，以及极角 \(\varphi\)。在球面坐标里，通常取\(0\le r<+\infty\)，\(0\le|\varphi|\le\pi\)，\(0\le|\theta|\le 2\pi\)。


\begin{center}
\begin{tikzpicture}[
    scale=1.5,
    x={(1cm,0cm)},
    y={(-0.4cm,-0.35cm)},
    z={(0cm,1cm)},
    line join=round,
    line cap=round
]
    % 坐标轴
    \draw[->] (-0.2,0,0) -- (3.2,0,0) node[below] {$x$};
    \draw[->] (0,-0.2,0) -- (0,2.8,0) node[right] {$y$};
    \draw[->] (0,0,0) -- (0,0,2.8) node[left] {$z$};
    \node[below left] at (0,0,0) {$O$};

    % 定义点 P 的坐标
    \def\rr{3.2}
    \def\phii{53}
    \def\thh{45}

    % 计算点 P 的直角坐标
    \pgfmathsetmacro\xx{\rr*sin(\phii)*cos(\thh)}
    \pgfmathsetmacro\yy{\rr*sin(\phii)*sin(\thh)}
    \pgfmathsetmacro\zz{\rr*cos(\phii)}

    % 点 P
    \coordinate (P) at (\xx,\yy,\zz);
    \fill (P) circle (1.2pt) node[above right] {$M(r,\varphi,\theta)$};

    % 投影点 P' 在 xOy 平面上
    \coordinate (Pxy) at (\xx,\yy,0);
    \fill (Pxy) circle (1.2pt) node[below right] {$M'(r,\varphi,0)$};

    % 连接
    \draw[dashed] (P) -- (Pxy);
    \draw[dashed] (0,0,0) -- (Pxy);

    % 标出 r
    \draw[thick] (0,0,0) -- (P);
    \node[above] at (2*\xx/3,2*\yy/3,2*\zz/3) {$r$};

    % 标出 phi 角（在 z--P 平面内）
    \draw[->] (0,0,0.4) arc (90-\phii:90:-0.4);
    \node[right] at (0.1,0,0.6) {$\varphi$};

    % 标出 theta 角
    \draw[->] (0.6,0,0) arc (0:\thh:0.6);
    \node[right] at (0.6,0.4,0) {$\theta$};

\end{tikzpicture}
\end{center}

点 \(M(x,y,z)\) 的球面坐标与直角坐标之间的关系是\(x = r\sin\varphi\cos\theta\)，\(y = r\sin\varphi\sin\theta\)，\(z = r\cos\varphi\)。与柱面坐标的关系是\(r = \sqrt{\rho^2+z^2}\)，\(\varphi = \arctan(\frac{\rho}{z})\)，\(\theta = \theta\)。反过来，如果已知球面坐标，令\(\rho=r\sin\varphi\)，\(z=r\cos\varphi\)，\(\theta=\theta\)，就得到柱面坐标。

有了坐标变换，下一步就要弄清楚球面坐标下的体积微元是什么。和二重积分极坐标、三重积分柱面坐标的思路如出一辙，我们把 \(r,\varphi,\theta\) 各自取一个微小的增量 \(\Delta r,\Delta\varphi,\Delta\theta\)，这一小块区域近似是一个弯曲的“长方体”。在 \(r\) 方向上的边长是 \(\Delta r\)，在 \(\varphi\) 方向上的“弧长”是 \(r\,\Delta\varphi\)，在 \(\theta\) 方向上的“弧长”是 \(r\sin\varphi\,\Delta\theta\)。当这些增量都趋向于零时，也就是取微分，那么\(\Delta r\)就是\(dr\)，\(r\,\Delta\varphi\)就是\(r\,d\varphi\)，\(r\sin\varphi\,\Delta\theta\)就是\(r\sin\varphi\,d\theta\)，这时候体积微元就是这三者的乘积，所以：
\[
dV = dr \cdot r\,d\varphi \cdot r\sin\varphi\,d\theta = r^2\sin\varphi\,dr\,d\varphi\,d\theta
\]

\begin{center}
\begin{tikzpicture}[
    scale=2,
    x={(1cm,0cm)},
    y={(-0.4cm,-0.35cm)},
    z={(0cm,1cm)},
    line join=round,
    line cap=round
]
    % 坐标轴
    \draw[->] (-0.2,0,0) -- (2.4,0,0) node[below] {$x$};
    \draw[->] (0,-0.2,0) -- (0,1.8,0) node[right] {$y$};
    \draw[->] (0,0,0) -- (0,0,2.4) node[left] {$z$};
    \node[above left] at (0,0,0) {$O$};

    % 定义点 P 的坐标
    \def\rr{2.8}
    \def\phii{45}
    \def\thh{30}
    \def\drr{0.8}
    \def\dthh{15}
    \def\dphii{10}

    % 计算点 P 的直角坐标
    \pgfmathsetmacro\xx{\rr*sin(\phii)*cos(\thh)}
    \pgfmathsetmacro\yy{\rr*sin(\phii)*sin(\thh)}
    \pgfmathsetmacro\zz{\rr*cos(\phii)}
    \pgfmathsetmacro\RR{\rr+\drr}

    % 点
    \coordinate (A) at (\xx,\yy,\zz);
    \coordinate (B) at ({\rr*sin(\phii+\dphii)*cos(\thh)},{\rr*sin(\phii+\dphii)*sin(\thh)},{\rr*cos(\phii+\dphii)});
    \coordinate (C) at ({\rr*sin(\phii)*cos(\thh+\dthh)},{\rr*sin(\phii)*sin(\thh+\dthh)},{\rr*cos(\phii)});
    \coordinate (D) at ({\rr*sin(\phii+\dphii)*cos(\thh+\dthh)},{\rr*sin(\phii+\dphii)*sin(\thh+\dthh)},{\rr*cos(\phii+\dphii)});
    \coordinate (A1) at ({\RR*sin(\phii)*cos(\thh)},{\RR*sin(\phii)*sin(\thh)},{\RR*cos(\phii)});
    \coordinate (B1) at ({\RR*sin(\phii+\dphii)*cos(\thh)},{\RR*sin(\phii+\dphii)*sin(\thh)},{\RR*cos(\phii+\dphii)});
    \coordinate (C1) at ({\RR*sin(\phii)*cos(\thh+\dthh)},{\RR*sin(\phii)*sin(\thh+\dthh)},{\RR*cos(\phii)});
    \coordinate (D1) at ({\RR*sin(\phii+\dphii)*cos(\thh+\dthh)},{\RR*sin(\phii+\dphii)*sin(\thh+\dthh)},{\RR*cos(\phii+\dphii)});

    \fill[gray!50,opacity=0.5]
  (C) -- (C1)
  plot[domain=\thh+\dthh:\thh, samples=40]
    ({\RR*sin(\phii)*cos(\x)}, {\RR*sin(\phii)*sin(\x)}, {\RR*cos(\phii)})
  -- plot[domain=\phii:\phii+\dphii, samples=40]
    ({\RR*sin(\x)*cos(\thh)}, {\RR*sin(\x)*sin(\thh)}, {\RR*cos(\x)})
  -- (B)
  -- plot[domain=\thh:\thh+\dthh, samples=40]
    ({\rr*sin(\phii+\dphii)*cos(\x)}, {\rr*sin(\phii+\dphii)*sin(\x)}, {\rr*cos(\phii+\dphii)})
  -- plot[domain=\phii+\dphii:\phii, samples=40]
    ({\rr*sin(\x)*cos(\thh+\dthh)}, {\rr*sin(\x)*sin(\thh+\dthh)}, {\rr*cos(\x)})
  -- cycle;


    % 投影点 P' 在 xOy 平面上
    \coordinate (Axy) at (\xx,\yy,0);
    \coordinate (D1xy) at ({\RR*sin(\phii+\dphii)*cos(\thh+\dthh)},{\RR*sin(\phii+\dphii)*sin(\thh+\dthh)},0);

    % 连接
    \draw[dashed] (A) -- (Axy);
    \draw[dashed] (0,0,0) -- (Axy);

    \draw[dashed] (D1) -- (D1xy);
    \draw[dashed] (0,0,0) -- (D1xy);

    % 标出 r
    \draw[dashed] (0,0,0) -- (A);
    \draw[dashed] (0,0,0) -- (B);
    \draw[thick,dashed] (A) -- (A1);
    \draw[thick] (B) -- (B1);
    \draw[thick] (C) -- (C1);
    \draw[thick] (D) -- (D1);

  \draw[thick,dashed] plot[domain=\phii:\phii+\dphii, samples=80]
  ({\rr*sin(\x)*cos(\thh)}, {\rr*sin(\x)*sin(\thh)}, {\rr*cos(\x)});

  \draw[thick,dashed] plot[domain=\thh:\thh+\dthh, samples=80]
  ({\rr*sin(\phii)*cos(\x)}, {\rr*sin(\phii)*sin(\x)}, {\rr*cos(\phii)});

  \draw[thick] plot[domain=\phii:\phii+\dphii, samples=80]
  ({\rr*sin(\x)*cos(\thh+\dthh)}, {\rr*sin(\x)*sin(\thh+\dthh)}, {\rr*cos(\x)});

\draw[thick] plot[domain=\thh:\thh+\dthh, samples=80]
  ({\rr*sin(\phii+\dphii)*cos(\x)}, {\rr*sin(\phii+\dphii)*sin(\x)}, {\rr*cos(\phii+\dphii)});

  % ---------- 外层弧线 ----------
% 经线弧：A1 -> B1 (等 θ = θ₀)
\draw[thick] plot[domain=\phii:\phii+\dphii, samples=80]
  ({\RR*sin(\x)*cos(\thh)}, {\RR*sin(\x)*sin(\thh)}, {\RR*cos(\x)});

% 纬线弧：A1 -> C1 (等 φ = φ₀)
\draw[thick] plot[domain=\thh:\thh+\dthh, samples=80]
  ({\RR*sin(\phii)*cos(\x)}, {\RR*sin(\phii)*sin(\x)}, {\RR*cos(\phii)});

% 经线弧：C1 -> D1 (等 θ = θ₀ + Δθ)
\draw[thick] plot[domain=\phii:\phii+\dphii, samples=80]
  ({\RR*sin(\x)*cos(\thh+\dthh)}, {\RR*sin(\x)*sin(\thh+\dthh)}, {\RR*cos(\x)});

% 纬线弧：B1 -> D1 (等 φ = φ₀ + Δφ)
\draw[thick] plot[domain=\thh:\thh+\dthh, samples=80]
  ({\RR*sin(\phii+\dphii)*cos(\x)}, {\RR*sin(\phii+\dphii)*sin(\x)}, {\RR*cos(\phii+\dphii)});

  \draw[->] ({0.8*cos(\thh)},{0.8*sin(\thh)},0) arc(\thh:(\thh+\dthh):0.8) node[midway,below right] {$d\theta$};

  \draw[->] 
  plot[domain=\phii:\phii+\dphii, samples=80]
    ({0.8*sin(\x)*cos(\thh)}, 
     {0.8*sin(\x)*sin(\thh)}, 
     {0.8*cos(\x)})
  node[below right] {$d\varphi$};

\pgfmathsetmacro{\thetaMid}{\thh + 0.5*\dthh}
\pgfmathsetmacro{\xMid}{\rr*sin(\phii)*cos(\thetaMid)}
\pgfmathsetmacro{\yMid}{\rr*sin(\phii)*sin(\thetaMid)}
\pgfmathsetmacro{\zMid}{\rr*cos(\phii)}

% 画引线（从弧中点向外偏移 0.6cm 的向量，可按视角调整）
\draw (\xMid, \yMid, \zMid)
  -- ++(-0.5, -1.5, 0) % ← 这里控制引线方向和长度
  node[above] {$r\sin\varphi\,d\theta$};

\pgfmathsetmacro{\phiMid}{\phii + 0.5*\dphii}
\pgfmathsetmacro{\xMidCD}{\rr*sin(\phiMid)*cos(\thh+\dthh)}
\pgfmathsetmacro{\yMidCD}{\rr*sin(\phiMid)*sin(\thh+\dthh)}
\pgfmathsetmacro{\zMidCD}{\rr*cos(\phiMid)}

\draw (\xMidCD, \yMidCD, \zMidCD)
  -- ++(-0.8, -0.8, 0)               % 偏移量可调
  node[above] {$r\,d\varphi$};


\draw ($(D)!0.5!(D1)$)
  -- ++(0.4, 1.2, 0)
  node[below] {$dr$};

\end{tikzpicture}
\end{center}


由此便得到了三重积分在球面坐标下的变换公式：

\begin{theorem}{三重积分球面坐标变换公式}{}
设函数 \(f(x,y,z)\) 在空间区域 \(\Omega\) 上可积，作球面坐标变换 \(x=r\sin\varphi\cos\theta,\;y=r\sin\varphi\sin\theta,\;z=r\cos\varphi\)，则
\[
\iiint_{\Omega} f(x,y,z)\,dxdydz
= \iiint_{\Omega} f(r\sin\varphi\cos\theta,\; r\sin\varphi\sin\theta,\; r\cos\varphi)\;
r^2\sin\varphi\,dr\,d\varphi\,d\theta
\]
\end{theorem}


要注意的是，球面坐标下的体积微元是 \(\textcolor{red}{r^2\sin\varphi}\,dr\,d\varphi\,d\theta\)，既不是 \(dr\,d\varphi\,d\theta\)，也不是 \(r\,dr\,d\varphi\,d\theta\)，而是\(\textcolor{red}{r^2\sin\varphi}\,dr\,d\varphi\,d\theta\)。

在积分次序上，球面坐标和柱面坐标一样，通常采用“先一后二”的思路。对大多数区域，我们是先固定 \(\varphi\) 和 \(\theta\)，从原点引出一条射线穿过区域，确定 \(r\) 的上下限，然后再对 \(\varphi,\theta\) 在某个平面区域上做二重积分。

用“画图穿线法”操作就是固定住 \(\varphi\) 和 \(\theta\)，然后从原点出发引一条射线，这条射线穿进区域的那一点就是积分下限 \(r=r_1(\varphi,\theta)\)，穿出区域的那一点就是积分上限 \(r=r_2(\varphi,\theta)\)，再让 \((\varphi,\theta)\) 在某个参数区域 \(\Delta\) 上变化，把所有这些射线上的累积量加起来，结果就是三重积分的值。所以我们有：
\begin{theorem}{球面坐标下的三重积分}{}
若空间闭区域 \(\Omega\) 可表示为
\[
\Omega = \{(r\sin\varphi\cos\theta,\; r\sin\varphi\sin\theta,\; r\cos\varphi) \mid
(\varphi,\theta)\in\Delta,\; r_1(\varphi,\theta)\le r\le r_2(\varphi,\theta)\}
\]
则
\[
\iiint_{\Omega} f(x,y,z)\,dxdydz
= \iiint_{\Omega} F(r,\varphi,\theta)\,r^2\sin\varphi\,dr\,d\varphi\,d\theta
\]
其中\(F(r,\varphi,\theta)=f(r\sin\varphi\cos\theta,\; r\sin\varphi\sin\theta,\; r\cos\varphi)\)。
\end{theorem}

至于 \(\varphi\) 和 \(\theta\) 的范围，需要根据积分区域的具体形状来确定。

\begin{myexample}
计算三重积分 \(\iiint_{\Omega} (x^2+y^2+z^2)\,dV\)，其中 \(\Omega\) 是由球面 \(x^2+y^2+z^2=R^2\) 围成的球体。

被积函数里出现了 \(x^2+y^2+z^2\)，整个区域又是一个标准的球体，这正是球面坐标大显身手的时候。作球面坐标变换，令 \(x=r\sin\varphi\cos\theta,\; y=r\sin\varphi\sin\theta,\; z=r\cos\varphi\)，则被积函数化为 \(r^2\)，体积微元是 \(r^2\sin\varphi\,dr\,d\varphi\,d\theta\)，区域 \(\Omega\) 就是 \(0\le r\le R\)，\(0\le\varphi\le\pi\)，\(0\le\theta\le 2\pi\)。

\begin{center}
    \begin{tikzpicture}
      \tdplotsetmaincoords{56}{135}
      
      \begin{scope}[tdplot_main_coords, scale=1.5]
        \draw[->, thick] (0,0,0) -- (2.5,0,0) node[anchor=north east] {$x$};
        \draw[->, thick] (0,0,0) -- (0,2.5,0) node[anchor=north west] {$y$};
        \draw[->, thick] (0,0,0) -- (0,0,2.5) node[anchor=south]      {$z$};
        \shade[ball color=blue!25] (0,0,0) circle (1.5cm);
      \end{scope}
    \end{tikzpicture}
\end{center}

所以：
\[
\iiint_{\Omega} (x^2+y^2+z^2)\,dV
= \iiint_{\Omega} r^2\cdot r^2\sin\varphi\,dr\,d\varphi\,d\theta
= \int_{0}^{2\pi} d\theta \int_{0}^{\pi}\sin\varphi \,d\varphi \int_{0}^{R} r^4\,dr
\]

先对 \(r\) 积分：
\[
\int_{0}^{R} r^4\,dr = \left[ \frac{r^5}{5} \right]_{0}^{R} = \frac{R^5}{5}
\]

再对 \(\varphi\) 积分：
\[
\int_{0}^{\pi} \frac{R^5}{5}\sin\varphi\,d\varphi = \frac{R^5}{5}\bigl[-\cos\varphi\bigr]_{0}^{\pi} = \frac{2R^5}{5}
\]

最后对 \(\theta\) 积分：
\[
\int_{0}^{2\pi}\frac{2R^5}{5}\,d\theta =\frac{2R^5}{5}\cdot 2\pi=\frac{4\pi R^5}{5}
\]

所以 \(\iiint_{\Omega} (x^2+y^2+z^2)\,dV = \dfrac{4\pi R^5}{5}\)。
\end{myexample}

\begin{myexample}
计算三重积分 \(\iiint_{\Omega} z\,dV\)，其中 \(\Omega\) 是由上半球面 \(x^2+y^2+z^2=1,\; z\ge 0\) 与圆锥面 \(z=\sqrt{x^2+y^2}\) 所围成的位于锥面下方的区域。

这个区域看起来有点复杂，但用球面坐标描述就非常自然。令 \(x=r\sin\varphi\cos\theta,\; y=r\sin\varphi\sin\theta,\; z=r\cos\varphi\)，则被积函数化为 \(r\cos\varphi\)，体积微元是 \(r^2\sin\varphi\,dr\,d\varphi\,d\theta\)。

先看看区域的边界。球面是 \(r=1\)。圆锥面 \(z=\sqrt{x^2+y^2}\) 化为球面坐标就是 \(r\cos\varphi = r\sin\varphi\)，即 \(\tan\varphi=1\)，也就是 \(\varphi=\frac{\pi}{4}\)。区域在锥面下方，即 \(z\le\sqrt{x^2+y^2}\)，对应 \(\cos\varphi\le\sin\varphi\)，所以 \(\varphi\) 从 \(\frac{\pi}{4}\) 到 \(\frac{\pi}{2}\)，也就是\(\varphi\in [\frac{\pi}{4},\frac{\pi}{2}]\)。\(\theta\) 可以绕区域一圈，所以 \(\theta\in [0,2\pi]\)。对于每一个固定的 \((\varphi,\theta)\)，射线从原点出发，先进入区域，穿进区域时是原点 \(r=0\)，穿出来碰到球面就是 \(r=1\)。所以 \(r\) 的范围是 \(0\le r\le 1\)。

\begin{center}
\begin{tikzpicture}[scale=2.5]
  \tdplotsetmaincoords{70}{120}
\begin{scope}[tdplot_main_coords]
  
  % ========== 基本常量 ==========
  \def\R{1}                     % 球半径
  \pgfmathsetmacro{\h}{1/sqrt(2)} % 交线高度
  \pgfmathsetmacro{\r}{\h}       % 交线半径

  
  
  % ========== 根据视角定义可见角度范围 ==========
  % 可见弧从 phi_start 到 phi_end（度）
  \def\phistart{-60}
  \def\phiend{120}
  % 不可见弧的结束角（360 + phistart 的等价正角度）
  \def\phihideend{300}  % 即 -40° + 360°
  
  % ========== 1. 不可见部分（虚线） ==========
  % 赤道圆（z=0）的不可见弧
  \draw[cyan, densely dashed, thick]
    ({\R*cos(\phiend)}, {\R*sin(\phiend)}, 0)
    arc[start angle=\phiend, end angle=\phihideend, radius=\R];
  
  \fill[red!80, opacity=0.72]
  plot[domain=0:360, smooth, variable=\t]
    ({\r*cos(\t)}, {\r*sin(\t)}, \h) -- cycle;

  
  % ========== 2. Ω 区域的可见表面（透明填充） ==========
  % 锥面可见部分（原点 → 左母线 → 交线可见弧 → 右母线 → 原点）
  \def\phicontourA{98.66}
  \def\phicontourB{-38.66}


  \fill[red!80, opacity=0.72]
    (0,0,0) --
    ({\r*cos(\phicontourA)}, {\r*sin(\phicontourA)}, \h) --
    plot[domain=\phicontourA:\phicontourB, smooth, variable=\phi]
      ({\r*cos(\phi)}, {\r*sin(\phi)}, \h) --
    cycle;

    
% 画两条轮廓母线（实线）
\draw[red!80, thick, dashed]
  ({\r*cos(\phicontourA)}, {\r*sin(\phicontourA)}, \h) -- (0,0,0)
  ({\r*cos(\phicontourB)}, {\r*sin(\phicontourB)}, \h) -- (0,0,0);
  
  % 球面可见部分（交线可见弧 + 赤道可见弧围成的区域）

  \pgfmathsetmacro{\latTop}{asin(\h)}   % 顶部纬度 (弧度)
  \pgfmathsetmacro{\latBottom}{0}       % 底部纬度

% 侧面填充
\fill[cyan, opacity=0.2]
    % 1. 底部纬线弧 (z=0, 半径 R=1)
    ({\R*cos(\phistart)}, {\R*sin(\phistart)}, 0)
    arc[start angle=\phistart, end angle=\phiend, radius=\R]
    
    % 2. 右侧经线弧：经度固定为 \phiend，纬度从 \latBottom 到 \latTop
    %    参数方程: (cosγ·cosφ, cosγ·sinφ, sinγ)
    -- plot[domain=\latBottom:\latTop, variable=\gamma]
        ({cos(\gamma)*cos(\phiend)}, {cos(\gamma)*sin(\phiend)}, {sin(\gamma)})
    
    % 3. 顶部纬线弧 (z=h, 半径 r, 反向从 \phiend 到 \phistart)
    arc[start angle=\phiend, end angle=\phistart, radius=\r]
    
    % 4. 左侧经线弧：经度 \phistart，纬度从 \latTop 回到 \latBottom
    -- plot[domain=\latTop:\latBottom, variable=\gamma]
        ({cos(\gamma)*cos(\phistart)}, {cos(\gamma)*sin(\phistart)}, {sin(\gamma)})
    -- cycle;
  
    % 右侧经线轮廓 (φ = φend)
\draw[cyan, thick]
  plot[domain=0:\latTop, variable=\gamma]
    ({cos(\gamma)*cos(\phiend)}, {cos(\gamma)*sin(\phiend)}, {sin(\gamma)});
% 左侧经线轮廓 (φ = φstart)
\draw[cyan, thick]
  plot[domain=0:\latTop, variable=\gamma]
    ({cos(\gamma)*cos(\phistart)}, {cos(\gamma)*sin(\phistart)}, {sin(\gamma)});
  % ========== 3. 可见轮廓与结构线（实线） ==========
  % 赤道圆可见弧
  \draw[cyan, thick]
    ({\R*cos(\phistart)}, {\R*sin(\phistart)}, 0)
    arc[start angle=\phistart, end angle=\phiend, radius=\R];
  
  % 交线圆可见弧
  \draw[red!80, thick] 
    plot[domain=0:360, smooth, samples=120, variable=\t] 
    ({\r*cos(\t)}, {\r*sin(\t)}, \h) -- cycle;
  
    \draw[dashed, thick] (0,0,0) -- (1,0,0);
  \draw[-stealth, thick] (1,0,0) -- (1.3,0,0) node[anchor=north east]{$x$};
  \draw[dashed, thick] (0,0,0) -- (0,1,0);
  \draw[-stealth, thick] (0,1,0) -- (0,1.3,0) node[anchor=north west]{$y$};
  \draw[dashed, thick] (0,0,0) -- (0,0,{1/sqrt(2)});
  \draw[-stealth, thick] (0,0,{1/sqrt(2)}) -- (0,0,1.3) node[anchor=south]{$z$};



  % ========== 4. 标签 ==========
  \node at (0.9, 0.15, 0.5) {\(\Omega\)};
  
\end{scope}
\end{tikzpicture}
\end{center}

于是：
\[
\iiint_{\Omega} z\,dV
= \iiint_{\Omega} r\cos\varphi\cdot r^2\sin\varphi\,dr\,d\varphi\,d\theta
= \int_{0}^{2\pi} d\theta \int_{\frac{\pi}{4}}^{\frac{\pi}{2}} \sin\varphi\,d\varphi \int_{0}^{1} r^3\cos\varphi\,dr
\]

先对 \(r\) 积分，\(\cos\varphi\) 与 \(r\) 无关：
\[
\int_{0}^{1} r^3\cos\varphi\,dr = \cos\varphi\left[ \frac{r^4}{4} \right]_{0}^{1} = \frac{\cos\varphi}{4}
\]

再对 \(\varphi\) 积分：
\[
\int_{\frac{\pi}{4}}^{\frac{\pi}{2}} \frac{\cos\varphi}{4} \sin\varphi\,d\varphi
= \frac{1}{4} \int_{\frac{\pi}{4}}^{\frac{\pi}{2}} \sin\varphi\cos\varphi\,d\varphi
= \frac{1}{8} \int_{\frac{\pi}{4}}^{\frac{\pi}{2}} \sin 2\varphi\,d\varphi
= \frac{1}{8} \left[ -\frac{\cos 2\varphi}{2} \right]_{\frac{\pi}{4}}^{\frac{\pi}{2}}
= \frac{1}{16}
\]

最后对 \(\theta\) 积分：
\[
\int_{0}^{2\pi} \frac{1}{16}\,d\theta = \frac{1}{16} \cdot 2\pi = \frac{\pi}{8}
\]

所以 \(\iiint_{\Omega} z\,dV = \dfrac{\pi}{8}\)。
\end{myexample}

总结一下，球面坐标主要是为球体和球壳一类问题准备的。当被积函数含有 \(x^2+y^2+z^2\)，或者区域边界是球面、锥面时，球面坐标往往比柱面坐标还要方便。使用时要牢记体积微元是 \(\textcolor{red}{r^2\sin\varphi}\,dr\,d\varphi\,d\theta\)。确定积分限时，用“画图穿线法”从原点引射线，穿入和穿出的位置分别是 \(r\) 的上下限，再根据区域在 \((\varphi,\theta)\) 上的投影确定角度范围，最后按 \(r\)、\(\varphi\)、\(\theta\) 的顺序化为累次积分即可。

至此，从二重积分到三重积分的几种常用坐标计算法就都介绍完了。不管面对什么形式的积分，只要抓住“分割、近似、求和、取极限”的核心思想，再结合区域的特点选择合适的坐标系，二重积分和三重积分就都不是难事。
% !TEX root = ../main.tex
\chapter{曲线积分与曲面积分}
到目前为止，我们都是在一片面积或一块体积上做积分，被积函数是密度函数，积分结果就是这片区域上的总质量。如果现在有一条弯曲的线段，已知它的线密度，怎么求总质量？你可能会觉得面积用二重积分，体积用三重积分，那么线段就该用定积分，这说不上错，但是如果是三维空间里的一条曲线呢？这时候要怎么使用定积分？更进一步的，如果是三维空间里的一块曲面呢？这时候又要怎么使用二重积分？曲线积分和曲面积分就是解决这些问题的。
\section{曲线积分}
曲线积分简单来说就是沿着一条曲线进行的积分，它把“求和”的操作转移到了一根弯曲的线上。它的核心思想和之前的积分一样，就是将曲线分割成无数个微小弧段，在每个弧段上取一个函数值乘以该弧段的某种微小度量，然后求和取极限。
\subsection{第一类曲线积分}
想象一根放在空间中的弯曲的细杆，如果已知线密度\(f(x,y,z)\)，那么要求整根细杆的总质量，就要对细杆做积分。怎么积分呢？仿照二重积分和三重积分的方式，我们就是对整个细杆区域进行积分。推广到一般情形，设空间内有一条光滑曲线弧段，已知函数\(f(x,y,z)\)，以该光滑曲线弧段为区域对积分函数\(f(x,y,z)\)积分，这种积分我们就称为\textbf{第一类曲线积分}。也叫\textbf{对弧长的曲线积分}。

具体的计算方法也是“分割、近似、求和、取极限”，如果\(f(x,y,z)\)是连续函数，我们就可以在光滑曲线弧段上取极小的一个弧段\(\Delta s\)，当这个小弧段趋于零时，我们可以把它看成是直线段，这就是弧微分\(ds\)。那么对弧微分\(ds\)积分，就得到第一类曲线积分的定义：
\begin{definition}{第一类曲线积分}{}
    设函数 \(f(x,y,z)\) 在空间光滑曲线弧段 \(C\) 上有界。将 \(C\) 任意分成 \(n\) 个小弧段
\[
\Delta s_1,\ \Delta s_2,\ \cdots,\ \Delta s_n
\]
每个小弧段既表示该弧段本身，也表示该弧段的弧长。在每个小弧段 \(\Delta s_i\) 上任取一点 \((\xi_i,\eta_i,\zeta_i)\)，作函数值 \(f(\xi_i,\eta_i,\zeta_i)\) 与小弧段弧长 \(\Delta s_i\) 的乘积 \(f(\xi_i,\eta_i,\zeta_i) \cdot \Delta s_i\)，并作出和
\[
S = \sum_{i=1}^{n} f(\xi_i,\eta_i,\zeta_i) \,\Delta s_i
\]

记 \(\lambda = \max\{\Delta s_1,\Delta s_2,\cdots,\Delta s_n\}\)。如果当 \(\lambda \to 0\) 时，\(S\) 的极限总存在，且与曲线 \(C\) 的分法及点 \((\xi_i,\eta_i,\zeta_i)\) 的取法无关，那么称 \(f(x,y,z)\) 在曲线弧段 \(C\) 上\textbf{可积}，并称这个极限 \(I\) 为函数 \(f(x,y,z)\) 在曲线弧段 \(C\) 上的\textbf{第一类曲线积分}（或\textbf{对弧长的曲线积分}），记作
\[
\int_{C} f(x,y,z)\,ds
\]
即：
\[
\int_{C} f(x,y,z)\,ds = \lim_{\lambda \to 0} \sum_{i=1}^{n} f(\xi_i,\eta_i,\zeta_i)\,\Delta s_i = I
\]

其中 \(f(x,y,z)\) 称为\textbf{被积函数}，\(f(x,y,z)\,ds\) 称为\textbf{被积表达式}，\(ds\) 称为\textbf{弧长微元}（或\textbf{弧微分}），\(x,y,z\) 称为\textbf{积分变量}，\(C\) 称为\textbf{积分曲线}（或\textbf{积分路径}），\(\int_C\) 称为\textbf{曲线积分号}。
\end{definition}
特别的，如果曲线弧段 \(C\) 是闭合曲线，那么函数在闭合曲线\(C\) 上对弧长的曲线积分记为\(\oint_C f(x,y,z)\,ds\)。

第一类曲线积分的性质和定积分是类似的，满足线性性质，可加性，保序性等性质，此处不再赘述。

这里我们主要讨论平面上的第一类曲线积分。由平面上的第一类曲线积分可以推广到空间的情形，具体过程我们省略不讲。

第一类曲线积分的物理意义我们已经说明了。更进一步的，如果被积函数换成电荷线密度，那么第一类曲线积分就表示沿着曲线的总电荷量。如果被积函数换成热流密度，那么第一类曲线积分就表示沿着曲线的总热流量。而在几何直观上，第一类曲线积分表示的是以\(C\)为准线，以\(f(x,y)\)为上边界的柱面面积。当被积函数为三元函数\(f(x,y,z)\)时，则在几何上无法直接表示为三维空间中的“面积”，可以理解为四维空间中“超柱面”的面积。

和多重积分一样，直接用定义计算第一类曲线积分并不现实，那么我们要如何计算第一类曲线积分呢？基本思路是通过曲线的参数方程，将弧长微元 \(ds\) 用参数的微分表示，从而把曲线积分转化为参数的定积分。

\begin{theorem}{第一类曲线积分的计算法}{}
设空间光滑曲线 \( C \) 的参数方程为  
\(
C(t) = \begin{cases}
x = x(t)\\
y = y(t)\\
z = z(t)
\end{cases}\)，\(t \in [\alpha, \beta]\)，被积函数为 \(f(x,y,z)\)，则弧长微元化为
\[
ds = |C\,'(t)|\,dt = \sqrt{[x'(t)]^2 + [y'(t)]^2 + [z'(t)]^2}\,dt
\]
积分化为
\[
\int_C f(x,y,z)\,ds = \int_{\alpha}^{\beta} f\big(x(t),y(t),z(t)\big) \cdot \sqrt{[x'(t)]^2 + [y'(t)]^2 + [z'(t)]^2}\,dt
\]

设平面光滑曲线 \( C \) 的参数方程为  
\(
C(t) = \begin{cases}
x = x(t)\\
y = y(t)
\end{cases}\)，\(t \in [\alpha, \beta]\)，被积函数为 \(f(x,y)\)，则弧长微元化为
\[
ds = |C\,'(t)|\,dt = \sqrt{[x'(t)]^2 + [y'(t)]^2}\,dt
\]
积分化为
\[
\int_C f(x,y)\,ds = \int_{\alpha}^{\beta} f\big(x(t),y(t)\big) \cdot \sqrt{[x'(t)]^2 + [y'(t)]^2}\,dt
\]
\end{theorem}

当\(C\)仅由\(y = y(x)\)或\(x = x(y)\)确定，则积分退化为：

\begin{theorem}{平面第一类曲线积分的计算法}{}
设平面光滑曲线 \( C:y=y(x) \)，\(x \in [a,b]\)，则弧长微元化为
\[
ds = \sqrt{1 + [y'(x)]^2}\,dx
\]
积分化为
\[
\int_C f(x,y)\,ds = \int_{a}^{b} f\big(x,y(x)\big) \cdot \sqrt{1 + [y'(x)]^2}\,dx
\]

设平面光滑曲线 \( C:x=x(y) \)，\(y \in [c,d]\)，则弧长微元化为
\[
ds = \sqrt{1 + [x'(y)]^2}\,dy
\]
积分化为
\[
\int_C f(x,y)\,ds = \int_{c}^{d} f\big(x(y),y\big) \cdot \sqrt{1 + [x'(y)]^2}\,dy
\]
\end{theorem}

具体的计算就是先求导，如果平面光滑曲线 \(C\) 由 \(y = y(x)\) 给出，那么就对 \(y(x)\) 求导，如果平面光滑曲线 \(C\) 由 \(x = x(y)\) 给出，那么就对 \(x(y)\) 求导，然后把弧微分\(ds\)换成关于\(dx\)或\(dy\)的表达式，最后代入公式转化为定积分计算。

我们看几个例子：

\begin{myexample}
计算\(\int_C \,ds\)，其中\(C\)是\(y=x\)由点\((0,0)\)到点\((1,1)\)的一段。

对\(y = x\)求导，得\(\frac{dy}{dx} = 1\)，所以弧长微元为：
\[
ds = \sqrt{1 + \left(\frac{dy}{dx}\right)^2}\,dx = \sqrt{1+1}\,dx = \sqrt{2}\,dx
\]

于是：
\[
\int_C ds = \int_0^1 \sqrt{2}\,dx = \sqrt{2}
\]

所以\(\int_C \,ds = \sqrt{2}\)。
\end{myexample}
这个例子说明，当被积函数为$1$时，对弧长的曲线积分就是曲线的长度。

\begin{myexample}
计算 \(\int_C x\,ds\)，其中 \(C\) 是抛物线 \(y = x^2\) 从点 \((0,0)\) 到点 \((1,1)\) 的一段。

对 \(y = x^2\) 求导得 \(y' = 2x\)，所以弧长微元为：
\[
ds = \sqrt{1 + (y')^2}\,dx = \sqrt{1 + 4x^2}\,dx
\]

于是：
\[
\int_C x\,ds = \int_0^1 x\sqrt{1+4x^2}\,dx
\]

令 \(u = 1+4x^2\)，则 \(du = 8x\,dx\)，即 \(x\,dx = \frac{du}{8}\)，当\(x=0\) 时 \(u=1\)，当\(x=1\) 时 \(u=5\)，代入得：
\[
\int_0^1 x\sqrt{1+4x^2}\,dx = \frac18 \int_1^5 \sqrt{u}\,du = \frac18 \cdot \frac23 u^{\frac{3}{2}}\Big|_1^5 = \frac{1}{12}(5\sqrt{5}-1)
\]

所以 \(\int_C x\,ds = \frac{5\sqrt{5}-1}{12}\)。
\end{myexample}

\begin{myexample}
    计算 \(\oint_C (x+y)\,ds\)，其中 \(C\) 是圆周 \(x^2+y^2=1\)。

设 \(x = \cos t\)，\(y = \sin t\)，\(t \in [0,2\pi]\)。求导得 \(x'(t) = -\sin t,\; y'(t) = \cos t\)，所以弧长微元为：
\[
ds = \sqrt{(-\sin t)^2 + (\cos t)^2}\,dt = dt
\]

于是：
\[
\oint_C (x+y)\,ds = \int_0^{2\pi} (\cos t + \sin t)\,dt
\]

计算得：
\[
 \int_0^{2\pi} (\cos t + \sin t)\,dt = \int_0^{2\pi} \cos t\,dt + \int_0^{2\pi} \sin t\,dt = [\sin t]_{0}^{2\pi} + [-\cos t]_{0}^{2\pi} = 0 + 0 = 0
\]

所以 \(\oint_C (x+y)\,ds = 0\)。

\end{myexample}

\begin{myexample}
    计算 \(\oint_C y\,ds\)，其中 \(C\) 是依次连接点 \((0,0), (2,0), (2,2), (0,2)\) 的正方形闭合边界。

将 \(C\) 分为四条边：
\[
C_1: (0,0)\to(2,0),\quad C_2: (2,0)\to(2,2),\quad C_3: (2,2)\to(0,2),\quad C_4: (0,2)\to(0,0)
\]

于是：
\[
\oint_C y\,ds = \int_{C_1} y\,ds + \int_{C_2} y\,ds + \int_{C_3} y\,ds + \int_{C_4} y\,ds
\]

计算\(\int_{C_1} y\,ds\)。\(C_1\)方程为 \(y=0\)，\(x\in[0,2]\)。求导得：\(\frac{dy}{dx}=0\)，所以：
\[
ds = \sqrt{1+\left(\frac{dy}{dx}\right)^2}\,dx = \sqrt{1+0}\,dx = dx
\]

积分：
\[
\int_{C_1} y\,ds = \int_{0}^{2} 0\,dx = 0
\]

计算\(\int_{C_2} y\,ds\)。\(C_2\)方程为\(x=2\)，\(y\in[0,2]\)。求导得：\(\frac{dx}{dy}=0\)，所以：
\[
ds = \sqrt{1+\left(\frac{dx}{dy}\right)^2}\,dy = \sqrt{1+0}\,dy = dy
\]

积分：
\[
\int_{C_2} y\,ds = \int_{0}^{2} y\,dy = \left[\frac{y^2}{2}\right]_{0}^{2} = 2
\]

计算\(\int_{C_3} y\,ds\)。\(C_3\)方程为\(y=2\)，\(x\in[0,2]\)。求导得：\(\frac{dy}{dx}=0\)。所以：
\[
ds = \sqrt{1+0}\,dx = dx
\]

积分：
\[
\int_{C_3} y\,ds = \int_{0}^{2} 2\,dx = 2\cdot 2 = 4
\]

计算\(\int_{C_4} y\,ds\)。\(C_4\)方程为\(x=0\)，\(y\in[0,2]\)。求导得：\(\frac{dx}{dy}=0\)。所以：
\[
ds = \sqrt{1+0}\,dy = dy
\]

积分：
\[
\int_{C_4} y\,ds = \int_{0}^{2} y\,dy = \left[\frac{y^2}{2}\right]_{0}^{2} = 2
\]

相加得：
\[
\oint_C y\,ds = 0 + 2 + 4 + 2 = 8
\]

所以 \(\oint_C y\,ds = 8\)。
\end{myexample}

\subsection{第二类曲线积分}
我们知道，第一类曲线积分的被积函数的函数值永远是数。而第二类曲线积分的被积函数有点不同，它的函数值是向量。所以在介绍第二类曲线积分之前，我们先来了解一下向量函数的概念：
\begin{definition}{向量函数}{}
    设有两个非空集合 \(D\) 与 \(Y\)，其中 \(Y\) 中的元素均为向量。若对于 \(D\) 中的每一个元素 \(x\)，按照某种确定的对应法则 \(\vec{f}\)，都有 \(Y\) 中唯一确定的向量 \(\vec{y}\) 与之对应，则称 \(\vec{f}\) 是定义在 \(D\) 上的一个\textbf{向量函数}，记作
\[
\vec{f}: D \to Y, \quad \vec{y} = \vec{f}(x), \quad x \in D
\]
其中 \(D\) 称为向量函数的\textbf{定义域}，全体函数值构成的集合 \(\{\vec{f}(x) \mid x \in D\}\) 称为向量函数的\textbf{值域}。
\end{definition}
我们只讨论三维空间里的一元向量函数。按照函数的定义，三维空间里的一元向量函数就意味着每一个自变量确定一个三维空间的向量。

简单来说，其实一元向量函数就是一个自变量，定义了\(n\)个分量函数，\(n\)是所在空间的维度，最后加上一个向量箭头。举个简单的例子，设一个三维向量函数 \(\vec{F}(t) = (t, t^2, t^3)\)，那么当 \(t=1\) 时，\(\vec{F}(1) = (1, 1, 1)\)，当 \(t=2\) 时，\(\vec{F}(2) = (2, 4, 8)\)。如果是二维向量函数 \(\vec{F}(t) = (t, t^2)\)，那么也是类似的，当 \(t=1\) 时，\(\vec{F}(1) = (1, 1)\)，诸如此类。

如果我们把之前一直使用的函数称为标量函数，那么向量函数就是标量函数的推广。它和标量函数的区别就在于，当自变量取不同的值的时候，标量函数的函数值是随着自变量变化而变化的数，而向量函数的函数值是随着自变量变化而变化的向量。

向量函数的性质和我们之前使用的函数类似，包括运算法则，极限，连续性和可导性等等，简述如下：

\begin{definition}{向量函数的极限}{}
设向量函数 \(\vec{f}(x)\) 在点 \(x_0\) 的某个去心邻域内有定义，\(\vec{L}\) 是一个常向量。若对于任意给定的 \(\epsilon > 0\)，总存在 \(\delta > 0\)，使得当  
\[
0 < |x - x_0| < \delta
\]  
时，一定有  
\[
|\vec{f}(x) - \vec{L}| < \epsilon
\]  
则称当 \(x\) 趋近于 \(x_0\) 时，向量函数 \(\vec{f}(x)\) 的\textbf{极限}为 \(\vec{L}\)，记作
\[
\lim_{x \to x_0} \vec{f}(x) = \vec{L}
\]
\end{definition}

\begin{definition}{向量函数的连续性}{}
设向量函数 \(\vec{f}(x)\) 在点 \(x_0\) 的某个邻域内有定义。若对任意 \(\epsilon > 0\)，存在 \(\delta > 0\)，使得当  
\[
|x - x_0| < \delta
\]  
时，有  
\[
|\vec{f}(x) - \vec{f}(x_0)| < \epsilon
\]  
则称向量函数 \(\vec{f}(x)\) 在点 \(x_0\) 处\textbf{连续}。
\end{definition}

\begin{definition}{向量函数的导数}{}
设向量函数 \(\vec{r}(t)\) 在点 \(t_0\) 的某个邻域内有定义。若极限  
\[
\lim_{\Delta t \to 0} \frac{\vec{r}(t_0 + \Delta t) - \vec{r}(t_0)}{\Delta t}
\]
存在，则称此极限值为向量函数 \(\vec{r}(t)\) 在点 \(t_0\) 处的\textbf{导数}（或\textbf{导向量}），记作  
\[
\vec{r}\,'(t_0), \quad \frac{d\vec{r}}{dt}\Big|_{t=t_0}, \quad \left.\frac{d\vec{r}}{dt}\right|_{t=t_0}
\]  
此时称 \(\vec{r}(t)\) 在点 \(t_0\) 处\textbf{可导}。
\end{definition}

\begin{theorem}{分量求导公式}{}
    若 \(\vec{r}(t) = (x(t), y(t), z(t))\)，则\(\vec{r}\,'(t) = ( x'(t),\, y'(t),\, z'(t) )\)。
\end{theorem}

\begin{theorem}{向量函数的求导法则}{}
\[
\begin{aligned}
&\text{(1)}\ \frac{d}{dt}\vec{C} = \vec{0} \\[4pt]
&\text{(2)}\ \frac{d}{dt}\big[c\,\vec{r}(t)\big] = c\,\vec{r}\,'(t) \\[4pt]
&\text{(3)}\ \frac{d}{dt}\big[\vec{r}_1(t) \pm \vec{r}_2(t)\big] = \vec{r}_1\,'(t) \pm \vec{r}_2\,'(t) \\[4pt]
&\text{(4)}\ \frac{d}{dt}\big[\varphi(t)\,\vec{r}(t)\big] = \varphi'(t)\,\vec{r}(t) + \varphi(t)\,\vec{r}\,'(t) \\[4pt]
&\text{(5)}\ \frac{d}{dt}\big[\vec{r}_1(t) \cdot \vec{r}_2(t)\big] = \vec{r}_1\,'(t) \cdot \vec{r}_2(t) + \vec{r}_1(t) \cdot \vec{r}_2\,'(t) \\[4pt]
&\text{(6)}\ \frac{d}{dt}\big[\vec{r}_1(t) \times \vec{r}_2(t)\big] = \vec{r}_1\,'(t) \times \vec{r}_2(t) + \vec{r}_1(t) \times \vec{r}_2\,'(t) \\[4pt]
&\text{(7)}\ \frac{d}{dt}\big[\vec{r}(\varphi(t))\big] = \vec{r}\,'(\varphi(t))\,\varphi'(t)
\end{aligned}
\]
\end{theorem}

有了对向量函数的基本认识，我们就可以定义第二类曲线积分了。

\begin{definition}{第二类曲线积分}{}
设向量函数 \(\vec{F}(x,y,z) = (P(x,y,z),\,Q(x,y,z),\,R(x,y,z))\) 在空间有向光滑曲线弧段 \(C\) 上有界。将 \(C\) 依其方向任意分成 \(n\) 个有向小弧段
\[
\overrightarrow{M_{0}M_{1}},\ \overrightarrow{M_{1}M_{2}},\ \cdots,\ \overrightarrow{M_{n-1}M_{n}}
\]
记每个有向小弧段在坐标轴上的投影为 \(\Delta x_i, \Delta y_i, \Delta z_i\)，并记其弧长为 \(\Delta s_i\)。在每个有向小弧段上任取一点 \((\xi_i,\eta_i,\zeta_i)\)，作和
\[
S = \sum_{i=1}^{n} \bigl[ P(\xi_i,\eta_i,\zeta_i)\,\Delta x_i + Q(\xi_i,\eta_i,\zeta_i)\,\Delta y_i + R(\xi_i,\eta_i,\zeta_i)\,\Delta z_i \bigr]
\]

记 \(\lambda = \max\{\Delta s_1,\Delta s_2,\cdots,\Delta s_n\}\)。如果当 \(\lambda \to 0\) 时，\(S\) 的极限总存在，且与曲线 \(C\) 的分法及点 \((\xi_i,\eta_i,\zeta_i)\) 的取法无关，那么称 \(\vec{F}(x,y,z)\) 在有向曲线弧段 \(C\) 上\textbf{可积}，并称这个极限为函数 \(\vec{F}(x,y,z)\) 在有向曲线弧段 \(C\) 上的\textbf{第二类曲线积分}（或\textbf{对坐标的曲线积分}），记作
\[
\int_{C} P(x,y,z)\,dx + Q(x,y,z)\,dy + R(x,y,z)\,dz
\]
或记作向量形式
\[
\int_{C} \vec{F}(x,y,z) \cdot d\vec{r}
\]
即：
\[
\int_{C} \vec{F}(x,y,z) \cdot d\vec{r} = \lim_{\lambda \to 0} \sum_{i=1}^{n} \bigl[ P(\xi_i,\eta_i,\zeta_i)\,\Delta x_i + Q(\xi_i,\eta_i,\zeta_i)\,\Delta y_i + R(\xi_i,\eta_i,\zeta_i)\,\Delta z_i \bigr] = I
\]

其中 \(P(x,y,z), Q(x,y,z), R(x,y,z)\) 称为\textbf{被积函数}，\(P\,dx + Q\,dy + R\,dz\)（或 \(\vec{F} \cdot d\vec{r}\)）称为\textbf{被积表达式}，\(d\vec{r} = (dx, dy, dz)\) 称为\textbf{有向弧微元}（或\textbf{弧微分向量}），\(x,y,z\) 称为\textbf{积分变量}，有向曲线 \(C\) 称为\textbf{积分路径}。
\end{definition}

这个定义似乎很难理解。想象一下你在无限大的海里划船，海里某个区域水流方向可能往北，也可能往南，水流速度可能很快，也可能很慢。向量函数就是整片海的“水流地图”，每个位置的水流方向和速度都标得清清楚楚。你划船时，只有水流沿船头方向的分量才会帮你前进，也就是对你做正功。侧向的水流只会让你偏移，对前进无效，所以不做功。沿船尾方向的分量则会阻碍你前进，对你做负功。你划船的路线可能是直线，可能是曲线，也有可能闭合的曲线。沿着你划船的路线，水流对你做的功就是第二类曲线积分。

换句话说，第二类曲线积分计算的是变力对物体做的功。我们知道功等于力与位移的点乘，力和位移都是矢量，所以第二类曲线积分是有方向的。

也就是说，如果将被积函数\(\vec{F}(x,y,z)\)看作是作用在物体上的变力，有向曲线 \(C\) 看作是物体的运动轨迹，那么\(\int_{C} \vec{F}(x,y,z) \cdot d\vec{r}\)就表示的是沿着物体的运动轨迹，这个变力对物体做的总功。

第二类曲线积分的定义也和之前的积分一脉相承。把整条路径切成很多有向小段，每一小段近似看成一段直线位移，而投影 \(\Delta x_i, \Delta y_i, \Delta z_i\)，就是每一小段在\(x\)，\(y\)，\(z\)三个方向上的分量。在每一小段里，力的大小和方向变化非常小，可以近似认为是恒定的。因此可以任取一点 \((\xi_i,\eta_i,\zeta_i)\)，作乘积\(P(\xi_i,\eta_i,\zeta_i)\,\Delta x_i + Q(\xi_i,\eta_i,\zeta_i)\,\Delta y_i + R(\xi_i,\eta_i,\zeta_i)\,\Delta z_i\)，如果我们把\(P\)，\(Q\)，\(R\)理解为力\(\vec{F}(x,y,z)\)在\(x\)，\(y\)，\(z\)三个方向上的分力，也就是\(\vec{F}(x,y,z)=P\vec{i}+Q\vec{j}+R\vec{k}\)，那么这个乘积其实就是力对物体做的功。我们把这个功累加，取极限，得到的就是总功。

第二类曲线积分有方向性，所以要尤其注意积分路径的方向。其他性质则和定积分类似。当积分路径是闭合曲线时，第二类曲线积分记作\(\oint_{C} \vec{F}(x,y,z) \cdot d\vec{r}\)。
\begin{theorem}{第二类曲线积分的方向性}{}
    若\(-C\)表示与\(C\)方向相反的同一条曲线，则有：
    \[\int_{-C} \vec{F}(x,y,z) \cdot d\vec{r}=\int_{C} -\vec{F}(x,y,z) \cdot d\vec{r}\]
\end{theorem}
计算第二类曲线积分，其实就是把力沿曲线做功转化为我们可以求解的定积分或重积分。顾名思义，第二类曲线积分又叫对坐标的曲线积分，所以我们主要对\(x\)，\(y\)，\(z\)三个坐标计算。

\begin{theorem}{第二类曲线积分的计算法}{}
设空间有向光滑曲线 \( C \) 的参数方程为  
\(
C(t) = \begin{cases}
x = x(t)\\
y = y(t)\\
z = z(t)
\end{cases}\)，\(t \in [\alpha, \beta]\)，其中 \(\alpha\) 到 \(\beta\) 的方向为对应曲线 \( C \) 的方向。被积向量函数为 \(\vec{F}(x,y,z) = (P(x,y,z), Q(x,y,z), R(x,y,z))\)，则有向弧微元化为
\[
d\vec{r} = C\,'(t)\,dt = (x'(t),\ y'(t),\ z'(t))\,dt
\]
积分化为
\[
\begin{aligned}
\int_C \vec{F} \cdot d\vec{r}
&= \int_C P\,dx + Q\,dy + R\,dz \\
&= \int_{\alpha}^{\beta} [ P(x(t),y(t),z(t))\,x'(t) \\
&\qquad + Q(x(t),y(t),z(t))\,y'(t) \\
&\qquad + R(x(t),y(t),z(t))\,z'(t) ]\,dt
\end{aligned}
\]
\end{theorem}


当光滑曲线 \( C \) 在平面上时，第二类曲线积分退化为：

\begin{theorem}{平面第二类曲线积分的计算法}{}
设平面有向光滑曲线 \( C \) 的参数方程为  
\(
C(t) = \begin{cases}
x = x(t)\\
y = y(t)
\end{cases}\)，\(t \in [\alpha, \beta]\)，其中 \(\alpha\) 到 \(\beta\) 的方向为对应曲线 \( C \) 的方向。被积向量函数为 \(\vec{F}(x,y) = (P(x,y), Q(x,y))\)，则有向弧微元化为
\[
d\vec{r} = C\,'(t)\,dt = (x'(t),\ y'(t))\,dt
\]
积分化为
\[
\int_C \vec{F} \cdot d\vec{r} = \int_C P\,dx + Q\,dy = \int_{\alpha}^{\beta} [ P(x(t),y(t))\,x'(t) + Q(x(t),y(t))\,y'(t) ]\,dt
\]
\end{theorem}

与第一类曲线积分的计算法不同，第二类曲线积分记作\(\int_C P\,dx + Q\,dy + R\,dz\)，这说明，第二类曲线积分是直接对\(x\)，\(y\)，\(z\)三个变量积分，而不是对弧长积分，这使得第二类曲线积分的计算略显简单，不过千万不能和第一类曲线积分混淆了。记住第一类曲线积分又叫对弧长的曲线积分，所以积分后面是\(ds\)，第二类曲线积分又叫对坐标的曲线积分，所以积分后面是\(dx\)，\(dy\)，\(dz\)。

看起来很复杂，但是第二类曲线积分是对\(dx\)，\(dy\)，\(dz\)积分，所以其实和定积分是类似的。基本思路就是对积分路径$C$求导，将所有\(dx\)，\(dy\)，\(dz\)都化为同一个变量。如果$C$是由参数$t$确定的，就换成$dt$，如果$C$是由\(x\)，\(y\)，\(z\)中某个变量确定的，就换成它的微分。统一成同一个变量之后就看最后对哪个变量积分，然后根据$C$的起点和终点确定上下限。例如$C$从\((0,0,0)\)到\((1,2,3)\)，如果最后是\(dx\)，对\(x\)积分，那就是从\(0\)积到\(1\)，如果是\(dy\)，对\(y\)积分，那就是从\(0\)积到\(2\)，以此类推。

我们来看几个例子：

\begin{myexample}
    计算\(\int_C y\,dx + x\,dy\)，其中 $C$ 是抛物线 $y = x^2$ 上从点 $(0,0)$ 到点 $(1,1)$ 的一段。

    对$y = x^2$求导得\(dy=2xdx\)。将$y = x^2$和\(dy=2xdx\)代入原积分：
    \[
    \int_C y\,dx + x\,dy=\int_C x^2\,dx + x\cdot 2xdx=\int_C 3x^2 \,dx
    \]

    此时原积分已经转化为对\(dx\)的积分，我们只需要再把它转化为定积分就可以了。

    由于$C$从点 $(0,0)$ 到点 $(1,1)$，因此原积分从\(0\)积到\(1\)：
    \[
    \int_C 3x^2 \,dx=\int_0^1 3x^2\,dx = [x^3]_0^1 = 1
    \]

    所以\(\int_C y\,dx + x\,dy=1\)。
\end{myexample}

整个步骤不需要什么代入弧长微元，只需要求微分，把所有的变量都统一用一个变量表达，再根据积分路径确定上下限，转换为定积分就可以直接计算了。

\begin{myexample}
    计算\(\int_C xy\,dx\)，其中 $C$ 是抛物线 $y = \sqrt{x}$ 上从点 $(1,1)$ 到点 $(0,0)$ 的一段。

    原来的$y = \sqrt{x}$求导后不好计算，我们不妨改写成$x = y^2$后再求导，得到\(dx=2ydy\)，将$x = y^2$和\(dx=2ydy\)代入原积分：
    \[
    \int_C xy\,dx=\int_C y^2\cdot y\cdot 2ydy=\int_C 2y^4dy
    \]

    $C$ 从点 $(1,1)$ 到点 $(0,0)$，原积分从\(1\)积到\(0\)：
    \[
    \int_C 2y^4dy=\int_{1}^{0} 2y^4\,dy=-\int_{0}^{1} 2y^4\,dy =-2\left[ \frac{y^5}{5} \right]_{0}^{1}= -\frac{2}{5}
    \]

    所以\(\int_C xy\,dx=-\frac{2}{5}\)。
\end{myexample}

这里要注意，$C$ 是从点 $(1,1)$ 到点 $(0,0)$，因此积分下限一定是$y=1$，积分上限一定是$y=0$，这也是第二类曲线积分方向性的体现。千万不能下意识就把小的放下限，大的放上限了。

\begin{myexample}
    计算\(\int_C z\,dx + x\,dy + y\,dz\)，其中 $C$ 为从点 $(0,0,0)$ 到点 $(1,1,1)$ 的直线段。

    这里$C$没有给出具体的方程，但是我们可以把它求出来。由于$C$ 是从点 $(0,0,0)$ 到点 $(1,1,1)$ 的直线段，所以$C$的方程可以用$x=y=z$表示，那么就有$dx=dy=dz$。代入$x=y=z$和$dx=dy=dz$，原积分化为：
\[
    \int_C z\,dx + x\,dy + y\,dz=\int_C x\,dx + x\,dx + x\,dx=\int_C 3x\,dx\]

    因此：
    \[
    \int_C 3x\,dx=\int_0^1 3x\,dx=3\left[ \frac{x^2}{2} \right]_{0}^{1}=\frac32
    \]

    所以\(\int_C z\,dx + x\,dy + y\,dz=\frac32\)。
\end{myexample}

\begin{myexample}
    计算\(\oint_C xy\,dy\)，其中 $C$ 为逆时针圆周 \(x^2+y^2=1\)。

    我们可以把$C$化成参数方程，令\(x = \cos t\)，\(y = \sin t\)，则\(0 \le t \le 2\pi\)，\(dy = \cos t\,dt\)，原积分化为：

    \[
\oint_C xy\,dy
= \oint_C \cos t \cdot \sin t \cdot \cos t\,dt
= \oint_C \cos^2 t \sin t\,dt
\]

$C$ 为逆时针方向，因此：
\[
\oint_C \cos^2 t \sin t\,dt=\int_0^{2\pi} \cos^2 t \sin t\,dt
\]

令 \(u = \cos t\)，则 \(du = -\sin t\,dt\)，当\(t=0\)时\(u=1\)，当\(t=2\pi\)时\(u=1\)，换元得：

\[
\int_0^{2\pi} \cos^2 t \sin t\,dt=-\int_1^1 u^2\,du=0
\]

所以\(\oint_C xy\,dy=0\)。
\end{myexample}

总结一下，第一类曲线积分是对积分路径的函数求导，然后再把结果代入公式替换掉弧微分\(ds\)，第二类曲线积分则把结果代入原积分，把积分化成只对一个变量积分的形式。要注意的是，第一类曲线积分没有方向性，积分上下限只与积分路径的端点有关。第二类曲线积分有方向性，要根据积分路径的起点和终点确定积分上下限。注意区分第一类曲线积分和第二类曲线积分。

\section{曲面积分}
解决了曲线积分，我们现在进入曲面积分。和曲线积分一样，曲面积分也分为两类。

\subsection{第一类曲面积分}
设空间内有一张光滑曲面片，它的面密度函数\(f(x,y,z)\)，以该光滑曲面片为区域对\(f(x,y,z)\)积分，这种积分我们就称为\textbf{第一类曲面积分}，也叫\textbf{对面积的曲面积分}。
\begin{definition}{第一类曲面积分}{}
    设函数 \(f(x,y,z)\) 在空间光滑曲面片 \(\Sigma\) 上有界。将 \(\Sigma\) 任意分成 \(n\) 个小曲面片
\[
\Delta S_1,\ \Delta S_2,\ \cdots,\ \Delta S_n
\]
每个小曲面片既表示该曲面片本身，也表示该曲面片的面积。在每个小曲面片 \(\Delta S_i\) 上任取一点 \((\xi_i,\eta_i,\zeta_i)\)，作函数值 \(f(\xi_i,\eta_i,\zeta_i)\) 与小曲面片面积 \(\Delta S_i\) 的乘积 \(f(\xi_i,\eta_i,\zeta_i) \cdot \Delta S_i\)，并作出和
\[
S = \sum_{i=1}^{n} f(\xi_i,\eta_i,\zeta_i) \,\Delta S_i
\]

记 \(\lambda = \max\{\Delta S_1,\Delta S_2,\cdots,\Delta S_n\}\)。如果当 \(\lambda \to 0\) 时，\(S\) 的极限总存在，且与曲面 \(\Sigma\) 的分法及点 \((\xi_i,\eta_i,\zeta_i)\) 的取法无关，那么称 \(f(x,y,z)\) 在曲面片 \(\Sigma\) 上\textbf{可积}，并称这个极限 \(I\) 为函数 \(f(x,y,z)\) 在曲面片 \(\Sigma\) 上的\textbf{第一类曲面积分}（或\textbf{对面积的曲面积分}），记作
\[
\iint_{\Sigma} f(x,y,z)\,dS
\]
即：
\[
\iint_{\Sigma} f(x,y,z)\,dS = \lim_{\lambda \to 0} \sum_{i=1}^{n} f(\xi_i,\eta_i,\zeta_i)\,\Delta S_i = I
\]

其中 \(f(x,y,z)\) 称为\textbf{被积函数}，\(f(x,y,z)\,dS\) 称为\textbf{被积表达式}，\(dS\) 称为\textbf{面积微元}（或\textbf{面积微分}），\(x,y,z\) 称为\textbf{积分变量}，\(\Sigma\) 称为\textbf{积分曲面}，\(\iint_{\Sigma}\) 称为\textbf{曲面积分号}。
\end{definition}

和第一类曲线积分类似，第一类曲面积分求的也是质量，只不过是把曲线换成了曲面，所以它具有与对弧长的曲线积分类似的性质，此处不再赘述。特殊的，当积分曲面是一张闭合曲面时，第一类曲面积分记作\(\oiint_{\Sigma} f(x,y,z)\,dS\)。

仿照第一类曲线积分，我们有：
\begin{theorem}{第一类曲面积分的计算法}{}
    设光滑曲面 $\Sigma$ 由 $z = z(x,y)$ 确定，$\Sigma$在$xOy$面上的投影区域为$D_{xy}$，则面积微元化为
\[
dS = \sqrt{1 + \left(\frac{\partial z}{\partial x}\right)^2 + \left(\frac{\partial z}{\partial y}\right)^2}\,dxdy
\]

积分化为
\[
\iint_{\Sigma} f(x,y,z)\,dS
= \iint_{D_{xy}} f\big(x,y,z(x,y)\big)\,
\sqrt{1 + \left(\frac{\partial z}{\partial x}\right)^2 + \left(\frac{\partial z}{\partial y}\right)^2}\,dxdy
\]

若积分由$x = x(y,z)$或$y = y(x,z)$确定，计算方法类似。
\end{theorem}

计算第一类曲面积分的基本思路就是化成二重积分。先把曲面投影到对应坐标面上，然后对曲面方程求偏导，替换掉面积微元，然后以投影区域为积分区域对被积函数求二重积分，得到的就是第一类曲面积分的结果。

我们来看几个例子：

\begin{myexample}
计算 \(\iint_{\Sigma} dS\)，其中 \(\Sigma\) 是由 \(x^2+y^2+z^2=1\) 被平面 \(z=0\) 所截的曲面的上半部分。

曲面方程可以写成 \(z=\sqrt{1-x^2-y^2}\)，由于\(\Sigma\)取上半部分，那么它在 \(xOy\) 平面上的投影区域\(D_{xy}\)就是\(x^2+y^2\le 1\)。

求偏导数得：
\[
z_x = \frac{-x}{\sqrt{1-x^2-y^2}},\qquad
z_y = \frac{-y}{\sqrt{1-x^2-y^2}}
\]

代入得：
\[
dS = \sqrt{1+z_x^2+z_y^2}\,dxdy
= \sqrt{1+\frac{x^2}{1-x^2-y^2}+\frac{y^2}{1-x^2-y^2}}\,dxdy
= \frac{1}{\sqrt{1-x^2-y^2}}\,dxdy
\]

于是：
\[
\iint_{\Sigma} dS = \iint_{D_{xy}} \frac{1}{\sqrt{1-x^2-y^2}}\,dxdy
\]

令 \(x=\rho\cos\theta,\; y=\rho\sin\theta\)，则 \(0\le \rho\le 1\)，\(0\le\theta\le2\pi\)，\(dxdy = \rho\,d\rho\,d\theta\)，原积分化为：
\[
\iint_{D_{xy}} \frac{1}{\sqrt{1-x^2-y^2}}\,dxdy = \int_0^{2\pi} d\theta \int_0^1 \frac{1}{\sqrt{1-\rho^2}}\,\rho\,d\rho
\]

令 \(u=1-\rho^2\)，\(du=-2\rho\,d\rho\)，当 \(\rho=0\) 时 \(u=1\)，当\(\rho=1\) 时 \(u=0\)，换元得：
\[
\int_0^1 \frac{\rho}{\sqrt{1-\rho^2}}\,d\rho = \int_1^0 \frac{1}{\sqrt{u}}\left(-\frac{du}{2}\right)
= \frac12\int_0^1 u^{-\frac12}\,du
= \frac12\cdot 2[\sqrt{u}]_0^1 = 1
\]

再对\(\theta\)积分：
\[\int_0^{2\pi} d\theta=2\pi\]

所以\(\iint_{\Sigma} dS = 2\pi\)。
\end{myexample}

实际上，由于被积函数为$1$，所以这正是这个曲面的面积。

\begin{myexample}
    计算 \(\iint_{\Sigma}(x^2+y^2)\,dS\)，其中 \(\Sigma\) 是锥面 \(z=\sqrt{x^2+y^2}\) 被平面 \(z=1\) 所截下的部分。

    曲面方程为 \(z=\sqrt{x^2+y^2}\)，它被平面 \(z=1\) 所截，那么它在 \(xOy\) 平面上的投影区域\(D_{xy}\)就是\(x^2+y^2\le 1\)。

    求偏导得：
    \[
\frac{\partial z}{\partial x}= \frac{x}{\sqrt{x^2+y^2}},\qquad
\frac{\partial z}{\partial y}= \frac{y}{\sqrt{x^2+y^2}}
\]

于是：
\[
dS = \sqrt{1+\left(\frac{\partial z}{\partial x}\right)^2+\left(\frac{\partial z}{\partial y}\right)^2}\,dxdy
= \sqrt{1+\frac{x^2}{x^2+y^2}+\frac{y^2}{x^2+y^2}}\,dxdy
= \sqrt{2}\,dxdy
\]

代入得：
\[
\iint_{\Sigma}(x^2+y^2)\,dS
= \iint_{D_{xy}}(x^2+y^2)\cdot\sqrt{2}\,dxdy
\]

令 \(x=\rho\cos\theta,\; y=\rho\sin\theta\)，则 \(0\le \rho\le 1\)，\(0\le\theta\le2\pi\)，\(dxdy = \rho\,d\rho\,d\theta\)，原积分化为：

\[
\iint_{D_{xy}}(x^2+y^2)\cdot\sqrt{2}\,dxdy
= \sqrt{2}\int_0^{2\pi}d\theta\int_0^1 \rho^2\cdot \rho\,d\rho
= \sqrt{2}\cdot 2\pi\cdot \frac{1}{4}
= \frac{\sqrt{2}\pi}{2}
\]

所以\(\iint_{\Sigma}(x^2+y^2)\,dS=\frac{\sqrt{2}}{2}\pi\)。
\end{myexample}

重点就是如何转换为二重积分，先将区域投影，再求偏导，然后替换面积微元，最后计算。注意二重积分的积分区域是投影区域。其实和第一类曲线积分的计算如出一辙。

\begin{myexample}
    计算 \(\iint_{\Sigma} z\,dS\)，其中 \(\Sigma\) 是球面 \(x^2+y^2+z^2=a^2\) 在第一卦限内的部分。

    在第一卦限内，球面方程可写成 \(z=\sqrt{a^2-x^2-y^2}\)，它在 \(xOy\) 面上的投影区域\(D_{xy}\)就是\(x^2+y^2\le a^2,\;x\ge0,\;y\ge0\)。

    求偏导得：
    \[
z_x = \frac{-x}{\sqrt{a^2-x^2-y^2}},\qquad
z_y = \frac{-y}{\sqrt{a^2-x^2-y^2}}
\]

于是：
\[
dS = \sqrt{1+z_x^2+z_y^2}\,dxdy
= \frac{a}{\sqrt{a^2-x^2-y^2}}\,dxdy
= \frac{a}{z}\,dxdy
\]

代入得：
\[
\iint_{\Sigma} z\,dS =\iint_{D_{xy}} z \cdot \frac{a}{z}\,dxdy = \iint_{D_{xy}} a \,dxdy
\]

\(a\)不是变量，与\(dxdy\)无关，因此：
\[
\iint_{D_{xy}} a\,dxdy = a\iint_{D_{xy}}dxdy = a\cdot \frac{\pi a^2}{4}=\frac{\pi}{4} a^3
\]

所以\(\iint_{\Sigma} z\,dS=\frac{\pi}{4} a^3\)。
\end{myexample}

\begin{myexample}
计算 \(\iint_{\Sigma} xyz \, dS\)，其中 \(\Sigma\) 是平面 \(x+y+z=1\) 在第一卦限的部分。

将平面写为 \(z=1-x-y\)，它在 \(xOy\) 平面上的投影区域\(D_{xy}\)为\(x\ge 0,\;y\ge 0,\;x+y\le 1\)。

求偏导数得：
\[
z_x = -1,\qquad z_y = -1
\]

代入得：
\[
dS = \sqrt{1+z_x^2+z_y^2}\,dxdy
= \sqrt{1+(-1)^2+(-1)^2}\,dxdy
= \sqrt{3}\,dxdy
\]

原积分化为：
\[
\iint_{\Sigma} xyz \, dS
= \iint_{D_{xy}} xy(1-x-y) \cdot \sqrt{3}\,dxdy
= \sqrt{3} \iint_{D_{xy}} (xy - x^2y - xy^2)\,dxdy
\]

将二重积分化为二次积分：
\[
\iint_{D_{xy}} (xy - x^2y - xy^2)\,dxdy
= \int_0^1 dx \int_0^{1-x} (xy - x^2y - xy^2)\,dy
\]

先对\(y\)积分：
\[
\begin{aligned}
    \int_0^{1-x} (xy - x^2y - xy^2)\,dy
    &= \int_0^{1-x} [ (x-x^2)y - xy^2 ]\,dy
    = \left[(x-x^2)\frac{y^2}{2} - x\frac{y^3}{3}\right]_0^{1-x}\\
    &= \frac{x-x^2}{2}(1-x)^2 - \frac{x}{3}(1-x)^3
    = x(1-x)^3\cdot(\frac12 - \frac13)\\
    &= \frac{1}{6}x(1-x)^3
\end{aligned}
\]

再对\(x\)积分：
\[
\begin{aligned}
    \int_0^1 \frac{1}{6}x(1-x)^3\,dx
    &=\frac16\int_0^1 (x - 3x^2 + 3x^3 - x^4)\,dx
    = \frac16 \cdot \left[ \frac{x^2}{2} - x^3 + \frac{3x^4}{4} - \frac{x^5}{5} \right]_0^1\\
    &= \frac16 \cdot(\frac12 - 1 + \frac34 - \frac15)
    = \frac16 \cdot\frac{1}{20}
    = \frac{1}{120}
\end{aligned}
\]

代回原积分，得：
\[
\iint_{\Sigma} xyz \, dS = \sqrt{3} \cdot \frac{1}{120} = \frac{\sqrt{3}}{120}
\]

所以 \(\iint_{\Sigma} xyz \, dS = \frac{\sqrt{3}}{120}\)。
\end{myexample}


\subsection{第二类曲面积分}
和第二类曲线积分一样，第二类曲面积分也是有方向的。不过曲面本身并没有方向，想要定义第二类曲面积分，还得先给曲面指定一个方向。

想象一张平面，它有上下两面。哪面是正面呢？我们一般规定一个法向量，法向量从哪一面穿出来，哪一面就是正面。例如对\(xOy\)面来说，如果法向量指向\(z\)轴正方向，那么就以上侧为正面，下侧为反面，反之亦然。这种规定对一般的曲面也是如此。由于光滑曲面在每一点都有法向量，而法向量有两个相反的方向，取哪一个方向，就决定了取曲面的哪一侧。

\begin{definition}{有向曲面}{}
    设\(S\)是一片光滑曲面。如果存在一个定义在整个\(S\)上的连续的单位法向量函数\(\vec{n}(P)\)，\(P\in S\)，则称\(S\)是\textbf{可定向的}。

    在选定了\(\vec{n}\)的指向后，\(S\)连同这个选定的指向，就称为一个\textbf{有向曲面}。
\end{definition}

绝大多数曲面都能按照类似的定义分出正反两面，这种曲面称为\textbf{双侧曲面}。因为这类曲面可以定义连续的单位法向量函数。但也有极少数曲面只有一侧，最有名的就是莫比乌斯带。拿一张长纸条，把一端扭转一百八十度后再和另一端连起来，就得到一条莫比乌斯带。沿着莫比乌斯带走一圈，就会回到出发点的“另一面”。在这个过程中，法向量的方向会反转。这样的曲面无法区分两侧，自然也无法定义有向曲面。这种不能定义方向的曲面我们不作讨论。

特别的，当有向曲面是闭合曲面时，我们通常取\textbf{外侧}为正面，此时法向量指向曲面所围区域的外部。

在有向曲面的基础上，接下来讨论有向投影面积。

在第二类曲线积分的定义中，我们也是先把各个小弧段投影到坐标轴，再分别对它们的分量计算。但是有向曲面的投影和一般的投影区域不同，有向曲面是有方向的，因此它在坐标平面上的投影也应该体现方向性。

\begin{definition}{有向投影面积}{}
    设\(S\)是一片光滑曲面，已指定单位法向量函数\(\vec{n}(P)\)。对于坐标平面，规定其正法向量：
    \[\vec{e_x}=(1,0,0),\quad \vec{e_y}=(0,1,0),\quad \vec{e_z}=(0,0,1)\]

    如果将\(S\)投影到\(xOy\)平面，所得的投影区域面积为$a$，那么：

    若\(\vec{n}(P)\cdot \vec{e_z}<0\)，则有向投影面积取\textbf{负值}，即$-a$。

    若\(\vec{n}(P)\cdot \vec{e_z}=0\)，则有向投影面积为\textbf{$0$}，此时$a=0$。

    若\(\vec{n}(P)\cdot \vec{e_z}>0\)，则有向投影面积取\textbf{正值}，即$a$。

    投影到其他坐标平面时定义类似。
\end{definition}

现在，我们就能仿照第二类曲线积分，定义第二类曲面积分了。

\begin{definition}{第二类曲面积分}{}
    设向量函数 \(\vec{F}(x,y,z) = (P(x,y,z),\,Q(x,y,z),\,R(x,y,z))\) 在有向光滑曲面片 \(\Sigma\) 上有界。将 \(\Sigma\) 任意分成 \(n\) 个小曲面片
\[
\Delta S_1,\ \Delta S_2,\ \cdots,\ \Delta S_n
\]
每个小曲面片既表示该曲面片本身，也表示该曲面片的面积。在每个小曲面片 \(\Delta S_i\) 上任取一点 \((\xi_i,\eta_i,\zeta_i)\)，并按 \(\Sigma\) 所取的一侧在该点取定单位法向量 \(\vec{n}_i\)。把小曲面片 \(\Delta S_i\) 分别投影到 \(yOz\)、\(zOx\)、\(xOy\) 三个坐标面上，得到三个有向投影面积，记为 \(\Delta S_{i,yz}\)、\(\Delta S_{i,zx}\)、\(\Delta S_{i,xy}\)。作和
\[
S = \sum_{i=1}^{n} \bigl[ P(\xi_i,\eta_i,\zeta_i)\,\Delta S_{i,yz} + Q(\xi_i,\eta_i,\zeta_i)\,\Delta S_{i,zx} + R(\xi_i,\eta_i,\zeta_i)\,\Delta S_{i,xy} \bigr]
\]

记 \(\lambda = \max\{\Delta S_1,\Delta S_2,\cdots,\Delta S_n\}\)。如果当 \(\lambda \to 0\) 时，\(S\) 的极限总存在，且与曲面 \(\Sigma\) 的分法及点 \((\xi_i,\eta_i,\zeta_i)\) 的取法无关，那么称 \(\vec{F}(x,y,z)\) 在有向曲面片 \(\Sigma\) 上\textbf{可积}，并称这个极限为函数 \(\vec{F}(x,y,z)\) 在有向曲面片 \(\Sigma\) 上的\textbf{第二类曲面积分}（或\textbf{对坐标的曲面积分}），记作
\[
\iint_{\Sigma} P(x,y,z)\,dy\,dz + Q(x,y,z)\,dz\,dx + R(x,y,z)\,dx\,dy
\]
或记作向量形式
\[
\iint_{\Sigma} \vec{F}(x,y,z) \cdot \vec{n}\,dS
\]
即：
\[
\begin{aligned}
    &\iint_{\Sigma} \vec{F}(x,y,z) \cdot \vec{n}\,dS\\ 
    &= \lim_{\lambda \to 0} \sum_{i=1}^{n} \bigl[ P(\xi_i,\eta_i,\zeta_i)\,\Delta S_{i,yz} + Q(\xi_i,\eta_i,\zeta_i)\,\Delta S_{i,zx} + R(\xi_i,\eta_i,\zeta_i)\,\Delta S_{i,xy} \bigr] 
    = I
\end{aligned}
\]

其中 \(P(x,y,z), Q(x,y,z), R(x,y,z)\) 称为\textbf{被积函数}，\(P\,dy\,dz + Q\,dz\,dx + R\,dx\,dy\)（或 \(\vec{F} \cdot \vec{n}\,dS\)）称为\textbf{被积表达式}，\(\vec{n}\,dS = (dy\,dz,\ dz\,dx,\ dx\,dy)\) 称为\textbf{有向面积微元}（或\textbf{面积微元向量}），\(x,y,z\) 称为\textbf{积分变量}，有向曲面 \(\Sigma\) 称为\textbf{积分曲面}，\(\iint_{\Sigma}\) 称为\textbf{曲面积分号}。
\end{definition}

这个定义看起来有点绕，我们来看看它到底在算什么。想象一条大河，河水在每一点的流速都不同，方向也不同。如果在大河里取无数多个点，每一点对应一个流速，这些点和流速就构成了一个向量函数 \(\vec{F}\)。现在我们拿一张渔网放进河里，某一点的水流可能流得很快，也可能流得很慢。水流可能会从网的正面穿过来，也可能从网的反面穿回去。我们规定，从正面穿过的水量为正，从反面穿回的水量为负，那么单位时间内穿过整张网的水的质量，就是总流量，也就是第二类曲面积分计算的结果。

我们一般认为，真正穿过网的水并不是整条水流，而只是水流沿网面法方向的那个分量。贴着网面流的水只会从网面上滑过去，并不会穿过网。所以每一小片对总流量的贡献，就是水流速度在法方向上的分量乘上这一小片的面积，也就是 \(\vec{F}\cdot\vec{n}\,dS\)。把这个贡献按三个坐标面分解，就得到了定义里的 \(P\,\Delta S_{yz} + Q\,\Delta S_{zx} + R\,\Delta S_{xy}\)。所以说，第二类曲面积分和之前的积分也是一脉相承的，只是把“函数值乘面积”换成了“向量沿法方向的分量乘面积”。

换句话说，第二类曲面积分计算的是向量函数穿过有向曲面的\textbf{通量}。如果 \(\vec{F}\) 是流体的速度函数，通量就是穿过曲面的流量。如果 \(\vec{F}\) 是电场强度，通量就是电通量。如果 \(\vec{F}\) 是磁感应强度，通量就是磁通量。总之，凡是“某种流穿过一张网”的场景，都离不开第二类曲面积分。

第二类曲面积分是有方向的，换了正反面，积分结果就要变号，这一点和第一类曲面积分不同，但是和第二类曲线积分是类似的。

\begin{theorem}{第二类曲面积分的方向性}{}
    若 \(-\Sigma\) 表示与 \(\Sigma\) 取相反一侧的同一张曲面，则有：
    \[\iint_{-\Sigma} \vec{F}(x,y,z) \cdot \vec{n}\,dS = \iint_{\Sigma} -\vec{F}(x,y,z) \cdot \vec{n}\,dS\]
\end{theorem}

第二类曲面积分的其他性质和第一类曲面积分类似，满足线性性质，可加性等性质，此处不再赘述。当积分曲面是闭合曲面时，第二类曲面积分记作 \(\oiint_{\Sigma} P\,dy\,dz + Q\,dz\,dx + R\,dx\,dy\)。

计算第二类曲面积分的基本思路和第一类曲面积分一样，都是化成二重积分。不过和第一类曲面积分相比，要注意两个问题。第一个问题是，第二类曲面积分是对坐标的积分，\(P\,dy\,dz\)，\(Q\,dz\,dx\)，\(R\,dx\,dy\) 三项分别对应曲面在 \(yOz\)，\(zOx\)，\(xOy\) 三个坐标面上的投影，所以题目要求哪一项，我们就把曲面方程改写成相应的形式。比如要求 \(R\,dx\,dy\)，就把曲面写成 \(z=z(x,y)\)，要求 \(P\,dy\,dz\)，就把曲面写成 \(x=x(y,z)\)。第二个问题就是符号。最后计算结果千万不能忘记判断方向，方向会影响结果的符号。

\begin{theorem}{第二类曲面积分的计算法}{}
    设光滑曲面 \(\Sigma\) 由 \(z = z(x,y)\) 确定，\(\Sigma\) 在 \(xOy\) 面上的投影区域为 \(D_{xy}\)。若 \(\Sigma\) 取正面，则
\[
\iint_{\Sigma} R(x,y,z)\,dx\,dy = \iint_{D_{xy}} R\big(x,y,z(x,y)\big)\,dx\,dy
\]
若 \(\Sigma\) 取反面，则
\[
\iint_{\Sigma} R(x,y,z)\,dx\,dy = -\iint_{D_{xy}} R\big(x,y,z(x,y)\big)\,dx\,dy
\]

若光滑曲面 \(\Sigma\) 由\(x = x(y,z)\)或\(y = y(z,x)\) 确定，计算方法类似。
\end{theorem}

如果觉得一项一项投影太麻烦，第二类曲面积分还有一个更统一的算法。设曲面 \(\Sigma\) 由 \(z = z(x,y)\) 确定并取上侧为正，则单位法向量为
\[
\vec{n} = \frac{(-z_x,\ -z_y,\ 1)}{\sqrt{1+z_x^2+z_y^2}}
\]

而第一类曲面积分中我们已经知道 \(dS = \sqrt{1+z_x^2+z_y^2}\,dx\,dy\)，于是：
\[
\vec{n}\,dS = (-z_x,\ -z_y,\ 1)\,dx\,dy
\]

代入向量形式，得：
\[
\iint_{\Sigma} \vec{F} \cdot \vec{n}\,dS = \iint_{D_{xy}} \bigl[ -P(x,y,z)\,z_x - Q(x,y,z)\,z_y + R(x,y,z) \bigr]\,dx\,dy
\]

这里的 \(z\) 都要代入 \(z = z(x,y)\)。这个公式把三个坐标面上的投影统一成了对 \(xOy\) 面的一个二重积分。不过我们一般还是用投影比较多。

我们来看几个例子：

\begin{myexample}
计算 \(\iint_{\Sigma} \,dx\,dy\)，其中 \(\Sigma\) 是球面 \(x^2+y^2+z^2=1\) 被平面 \(z=0\) 所截的曲面的上半部分，取上侧。

曲面方程可以写成 \(z=\sqrt{1-x^2-y^2}\)，由于 \(\Sigma\) 取上半部分，那么它在 \(xOy\) 平面上的投影区域\(D_{xy}\)就是\(x^2+y^2\le 1\)。\(\Sigma\) 取上侧，符号取正，于是：
\[
\iint_{\Sigma} \,dx\,dy = \iint_{D_{xy}} \,dx\,dy
\]

因为被积函数是 \(1\)，这个二重积分就是投影区域 \(D_{xy}\) 的面积。\(D_{xy}\) 是半径为 \(1\) 的圆，面积为 \(\pi\)，所以：
\[
\iint_{\Sigma} \,dx\,dy = \pi
\]
\end{myexample}

上一节我们算过，同一张上半球面的第一类曲面积分 \(\iint_{\Sigma} dS = 2\pi\)，而现在的结果是 \(\iint_{\Sigma} \,dx\,dy = \pi\)，这正是因为第一类曲面积分积的是曲面本身，第二类曲面积分积的是投影。事实上，当被积函数为 \(1\) 时，第二类曲面积分 \(\iint_{\Sigma} dx\,dy\) 就等于积分曲面在 \(xOy\) 面上的投影区域的面积。

\begin{myexample}
计算 \(\iint_{\Sigma} z\,dx\,dy\)，其中 \(\Sigma\) 是锥面 \(z=\sqrt{x^2+y^2}\) 被平面 \(z=1\) 所截下的部分，取下侧。

曲面方程为 \(z=\sqrt{x^2+y^2}\)，它被平面 \(z=1\) 所截，那么它在 \(xOy\) 平面上的投影区域\(D_{xy}\)就是\(x^2+y^2\le 1\)。\(\Sigma\) 取下侧，符号取负，于是：
\[
\iint_{\Sigma} z\,dx\,dy = -\iint_{D_{xy}} \sqrt{x^2+y^2}\,dx\,dy
\]

令 \(x=\rho\cos\theta,\; y=\rho\sin\theta\)，则 \(0\le \rho\le 1\)，\(0\le\theta\le2\pi\)，\(dx\,dy = \rho\,d\rho\,d\theta\)，原积分化为：
\[
-\iint_{D_{xy}} \sqrt{x^2+y^2}\,dx\,dy = -\int_0^{2\pi} d\theta \int_0^1 \rho\cdot \rho\,d\rho
= -2\pi\cdot\frac13 = -\frac{2\pi}{3}
\]

所以\(\iint_{\Sigma} z\,dx\,dy = -\frac{2\pi}{3}\)。
\end{myexample}

注意这里\(\Sigma\)取下侧，因此符号取负。如果我们把取下侧改成取上侧，答案就会从负的 \(\frac{2\pi}{3}\) 变成正的 \(\frac{2\pi}{3}\)。这就是第二类曲面积分的方向性的体现。

\begin{myexample}
计算 \(\iint_{\Sigma} x\,dy\,dz + y\,dz\,dx + z\,dx\,dy\)，其中 \(\Sigma\) 是平面 \(x+y+z=1\) 在第一卦限的部分，取上侧。

\(\Sigma\) 取上侧时，单位法向量为 \(\vec{n}=\frac{1}{\sqrt{3}}(1,1,1)\)，三个分量都为正，所以三个坐标面上的投影都取正号。我们把三项分别算出来。

先算 \(z\,dx\,dy\) 项。把平面写成 \(z=1-x-y\)，它在 \(xOy\) 面上的投影区域\(D_{xy}\)为\(x\ge 0,\;y\ge 0,\;x+y\le 1\)。于是：
\[
\iint_{\Sigma} z\,dx\,dy = \iint_{D_{xy}} (1-x-y)\,dx\,dy
= \int_0^1 dx \int_0^{1-x} (1-x-y)\,dy
= \int_0^1 \frac{(1-x)^2}{2}\,dx = \frac16
\]

再算 \(x\,dy\,dz\) 项。把平面写成 \(x=1-y-z\)，它在 \(yOz\) 面上的投影区域\(D_{yz}\)为\(y\ge 0,\;z\ge 0,\;y+z\le 1\)。于是：
\[
\iint_{\Sigma} x\,dy\,dz = \iint_{D_{yz}} (1-y-z)\,dy\,dz = \frac16
\]

同理，\(y\,dz\,dx\) 项也是 \(\frac16\)。相加得：
\[
\iint_{\Sigma} x\,dy\,dz + y\,dz\,dx + z\,dx\,dy = \frac16+\frac16+\frac16 = \frac12
\]

所以 \(\iint_{\Sigma} x\,dy\,dz + y\,dz\,dx + z\,dx\,dy = \frac12\)。

其实还有第二种方法。既然 \(\Sigma\) 取上侧，单位法向量 \(\vec{n}=\frac{1}{\sqrt{3}}(1,1,1)\)，被积向量函数 \(\vec{F}=(x,y,z)\)，而在平面 \(x+y+z=1\) 上，\(\vec{F}\cdot\vec{n}=\frac{x+y+z}{\sqrt{3}}=\frac{1}{\sqrt{3}}\)。我们之前算过这张平面的面积微元是 \(dS=\sqrt{3}\,dx\,dy\)，于是：
\[
\iint_{\Sigma} \vec{F}\cdot\vec{n}\,dS = \iint_{D_{xy}} \frac{1}{\sqrt{3}}\cdot\sqrt{3}\,dx\,dy = \iint_{D_{xy}} dx\,dy = \frac12
\]

这里的 \(\iint_{D_{xy}} dx\,dy\) 就是平面在第一卦限的三角形在 \(xOy\) 面上的投影面积，所以就是 \(\frac12\)。
\end{myexample}

总结一下，第一类曲面积分的被积函数是数，积的是曲面本身的面积微元 \(dS\)，没有方向性，算的是质量。第二类曲面积分的被积函数是向量，积的是曲面在三个坐标面上的有向投影，有方向性，算的是通量。这和第一类曲线积分与第二类曲线积分的区别如出一辙。一个对面积，一个对坐标。一个没有方向，一个有方向。两者的计算最后都化成二重积分，区别只是面积微元换成了 \(dx\,dy\)、\(dy\,dz\)、\(dz\,dx\)，并且多了正负号。注意区分第一类曲面积分和第二类曲面积分。

\section{格林公式}

前面我们分别介绍了曲线积分和曲面积分。现在，我们要探索一下，闭曲线上的第二类曲线积分和它所围成的平面区域上的二重积分之间会不会有什么联系？答案是肯定的。那就是格林公式。

\subsection{格林公式}

要理解格林公式，我们不妨先回到之前的牛顿-莱布尼茨公式：
\[
\int_a^b f'(x)\,dx = f(b) - f(a)
\]

它说的是，函数在区间内部的“累计变化”，其实完全由区间两个端点上的函数值决定。换句话说，内部的信息，被边界的信息包住了。格林公式正是这句话的二维版本。闭曲线围成的区域内部的二重积分，等于边界上的曲线积分。一维里是两端相减，二维里就变成沿闭合边界走一圈。

设 \(D\) 是一块由分段光滑的简单闭曲线围成的平面区域，\(L\) 是它的边界。我们规定，\(D\) 的正向边界是指人沿着边界走一圈时，区域 \(D\) 始终在他的左手边。在通常的坐标系里，这其实是逆时针方向。有了这个正方向，我们就可以写出格林公式了。

\begin{theorem}{格林公式}{}
设闭区域 \(D\) 由分段光滑的简单闭曲线 \(L\) 围成，函数 \(P(x,y)\)、\(Q(x,y)\) 在 \(D\) 上具有一阶连续偏导数，则
\[
\oint_L P\,dx + Q\,dy = \iint_D \left(\frac{\partial Q}{\partial x} - \frac{\partial P}{\partial y}\right) dx\,dy
\]
其中 \(L\) 取 \(D\) 的正向边界。
\end{theorem}

我们简要证明一下。

先只看 \(P\) 这一项，也就是验证\(\oint_L P\,dx = -\iint_D \frac{\partial P}{\partial y}\,dx\,dy\)。

用平行于 \(y\) 轴的直线穿入区域 \(D\)。我们设区域是\(X\)型区域，它的下边界是 \(y=y_1(x)\)，上边界是 \(y=y_2(x)\)，$x\in[a,b]$。边界 $L$ 被直线分为四段，我们记下边界为 $L_1$，上边界为$L_2$。在竖直边界上 $x$ 不变，故 $dx=0$，积分 $\int P\,dx=0$。于是曲线积分仅剩上下两段：
\[\oint_L P\,dx = \int_{L_1} P\,dx + \int_{L_2} P\,dx
= \int_a^b P(x,y_1(x))\,dx + \int_b^a P(x,y_2(x))\,dx\]

另一方面，二重积分可化为先对 $y$ 后对 $x$ 的累次积分：
\[
-\iint_D \frac{\partial P}{\partial y}\,dxdy
= -\int_a^b \left( \int_{y_1(x)}^{y_2(x)} \frac{\partial P}{\partial y}\,dy \right)dx
\]

由牛顿-莱布尼茨公式，内层积分为：
\[
\int_{y_1(x)}^{y_2(x)} \frac{\partial P}{\partial y}\,dy
= P(x,y_2(x)) - P(x,y_1(x))
\]

代入二重积分得：
\[-\iint_D \frac{\partial P}{\partial y}\,dxdy
= \int_a^b P(x,y_1(x))\,dx - \int_a^b P(x,y_2(x))\,dx\]

我们知道\(\oint_L P\,dx = \int_a^b P(x,y_1(x))\,dx + \int_b^a P(x,y_2(x))\,dx\)，即：
\[
\oint_L P\,dx = -\iint_D \frac{\partial P}{\partial y}\,dxdy
\]

同理，对 \(Q\) 那一项，用平行于 \(x\) 轴的直线穿入区域 \(D\)，得到：
\[
\oint_L Q\,dy = \iint_D \frac{\partial Q}{\partial x}\,dx\,dy
\]

两式相加，就是格林公式。

格林公式的左边 \(\oint_L P\,dx+Q\,dy\)，是向量函数 \(\vec{F}=(P,Q)\) 沿闭曲线一圈的积分，我们把它称为\textbf{环流量}。如果 \(\vec{F}\) 是水流，环流量刻画的就是水流绕这条曲线“转”了多少。公式右边的 \(\frac{\partial Q}{\partial x}-\frac{\partial P}{\partial y}\)，刻画的是场在一点附近的旋转程度，我们叫它\textbf{二维旋度}。于是格林公式的物理意义就是\textbf{环流量等于区域内所有点旋度的总和。}

格林公式还有一个特别实用的副产品。在公式里取 \(P=-\frac{y}{2}\)，\(Q=\frac{x}{2}\)，右边的被积函数变成 \(\frac12+\frac12=1\)，于是：
\[
\iint_D dx\,dy = \frac12 \oint_L x\,dy - y\,dx
\]

这就把一块区域的面积，化成了沿它边界的曲线积分。有时候区域的边界形状复杂，直接算面积要费好大劲，但沿边界走一圈的曲线积分反而好算，这个公式就派上了大用场。古代的测绘师发明过一种叫“求积仪”的机械，沿着地块的边界走一圈，表盘上的读数就是面积，靠的就是这个原理。

我们看几个例子感受一下：

\begin{myexample}
    计算 \(\oint_L xy\,dx\)，其中 \(L\) 是抛物线 \(y=x^2\) 与直线 \(y=x\) 所围成的闭区域 \(D\) 的正向边界。

    这个闭区域夹在两条曲线之间，边界由两段组成。\(L_1\) 是抛物线 \(y=x^2\) 上从 \((0,0)\) 到 \((1,1)\) 的一段，\(L_2\) 是直线 \(y=x\) 上从 \((1,1)\) 回到 \((0,0)\) 的一段。

    沿 \(L_1\)，把 \(y=x^2\) 代进去：
    \[
    \int_{L_1} xy\,dx = \int_0^1 x\cdot x^2\,dx = \int_0^1 x^3\,dx = \frac14
    \]

    沿 \(L_2\)，把 \(y=x\) 代进去，注意方向是从 \((1,1)\) 到 \((0,0)\)：
    \[
    \int_{L_2} xy\,dx = \int_1^0 x\cdot x\,dx = -\int_0^1 x^2\,dx = -\frac13
    \]

    两者相加，直接算得：
    \[
    \oint_L xy\,dx = \frac14 - \frac13 = -\frac{1}{12}
    \]

    我们再用格林公式算一下：
    \[
    \oint_L xy\,dx = \iint_D (0-x)\,dx\,dy = -\iint_D x\,dx\,dy
    \]
    
    这里 \(P=xy\)，\(Q=0\)。

    在区域 \(D\) 里，\(x\) 从 \(0\) 到 \(1\)，\(y\) 从 \(x^2\) 到 \(x\)，所以：
    \[
    -\iint_D x\,dx\,dy = -\int_0^1\int_{x^2}^{x} x\,dy\,dx = -\int_0^1 x(x-x^2)\,dx = -\left(\frac13-\frac14\right) = -\frac1{12}
    \]
\end{myexample}

格林公式适合处理的是闭曲线积分，那如果不是闭曲线，可以用格林公式吗？其实是可以的。

\begin{myexample}
    计算 \(\int_C y\,dx + x^2\,dy\)，其中 \(C\) 是抛物线 \(y=1-x^2\) 上从 \((1,0)\) 到 \((0,1)\) 的一段。

    这条曲线不是闭的，缺口就在两个端点之间。我们补上两段。从 \((0,1)\) 沿 \(y\) 轴回到 \((0,0)\)，再从 \((0,0)\) 沿 \(x\) 轴回到 \((1,0)\)。这样一来，原来的抛物线加上这两条轴线段，正好围成第一象限里由 \(y=1-x^2\)、\(x\) 轴、\(y\) 轴围成的闭区域 \(D\)，而且走的方向是正向。设补上的两条线段合起来叫 \(L\)。

    先对闭曲线 \(C+L\) 用格林公式。这里 \(P=y\)，\(Q=x^2\)，于是：
    \[
    \frac{\partial Q}{\partial x}-\frac{\partial P}{\partial y} = 2x-1
    \]

    在区域 \(D\) 里，\(x\) 从 \(0\) 到 \(1\)，\(y\) 从 \(0\) 到 \(1-x^2\)，所以：
    \[
    \oint_{C+L} P\,dx+Q\,dy = \iint_D (2x-1)\,dx\,dy = \int_0^1\int_0^{1-x^2} (2x-1)\,dy\,dx = \int_0^1 (2x-1)(1-x^2)\,dx = -\frac16
    \]

    再看补上的那两段。沿 \(y\) 轴的一段上 \(x=0\)，被积表达式只剩 \(Q\,dy=x^2\,dy=0\)。沿 \(x\) 轴的一段上 \(y=0\)，只剩 \(P\,dx=y\,dx=0\)。所以：
    \[
    \int_L y\,dx + x^2\,dy = 0
    \]

    于是：
    \[
    \int_C y\,dx + x^2\,dy = \oint_{C+L} P\,dx+Q\,dy - 0 = -\frac16
    \]
\end{myexample}

这个例子说明了，如果曲线不是闭曲线，我们可以用“补线法”把曲线补成闭曲线，再用格林公式，最后记得减去补线的积分就可以了。

\begin{myexample}
计算\(\oint_L \frac{y\,dx - x\,dy}{2(x^2 + y^2)}\)，其中 \(L\) 是圆周 \((x-1)^2 + y^2 = 2\) 的正向。
 
令
\(P = \frac{y}{2(x^2 + y^2)}\)，\(Q = -\frac{x}{2(x^2 + y^2)}\)，当 \((x,y) \neq (0,0)\) 时，计算偏导数：

\[
\frac{\partial Q}{\partial x} = \frac{x^2 - y^2}{2(x^2 + y^2)^2},\qquad
\frac{\partial P}{\partial y} = \frac{x^2 - y^2}{2(x^2 + y^2)^2}
\]

所以：
\[
\frac{\partial Q}{\partial x} - \frac{\partial P}{\partial y} = 0
\]

但闭曲线 \(L\) 所围区域包含原点 \((0,0)\)，而 \(P,Q\) 在原点无定义，格林公式要求函数 \(P(x,y)\)、\(Q(x,y)\) 在 \(D\) 上具有一阶连续偏导数，这不满足使用格林公式的条件。我们取充分小的正数 \(\epsilon\)，作以原点为圆心，\(\epsilon\) 为半径的圆周 \(L_\epsilon\)，使其完全包含在 \(L\) 内部。记 \(L\) 与 \(L_\epsilon\) 之间的区域为 \(D'\)，在 \(D'\) 上可用格林公式：  
\[
\oint_{L \cup L_\epsilon} P\,dx + Q\,dy = \iint_{D'} \left(\frac{\partial Q}{\partial x} - \frac{\partial P}{\partial y}\right) dx\,dy = 0
\]

其中\(L \cup L_\epsilon\)表示先沿\(L\)逆时针走一圈，再沿\(L_\epsilon\)顺时针走一圈。

因此：
\[
\oint_L P\,dx + Q\,dy = -\oint_{L_\epsilon} P\,dx + Q\,dy = \oint_{L_\epsilon^+} \frac{y\,dx - x\,dy}{2(x^2 + y^2)}
\]

其中 \(L_\epsilon^+\) 表示以原点为心、半径为 \(\epsilon\) 的逆时针圆周。

在 \(L_\epsilon^+\) 上，参数化为 \(x = \epsilon\cos\theta\)，\(y = \epsilon\sin\theta\)，\(\theta\in[0,2\pi]\)，则：
\[
y\,dx - x\,dy = \epsilon\sin\theta\,(-\epsilon\sin\theta\,d\theta) - \epsilon\cos\theta\,(\epsilon\cos\theta\,d\theta) = -\epsilon^2\,d\theta
\]

分母 \(2(x^2 + y^2) = 2\epsilon^2\)，于是：
\[
\frac{y\,dx - x\,dy}{2(x^2 + y^2)} = -\frac12\,d\theta
\]

这里\(\epsilon\)被消去了。这说明这个积分跟小圆的半径无关。所以无论我们挖的洞多小，结果都一样。代入得：
\[
\oint_{L_\epsilon^+} = \int_0^{2\pi} \left(-\frac12\right) d\theta = -\pi
\]

所以\(\oint_L \frac{y\,dx - x\,dy}{2(x^2 + y^2)} = -\pi
\)。

\end{myexample}

这种方法叫“挖洞法”。遇到不能使用格林公式的区域，就先挖去，后面再补上。

总结一下，格林公式把平面闭曲线上的第二类曲线积分，化成了曲线所围区域上的二重积分，它是牛顿-莱布尼茨公式的二维版本，体现的正是“内部的变化由边界决定”这个思想。用的时候要注意，边界要取正向，也就是区域始终在左手边的逆时针方向，方向反了，结果就要变号。曲线不是闭曲线时可以用“补线法”，先补上一条好算的曲线再用格林公式，最后记得把补上的那段的积分减掉。如果 \(P\)，\(Q\) 在区域内某点没有定义，格林公式就用不了，这时可以用“挖洞法”，先挖去该点周围的小圆，再在小圆与边界围成的区域上用格林公式，而小圆上的积分往往跟圆的大小无关。


\subsection{与路径无关的第二类曲线积分}

回忆第二类曲线积分的物理意义，它计算的是变力做的功。然而有些力做功不考虑路径，比如重力，把物体从 \(A\) 点搬到 \(B\) 点，无论走哪条路，重力做的功都一样，只跟 \(A\)、\(B\) 两点的高度有关。从积分角度来说，这就是第二类曲线积分与路径无关。那么在平面上，什么样的第二类曲线积分与路径无关呢？

先把问题转化一下。第二类曲线积分在区域里与路径无关，等价于沿区域里任何一条闭曲线的积分为零。因为如果两条不同的路径 \(L_1\)，\(L_2\) 从 \(A\) 到 \(B\) 积分值相同，把它们首尾相接拼成一条闭曲线，沿着\(L_1\)从 \(A\) 到 \(B\) 积分，再沿着\(L_2\)从 \(B\) 到 \(A\) 积分，两者相加正好得零。反过来，如果沿任意闭曲线积分为零，那么任意两条同起点同终点的路径正好围出一条闭曲线，积分当然也为零。

这样一来，问题就变成了，什么时候沿任意闭曲线的积分都为零？格林公式告诉我们，沿闭曲线 \(L\) 的积分等于它围成区域上的 \(\iint_D\left(\frac{\partial Q}{\partial x}-\frac{\partial P}{\partial y}\right)dx\,dy\)。如果区域里处处有 \(\frac{\partial Q}{\partial x}=\frac{\partial P}{\partial y}\)，这样二重积分的被积函数为零，那么这个第二类曲线积分自然恒为零。

有了初步的认识，我们现在先介绍单连通区域：
\begin{definition}{单连通区域}{}
    设 $D$ 是一个在平面上的区域，如果对于 $D$ 内的任意一条简单闭曲线 $L$，$L$ 所围成的内部区域都完全包含于 $D$，则称 $D$ 为\textbf{单连通区域}。
\end{definition}

直观上理解，就是说在平面上的一个区域，区域内任意一条简单闭曲线，都可以在区域内连续地收缩成一个点，且收缩过程中始终不超出该区域，那么这个区域就是单连通区域。更形象的说，“没有洞”的区域就是单连通区域。因为“有洞”的区域可以在区域内找到一条包围洞的闭曲线，它的内部不在区域内，直接违反了单连通定义中“所有闭曲线的内部都必须在区域内”的要求。

为什么要强调单连通？我们来看一个例子。

设在挖去原点的平面上有一个向量函数\(\vec{F} = \left(-\frac{y}{x^2+y^2},\ \frac{x}{x^2+y^2}\right)\)，闭曲线 \(L:x^2+y^2=1\)，除原点外，处处有 \(\frac{\partial Q}{\partial x}=\frac{\partial P}{\partial y}\)。令 \(x=\cos t\)，\(y=\sin t\)，分母恒为 \(1\)，沿\(L\)逆时针积分：
\[
\oint_L \vec{F}\cdot d\vec{r} = \oint_L P\,dx+Q\,dy = \int_0^{2\pi} ((-\sin t)(-\sin t) + \cos t\cos t)dt = \int_0^{2\pi} dt = 2\pi
\]

然而刚刚我们知道，如果\(\frac{\partial Q}{\partial x}=\frac{\partial P}{\partial y}\)，那么积分与路径无关，所以结果应该是\(0\)。也就是：
\[
\oint_L \vec{F}\cdot d\vec{r} = \oint_L P\,dx+Q\,dy = \iint_D\left(\frac{\partial Q}{\partial x}-\frac{\partial P}{\partial y}\right)dx\,dy = \iint_D 0 \,dx\,dy = 0
\]

这正是因为我们去掉了原点，导致区域不再是单连通区域，因此条件 \(\frac{\partial Q}{\partial x}=\frac{\partial P}{\partial y}\)不再适用。

把上面的讨论整理一下，就得到第二类曲线积分与路径无关的充要条件：

\begin{theorem}{第二类曲线积分与路径无关的充要条件}{}
设 \(D\) 是平面上的单连通区域，函数 \(P(x,y)\)、\(Q(x,y)\) 在 \(D\) 内具有一阶连续偏导数，则下列四个条件等价：
\[
\begin{aligned}
&\text{(1)}\ \text{沿 }D\text{ 内任意闭曲线 }L\text{，}\oint_L P\,dx+Q\,dy = 0 \\[6pt]
&\text{(2)}\ \text{曲线积分 }\int_L P\,dx+Q\,dy\text{ 在 }D\text{ 内与路径无关，只与起点和终点有关} \\[6pt]
&\text{(3)}\ \text{存在可微函数 }u(x,y)\text{，使得 }du = P\,dx+Q\,dy\text{，即 }\frac{\partial u}{\partial x}=P\text{，}\frac{\partial u}{\partial y}=Q \\[6pt]
&\text{(4)}\ \text{在 }D\text{ 内处处有 }\frac{\partial Q}{\partial x} = \frac{\partial P}{\partial y}
\end{aligned}
\]
这四个条件是第二类曲线积分与路径无关的充要条件。
\end{theorem}

一旦找到函数 \(u(x,y)\)，曲线积分就变成两端点上函数之差，这正是牛顿-莱布尼茨公式的二维翻版：
\[
\int_A^B P\,dx + Q\,dy = u(B) - u(A)
\]

\begin{myexample}
    计算 \(\int_L (3x^2y+1)\,dx + x^3\,dy\)，其中 \(L\) 是从 \((0,0)\) 到 \((1,2)\) 的直线\(y=2x\)。

    这里 \(P=3x^2y+1\)，\(Q=x^3\)，于是：
    \[
    \frac{\partial P}{\partial y} = 3x^2,\qquad \frac{\partial Q}{\partial x} = 3x^2
    \]

    偏导数相等，所以积分与路径无关。既然积分与路径无关，我们当然挑一条最好走的路。先从 \((0,0)\) 水平走到 \((1,0)\)，再竖直走到 \((1,2)\)。

    第一段上 \(y=0\)，\(dy=0\)，被积表达式只剩 \(P\,dx=(0+1)dx=dx\)，所以：
    \[
    \int_{(0,0)}^{(1,0)} dx = 1
    \]

    第二段上 \(x=1\)，\(dx=0\)，被积表达式只剩 \(Q\,dy=1^3\,dy=dy\)，所以：
    \[
    \int_{(1,0)}^{(1,2)} dy = 2
    \]

    加起来就是：
    \[
    \int_L (3x^2y+1)\,dx + x^3\,dy = 1+2 = 3
    \]
\end{myexample}

\begin{myexample}
    计算 \(\int_L 2xy\,dx + x^2\,dy\)，其中 \(L\) 是从 \((0,0)\) 到 \((2,1)\) 的抛物线 \(y=\frac{x^2}{4}\)。

    算偏导：\(\frac{\partial P}{\partial y}=2x\)，\(\frac{\partial Q}{\partial x}=2x\)。所以积分与路径无关。

    我们沿抛物线 \(y=\frac{x^2}{4}\) 走，\(dy=\frac{x}{2}\,dx\)，于是：
    \[
    \int_L 2xy\,dx + x^2\,dy = \int_0^2 \left(2x\cdot\frac{x^2}{4} + x^2\cdot\frac{x}{2}\right)dx = \int_0^2 \left(\frac{x^3}{2}+\frac{x^3}{2}\right)dx = \int_0^2 x^3\,dx = 4
    \]

    积分与路径无关，我们也可以沿折线 \((0,0)\to(2,0)\to(2,1)\) 走，沿\((0,0)\to(2,0)\)，\(y=0\)，\(dy=0\)，只剩：
    \[
    \int_0^2 2x \cdot0\,dx + x^2 \cdot 0=0 
    \]
    
    沿\((2,0)\to(2,1)\)，\(x=2\)，\(dx=0\)，只剩：
    \[
    2 \cdot 2y \cdot 0 + \int_0^1 2^2 \, dy=4
    \]

    相加得：
    \[
    \int_L 2xy\,dx + x^2\,dy = 0 + 4 =4
    \]
\end{myexample}

\begin{myexample}
    计算 \(\int_L (2xy-y^2)\,dx + (x^2-2xy)\,dy\) 其中\(L\)是 \((0,0)\) 到 \((2,1)\) 的曲线\(y=\sin(\frac{\pi}{2}\sin(\frac{\pi x}{4}))\)。

    这里 \(P=2xy-y^2\)，\(Q=x^2-2xy\)。求偏导：
    \[
    \frac{\partial P}{\partial y} = 2x-2y,\qquad \frac{\partial Q}{\partial x} = 2x-2y
    \]

    偏导数相等，所以积分与路径无关。所以我们可以完全不用理会原来复杂的三角函数积分路径，重新选择一条简单的折线路径来代替原曲线。

    取路径\(C_1\)，从 \((0,0)\) 到 \((2,0)\)，\(y=0\)，\(dy=0\)，此时原积分为：
    \[
\int_{C_1} (2xy-y^2)\,dx + (x^2-2xy)\,dy = \int_0^2 (2x\cdot0 - 0^2)\,dx + (x^2-0)\cdot0 = 0
\]

    取路径\(C_2\),从 \((2,0)\) 到 \((2,1)\)，\(x=2\)，\(dx=0\)，此时原积分为：
    \[
    \int_{C_2} (2xy-y^2)\,dx + (x^2-2xy)\,dy = \int_0^1 (2^2 - 2\cdot2\cdot y)\,dy = \int_0^1 (4-4y)\,dy = 2
    \]

相加得：
\[\int_L (2xy-y^2)\,dx + (x^2-2xy)\,dy=\int_{C_1}+\int_{C_2}=0+2=2\]

\end{myexample}

总结一下，格林公式把平面闭曲线上的第二类曲线积分，化成了它所围区域上的二重积分。它是牛顿-莱布尼茨公式从“一段区间”到“一块区域”的推广。

\section{高斯公式}
格林公式把平面闭曲线上的积分转化成了二重积分，那么，在空间中，闭曲面上的第二类曲面积分，能不能类似地化成其他积分？答案是可以的。那就是高斯公式。
\subsection{高斯公式}
高斯公式说的是空间闭区域上的三重积分与其边界曲面上的曲面积分之间的关系。和格林公式一样，它也是“内部的变化由边界决定”思想的体现。

\begin{theorem}{高斯公式}{}
设空间闭区域 \(\Omega\) 由分片光滑的闭曲面 \(\Sigma\) 围成，函数 \(P(x,y,z)\)，\(Q(x,y,z)\)，\(R(x,y,z)\) 在 \(\Omega\) 上具有一阶连续偏导数，则
\[
\oiint_{\Sigma} P\,dy\,dz + Q\,dz\,dx + R\,dx\,dy = \iiint_{\Omega} \left( \frac{\partial P}{\partial x} + \frac{\partial Q}{\partial y} + \frac{\partial R}{\partial z} \right) dV
\]
其中 \(\Sigma\) 取外侧。
\end{theorem}

高斯公式的物理意义非常直观。我们已经知道，第二类曲面积分表示的是向量函数穿过有向曲面的通量。而被积函数 \(\frac{\partial P}{\partial x}+\frac{\partial Q}{\partial y}+\frac{\partial R}{\partial z}\) 称为向量函数 \(\vec{F}\) 的\textbf{散度}，记作 \(\operatorname{div}\vec{F}\)，所以高斯公式讲的就是\textbf{区域内部产生的量等于从该区域边界流出的量}。

接下来我们通过几个例子感受一下：

\begin{myexample}
计算 \(\oiint_{\Sigma} x\,dy\,dz + y\,dz\,dx + z\,dx\,dy\)，其中 \(\Sigma\) 是由上半球面 \(x^2+y^2+z^2=1,\ z\ge 0\) 与底面 \(z=0,\ x^2+y^2\le 1\) 所围成的封闭曲面，取外侧。

这里 \(P=x\)，\(Q=y\)，\(R=z\)。偏导数分别是
\[
\frac{\partial P}{\partial x}=1,\quad \frac{\partial Q}{\partial y}=1,\quad \frac{\partial R}{\partial z}=1
\]

代入高斯公式，得：
\[
\oiint_{\Sigma} x\,dy\,dz + y\,dz\,dx + z\,dx\,dy = \iiint_{\Omega} (1+1+1)\,dV = 3 \iiint_{\Omega} dV
\]

因为\(\Omega\) 是上半球体，体积为 \(\frac{2}{3}\pi\)，因此：
\[
3 \iiint_{\Omega} dV = 3 \times \frac{2}{3}\pi = 2\pi
\]

所以 \(\oiint_{\Sigma} x\,dy\,dz + y\,dz\,dx + z\,dx\,dy = 2\pi\)。
\end{myexample}

如果把这个闭曲面分成上半球面和底面分别计算，计算量会大很多，用高斯公式一步就完成了。注意这里一定要把底面算进去，因为高斯公式要求积分曲面是封闭的。少了底面，就不再是闭曲面，也就不能直接套用高斯公式了。

和格林公式一样，如果区域不是闭合曲面，我们可以用“补面法”补成闭曲面再用高斯公式计算，最后再减去补面的积分就可以了。

\begin{myexample}
    计算\(\iint_{\Sigma} \frac{x\,dy\,dz + y\,dz\,dx + z\,dx\,dy}{(x^2+y^2+z^2)^{\frac32}}\)，其中曲面 $\Sigma$ 是上半锥面 $z=\sqrt{x^2+y^2}$ 被平面 $z=1$ 与 $z=2$ 截出的部分，取上侧。

    $\Sigma$本身不是封闭的，可以看成是一段开口的“喇叭筒”，上下各有一个圆口。我们用两个水平圆盘把它封住，上底面 $\Sigma_1: z=2,\ x^2+y^2\le 4$，取上侧，下底面 $\Sigma_2: z=1,\ x^2+y^2\le 1$，取下侧。记闭曲面 $\Sigma' = \Sigma \cup \Sigma_1 \cup \Sigma_2$，$\Sigma'$ 取上侧，包围的空间区域记为 $\Omega$。$\Omega$ 是圆台体，上底半径 $2$，下底半径 $1$，高 $1$。

    不妨设：
\[
P=\frac{x}{(x^2+y^2+z^2)^{\frac32}},\quad
Q=\frac{y}{(x^2+y^2+z^2)^{\frac32}},\quad
R=\frac{z}{(x^2+y^2+z^2)^{\frac32}}
\]

求偏导：
\[
\frac{\partial P}{\partial x} = \frac{(y^2+z^2-2x^2)}{(x^2+y^2+z^2)^{\frac52}},\qquad
\frac{\partial Q}{\partial y} = \frac{(x^2+z^2-2y^2)}{(x^2+y^2+z^2)^{\frac52}},\qquad
\frac{\partial R}{\partial z} = \frac{(x^2+y^2-2z^2)}{(x^2+y^2+z^2)^{\frac52}}
\]

三个偏导相加：
\[
\frac{\partial P}{\partial x} + \frac{\partial Q}{\partial y} + \frac{\partial R}{\partial z} = 0
\]

由高斯公式：
\[
\oiint_{\Sigma'} P\,dy\,dz + Q\,dz\,dx + R\,dx\,dy = \iiint_{\Omega} 0\,dV = 0
\]

我们求的是原来\(\Sigma\)的积分，所以要减去补面的积分，即：
\[
\iint_{\Sigma} = 0 -\iint_{\Sigma_1} - \iint_{\Sigma_2} = -(\iint_{\Sigma_1} + \iint_{\Sigma_2})
\]

先算 $\Sigma_1$ 上的积分。$\Sigma_1$ 的方程为 $z=2$，法向量向上，故 $dy\,dz=0,\ dz\,dx=0$，只留 $dx\,dy$ 项，所以：
\[
\iint_{\Sigma_1} \frac{z\,dx\,dy}{(x^2+y^2+z^2)^{\frac32}}
= \iint_{D_1} \frac{2\,dx\,dy}{(x^2+y^2+4)^{\frac32}}
\]

其中 $D_1: x^2+y^2\le 4$。化为极坐标，即：
\[
\iint_{D_1} \frac{2}{(\rho^2+4)^{\frac32}}\cdot \rho\,d\rho\,d\theta
= 2\pi \int_0^2 \frac{2\rho}{(\rho^2+4)^{\frac32}}\,d\rho
\]

令 $u=\rho^2+4$，$du=2\rho\,d\rho$，上下限 $u$ 从 $4$ 到 $8$，原积分化为：
\[2\pi \int_0^2 \frac{2\rho}{(\rho^2+4)^{\frac32}}\,d\rho
= 2\pi \int_4^8 u^{-\frac32}\,du
= 2\pi \cdot \left[ -2u^{-\frac12} \right]_4^8
= 2\pi ( 1 - \frac{1}{\sqrt{2}} )
\]

再算 $\Sigma_2$ 上的积分。$\Sigma_2$ 的方程为 $z=1$，取下侧，$dx\,dy$ 项取负号。即：
\[
\iint_{\Sigma_2} \frac{z\,dx\,dy}{(x^2+y^2+z^2)^{\frac{3}{2}}}
= -\iint_{D_2} \frac{1\,dx\,dy}{(x^2+y^2+1)^{\frac{3}{2}}}
\]

其中 $D_2: x^2+y^2\le 1$。化为极坐标，即：
\[
- \int_0^{2\pi} d\theta \int_0^1 \frac{\rho}{(\rho^2+1)^{\frac{3}{2}}}\,d\rho
= -2\pi \int_0^1 \frac{\rho}{(\rho^2+1)^{\frac{3}{2}}}\,d\rho
\]

令 $u=\rho^2+1$，$du=2\rho\,d\rho$，$u$ 从 $1$ 到 $2$，原积分化为：
\[-2\pi \int_0^1 \frac{\rho}{(\rho^2+1)^{\frac{3}{2}}}\,d\rho
= -2\pi \int_1^2 \frac{1}{2}u^{-\frac{3}{2}}\,du
= -\pi \cdot \left[ -2u^{-\frac{1}{2}} \right]_1^2
= -\pi (2 - \sqrt{2})
\]

代入得：
\[
-(\iint_{\Sigma_1} + \iint_{\Sigma_2})=-(2\pi ( 1 - \frac{1}{\sqrt{2}} )-\pi (2 - \sqrt{2}))=0\]

所以\(\iint_{\Sigma} \frac{x\,dy\,dz + y\,dz\,dx + z\,dx\,dy}{(x^2+y^2+z^2)^{\frac32}}=0\)。
\end{myexample}

同样的，我们也有“挖洞法”：

\begin{myexample}
    计算\(\oiint_{\Sigma} \frac{x\,dy\,dz + y\,dz\,dx + z\,dx\,dy}{(x^2+y^2+z^2)^{\frac32}}\)，其中 $\Sigma$ 是球面 $x^2+y^2+z^2=1$，取外侧。

    被积函数在原点 $(0,0,0)$ 处无定义，而原点恰好位于 $\Sigma$ 所围球体的内部，以原点为圆心作充分小的球面 $\Gamma_\epsilon: x^2+y^2+z^2=\epsilon^2$，取内侧，它与 $\Sigma$ 一起围成一个不包含原点的闭区域，记为 $\Omega_\epsilon$。将 $\Gamma_\epsilon$ 取内侧时记作 $\Gamma_\epsilon^+$，取外侧时记作 $\Gamma_\epsilon^-$

    被积函数和上个例子一样，所以由高斯公式：
    \[
    \oiint_{\Sigma} P\,dy\,dz + Q\,dz\,dx + R\,dx\,dy
    + \oiint_{\Gamma_\epsilon^+} P\,dy\,dz + Q\,dz\,dx + R\,dx\,dy
    = \iiint_{\Omega_\epsilon} 0\,dV = 0
    \]

    因此：
    \[
\begin{aligned}
        \oiint_{\Sigma} P\,dy\,dz + Q\,dz\,dx + R\,dx\,dy
        &= - \oiint_{\Gamma_\epsilon^+} P\,dy\,dz + Q\,dz\,dx + R\,dx\,dy\\
        &= \oiint_{\Gamma_\epsilon^-} P\,dy\,dz + Q\,dz\,dx + R\,dx\,dy
\end{aligned}
    \]

    计算 $\Gamma_\epsilon^-$ 上的积分。在 $\Gamma_\epsilon$上，$x^2+y^2+z^2=\epsilon^2$，面积元投影满足
    $dy\,dz = \frac{x}{\epsilon}\,dS,\ dz\,dx = \frac{y}{\epsilon}\,dS,\ dx\,dy = \frac{z}{\epsilon}\,dS$，所以：
    \[
    \frac{x\,dy\,dz + y\,dz\,dx + z\,dx\,dy}{(x^2+y^2+z^2)^{\frac32}}
    = \frac{x\cdot\frac{x}{\epsilon} + y\cdot\frac{y}{\epsilon} + z\cdot\frac{z}{\epsilon}}{\epsilon^3}\,dS
    = \frac{1}{\epsilon^2}\,dS
    \]

    于是：
    \[
    \oiint_{\Gamma_\epsilon^-} \frac{1}{\epsilon^2}\,dS
    = \frac{1}{\epsilon^2}\cdot 4\pi\epsilon^2 = 4\pi
    \]

    所以\(\oiint_{\Sigma} \frac{x\,dy\,dz + y\,dz\,dx + z\,dx\,dy}{(x^2+y^2+z^2)^{\frac32}} = 4\pi\)。
\end{myexample}

\subsection{通量与散度}

我们已经知道，第二类曲面积分计算的是向量函数穿过有向曲面的通量。现在我们详细介绍一下。

\begin{definition}{通量}
设向量函数 \(\vec{F}(x,y,z)\) 定义在空间区域上，\(\Sigma\) 是一张有向光滑曲面，则 \(\vec{F}\) 穿过 \(\Sigma\) 的\textbf{通量}为
\[
\Phi = \iint_{\Sigma} \vec{F} \cdot \vec{n}\,dS = \iint_{\Sigma} P\,dy\,dz + Q\,dz\,dx + R\,dx\,dy
\]
其中 \(\vec{n}\) 是 \(\Sigma\) 上选定一侧的单位法向量。
\end{definition}

通俗地讲，如果把 \(\vec{F}\) 看作水流的速度，通量就是单位时间内水流穿过这张曲面的水量。如果 \(\vec{F}\) 是电场强度，通量就是电通量。总之，通量刻画的是某种“流”穿过一张曲面的总量，穿出为正，穿入为负。

那么，在这张曲面所包围的区域内部，到底是“往外冒水”的多，还是“往里吸水”的多？这就引出了散度的概念。

\begin{definition}{散度}
设向量函数 \(\vec{F}(x,y,z) = (P,Q,R)\) 在点 \((x,y,z)\) 处具有一阶连续偏导数，则称
\[
\operatorname{div} \vec{F} = \frac{\partial P}{\partial x} + \frac{\partial Q}{\partial y} + \frac{\partial R}{\partial z}
\]
为向量函数 \(\vec{F}\) 在点 \((x,y,z)\) 处的\textbf{散度}。
\end{definition}

散度是一个标量。如果在某一点 \(\operatorname{div}\vec{F}>0\)，说明这一点是“\textbf{源}”，向量从这一点向外发散，好比水龙头往外冒水。如果 \(\operatorname{div}\vec{F}<0\)，说明这一点是“\textbf{汇}”，向量向这一点汇聚，好比下水道口往里吸水。如果 \(\operatorname{div}\vec{F}=0\)，则说明这一点既没有发散也没有汇聚。

有了散度的概念，高斯公式就可以写成简洁的向量形式：
\[
\oiint_{\Sigma} \vec{F} \cdot \vec{n}\,dS = \iiint_{\Omega} \operatorname{div} \vec{F}\,dV
\]

这个形式表达的就是高斯公式的物理意义：\textbf{闭曲面上的通量等于曲面所围区域内各点散度的总和。}所以高斯公式又称散度定理。比如你要统计一条马路一小时内的净流出车流，你可以守在这条路口数通过了多少车，也可以数这条路上某时刻有多少车，两者结果是相等的。

我们看一个散度的例子：
\begin{myexample}
求 \(\vec{F} = (x^2y,\, y^2z,\, z^2x)\) 在点 \((1,1,1)\) 处的散度。

这里 \(P=x^2y\)，\(Q=y^2z\)，\(R=z^2x\)。先求偏导：
\[
\frac{\partial P}{\partial x}=2xy,\qquad
\frac{\partial Q}{\partial y}=2yz,\qquad
\frac{\partial R}{\partial z}=2zx
\]

所以散度为：
\[
\operatorname{div} \vec{F} = 2xy + 2yz + 2zx = 2(xy+yz+zx)
\]

代入点 \((1,1,1)\)，得：
\[
\operatorname{div} \vec{F}(1,1,1) = 2(1+1+1) = 6
\]
\end{myexample}

我们看一下如何判断是发散还是汇聚。
\begin{myexample}
设 \(\vec{F} = (3x^2, -2y, z^3)\)，求 \(\operatorname{div} \vec{F}\)，并判断原点 \((0,0,0)\) 是源还是汇。

求偏导：
\[
\frac{\partial P}{\partial x}=6x,\quad 
\frac{\partial Q}{\partial y}=-2,\quad 
\frac{\partial R}{\partial z}=3z^2
\]

于是：
\[
\operatorname{div} \vec{F} = 6x - 2 + 3z^2
\]

在原点处，\(x=0\)，\(z=0\)，所以 \(\operatorname{div} \vec{F}(0,0,0) = -2 < 0\)，因此原点是汇。
\end{myexample}

\subsection{与曲面无关的第二类曲面积分}

前面我们看到，第二类曲线积分在某种条件下与路径无关。那么，在某种条件下，曲面积分的值是否可以与曲面无关呢？也就是说，曲面积分的值是否可以只由曲面的边界曲线决定？

我们不妨先设想两片光滑曲面 \(\Sigma_1\) 与 \(\Sigma_2\)，它们有同一条边界曲线 \(\Gamma\)，并且彼此不交错。把 \(\Sigma_2\) 反向之后与 \(\Sigma_1\) 拼在一起，就构成一张分片光滑的闭曲面，记作 \(\Sigma\)。如果向量函数 \(\vec{F}=(P,Q,R)\) 的散度在 \(\Sigma\) 所包围的区域上处处为零，那么由高斯公式立即得到：
\[
\oiint_{\Sigma} \vec{F}\cdot\vec{n}\,dS = \iiint_{\Omega} 0\,dV = 0
\]

由于 \(\Sigma\) 是由 \(\Sigma_1\) 与反向的 \(\Sigma_2\) 拼成的，上式意味着：
\[
\iint_{\Sigma_1} \vec{F}\cdot\vec{n}\,dS
=
\iint_{\Sigma_2} \vec{F}\cdot\vec{n}\,dS
\]

也就是说，只要散度处处为零，穿过同一边界曲线的两张曲面的通量就相等，曲面积分不再依赖于曲面的具体形状。

不过，上面的推理还隐含了一个要求，由 \(\Sigma_1\) 和反向 \(\Sigma_2\) 拼成的闭曲面所包围的区域，必须完全落在向量函数有定义且满足散度为零的区域内。如果区域中有“空洞”，这个条件就可能被破坏。因此，需要引入空间二维单连通区域的概念。

\begin{definition}{空间二维单连通区域}{}
设 \(\Omega\) 是空间中的一个区域。如果对于 \(\Omega\) 内任意一张分片光滑的简单闭曲面 \(\Sigma\)，\(\Sigma\) 所包围的空间区域都完全包含于 \(\Omega\)，则称 \(\Omega\) 为\textbf{空间二维单连通区域}。
\end{definition}

和单连通区域一样，如果在\(\Omega\)内任意一个简单闭合曲面都能在\(\Omega\)内部连续收缩为一个点，且始终不离开\(\Omega\)，那么\(\Omega\)就是二维单连通区域。直观地说，空间二维单连通区域就是“没有空洞”的区域。例如实心球体，实心立方体等都是空间二维单连通的。而一个实心球体内部挖去一个小球后，由于出现了空腔，就不再是空间二维单连通的了。在讨论曲面积分与曲面无关时，这一条件保证了我们可以对两个同边界的曲面拼成的闭曲面所围区域使用高斯公式。

总结一下就得到：

\begin{theorem}{第二类曲面积分与曲面无关的条件}{}
设 \(\Omega\) 是空间二维单连通区域，函数 \(P(x,y,z)\)、\(Q(x,y,z)\)、\(R(x,y,z)\) 在 \(\Omega\) 内具有一阶连续偏导数。则第二类曲面积分
\[
\iint_{\Sigma} P\,dy\,dz + Q\,dz\,dx + R\,dx\,dy
\]
在 \(\Omega\) 内与曲面 \(\Sigma\) 无关的\textbf{充分必要条件}是在\(\Omega\) 内处处散度为\(0\)，即在 \(\Omega\) 内处处有
\[
\frac{\partial P}{\partial x} + \frac{\partial Q}{\partial y} + \frac{\partial R}{\partial z} = 0
\]
\end{theorem}

这个定理说明，散度为零的\(\vec{F}=(P,Q,R)\)中间没有新的“源”或“汇”，流入多少就流出多少。

\begin{myexample}
计算 \(\iint_{\Sigma} 2xz\,dy\,dz + 2yz\,dz\,dx + (x^2+y^2-2z^2)\,dx\,dy\)，其中 \(\Sigma\) 是上半球面 \(z=\sqrt{1-x^2-y^2}\) 的上侧。

记 \(P=2xz\)，\(Q=2yz\)，\(R=x^2+y^2-2z^2\)，求散度：
\[
\frac{\partial P}{\partial x}+\frac{\partial Q}{\partial y}+\frac{\partial R}{\partial z}
= 2z+2z-4z = 0
\]

整个空间是空间二维单连通区域，因此该曲面积分与曲面无关，只依赖于边界曲线。\(\Sigma\) 的边界是圆周 \(x^2+y^2=1,\ z=0\)。我们取同边界的平面圆盘 \(D:\ x^2+y^2\le 1,\ z=0\)，取上侧。在 \(D\) 上，\(z=0\)，法向量为 \((0,0,1)\)，于是：
\[
dy\,dz=0,\qquad dz\,dx=0,\qquad dx\,dy=dS
\]

被积表达式化为：
\[
2xz\,dy\,dz + 2yz\,dz\,dx + (x^2+y^2-2z^2)\,dx\,dy
= (x^2+y^2)\,dx\,dy
\]

所以：
\[
\iint_{\Sigma} 2xz\,dy\,dz + 2yz\,dz\,dx + (x^2+y^2-2z^2)\,dx\,dy
= \iint_{D} (x^2+y^2)\,dx\,dy
\]

用极坐标计算：
\[
\iint_{D} (x^2+y^2)\,dx\,dy
= \int_0^{2\pi} d\theta \int_0^1 \rho^2\cdot \rho\,d\rho
= 2\pi \cdot \frac14 = \frac{\pi}{2}
\]

所以\(\iint_{\Sigma} 2xz\,dy\,dz + 2yz\,dz\,dx + (x^2+y^2-2z^2)\,dx\,dy = \frac{\pi}{2}\)。
\end{myexample}
如果直接计算上半球面上的积分会比较麻烦，但利用积分与曲面无关的性质，我们只需换成一个容易计算的曲面，就得到了结果。

\section{斯托克斯公式}
前面格林公式告诉我们，平面闭曲线上的第二类曲线积分，可以化成它所围区域上的二重积分。现在我们把这个问题推广到空间中。空间中的闭曲线本身不能直接围出一块平面区域，但它可以作为某张曲面的边界。那么，沿一张曲面边界闭曲线上的第二类曲线积分，能不能化成这张曲面上的第二类曲面积分？这就是斯托克斯公式。
\subsection{斯托克斯公式}
斯托克斯公式把空间闭曲线上的曲线积分化为曲面积分。在介绍它之前，我们先讲一下右手法则。
\begin{definition}{右手法则}{}
设 $\Gamma$ 是空间中的分段光滑有向闭曲线，$\Sigma$ 是以 $\Gamma$ 为边界的分片光滑有向曲面。若右手四指沿 $\Gamma$ 的正向弯曲时，拇指所指的方向与 $\Sigma$ 上选定的法向量指向一致，则称 $\Gamma$ 的正向与 $\Sigma$ 的一侧符合\textbf{右手法则}。
\end{definition}

直观地说，当你用右手握住边界曲线，手指顺着曲线方向弯曲，拇指所指的方向就是曲面正法向的方向。右手法则明确了空间中有向曲线与有向曲面之间的方向关系。

现在我们就能给出斯托克斯公式了。

\begin{theorem}{斯托克斯公式}{}
设 $\Gamma$ 为空间中的分段光滑有向闭曲线，$\Sigma$ 是以 $\Gamma$ 为边界的分片光滑有向曲面，$\Gamma$ 的正向与 $\Sigma$ 的一侧符合右手法则。若函数 $P(x,y,z)$，$Q(x,y,z)$，$R(x,y,z)$ 在包含 $\Sigma$ 的空间区域上具有一阶连续偏导数，则
\[
\begin{aligned}
    &\oint_{\Gamma} P\,dx + Q\,dy + R\,dz
    =
    \begin{vmatrix}
    dy\,dz & dz\,dx & dx\,dy \\[4pt]
    \dfrac{\partial}{\partial x} & \dfrac{\partial}{\partial y} & \dfrac{\partial}{\partial z} \\[6pt]
    P & Q & R
    \end{vmatrix}\\
    &=
    \iint_{\Sigma}
    \left(\frac{\partial R}{\partial y}-\frac{\partial Q}{\partial z}\right)dy\,dz
    +\left(\frac{\partial P}{\partial z}-\frac{\partial R}{\partial x}\right)dz\,dx
    +\left(\frac{\partial Q}{\partial x}-\frac{\partial P}{\partial y}\right)dx\,dy
\end{aligned}
\]
\end{theorem}

公式中如果$\Gamma$ 的方向与 $\Sigma$ 的侧取相反的一侧，右手法则就不适用，公式右边的积分结果就要变号。这一点和第二类曲面积分的方向性是一致的。

这个公式的证明思路与格林公式类似，我们这里只简要分析证明思路，不再展开详细计算。

以曲面 $\Sigma$ 由 $z=z(x,y)$ 给出并取上侧为例，先把曲面积分中的 $dy\,dz$、$dz\,dx$、$dx\,dy$ 各项用 $z_x,z_y$ 表示成对 $xOy$ 面上投影区域 $D_{xy}$ 的二重积分，再对所得的二重积分应用格林公式，最后整理便得到沿 $\Gamma$ 的曲线积分。这样就把曲面积分转换为了曲线积分。

直观地看，斯托克斯公式是格林公式在空间中的推广。当曲面 $\Sigma$ 退化到 $xOy$ 面上的一块平面区域时，$dz=0$，斯托克斯公式就回到了格林公式。因此，格林公式可以看成斯托克斯公式的特殊情形。

我们看几个例子感受一下：

\begin{myexample}
计算 $\oint_{\Gamma} y\,dx + z\,dy + x\,dz$，其中 $\Gamma$ 是平面 $z=1$ 上的圆周 $x^2+y^2=1$，取逆时针方向。

先直接计算。令$x=\cos t$，$y=\sin t$，$z=1$，$t\in[0,2\pi]$，则\(dx=-\sin t\,dt\)，\(dy=\cos t\,dt\)，\(dz=0\)，代入得：
\[
\oint_{\Gamma} y\,dx+z\,dy+x\,dz
=
\int_0^{2\pi}\left[\sin t(-\sin t)+1\cdot\cos t+\cos t\cdot 0\right]dt
= -\pi
\]

我们也可以用斯托克斯公式计算。设 $\Sigma$ 为圆盘 $z=1$，$x^2+y^2\le 1$，取上侧，法向量 $\vec{n}=(0,0,1)$。这里 $P=y$，$Q=z$，$R=x$，求偏导得：
\[
\begin{aligned}
\frac{\partial R}{\partial y}-\frac{\partial Q}{\partial z}&=0-1=-1\\
\frac{\partial P}{\partial z}-\frac{\partial R}{\partial x}&=0-1=-1\\
\frac{\partial Q}{\partial x}-\frac{\partial P}{\partial y}&=0-1=-1
\end{aligned}
\]

代入斯托克斯公式：
\[
\oint_{\Gamma} y\,dx+z\,dy+x\,dz
=
\iint_{\Sigma}(-1,-1,-1)\cdot(0,0,1)\,dS
=
-\iint_{\Sigma}dS
= -\pi
\]

两者结果一致。
\end{myexample}

\begin{myexample}
计算 $\oint_{\Gamma}(z-y)\,dx+(x-z)\,dy+(y-x)\,dz$，其中 $\Gamma$ 是平面 $x+y+z=1$ 在第一卦限内三角形的边界，取从 $z$ 轴正方向看为逆时针的方向。

由斯托克斯公式，令 $\Sigma$ 为平面 $x+y+z=1$ 在第一卦限内的三角形，取上侧。其单位法向量为\(\vec{n}=\frac{1}{\sqrt{3}}(1,1,1)\)。这里 $P=z-y$，$Q=x-z$，$R=y-x$，于是：
\[
\begin{aligned}
\frac{\partial R}{\partial y}-\frac{\partial Q}{\partial z}&=1-(-1)=2\\
\frac{\partial P}{\partial z}-\frac{\partial R}{\partial x}&=1-(-1)=2\\
\frac{\partial Q}{\partial x}-\frac{\partial P}{\partial y}&=1-(-1)=2
\end{aligned}
\]

代入斯托克斯公式：
\[
\oint_{\Gamma}(z-y)\,dx+(x-z)\,dy+(y-x)\,dz
=
\iint_{\Sigma}(2,2,2)\cdot\frac{1}{\sqrt{3}}(1,1,1)\,dS
=
\iint_{\Sigma}2\sqrt{3}\,dS
\]

则积分为\(2\sqrt{3}\)的三角形面积。三角形 $\Sigma$ 的三个顶点为 $(1,0,0)$，$(0,1,0)$，$(0,0,1)$，是边长为 $\sqrt{2}$ 的等边三角形，面积为 $\frac{\sqrt{3}}{2}$。所以：
\[
\iint_{\Sigma}2\sqrt{3}\,dS
=
2\sqrt{3}\cdot\frac{\sqrt{3}}{2}
=
3
\]

所以 $\oint_{\Gamma}(z-y)\,dx+(x-z)\,dy+(y-x)\,dz=3$。
\end{myexample}

\begin{myexample}
计算 $\oint_{\Gamma} (y^2+z^2)\,dx+(z^2+x^2)\,dy+(x^2+y^2)\,dz$，其中 $\Gamma$ 是圆周 $x^2+y^2=1$，$z=1$，取从 $z$ 轴正方向看为逆时针的方向。

这里 $P=y^2+z^2$，$Q=z^2+x^2$，$R=x^2+y^2$。先求偏导：
\[
\begin{aligned}
\frac{\partial R}{\partial y}-\frac{\partial Q}{\partial z}&=2y-2z\\
\frac{\partial P}{\partial z}-\frac{\partial R}{\partial x}&=2z-2x\\
\frac{\partial Q}{\partial x}-\frac{\partial P}{\partial y}&=2x-2y
\end{aligned}
\]

令 $\Sigma$ 为圆盘 $z=1$，$x^2+y^2\le 1$，取上侧，单位法向量 $\vec{n}=(0,0,1)$。于是\(dydz=0\)，\(dzdx=0\)，\(dxdy=dS\)，代入斯托克斯公式得：
\[
\begin{aligned}
\oint_{\Gamma} P\,dx+Q\,dy+R\,dz
&=
\iint_{\Sigma}
\left[(2y-2z)dy\,dz+(2z-2x)dz\,dx+(2x-2y)dx\,dy\right] \\
&=
\iint_{D_{xy}} (2x-2y)\,dx\,dy
\end{aligned}
\]

其中 $D_{xy}$ 为圆盘 $x^2+y^2\le 1$。

令 $x=\rho\cos\theta$，$y=\rho\sin\theta$，则 $0\le\rho\le 1$，$0\le\theta\le2\pi$，$dx\,dy=\rho\,d\rho\,d\theta$。于是：
\[
\begin{aligned}
\iint_{D_{xy}} (2x-2y)\,dx\,dy
&=
\int_0^{2\pi}\int_0^1
\left[2\rho\cos\theta-2\rho\sin\theta\right]\rho\,d\rho\,d\theta \\
&=
\int_0^{2\pi}\cos\theta\,d\theta\int_0^1 2\rho^2\,d\rho
-\int_0^{2\pi}\sin\theta\,d\theta\int_0^1 2\rho^2\,d\rho \\
&=
0 - 0
= 0
\end{aligned}
\]

所以 $\oint_{\Gamma} (y^2+z^2)\,dx+(z^2+x^2)\,dy+(x^2+y^2)\,dz = 0$。
\end{myexample}

斯托克斯公式把空间闭曲线上的第二类曲线积分，转化成了以该曲线为边界的曲面上的第二类曲面积分。和格林公式一样，它体现的也是“边界上的积分由内部决定”的思想，只不过从平面区域推广到了空间曲面。

\subsection{环流量与旋度}

在第二类曲线积分中，我们已经知道曲线积分可以表示变力沿路径做的功。如果积分路径是一条闭曲线，这个积分就有了专门的名称。

\begin{definition}{环流量}{}
设向量函数 $\vec{F}(x,y,z)=(P,Q,R)$ 定义在空间区域上，$C$ 是一条有向光滑闭曲线，则称
\[\Gamma =
\oint_{C}\vec{F}\cdot d\vec{r}
=
\oint_{C}P\,dx+Q\,dy+R\,dz
\]
为向量函数 $\vec{F}$ 沿闭曲线 $C$ 的\textbf{环流量}。
\end{definition}

通俗地讲，如果 $\vec{F}$ 是水流的速度，环流量刻画的就是水流绕这条闭曲线“转”了多少。如果 $\vec{F}$ 表示的是力，环流量就是沿闭曲线运动一周时力所做的功。环流量越大，说明流体绕该闭曲线的环流越明显。

那么，这种“打转”在每一点附近到底有多强？这就引出了旋度。

\begin{definition}{旋度}{}
设向量函数 $\vec{F}(x,y,z)=(P,Q,R)$ 在点 $(x,y,z)$ 处具有一阶连续偏导数，则称向量
\[
\operatorname{rot}\vec{F}
=
\begin{vmatrix}
\vec{i} & \vec{j} & \vec{k} \\
\dfrac{\partial}{\partial x} & \dfrac{\partial}{\partial y} & \dfrac{\partial}{\partial z} \\
P & Q & R
\end{vmatrix}
=
\left(
\frac{\partial R}{\partial y}-\frac{\partial Q}{\partial z},
\ \frac{\partial P}{\partial z}-\frac{\partial R}{\partial x},
\ \frac{\partial Q}{\partial x}-\frac{\partial P}{\partial y}
\right)
\]
为向量函数 $\vec{F}$ 在点 $(x,y,z)$ 处的\textbf{旋度}。
\end{definition}


散度是一个标量，而旋度是一个向量。散度表示的是“这一点附近有没有源或汇”，旋度表示的是“这一点附近有没有旋转”。如果在某一点 $\operatorname{rot}\vec{F}\neq\vec{0}$，说明向量函数在该点附近有旋转趋势，旋度的方向就是旋转轴的方向，旋度的大小刻画旋转的强弱。如果处处有 $\operatorname{rot}\vec{F}=\vec{0}$，则说明区域内没有旋转。

我们看一个简单而典型的例子，它说明了旋度与旋转的关系。

\begin{myexample}
设刚体绕 $z$ 轴以恒定角速度 $\omega$ 旋转，其速度为\(
\vec{v}=(-\omega y,\ \omega x,\ 0)\)，求 $\operatorname{rot}\vec{v}$。

这里 $P=-\omega y$，$Q=\omega x$，$R=0$。求偏导：
\[
\frac{\partial R}{\partial y}=0,\quad \frac{\partial Q}{\partial z}=0,\quad
\frac{\partial P}{\partial z}=0,\quad \frac{\partial R}{\partial x}=0,\quad
\frac{\partial Q}{\partial x}=\omega,\quad \frac{\partial P}{\partial y}=-\omega
\]

于是：
\[
\operatorname{rot}\vec{v}
=
(0-0,\ 0-0,\ \omega-(-\omega))
=
(0,0,2\omega)
\]

这说明刚体旋转的旋度方向沿转轴方向，大小等于角速度的两倍。
\end{myexample}

有了旋度概念，斯托克斯公式就可以写成另一种形式：
\[
\oint_{\Gamma}\vec{F}\cdot d\vec{r}
=
\iint_{\Sigma}(\operatorname{rot}\vec{F})\cdot\vec{n}\,dS
\]

这个形式表达的正是斯托克斯公式的物理意义：\textbf{沿闭曲线的环流量等于旋度穿过以该曲线为边界的曲面的通量。}所以斯托克斯公式又被称为旋度定理。

我们看一个计算旋度的例子：

\begin{myexample}
求向量函数 $\vec{F}=(xy,\ yz,\ zx)$ 在点 $(1,1,1)$ 处的旋度。

这里 $P=xy$，$Q=yz$，$R=zx$。先求偏导：
\[
\frac{\partial R}{\partial y}=0,\quad \frac{\partial Q}{\partial z}=y,\quad
\frac{\partial P}{\partial z}=0,\quad \frac{\partial R}{\partial x}=z,\quad
\frac{\partial Q}{\partial x}=0,\quad \frac{\partial P}{\partial y}=x
\]

于是：
\[
\operatorname{rot}\vec{F}
=
\left(\frac{\partial R}{\partial y}-\frac{\partial Q}{\partial z},
\frac{\partial P}{\partial z}-\frac{\partial R}{\partial x},
\frac{\partial Q}{\partial x}-\frac{\partial P}{\partial y}\right)
=
(-y,\ -z,\ -x)
\]

代入点 $(1,1,1)$，得：
\[
\operatorname{rot}\vec{F}(1,1,1)=(-1,-1,-1)
\]

或者直接用行列式计算：
\[
\begin{aligned}
    \operatorname{rot}\vec{F}
    &=
    \begin{vmatrix}
    \vec{i} & \vec{j} & \vec{k} \\
    \dfrac{\partial}{\partial x} & \dfrac{\partial}{\partial y} & \dfrac{\partial}{\partial z} \\
    P & Q & R
    \end{vmatrix}
    =
    \vec{i}
    \begin{vmatrix}
    \frac{\partial}{\partial y} & \frac{\partial}{\partial z} \\
    yz & zx
    \end{vmatrix}
    -
    \vec{j}
    \begin{vmatrix}
    \frac{\partial}{\partial x} & \frac{\partial}{\partial z} \\
    xy & zx
    \end{vmatrix}
    +
    \vec{k}
    \begin{vmatrix}
    \frac{\partial}{\partial x} & \frac{\partial}{\partial y} \\
    xy & yz
    \end{vmatrix}\\
    &=
    \frac{\partial}{\partial y}(zx)
-
\frac{\partial}{\partial z}(yz)-(\frac{\partial}{\partial x}(zx)
-
\frac{\partial}{\partial z}(xy))+\frac{\partial}{\partial x}(yz)
-
\frac{\partial}{\partial y}(xy)\\
&=
-y\,\vec{i}
-z\,\vec{j}
-x\,\vec{k}
\end{aligned}
\]

代入点 \((1,1,1)\)得到：
\[
\operatorname{rot}\vec{F}(1,1,1)
=
(-1, -1, -1)
\]
\end{myexample}

\begin{myexample}
设 $\vec{F}=(y,\ x,\ 0)$，求 $\operatorname{rot}\vec{F}$，并判断是否无旋。

这里 $P=y$，$Q=x$，$R=0$。计算偏导：
\[
\frac{\partial R}{\partial y}=0,\quad \frac{\partial Q}{\partial z}=0,\quad
\frac{\partial P}{\partial z}=0,\quad \frac{\partial R}{\partial x}=0,\quad
\frac{\partial Q}{\partial x}=1,\quad \frac{\partial P}{\partial y}=1
\]

所以：
\[
\operatorname{rot}\vec{F}
=
(0-0,\ 0-0,\ 1-1)
=
(0,0,0)
\]

因此 $\vec{F}=(y,x,0)$ 无旋。
\end{myexample}

总结一下，格林公式可以把二重积分化为平面闭曲线上的第二类曲线积分，高斯公式可以把三重积分化为闭曲面上的第二类曲面积分，斯托克斯公式可以把第二类曲面积分化为空间闭曲线上的第二类曲线积分。它们和牛顿-莱布尼茨公式共同说明了区域内部的整体变化可以由其边界上的积分表示出来。
% !TEX root = ../main.tex
\chapter{广义斯托克斯定理}
上一章我们分别得到了格林公式，高斯公式和斯托克斯公式。从形式上看，这三个公式各不相同，但它们有一个共同的思想，就是边界上的积分等于内部某种“导数”的积分。这个思想在牛顿-莱布尼茨公式中已经出现了。区间上的导数积分等于两端点函数值之差。如果把端点看作区间的“边界”，那么它说的正是“内部导数的累计等于边界上的函数值之差”。格林公式，高斯公式和斯托克斯公式都是这个思想向高维或不同几何对象推广的产物。

那么，能否有一条统一的公式，把格林公式，高斯公式和斯托克斯公式全部包含进去？我们一起探索一下。
\section{场论基础}
我们在多元函数微分学遇到了梯度，在曲线积分与曲面积分遇到了散度和旋度。现在，我们从“场”的角度再次认识这些概念。
\subsection{标量场与向量场}
前面我们研究的函数，可以理解为自变量是点的坐标，因变量是一个数。比如二元函数 $z = f(x, y)$ 就是平面上每一个点对应一个实数，三元函数$w = f(x, y, z)$ 就是空间中每一个点对应一个实数。在物理和工程中，这种“给每个点分配一个数”的模式极其常见。比如房间里每一点有一个温度，大气中每一点有一个气压，电场中每一点有一个电势。我们把这种对象叫做\textbf{标量场}。

\begin{definition}{标量场}{}
设 $D$ 是空间中的一个区域。若对于 $D$ 中的每一个点 $P(x, y, z)$，都有唯一确定的标量$\varphi(x, y, z)$ 与之对应，则称 $\varphi$ 是定义在 $D$ 上的一个\textbf{标量场}。
\end{definition}

简单来说，标量场就是一个三元函数 $\varphi(x, y, z)$，只不过我们换了一个名字，强调它描述的是空间中每一点的某种“量”。温度场 $T(x, y, z)$，密度场 $\rho(x, y, z)$，电势场 $V(x, y, z)$ 等等，它们都可以算作标量场。本质上，标量场其实是函数的一种。

标量场有一个非常直观的几何工具，就是\textbf{等值面}。把所有函数值相等的点连起来，就得到一张曲面。

\begin{definition}{等值面}{}
标量场 $\varphi(x, y, z)$ 中，满足 $\varphi(x, y, z) = C$（$C$ 为常数）的点构成的曲面，称为标量场的一个\textbf{等值面}。
\end{definition}

等值面就像地图上的等高线。想象一座山，海拔相同的点连成一条闭合曲线，就是等高线。推广到三维空间，海拔相同的点连成一张曲面，就是等值面。比如温度场 $T(x, y, z) = 100$ 的等值面，就是所有温度恰好为 $100$ 度的点组成的曲面。取不同的 $C$ 值，就得到不同的一族等值面。

对于二维标量场 $\varphi(x, y)$，等值面退化为\textbf{等值线}，即满足 $\varphi(x, y) = C$ 的曲线。

等值面的法方向和标量场的变化率密切相关。我们之前学过梯度，它垂直于等值面，指向函数值增加最快的方向。

标量场描述的是“每个点一个数”，但很多物理量不仅有大小，还有方向。比如风速，你不仅要知道风有多大，还要知道风往哪吹。这种描述的是“给每个点分配一个向量”的对象，就是\textbf{向量场}。

\begin{definition}{向量场}{}
设 $D$ 是空间中的一个区域。若对于 $D$ 中的每一个点 $P(x, y, z)$，都有唯一确定的向量
\[
\vec{F}(x, y, z) = (P(x, y, z),\; Q(x, y, z),\; R(x, y, z))
\]
与之对应，则称 $\vec{F}$ 是定义在 $D$ 上的一个\textbf{向量场}。其中 $P$、$Q$、$R$ 是三个标量函数，分别称为向量场的 $x$、$y$、$z$ 分量。
\end{definition}

向量场在每一点给出一个向量。想象一下，你在一片草地上，每一根草都指向某个方向，倒向某个角度。本质上，向量场是向量函数的一种。

为了直观地表示向量场，我们可以在区域内的若干点处画出对应的向量箭头。箭头的方向表示向量场在该点的方向，箭头的长度表示向量场在该点的大小。这种图就是\textbf{向量场的图示}。例如涡流。

\begin{center}
  \begin{tikzpicture}
    % 坐标轴
    \draw[->] (-2.8,0) -- (2.8,0) node[right] {$x$};
    \draw[->] (0,-2.8) -- (0,2.8) node[above] {$y$};

    % 向量场 F = (-y, x, 0) 在 xOy 平面上的投影
    \foreach \x in {-2,-1,0,1,2} {
      \foreach \y in {-2,-1,0,1,2} {
        \pgfmathsetmacro{\vx}{-\y}
        \pgfmathsetmacro{\vy}{\x}
        \pgfmathsetmacro{\len}{sqrt(\vx*\vx + \vy*\vy)}
        \ifdim\len pt > 0.1pt
          \pgfmathsetmacro{\nx}{\vx/\len}
          \pgfmathsetmacro{\ny}{\vy/\len}
          \draw[->, blue!70!black, thick] (\x,\y) -- ({\x+0.55*\nx},{\y+0.55*\ny});
        \fi
      }
    }
  \end{tikzpicture}
\end{center}

生活中，风场可以看成一个向量场，箭头方向表示风向，箭头大小表示风速。水流也是一个典型的例子，例如海洋里每一点的水流方向与快慢构成一个向量场。物理中，我们熟悉的电场其实就是向量场，箭头方向表示电场方向，箭头大小表示电场强度大小，类似还有磁场，力场等等。

标量场和向量场是场论的两个基本对象。

\subsection{场的性质}
假设我们现在知道了一个场，那么我们要怎么描述这个场呢？比如温度场，有些地方热，有些地方冷，我们要怎么描述温度场的这种变化呢？再比如有一个风场，如果这个风场是一个龙卷风，那么在数学上要怎么描述这个风场呢？或者说负电荷的电场，所有电场方向都指向负电荷中心，这又要怎么描述呢？

我们之前学的梯度描述的是函数值变化最快的方向和快慢。如果把它作用在标量场上，就得到：

\begin{definition}{场的梯度}{}
  设标量场 \(f(x,y,z)\) 具有一阶连续偏导数，则称
  \[
  \operatorname{grad} f = \nabla f = \left(\frac{\partial f}{\partial x},\; \frac{\partial f}{\partial y},\; \frac{\partial f}{\partial z}\right)
  \]
  为标量场 \(f\) 的\textbf{梯度}。
\end{definition}

我们知道，梯度是一个向量，它指向函数值增加最快的方向，其模长等于沿该方向的方向导数。因此梯度刻画了标量场在一点附近的“变化趋势”。

如果我们把温度场理解为标量场 \(f(x,y,z)\)，那么沿着梯度的方向就是温度升高最快的方向。反过来，温度降低最快的方向就是负梯度方向。这样我们就用梯度把温度场的冷热变化描述清楚了。

梯度描述的是标量场的变化，但向量场本身既有大小又有方向，它的变化规律更复杂。比如龙卷风的风速场，既有“向外散开”的趋势，也有“绕中心旋转”的趋势。这就需要散度和旋度了。

散度刻画的是向量场在一点处的“发散”或“汇聚”程度。

\begin{definition}{场的散度}{}
  设向量场 \(\vec{F}(x,y,z)=(P(x,y,z),\;Q(x,y,z),\;R(x,y,z))\) 具有一阶连续偏导数，则称
  \[
  \operatorname{div}\vec{F} = \nabla \cdot \vec{F} = \frac{\partial P}{\partial x} + \frac{\partial Q}{\partial y} + \frac{\partial R}{\partial z}
  \]
  为向量场 \(\vec{F}\) 的\textbf{散度}。
\end{definition}

散度是一个\textbf{标量}。它把向量场的三个分量分别沿自身方向求偏导后相加，也就是 \(\nabla\) 与 \(\vec{F}\) 做“点积”。

我们知道，如果把 \(\vec{F}\) 理解为流体的速度场，则 \(\operatorname{div}\vec{F}>0\) 表示该点有流体向外流出，也就是“源”。\(\operatorname{div}\vec{F}<0\) 表示该点有流体向内流入，也就是“汇”。\(\operatorname{div}\vec{F}=0\) 表示该点无源无汇，流入与流出平衡。这样我们就用散度把向量场在某一点是向外发散还是向内汇聚的程度描述清楚了。

再来看旋度。旋度刻画的是向量场在一点处的“旋转”趋势。

\begin{definition}{场的旋度}{}
  设向量场 \(\vec{F}(x,y,z)=(P(x,y,z),\;Q(x,y,z),\;R(x,y,z))\) 具有一阶连续偏导数，则称
  \[
  \operatorname{rot}\vec{F} = \nabla \times \vec{F}
  =
  \begin{vmatrix}
  \vec{i} & \vec{j} & \vec{k} \\[2pt]
  \dfrac{\partial}{\partial x} & \dfrac{\partial}{\partial y} & \dfrac{\partial}{\partial z} \\[6pt]
  P & Q & R
  \end{vmatrix}
  \]
  为向量场 \(\vec{F}\) 的\textbf{旋度}。
\end{definition}

旋度是一个\textbf{向量}。它的方向表示旋转轴的方向，模长表示旋转强度。它把\(\nabla\) 与 \(\vec{F}\) 做“叉积”，计算时和普通向量叉积完全一样，只是把“分量”换成了偏导数。仔细观察就会发现，行列式的第二行就是Nabla算子。

如果把 \(\vec{F}\) 理解为流体的速度场，那么它的方向表示流体在该点旋转的转轴方向，它的模长表示旋转的强度。\(\operatorname{rot}\vec{F} \ne 0\)，表示该点附近的流体有旋转趋势或正在旋转。\(\operatorname{rot}\vec{F} = 0\)，表示该点没有旋转趋势，流体可能只是弯曲流动而没有自身打转。这样我们就用旋度把向量场在某一点旋转的趋势描述清楚了。

利用散度和旋度，可以对向量场进行分类。

\begin{definition}{场的分类}{}
  设向量场 \(\vec{F}\) 具有连续的一阶偏导数。

  若 \(\operatorname{div}\vec{F}=0\)，则称 \(\vec{F}\) 为\textbf{无源场}。

  若 \(\operatorname{rot}\vec{F}=\vec{0}\)，则称 \(\vec{F}\) 为\textbf{无旋场}。

  若同时有 \(\operatorname{div}\vec{F}=0\) 且 \(\operatorname{rot}\vec{F}=\vec{0}\)，则称 \(\vec{F}\) 为\textbf{调和场}。
\end{definition}

其中无旋场中有一类最特殊的，那就是\textbf{保守场}。

\begin{definition}{保守场}{}
  若向量场 \(\vec{F}\) 是某个标量场 \(f\) 的梯度，即 \(\vec{F}=\nabla f\)，则称 \(\vec{F}\) 为\textbf{保守场}，\(f\) 称为 \(\vec{F}\) 的\textbf{势函数}。
\end{definition}

保守场在物理学中极为常见。重力场，静电场都是保守场。保守场的一个重要特征是第二类曲线积分与路径无关。结合之前对于曲线积分与路径无关的结论，我们有以下等价条件。

\begin{theorem}{保守场的等价条件}{}
  设 \(\vec{F}\) 在单连通区域 \(\Omega\) 上具有一阶连续偏导数，则以下条件等价：
\[
\begin{aligned}
&\text{(1)}\ \vec{F}\text{ 是保守场，即存在势函数 }f\text{，使得 }\vec{F}=\nabla f \\[6pt]
&\text{(2)}\ \vec{F}\text{ 是无旋场，即 }\operatorname{rot}\vec{F}=\vec{0} \\[6pt]
&\text{(3)}\ \vec{F}\text{ 的第二类曲线积分在 }\Omega\text{ 内与路径无关} \\[6pt]
&\text{(4)}\ \vec{F}\text{ 沿 }\Omega\text{ 内任意闭合曲线的第二类曲线积分为零}
\end{aligned}
\]
\end{theorem}

注意观察梯度、散度、旋度的形式，它们似乎都可以用带有Nabla算子的式子表达，这是否暗示着它们之间可能存在某种联系？答案是肯定的。我们把这些运算做复合，就能得到场论最基本的恒等式。

\begin{theorem}{场论基本恒等式}{}
  设 \(f\) 是标量场，\(\vec{F}\) 是向量场，且各阶偏导数连续，则

  梯度的旋度恒为零：
  \[
  \operatorname{rot}(\nabla f)=\nabla \times (\nabla f)=\vec{0}
  \]

  旋度的散度恒为零：
  \[
  \operatorname{div}(\operatorname{rot}\vec{F})=\nabla \cdot (\nabla \times \vec{F})=0
  \]

  梯度的散度为拉普拉斯算子：
  \[
  \operatorname{div}(\nabla f)=\nabla \cdot (\nabla f)
  =\nabla^2 f
  =\frac{\partial^2 f}{\partial x^2}+\frac{\partial^2 f}{\partial y^2}+\frac{\partial^2 f}{\partial z^2}
  \]

其中\(\nabla^2=\frac{\partial^2}{\partial x^2}+\frac{\partial^2}{\partial y^2}+\frac{\partial^2}{\partial z^2}\)称为\textbf{拉普拉斯算子}。满足 \(\nabla^2\varphi=0\) 的函数称为\textbf{调和函数}，对应的方程 \(\nabla^2\varphi=0\) 称为\textbf{拉普拉斯方程}。
\end{theorem}

梯度的散度表示标量场在某点相对于周围环境的“凸起”或“凹陷”程度。如果在某点\(\nabla^2 f>0\)，那么从该点出发，往四周走 \(f\)都会变大，说明这个点是“低谷”。如果在某点\(\nabla^2 f<0\)，那么从该点出发，往四周走 \(f\)都会变小，说明这个点是“高峰”。如果在某点\(\nabla^2 f=0\)，那么该点与周围持平。这和我们判断函数的极值是类似的。

“梯度的旋度为零”意味着，如果一个向量场是某个标量场的梯度，那么它一定是无旋的。“旋度的散度为零”意味着，如果一个向量场是某个向量场的旋度，那么它一定是无源的。这两个恒等式表示的逆命题也是成立的，在一定的区域条件下，无旋场必为梯度场，无源场必为旋度场。

\begin{myexample}
  判断向量场 \(\vec{F}=(2xy,\;x^2,\;1)\) 是否为保守场，是否为调和场。

  计算旋度：
  \[
  \operatorname{rot}\vec{F}
  =
  \begin{vmatrix}
  \vec{i} & \vec{j} & \vec{k} \\[2pt]
  \dfrac{\partial}{\partial x} & \dfrac{\partial}{\partial y} & \dfrac{\partial}{\partial z} \\[6pt]
  2xy & x^2 & 1
  \end{vmatrix}
  =
  (0-0,\;0-0,\;2x-2x)
  =
  (0,\;0,\;0)
  \]

  旋度为零，所以 \(\vec{F}\) 是保守场。

  计算散度：
  \[
  \operatorname{div}\vec{F}
  =
  \nabla \cdot \vec{F}
  =
  \frac{\partial (2xy)}{\partial x} + \frac{\partial x^2}{\partial y} + \frac{\partial (1)}{\partial z}
  =
  2y
  \]

  散度不为零，所以 \(\vec{F}\) 不是调和场。

\end{myexample}

总结一下，梯度把标量场变为向量场，散度把向量场变为标量场，旋度把向量场变为向量场。梯度描述标量场的变化，散度描述向量场的发散与汇聚，旋度描述向量场的旋转。三者的复合产生了场论的基本恒等式，也为无源场，无旋场，保守场提供了判别方法。

\section{微分形式与外微分}


\subsection{微分形式}
前面我们分别讨论了格林公式、高斯公式和斯托克斯公式，它们的被积表达式看起来各不相同，但它们有一个共同的特点，就是由\(dx\)，\(dy\)，\(dz\) 这些基本微分按照一定方式组合起来的。它们在结构上又有相似的地方，即区域的边界上的积分等于区域内部某种“导数”的积分。但是要把这种相似性说清楚，需要解决一个语言问题。于是我们自然产生一个想法，能不能把这些\(dx\)，\(dy\)，\(dz\) 这些基本微分统一成同一种记号？如果能，那么这三个公式能不能写成一个公式？微分形式就是用来做这件事的。

我们先从最简单的对象开始。

一个普通的光滑函数 $\varphi(x,y,z)$ 不含任何微分，称为 $0$-形式。

\begin{definition}{0-形式}{}
设 $\varphi(x,y,z)$ 是光滑函数，则称 $\varphi$ 为一个\textbf{$0$-形式}。
\end{definition}

$0$-形式本身不需要积分，它只是函数值。

接下来看线积分。第二类曲线积分的被积表达式是\(P\,dx+Q\,dy+R\,dz\)，这里的 \(dx\)、\(dy\)、\(dz\) 在积分中总是与某个分量函数搭配出现，整体表示沿曲线的某种作用量。我们把这种含有一个微分的表达式称为 $1$-形式。

\begin{definition}{1-形式}{}
设 \(P(x,y,z)\)、\(Q(x,y,z)\)、\(R(x,y,z)\) 是空间区域上的光滑函数，则称
\[
\omega=P\,dx+Q\,dy+R\,dz
\]
为一个\textbf{$1$-形式}。
\end{definition}

在微分形式的语言里，\(dx\)，\(dy\)，\(dz\) 不再表示某个变量的增量，而是三个基本符号，或者说三个基本“基底”。$1$-形式就是这三个基底的组合，系数是函数 \(P\)、\(Q\)、\(R\)。例如，力场 \(\vec{F}=(P,Q,R)\) 沿曲线做功时，被积表达式\(\vec{F}\cdot d\vec{r}=P\,dx+Q\,dy+R\,dz\)就是一个典型的$1$-形式。

再看面积分。第二类曲面积分的被积表达式含有两个微分，即 $dy\,dz$、$dz\,dx$、$dx\,dy$。这种表达式显然不能写成$1$-形式，因为$1$-形式只含有一个基本微分。我们需要一种由两个基本微分相乘得到的新对象。但这里的“乘积”不能是普通的数的乘法。因为在曲面积分中，\(dx\,dy\) 表示一个有方向的面积元素，交换顺序应当改变方向，即 \(dx\,dy=-dy\,dx\)。如果把这种乘积看成普通乘法，交换顺序不变，就会丢掉方向性。因此，我们引入一种新的乘积，称为\textbf{楔积}，它最重要的性质是反对称。

\begin{definition}{楔积的运算规则}{}
  在基本微分 \(dx\)，\(dy\)，\(dz\) 之间定义一种乘积，记作 \(dx\wedge dy\)，称为\textbf{楔积}。它满足：
\[
dx\wedge dx=0,\quad dy\wedge dy=0,\quad dz\wedge dz=0
\]
\[
dx\wedge dy=-dy\wedge dx,\quad dx\wedge dz=-dz\wedge dx,\quad dy\wedge dz=-dz\wedge dy
\]
\end{definition}

为什么要有这样一条规则？这和向量叉乘的性质完全类似。交换两个方向，符号改变。$dx\wedge dy$ 可以理解为 $xOy$ 平面上的有向面积元素，$dy\wedge dx$ 则是相反方向的面积元素。换句话说，楔积记住了面积的“方向”。如果两个微分相同，比如 $dx\wedge dx$，就相当于同一个方向和自己交换，因为 $dx\wedge dx=-dx\wedge dx$，所以$dx\wedge dx$只能等于零。

有了楔积，就可以定义 $2$-形式。

\begin{definition}{2-形式}{}
设 \(P(x,y,z)\)、\(Q(x,y,z)\)、\(R(x,y,z)\) 是空间区域上的光滑函数，则称
\[
\omega=P\,dy\wedge dz+Q\,dz\wedge dx+R\,dx\wedge dy
\]
为一个\textbf{$2$-形式}。
\end{definition}

这里的顺序不是随意写的。我们采用循环顺序 \(dy\wedge dz\)、\(dz\wedge dx\)、\(dx\wedge dy\)，正是为了与第二类曲面积分中的有向面积元素 \(dy\,dz\)、\(dz\,dx\)、\(dx\,dy\) 对应。如果写成 \(dx\wedge dz\)，根据反对称性会多出一个负号。

若 $\vec F=(P,Q,R)$ 是向量场，它穿过曲面 $\Sigma$ 的通量为\(\iint_\Sigma P\,dy\,dz+Q\,dz\,dx+R\,dx\,dy
\)。在微分形式中，这个被积表达式就是 $2$-形式。$dy\wedge dz$，$dz\wedge dx$，$dx\wedge dy$ 分别对应三个坐标平面上的有向面积元素。

再进一步，三重积分中的被积表达式含有三个基本微分 \(dx\,dy\,dz\)。同样，我们把它写成 \(dx\wedge dy\wedge dz\)，称为 $3$-形式。

\begin{definition}{3-形式}{}
设 $f(x,y,z)$ 为光滑函数，则称
\[
\omega=f\,dx\wedge dy\wedge dz
\]
为一个\textbf{$3$-形式}。
\end{definition}

$3$-形式对应空间区域上的积分，$dx\wedge dy\wedge dz$ 表示有向体积元素。在三维空间中，最高只能到 $3$-形式，因为只有三个独立的方向。如果再写 $dx\wedge dy\wedge dz\wedge dx$，由于前面已经有一个 $dx$，根据反对称性，整个式子必然为零。因此三维空间中的微分形式到$3$-形式为止，没有$4$-形式或更高阶形式。

现在我们可以把前面各章中不同维度的积分统一起来了。$0$-形式对应点上的函数值，$1$-形式对应第二类曲线积分，$2$-形式对应第二类曲面积分，$3$-形式对应三重积分。仔细观察就会发现，微分形式的“阶数”正好等于积分区域的维数。点是 $0$ 维的，曲线是 $1$ 维的，曲面是 $2$ 维的，空间区域是 $3$ 维的。至此，我们所学的积分算是初步统一了。

\subsection{外微分}
格林公式、高斯公式和斯托克斯公式都是把区域内部的积分与边界上的积分联系起来，而边界的维数恰好比内部少一维。换句话说，边界上出现的微分形式通常比内部少一阶。例如斯托克斯公式中，曲线边界上出现的是$1$-形式，曲面上出现的是$2$-形式。高斯公式中，封闭曲面上出现的是$2$-形式，体内出现的是$3$-形式。这正好说明了区域的边界上的积分等于区域内部某种“导数”的积分，因为我们知道求导数总是会降阶的。那么这种“导数”究竟是什么呢？

前面我们已经把各种积分中的被积表达式统一成了微分形式，接下来就要把格林公式、高斯公式和斯托克斯公式中的“导数”统一起来了。格林公式中出现的是 \(\frac{\partial Q}{\partial x}-\frac{\partial P}{\partial y}\)，斯托克斯公式中出现的是旋度，高斯公式中出现的是散度。和微分形式一样，它们都应当理解为对微分形式的一种微分运算。这个运算称为\textbf{外微分}。

\begin{definition}{外微分}{}
设 \(f,P,Q,R\) 为光滑函数，定义外微分算子 \(d\) 如下：

\(0\)-形式 \(f\) 的外微分为：
  \[
  df=\frac{\partial f}{\partial x}\,dx+\frac{\partial f}{\partial y}\,dy+\frac{\partial f}{\partial z}\,dz
  \]

  \(1\)-形式 \(\omega=P\,dx+Q\,dy+R\,dz\) 的外微分为：
  \[
  d\omega=
  \left(\frac{\partial R}{\partial y}-\frac{\partial Q}{\partial z}\right)dy\wedge dz
  +\left(\frac{\partial P}{\partial z}-\frac{\partial R}{\partial x}\right)dz\wedge dx
  +\left(\frac{\partial Q}{\partial x}-\frac{\partial P}{\partial y}\right)dx\wedge dy
  \]

  \(2\)-形式 \(\omega=P\,dy\wedge dz+Q\,dz\wedge dx+R\,dx\wedge dy\) 的外微分为：
  \[
  d\omega=
  \left(\frac{\partial P}{\partial x}+\frac{\partial Q}{\partial y}+\frac{\partial R}{\partial z}\right)
  dx\wedge dy\wedge dz
  \]

  \(3\)-形式 \(\omega=f\,dx\wedge dy\wedge dz\) 的外微分为：
  \[
  d\omega=0
  \]
\end{definition}

外微分不是随便定义的，我们希望有这样一个算子，它是普通微分的推广，满足庞加莱引理和莱布尼茨法则，保持方向性，并且满足这些条件的算子是唯一的。

先介绍庞加莱引理和莱布尼茨法则：

\begin{theorem}{庞加莱引理}{}
  设 \(\omega\) 是任意微分形式，且系数具有二阶连续偏导数，则\(d(d\omega)=0\)。

\end{theorem}
庞加莱引理其实是“边界的边界为空”这一几何事实在代数上的表现。

\begin{theorem}{莱布尼茨法则}{}
  设 \(\alpha\) 是 \(p\)-形式，\(\beta\) 是 \(q\)-形式，且它们的系数光滑，则
\[
d(\alpha\wedge\beta)=d\alpha\wedge\beta+(-1)^p\alpha\wedge d\beta
\]
\end{theorem}
莱布尼茨法则是保证外微分与楔积结构相容且满足庞加莱引理的必要条件。

仔细观察不同微分形式的外微分，把它们和刚刚的梯度，散度，旋度作对比，我们可以看到一个很整齐的对应关系。\(df\) 的三个系数正是 \(\nabla f\) 的三个分量。\(1\)-形式外微分后的三个系数正是 \(\nabla\times\vec{F}\) 的三个分量。\(2\)-形式外微分后的系数正是 \(\nabla\cdot\vec{F}\)。换句话说，\(0\)-形式的外微分是梯度，\(1\)-形式的外微分是旋度，\(2\)-形式的外微分是散度。

这说明，\(0\)-形式的外微分给出梯度对应的 \(1\)-形式，\(1\)-形式的外微分给出旋度对应的 \(2\)-形式，\(2\)-形式的外微分给出散度对应的 \(3\)-形式。而因为没有\(4\)-形式，所以\(3\)-形式的外微分就是\(0\)。也就是说，外微分把梯度，散度，旋度统一到了同一个运算框架中。

同时，按照这样的定义的外微分还满足我们刚刚提到的庞加莱引理和莱布尼茨法则。或者说，在我们对外微分的定义的要求下，只有这样定义，才能满足我们的需求。也正是因为这样定义，把微积分中不同的运算联系起来了。

由庞加莱引理，我们甚至能证明两条场论基本恒等式。

\begin{myexample}
  证明：梯度的旋度恒为零。

  取一个 \(0\)-形式 \(f\)，先求一次外微分，得到 \(df\)。它变成\(1\)-形式，对应梯度场 \(\nabla f\)。再求一次外微分，得到 \(d(df)\)，它变成\(2\)-形式，对应旋度场 \(\nabla\times(\nabla f)\)。

由庞加莱引理，可得\(d(df)=0\)。

所以\(\nabla\times(\nabla f)=\vec{0}\)。即梯度的旋度恒为零。
\end{myexample}

\begin{myexample}
  证明：旋度的散度恒为零。

  取一个 \(1\)-形式 \(\omega\)，先求一次外微分，得到 \(d\omega\)，它变成\(2\)-形式，对应旋度场 \(\nabla\times\vec{F}\)。再求一次外微分，得到 \(d(d\omega)\)，它变成\(3\)-形式，对应散度场 \(\nabla\cdot(\nabla\times\vec{F})\)。

  由庞加莱引理，可得\(d(d\omega)=0\)。

所以\(\nabla\cdot(\nabla\times\vec{F})=0\)。即旋度的散度恒为零。
\end{myexample}

我们来看两个例子，具体求一下外微分。

\begin{myexample}
设 \(\omega=xy\,dx+yz\,dy+zx\,dz\)，求 \(d\omega\)。

这里 \(P=xy,\ Q=yz,\ R=zx\)，代入 \(1\)-形式外微分展开式，得到：
\[
d\omega=\left(\frac{\partial R}{\partial y}-\frac{\partial Q}{\partial z}\right)dy\wedge dz
+\left(\frac{\partial P}{\partial z}-\frac{\partial R}{\partial x}\right)dz\wedge dx
+\left(\frac{\partial Q}{\partial x}-\frac{\partial P}{\partial y}\right)dx\wedge dy
\]

计算偏导数，得到：
\[
d\omega=-y\,dy\wedge dz-z\,dz\wedge dx-x\,dx\wedge dy
\]
\end{myexample}

\begin{myexample}
设 \(\omega=xy\,dy\wedge dz+yz\,dz\wedge dx+zx\,dx\wedge dy\)，求 \(d\omega\)。

这里 \(P=xy,\ Q=yz,\ R=zx\)，代入 \(2\)-形式外微分展开式，得到：
\[
d\omega=\left(\frac{\partial (xy)}{\partial x}
+\frac{\partial (yz)}{\partial y}
+\frac{\partial (zx)}{\partial z}\right)dx\wedge dy\wedge dz
\]

计算偏导数，得到：
\[
d\omega=(y+z+x)\,dx\wedge dy\wedge dz
\]
\end{myexample}

实际上，求不同形式的外微分，其实就是分别求梯度，散度，旋度。

有了外微分之后，格林公式，高斯公式和斯托克斯公式中的“内部导数”就有了统一的表达。
\section{广义斯托克斯定理}
在引入微分形式和外微分之后，我们可以开始尝试把格林公式、高斯公式和斯托克斯公式统一成一个公式了。

\subsection{斯托克斯定理的统一形式}
前面我们已经看到，格林公式、高斯公式和斯托克斯公式虽然研究的对象不同，但都可以写成“边界上的积分等于内部某种微分运算的积分”。在引入微分形式和外微分之后，这种共同结构可以用一个非常简洁的公式表达出来。

先说明边界与定向。设 \(M\) 是空间中的一段曲线、一块曲面或一个空间区域，把 \(M\) 的边界记作 \(\partial M\)，其中\(\partial\)称为边界算子。这里的边界不是简单的一个集合，还需要带上与 \(M\) 相容的定向。

\begin{definition}{边界与定向}{}
  我们规定，区域\(M\)的边界记作 \(\partial M\)，其方向为：
  
  若 \(M=[a,b]\) 是区间，则 \(\partial M=\{b\}-\{a\}\)，即右端点取正号，左端点取负号。

  若 \(M\) 是平面区域，则 \(\partial M\) 取正向边界曲线。

  若 \(M\) 是曲面，则 \(\partial M\) 是其边界曲线，方向由曲面的定向按右手法则确定。

  若 \(M\) 是空间区域，则 \(\partial M\) 取外侧边界曲面。
\end{definition}

根据微分形式和外微分，我们就可以写出广义斯托克斯定理了。

\begin{theorem}{广义斯托克斯定理}{}
设 \(M\) 是三维空间中具有分段光滑边界的 \((k+1)\) 维有向区域，若 \(\omega\) 是定义在 \(M\) 及其边界上的 \(k\)-形式，且系数具有一阶连续偏导数，则
\[
\int_{\partial M}\omega=\int_M d\omega
\]
\end{theorem}

这个公式的含义是微分形式 \(\omega\) 在边界上的积分，等于它的外微分 \(d\omega\) 在区域内部的积分。外微分把 \(0\)-形式、\(1\)-形式、\(2\)-形式分别与梯度、旋度、散度对应起来，因此格林公式、高斯公式和斯托克斯公式都可以看成这个公式在不同阶数下的具体表现。

这条定理的左端是 \(k\)-形式在 \(k\) 维边界上的积分，右端是外微分 \(d\omega\) 在 \(k+1\) 维区域上的积分。外微分把 \(k\)-形式提升为 \(k+1\)-形式，而边界算子 \(\partial\) 把 \(k+1\) 维区域降低为 \(k\) 维对象，二者的次数变化正好相反。

从形式上看，广义斯托克斯定理也解释了为什么外微分要定义成前面那样。\(d\omega\) 的阶数比 \(\omega\) 高一阶，而 \(M\) 的维数比 \(\partial M\) 高一维，所以等式两边的积分号正好匹配。反过来，如果外微分的定义稍有变化，就无法同时容纳格林公式、高斯公式和斯托克斯公式。

要证明它，依旧是“分割，近似，求和，取极限”的思想。把区域 \(M\) 剖分成许多小网格，在每个小网格上应用微积分基本定理，再把各网格上的结果相加。此时相邻网格的公共边界上的积分相互抵消，最后只剩下 \(\partial M\) 上的积分。广义斯托克斯定理正是这一证明过程的抽象写法。

于是，微积分中原本分散的几条积分公式，就这样被统一了，而且统一的形式极其优雅。外微分 \(d\) 负责描述局部变化，边界算子 \(\partial\) 负责划定整体轮廓。我们从最朴素的牛顿-莱布尼茨公式出发，经过格林公式、高斯公式和斯托克斯公式，最终在微分形式与外微分的语言下，看到了它们共同的源头。广义斯托克斯定理把区间上的微积分基本定理、平面上的格林公式、空间中的斯托克斯公式和高斯公式全部纳入其中。

广义斯托克斯定理告诉我们，边界上的量，恰好等于内部某种微分量的累计。它不依赖具体的维数，也不在意积分的具体名称，只用一个公式就概括了微积分中最重要的积分定理。我们之前学习的积分公式，其实并不是孤立的，而是同一座山峰在不同高度看到的不同侧面。

其实微积分一直以来都在诉说着同一件事，那就是边界与内部的关系、局部与整体的关系、微分与积分的关系，它们从来就是不可分割的。
\subsection{斯托克斯定理的不同形式}
下面把广义斯托克斯定理逐一带入 \(k=0,1,2\)，看它如何回到我们熟悉的各种积分公式。

\begin{myexample}
  设 \(M=[a,b]\)，\(\omega=f(x)\) 是 \(0\)-形式，代入广义斯托克斯定理，求具体形式。
  
  由于\(M=[a,b]\)，\(\omega=f(x)\) 是 \(0\)-形式，此时：
\[
\partial M=\{b\}-\{a\},\qquad d\omega=f'(x)\,dx
\]

于是广义斯托克斯定理给出：
\[
\int_{\partial M}\omega=f(b)-f(a),
\qquad
\int_M d\omega=\int_a^b f'(x)\,dx
\]

因此：
\[
\int_a^b f'(x)\,dx=f(b)-f(a)
\]

这正是牛顿-莱布尼茨公式。由此可见，微积分基本定理也可以纳入广义斯托克斯定理的框架。
\end{myexample}

\begin{myexample}
  设 \(M=D\) 是平面有界闭区域，\(\omega=P\,dx+Q\,dy\) 是 \(1\)-形式，代入广义斯托克斯定理，求具体形式。
  
  由于\(M=D\) 是平面有界闭区域，\(\omega=P\,dx+Q\,dy\) 是 \(1\)-形式，其外微分为：
\[
d\omega
=
\left(\frac{\partial Q}{\partial x}-\frac{\partial P}{\partial y}\right)dx\wedge dy
\]

此时 \(\partial D\) 是 \(D\) 的正向边界曲线 \(L\)，所以：
\[
\int_{\partial D}\omega=\oint_L P\,dx+Q\,dy
\]

由微分形式：
\[
\int_D d\omega
=
\iint_D
\left(\frac{\partial Q}{\partial x}-\frac{\partial P}{\partial y}\right)dx\,dy
\]

由广义斯托克斯定理得：
\[
\oint_L P\,dx+Q\,dy
=
\iint_D
\left(\frac{\partial Q}{\partial x}-\frac{\partial P}{\partial y}\right)dx\,dy
\]

这正是格林公式。
\end{myexample}


\begin{myexample}
  设 \(M=\Sigma\) 是空间中的有向曲面，\(\omega=P\,dx+Q\,dy+R\,dz\) 是 \(1\)-形式，代入广义斯托克斯定理，求具体形式。
  
  由于\(M=\Sigma\) 是空间中的有向曲面，\(\omega=P\,dx+Q\,dy+R\,dz\) 是 \(1\)-形式，其外微分为：
\[
d\omega
=
\left(\frac{\partial R}{\partial y}-\frac{\partial Q}{\partial z}\right)dy\wedge dz
+
\left(\frac{\partial P}{\partial z}-\frac{\partial R}{\partial x}\right)dz\wedge dx
+
\left(\frac{\partial Q}{\partial x}-\frac{\partial P}{\partial y}\right)dx\wedge dy
\]

设 \(\partial\Sigma=\Gamma\) 是 \(\Sigma\) 的有向边界曲线，则由广义斯托克斯定理得：
\[
\int_\Gamma P\,dx+Q\,dy+R\,dz
=
\iint_\Sigma
\left[
\left(\frac{\partial R}{\partial y}-\frac{\partial Q}{\partial z}\right)dy\,dz
+
\left(\frac{\partial P}{\partial z}-\frac{\partial R}{\partial x}\right)dz\,dx
+
\left(\frac{\partial Q}{\partial x}-\frac{\partial P}{\partial y}\right)dx\,dy
\right]
\]

若记 \(\vec F=(P,Q,R)\)，右端正是旋度的通量：
\[
\iint_\Sigma
\operatorname{rot}\vec F\cdot \vec n\,dS
\]

因此上式就是斯托克斯公式。
\end{myexample}

\begin{myexample}
  设 \(M=V\) 是空间中的有界闭区域，\(\omega=P\,dy\wedge dz+Q\,dz\wedge dx+R\,dx\wedge dy\) 是 \(2\)-形式，代入广义斯托克斯定理，求具体形式。
  
  由于 \(M=V\) 是空间中的有界闭区域，\(\omega=P\,dy\wedge dz+Q\,dz\wedge dx+R\,dx\wedge dy\) 是 \(2\)-形式，其外微分为：
\[
d\omega
=
\left(\frac{\partial P}{\partial x}
+\frac{\partial Q}{\partial y}
+\frac{\partial R}{\partial z}\right)dx\wedge dy\wedge dz
\]

此时 \(\partial V\) 是 \(V\) 的外侧边界曲面 \(S\)，于是：
\[
\int_{\partial V}\omega
=
\oiint_S P\,dy\,dz+Q\,dz\,dx+R\,dx\,dy
\]

由微分形式：
\[
\int_V d\omega
=
\iiint_V
\left(\frac{\partial P}{\partial x}
+\frac{\partial Q}{\partial y}
+\frac{\partial R}{\partial z}\right)dx\,dy\,dz
\]

由广义斯托克斯定理得：
\[
\oiint_S P\,dy\,dz+Q\,dz\,dx+R\,dx\,dy
=
\iiint_V
\left(\frac{\partial P}{\partial x}
+\frac{\partial Q}{\partial y}
+\frac{\partial R}{\partial z}\right)dx\,dy\,dz
\]

这正是高斯公式。

\end{myexample}


由此可见，牛顿-莱布尼茨公式、格林公式、斯托克斯公式和高斯公式并不是四个彼此孤立的结果，而是同一个广义斯托克斯定理在不同维数下的具体形式。微分形式和外微分提供了一套统一的语言，使得边界与内部、积分与微分之间的关系得到了最简洁的表达。


\backmatter
\chapter{后记}

本书是从“极限”开始的。

那时我们问，当一点无限接近另一点时，函数值会怎样？这个问题后来有了一个朴素的名字——极限。再后来，这个问题发展成导数，发展成微分，发展成积分，发展成多元函数、曲线积分、曲面积分，最后走到了广义斯托克斯定理。回头看，起点和终点其实说的是同一件事。

格林公式、高斯公式、斯托克斯公式，曾经像三座各自独立的山峰。它们一个联系平面区域与边界曲线，一个联系空间区域与边界曲面，一个联系曲面与其边界。但当它们被统一成广义斯托克斯公式时，三座山峰忽然连成了一片山脉。牛顿-莱布尼茨公式是它的一维形态，格林公式是它的平面形态，斯托克斯公式是它的曲面形态，高斯公式是它的空间形态。它揭示了一种统一性，就是局部变化能够决定整体行为，边界上的量恰好等于内部某种微分量的累计。

微积分不只是一种算数方法，它告诉我们，复杂的变化可以从局部加以理解，而局部的规则又会在整体边界上显现出来。科学中的许多基本定律，本质上都是这一结构的翻版，比如一个区域内的变化率，等于穿过其边界向外流出的量。一个系统内部的守恒，表现为其边界上的净通量。

数学中最值得记住的，不是某一条定理的具体形式，而是它背后那个简单而统一的思想。格林公式、高斯公式、斯托克斯公式，如果只当作三个工具，可能会很快遗忘。但如果看见它们是同一公式在不同维度上的投影，它们便会长久地留在心中。

当然，本书并不完美。有些证明为了可读性做了简化，有些内容可以更深入，有些例题还可以更丰富，甚至某些描述或者计算还有纰漏。但这些遗憾，或许正说明写作本身也像微积分一样：总有一些变化尚未完全展开，总有一些边界尚未最终封闭。它不完美，但它尽力把一条清晰的路径呈现了出来。如果你能在最后一章回头看见我们的来时路，那么本书的目的就达到了。

最后，经历了130小时的编写，我想感谢DeepSeek，没有AI的帮助，本书不可能出现。

我想说，如果你看到这里，只能记住一个公式，那么请记住它：\(\int_{\partial M}\omega=\int_M d\omega\)。

如果只能记住一个思想，那么请带走它：边界发生的一切，由内部决定。

这或许就是微积分带给我们的最重要的一件事。

\end{document}
